\documentclass[11pt,a4paper]{article}
\usepackage[utf8]{inputenc}
\usepackage[vietnamese,english]{babel}
\usepackage{amsmath,amssymb}
\usepackage{geometry}
\geometry{a4paper, margin=2cm}
\usepackage{enumitem}
\usepackage{xcolor}
\usepackage{hyperref}

\hypersetup{
    colorlinks=true,
    linkcolor=blue,
    urlcolor=blue
}

\title{\textbf{CSI106 - Foundations of Computer Science}\\\Large Full Question Bank with Answers and Explanations}
\author{Converted via image-to-text-local}
\date{\today}

\begin{document}
\maketitle
\tableofcontents

\vspace{0.8cm}
\hrule
\vspace{0.5cm}


\newpage
\section*{1. Computer Organization}
\addcontentsline{toc}{section}{1. Computer Organization}
\hrule\vspace{0.5cm}

\begin{minipage}{\textwidth}
\subsection*{Question 1 \quad \footnotesize{[0001\_q001\_Question\_1.png]}}
\noindent \textbf{The personnel who deals with the computer \& its management put together are called}

\begin{enumerate}[label=\Alph*.]
    \item Software
    \item \textbf{Humanware} \hfill \textcolor{blue}{\textbf{[Correct]}}
    \item Firmware
    \item Hardware
    \item None of the others
\end{enumerate}

\noindent \textbf{Correct Answer:} \textcolor{blue}{\textbf{B}} \\[2pt]
\noindent \textbf{Explanation:} Mô hình Turing là nền tảng lý thuyết của máy tính. Alan Turing đề xuất máy Turing (TM) như một mô hình toán học có khả năng thực hiện bất kỳ thuật toán nào. Máy Turing gồm: bộ điều khiển hữu hạn trạng thái, băng nhớ vô hạn, và đầu đọc/ghi. Von Neumann sau đó đã phát triển mô hình này thành kiến trúc máy tính lưu trữ chương trình (stored-program concept).

\par\vspace{0.3cm}
\hrule
\vspace{0.3cm}
\end{minipage}

\begin{minipage}{\textwidth}
\subsection*{Question 2 \quad \footnotesize{[0002\_q002\_Question\_2.png]}}
\noindent \textbf{FORTRAN and COBOL are examples of}

\begin{enumerate}[label=\Alph*.]
    \item Hardware
    \item Operating systems
    \item \textbf{Computer languages} \hfill \textcolor{blue}{\textbf{[Correct]}}
    \item Algorithms
\end{enumerate}

\noindent \textbf{Correct Answer:} \textcolor{blue}{\textbf{C}} \\[2pt]
\noindent \textbf{Explanation:} Kiến trúc Von Neumann gồm 4 thành phần chính: CPU (Central Processing Unit - đơn vị xử lý trung tâm), bộ nhớ (memory), bus kết nối, và thiết bị I/O (input/output). CPU chứa ALU (Arithmetic Logic Unit - đơn vị tính toán số học và logic) và CU (Control Unit - đơn vị điều khiển). Bộ nhớ lưu trữ cả dữ liệu và chương trình.

\par\vspace{0.3cm}
\hrule
\vspace{0.3cm}
\end{minipage}

\begin{minipage}{\textwidth}
\subsection*{Question 3 \quad \footnotesize{[0003\_q003\_Question\_3.png]}}
\noindent \textbf{is a stand-alone storage location that holds data temporarily.}

\begin{enumerate}[label=\Alph*.]
    \item An ALU
    \item \textbf{A register} \hfill \textcolor{blue}{\textbf{[Correct]}}
    \item A control unit
    \item A tape drive
\end{enumerate}

\noindent \textbf{Correct Answer:} \textcolor{blue}{\textbf{B}} \\[2pt]
\noindent \textbf{Explanation:} Trong kiến trúc máy tính, CPU (Central Processing Unit) là 'bộ não' của máy tính. CPU thực thi các lệnh (instructions) được lưu trong bộ nhớ. Chu kỳ thực thi lệnh (machine cycle) gồm: Fetch (lấy lệnh từ bộ nhớ) → Decode (giải mã lệnh) → Execute (thực thi) → Store (lưu kết quả).

\par\vspace{0.3cm}
\hrule
\vspace{0.3cm}
\end{minipage}

\begin{minipage}{\textwidth}
\subsection*{Question 4 \quad \footnotesize{[0004\_q004\_Question\_4.png]}}
\noindent \textbf{What is the 80-20 rule in cache?}

\begin{enumerate}[label=\Alph*.]
    \item \textbf{Most computers typically spend 80 percent of their time accessing only 20 percent of the data} \hfill \textcolor{blue}{\textbf{[Correct]}}
    \item Cache memory at any time contains a copy of a portion of main memory.
    \item It can hold data 20 percent to make access faster at least 80 percent of the time.
    \item The CPU need to checks the cache first.
\end{enumerate}

\noindent \textbf{Correct Answer:} \textcolor{blue}{\textbf{A}} \\[2pt]
\noindent \textbf{Explanation:} Bộ nhớ máy tính được tổ chức theo thứ bậc (memory hierarchy): thanh ghi (registers) trong CPU → bộ nhớ cache → RAM (bộ nhớ chính) → bộ nhớ phụ (ổ cứng, SSD). Tốc độ truy xuất giảm dần nhưng dung lượng tăng dần từ trên xuống dưới. Cache là bộ nhớ nhỏ nhưng cực nhanh, lưu dữ liệu thường dùng.

\par\vspace{0.3cm}
\hrule
\vspace{0.3cm}
\end{minipage}

\begin{minipage}{\textwidth}
\subsection*{Question 5 \quad \footnotesize{[0005\_q005\_Question\_5.png]}}
\noindent \textbf{A monitor's \_\_ is the distance between the holes in the mask behind the screen. This helps determine how sharp the dots appear.}

\begin{enumerate}[label=\Alph*.]
    \item Refresh rate
    \item \textbf{Dot pitch} \hfill \textcolor{blue}{\textbf{[Correct]}}
    \item Resolution
    \item Color depth
    \item None of the others
\end{enumerate}

\noindent \textbf{Correct Answer:} \textcolor{blue}{\textbf{B}} \\[2pt]
\noindent \textbf{Explanation:} Bus hệ thống là tập hợp các đường dây vật lý kết nối các thành phần của máy tính. Có 3 loại bus chính: data bus (truyền dữ liệu), address bus (truyền địa chỉ), và control bus (truyền tín hiệu điều khiển). Bus được đo bằng độ rộng (số bit truyền đồng thời) và tốc độ (MHz).

\par\vspace{0.3cm}
\hrule
\vspace{0.3cm}
\end{minipage}

\begin{minipage}{\textwidth}
\subsection*{Question 6 \quad \footnotesize{[0006\_q006\_Question\_6.png]}}
\noindent \textbf{What is taxonomy of computer organization?}

\begin{enumerate}[label=\Alph*.]
    \item SIMD - Single Instruction-stream, Multiple Data-stream
    \item MISD - Multiple Instruction-stream, Single Data-stream
    \item MIMD - Multiple Instruction-stream, Multiple Data-stream
    \item \textbf{All of others} \hfill \textcolor{blue}{\textbf{[Correct]}}
\end{enumerate}

\noindent \textbf{Correct Answer:} \textcolor{blue}{\textbf{D}} \\[2pt]
\noindent \textbf{Explanation:} Thiết bị I/O (Input/Output) gồm thiết bị nhập (bàn phím, chuột, máy quét) và thiết bị xuất (màn hình, máy in, loa). Các thiết bị này kết nối với CPU qua các cổng I/O hoặc bus. Trong mô hình Von Neumann, I/O được quản lý bởi CU thông qua các lệnh I/O đặc biệt.

\par\vspace{0.3cm}
\hrule
\vspace{0.3cm}
\end{minipage}

\begin{minipage}{\textwidth}
\subsection*{Question 7 \quad \footnotesize{[0007\_q007\_Question\_7.png]}}
\noindent \textbf{Instructions and memory addresses are represented by}

\begin{enumerate}[label=\Alph*.]
    \item Character codes
    \item \textbf{Binary codes} \hfill \textcolor{blue}{\textbf{[Correct]}}
    \item Binary word
    \item Parity bit
    \item None of the others
\end{enumerate}

\noindent \textbf{Correct Answer:} \textcolor{blue}{\textbf{B}} \\[2pt]
\noindent \textbf{Explanation:} ALU (Arithmetic Logic Unit) là phần của CPU thực hiện các phép tính số học (cộng, trừ, nhân, chia) và các phép logic (AND, OR, NOT, XOR). ALU nhận dữ liệu từ các thanh ghi, xử lý và trả kết quả về thanh ghi accumulator (ACC).

\par\vspace{0.3cm}
\hrule
\vspace{0.3cm}
\end{minipage}

\begin{minipage}{\textwidth}
\subsection*{Question 8 \quad \footnotesize{[0008\_q008\_Question\_8.png]}}
\noindent \textbf{The base of the hexadecimal number system is}

\begin{enumerate}[label=\Alph*.]
    \item 2
    \item 8
    \item 10
    \item \textbf{16} \hfill \textcolor{blue}{\textbf{[Correct]}}
\end{enumerate}

\noindent \textbf{Correct Answer:} \textcolor{blue}{\textbf{D}} \\[2pt]
\noindent \textbf{Explanation:} Thanh ghi (register) là bộ nhớ nhỏ, cực nhanh nằm trong CPU. Các thanh ghi quan trọng: MAR (Memory Address Register - chứa địa chỉ bộ nhớ cần truy xuất), MDR (Memory Data Register - chứa dữ liệu đọc/ghi), PC (Program Counter - chứa địa chỉ lệnh tiếp theo), IR (Instruction Register - chứa lệnh đang thực thi).

\par\vspace{0.3cm}
\hrule
\vspace{0.3cm}
\end{minipage}

\begin{minipage}{\textwidth}
\subsection*{Question 9 \quad \footnotesize{[0009\_q009\_Question\_9.png]}}
\noindent \textbf{Which of the following is equivalent to 12 in decimal?}

\begin{enumerate}[label=\Alph*.]
    \item (1110)2
    \item \textbf{(C)16} \hfill \textcolor{blue}{\textbf{[Correct]}}
    \item (15)8
    \item None of the others
\end{enumerate}

\noindent \textbf{Correct Answer:} \textcolor{blue}{\textbf{B}} \\[2pt]
\noindent \textbf{Explanation:} Chu kỳ lệnh (instruction cycle) là quy trình CPU thực thi một lệnh: (1) Fetch: CPU lấy lệnh từ bộ nhớ tại địa chỉ trong PC vào IR, (2) Decode: CU giải mã lệnh trong IR, (3) Execute: ALU thực thi lệnh, (4) PC tăng lên để trỏ đến lệnh tiếp theo. Mỗi bước được đồng bộ bởi xung clock.

\par\vspace{0.3cm}
\hrule
\vspace{0.3cm}
\end{minipage}

\begin{minipage}{\textwidth}
\subsection*{Question 10 \quad \footnotesize{[0010\_q010\_Question\_10.png]}}
\noindent \textbf{The maximum value of an octal integer with K=2 digits is Nmax = ?}

\begin{enumerate}[label=\Alph*.]
    \item \textbf{63} \hfill \textcolor{blue}{\textbf{[Correct]}}
    \item 64
    \item 65
    \item 66
    \item None of the others
\end{enumerate}

\noindent \textbf{Correct Answer:} \textcolor{blue}{\textbf{A}} \\[2pt]
\noindent \textbf{Explanation:} Clock CPU là bộ tạo xung nhịp đồng bộ hóa mọi hoạt động của CPU. Tốc độ clock đo bằng Hz (số chu kỳ/giây). Mỗi lệnh cần một số chu kỳ clock nhất định (CPI - Cycles Per Instruction). Hiệu năng CPU = (Số lệnh × CPI) / Tốc độ clock.

\par\vspace{0.3cm}
\hrule
\vspace{0.3cm}
\end{minipage}

\begin{minipage}{\textwidth}
\subsection*{Question 11 \quad \footnotesize{[0011\_q011\_Question\_11.png]}}
\noindent \textbf{Conversion of binary number 1010101000010111 to hexadecimal number is}

\begin{enumerate}[label=\Alph*.]
    \item D8F9
    \item A8B9
    \item \textbf{AA17} \hfill \textcolor{blue}{\textbf{[Correct]}}
    \item D9F8
    \item None of the others
\end{enumerate}

\noindent \textbf{Correct Answer:} \textcolor{blue}{\textbf{C}} \\[2pt]
\noindent \textbf{Explanation:} RISC (Reduced Instruction Set Computer) là kiến trúc CPU với tập lệnh nhỏ gọn, mỗi lệnh thực thi trong 1 chu kỳ clock. CISC (Complex Instruction Set Computer) có tập lệnh phức tạp, một lệnh có thể làm nhiều việc nhưng cần nhiều chu kỳ. Ví dụ: ARM là RISC, Intel x86 là CISC.

\par\vspace{0.3cm}
\hrule
\vspace{0.3cm}
\end{minipage}

\begin{minipage}{\textwidth}
\subsection*{Question 12 \quad \footnotesize{[0012\_q012\_Question\_12.png]}}
\noindent \textbf{The precision of the fractional part of a number stored in a computer is defined by the}

\begin{enumerate}[label=\Alph*.]
    \item Sign
    \item Exponent
    \item \textbf{Mantissa} \hfill \textcolor{blue}{\textbf{[Correct]}}
    \item Last digit
\end{enumerate}

\noindent \textbf{Correct Answer:} \textcolor{blue}{\textbf{C}} \\[2pt]
\noindent \textbf{Explanation:} Bộ nhớ cache là bộ nhớ đệm nhỏ nhưng cực nhanh, nằm giữa CPU và RAM. Cache lưu bản sao của dữ liệu thường xuyên truy cập từ RAM, giúp giảm thời gian chờ CPU. Nguyên lý locality (tính cục bộ): dữ liệu vừa dùng có khả năng cao sẽ được dùng lại (temporal locality) hoặc dữ liệu lân cận sẽ được dùng (spatial locality).

\par\vspace{0.3cm}
\hrule
\vspace{0.3cm}
\end{minipage}

\begin{minipage}{\textwidth}
\subsection*{Question 13 \quad \footnotesize{[0013\_q013\_Question\_13.png]}}
\noindent \textbf{Which of the following is the coding of data so that it can't be easily understood if intercepted?}

\begin{enumerate}[label=\Alph*.]
    \item Barcode
    \item Decoder
    \item \textbf{Encryption} \hfill \textcolor{blue}{\textbf{[Correct]}}
    \item Memonics
    \item None of the others
\end{enumerate}

\noindent \textbf{Correct Answer:} \textcolor{blue}{\textbf{C}} \\[2pt]
\noindent \textbf{Explanation:} Trong kiến trúc Harvard, bộ nhớ lệnh và bộ nhớ dữ liệu được tách biệt, mỗi loại có bus riêng. Kiến trúc Von Neumann dùng chung một bộ nhớ cho cả lệnh và dữ liệu. Harvard nhanh hơn vì có thể đọc lệnh và dữ liệu đồng thời, nhưng phức tạp hơn.

\par\vspace{0.3cm}
\hrule
\vspace{0.3cm}
\end{minipage}

\begin{minipage}{\textwidth}
\subsection*{Question 14 \quad \footnotesize{[0014\_q014\_Question\_14.png]}}
\noindent \textbf{If we call the bit depth or number of bits per sample B, the number of samples per second S, so how to calculate BIT RATE (R)?}

\begin{enumerate}[label=\Alph*.]
    \item B × 2S
    \item 2B ×S
    \item \textbf{B × S} \hfill \textcolor{blue}{\textbf{[Correct]}}
    \item B2 × S2
\end{enumerate}

\noindent \textbf{Correct Answer:} \textcolor{blue}{\textbf{C}} \\[2pt]
\noindent \textbf{Explanation:} Pipelining (xử lý ống dẫn) là kỹ thuật CPU thực thi nhiều lệnh đồng thời ở các giai đoạn khác nhau. Ví dụ: trong khi lệnh 1 đang Execute, lệnh 2 được Decode, lệnh 3 được Fetch. Pipelining làm tăng throughput (thông lượng) của CPU nhưng có thể gặp hazard (xung đột dữ liệu/điều khiển).

\par\vspace{0.3cm}
\hrule
\vspace{0.3cm}
\end{minipage}

\begin{minipage}{\textwidth}
\subsection*{Question 15 \quad \footnotesize{[0015\_q015\_Question\_15.png]}}
\noindent \textbf{(which stands for picture elements). In the storing images, the samples are called}

\begin{enumerate}[label=\Alph*.]
    \item Bit map
    \item \textbf{Pixels} \hfill \textcolor{blue}{\textbf{[Correct]}}
    \item Resolution
    \item Color depth
    \item None of others
\end{enumerate}

\noindent \textbf{Correct Answer:} \textcolor{blue}{\textbf{B}} \\[2pt]
\noindent \textbf{Explanation:} DMA (Direct Memory Access) là cơ chế cho phép thiết bị I/O truy xuất trực tiếp bộ nhớ mà không cần CPU xử lý từng byte. Bộ điều khiển DMA quản lý việc truyền dữ liệu, giải phóng CPU để làm việc khác. DMA rất quan trọng cho các thiết bị tốc độ cao như ổ cứng, card mạng.

\par\vspace{0.3cm}
\hrule
\vspace{0.3cm}
\end{minipage}

\begin{minipage}{\textwidth}
\subsection*{Question 16 \quad \footnotesize{[0016\_q016\_Question\_16.png]}}
\noindent \textbf{Use a arithmetic right shift operation on the bit pattern 1001 1000}

\begin{enumerate}[label=\Alph*.]
    \item 0100 1100
    \item \textbf{1100 1100} \hfill \textcolor{blue}{\textbf{[Correct]}}
    \item 1101 1001
    \item 1001 1000
\end{enumerate}

\noindent \textbf{Correct Answer:} \textcolor{blue}{\textbf{B}} \\[2pt]
\noindent \textbf{Explanation:} Interrupt (ngắt) là tín hiệu từ thiết bị I/O báo cho CPU rằng có sự kiện cần xử lý. Khi nhận interrupt, CPU tạm dừng chương trình hiện tại, lưu trạng thái vào stack, xử lý interrupt handler, sau đó quay lại chương trình. Interrupt giúp CPU không phải chờ I/O (polling).

\par\vspace{0.3cm}
\hrule
\vspace{0.3cm}
\end{minipage}

\begin{minipage}{\textwidth}
\subsection*{Question 17 \quad \footnotesize{[0017\_q017\_Question\_17.png]}}
\noindent \textbf{is usually privately owned, limited in size and connects some hosts in a single office, building, or campus.}

\begin{enumerate}[label=\Alph*.]
    \item \textbf{LAN} \hfill \textcolor{blue}{\textbf{[Correct]}}
    \item WAN
    \item Internet
    \item Network
\end{enumerate}

\noindent \textbf{Correct Answer:} \textcolor{blue}{\textbf{A}} \\[2pt]
\noindent \textbf{Explanation:} Bộ nhớ chính (RAM - Random Access Memory) là bộ nhớ khả biến (volatile), mất dữ liệu khi mất điện. RAM được tổ chức thành các ô nhớ (cells), mỗi ô có địa chỉ duy nhất. Thời gian truy xuất RAM không phụ thuộc vào vị trí địa chỉ (random access). SRAM (static) nhanh hơn DRAM (dynamic).

\par\vspace{0.3cm}
\hrule
\vspace{0.3cm}
\end{minipage}

\begin{minipage}{\textwidth}
\subsection*{Question 18 \quad \footnotesize{[0018\_q018\_Question\_18.png]}}
\noindent \textbf{In physical layer, a \_can be broadly defined as anything that can carry information from a source to a destination.}

\begin{enumerate}[label=\Alph*.]
    \item Digital-to-digital conversion
    \item Analog-to-digital conversion
    \item Digital-to-analog conversion
    \item Analog-to-analog conversion
    \item \textbf{Transmission medium} \hfill \textcolor{blue}{\textbf{[Correct]}}
\end{enumerate}

\noindent \textbf{Correct Answer:} \textcolor{blue}{\textbf{E}} \\[2pt]
\noindent \textbf{Explanation:} ROM (Read-Only Memory) là bộ nhớ chỉ đọc, không mất dữ liệu khi mất điện (non-volatile). ROM chứa firmware - chương trình khởi động (bootloader) và BIOS. Các loại ROM: PROM (lập trình một lần), EPROM (xóa bằng tia UV), EEPROM (xóa bằng điện), Flash memory (có thể xóa/ghi nhiều lần).

\par\vspace{0.3cm}
\hrule
\vspace{0.3cm}
\end{minipage}

\begin{minipage}{\textwidth}
\subsection*{Question 19 \quad \footnotesize{[0019\_q019\_Question\_19.png]}}
\noindent \textbf{The layer of the TCP/IP protocol suite is responsible for source-to-destination delivery of the entire message.}

\begin{enumerate}[label=\Alph*.]
    \item \textbf{Transport} \hfill \textcolor{blue}{\textbf{[Correct]}}
    \item Network
    \item Data-link
    \item Session
\end{enumerate}

\noindent \textbf{Correct Answer:} \textcolor{blue}{\textbf{A}} \\[2pt]
\noindent \textbf{Explanation:} Đơn vị đo thông tin nhỏ nhất là bit (binary digit). 1 byte = 8 bits. 1 KB (Kilobyte) = 2\textasciicircum{}10 = 1024 bytes. 1 MB = 1024 KB. 1 GB = 1024 MB. 1 TB = 1024 GB. Tốc độ xử lý đo bằng Hz (số chu kỳ/giây): 1 MHz = 10\textasciicircum{}6 Hz, 1 GHz = 10\textasciicircum{}9 Hz.

\par\vspace{0.3cm}
\hrule
\vspace{0.3cm}
\end{minipage}

\begin{minipage}{\textwidth}
\subsection*{Question 20 \quad \footnotesize{[0020\_q020\_Question\_20.png]}}
\noindent \textbf{protocol (s) is one of the protocols in the transport layer.}

\begin{enumerate}[label=\Alph*.]
    \item Only TCP
    \item Only UDP
    \item Only SCTP
    \item \textbf{TCP, UDP, and SCTP} \hfill \textcolor{blue}{\textbf{[Correct]}}
\end{enumerate}

\noindent \textbf{Correct Answer:} \textcolor{blue}{\textbf{D}} \\[2pt]
\noindent \textbf{Explanation:} Máy tính Turing (Turing Machine - TM) là mô hình toán học gồm: (1) băng nhớ vô hạn chia thành các ô, (2) đầu đọc/ghi di chuyển trái/phải, (3) bộ điều khiển với tập trạng thái hữu hạn. TM có thể mô phỏng mọi thuật toán. Máy Turing phổ dụng (UTM) có thể mô phỏng mọi TM khác.

\par\vspace{0.3cm}
\hrule
\vspace{0.3cm}
\end{minipage}

\begin{minipage}{\textwidth}
\subsection*{Question 21 \quad \footnotesize{[0021\_q021\_Question\_21.png]}}
\noindent \textbf{In a \_\_, the directory system listing of the peers and what they offer uses the client-server paradigm, but the storing and downloading of the files are done using the peer-to-peer paradigm.}

\begin{enumerate}[label=\Alph*.]
    \item WWW
    \item \textbf{Centralized network} \hfill \textcolor{blue}{\textbf{[Correct]}}
    \item Hybrid network
    \item Decentralized network
\end{enumerate}

\noindent \textbf{Correct Answer:} \textcolor{blue}{\textbf{B}} \\[2pt]
\noindent \textbf{Explanation:} Chu kỳ Fetch-Decode-Execute: (1) Fetch: CPU lấy lệnh từ bộ nhớ tại địa chỉ PC → IR, (2) Decode: CU giải mã IR để xác định thao tác, (3) Execute: ALU thực thi, kết quả lưu vào ACC hoặc bộ nhớ. Đây là quy trình cơ bản của mọi CPU Von Neumann.

\par\vspace{0.3cm}
\hrule
\vspace{0.3cm}
\end{minipage}

\begin{minipage}{\textwidth}
\subsection*{Question 22 \quad \footnotesize{[0022\_q022\_Question\_22.png]}}
\noindent \textbf{The major design goals of an operating system:}

\begin{enumerate}[label=\Alph*.]
    \item \textbf{Efficient use of hardware and easy use of resources} \hfill \textcolor{blue}{\textbf{[Correct]}}
    \item \textbf{Control subsystems and manage memory} \hfill \textcolor{blue}{\textbf{[Correct]}}
    \item \textbf{Access I/O device and CPU} \hfill \textcolor{blue}{\textbf{[Correct]}}
    \item \textbf{Provide UI} \hfill \textcolor{blue}{\textbf{[Correct]}}
    \item Operation on data
    \item Text editor
\end{enumerate}

\noindent \textbf{Correct Answer:} \textcolor{blue}{\textbf{A, B, C, D}} \\[2pt]
\noindent \textbf{Explanation:} Các thành phần trong mô hình Von Neumann: (A) CPU gồm ALU và CU, (B) Bộ nhớ chứa dữ liệu và chương trình, (C) Bus kết nối các thành phần, (D) Thiết bị I/O giao tiếp với bên ngoài. Mô hình này là nền tảng của hầu hết máy tính ngày nay.

\par\vspace{0.3cm}
\hrule
\vspace{0.3cm}
\end{minipage}

\begin{minipage}{\textwidth}
\subsection*{Question 23 \quad \footnotesize{[0023\_q023\_Question\_23.png]}}
\noindent \textbf{A multi programming system is one that can}

\begin{enumerate}[label=\Alph*.]
    \item Run very fast
    \item \textbf{Share hardware resources with many programs simultaneously} \hfill \textcolor{blue}{\textbf{[Correct]}}
    \item Compute many programs simultaneously
    \item Use many operating systems
    \item None of the others
\end{enumerate}

\noindent \textbf{Correct Answer:} \textcolor{blue}{\textbf{B}} \\[2pt]
\noindent \textbf{Explanation:} Chu kỳ xung nhịp (clock cycle) là đơn vị thời gian cơ bản của CPU. Tần số clock càng cao, CPU càng thực thi nhiều lệnh mỗi giây. Tuy nhiên, hiệu năng còn phụ thuộc vào số chu kỳ mỗi lệnh (CPI) và số lệnh trong chương trình. Hiệu năng = 1 / (Số lệnh × CPI × Chu kỳ clock).

\par\vspace{0.3cm}
\hrule
\vspace{0.3cm}
\end{minipage}

\begin{minipage}{\textwidth}
\subsection*{Question 24 \quad \footnotesize{[0024\_q024\_Question\_24.png]}}
\noindent \textbf{In , only one program can reside in memory for execution.}

\begin{enumerate}[label=\Alph*.]
    \item \textbf{Mono-programming} \hfill \textcolor{blue}{\textbf{[Correct]}}
    \item Multi-programming
    \item Partitioning
    \item Paging
\end{enumerate}

\noindent \textbf{Correct Answer:} \textcolor{blue}{\textbf{A}} \\[2pt]
\noindent \textbf{Explanation:} Control Unit (CU) là bộ phận điều khiển của CPU, có nhiệm vụ: giải mã lệnh, tạo tín hiệu điều khiển cho ALU, bộ nhớ và I/O, quản lý thứ tự thực thi lệnh, xử lý ngắt. CU hoạt động theo chu kỳ xung nhịp và là 'nhạc trưởng' của CPU.

\par\vspace{0.3cm}
\hrule
\vspace{0.3cm}
\end{minipage}

\begin{minipage}{\textwidth}
\subsection*{Question 25 \quad \footnotesize{[0025\_q025\_Question\_25.png]}}
\noindent \textbf{The scheduler creates a process from a job and changes a process back to a job.}

\begin{enumerate}[label=\Alph*.]
    \item \textbf{Job} \hfill \textcolor{blue}{\textbf{[Correct]}}
    \item Process
    \item Virtual
    \item Queue
\end{enumerate}

\noindent \textbf{Correct Answer:} \textcolor{blue}{\textbf{A}} \\[2pt]
\noindent \textbf{Explanation:} Bộ nhớ máy tính có thứ bậc (hierarchy): Thanh ghi (register) - nhanh nhất, nhỏ nhất; Cache L1, L2, L3; RAM (bộ nhớ chính); SSD/HDD (bộ nhớ phụ). Nguyên lý locality giúp cache hoạt động hiệu quả: dữ liệu thường dùng sẽ được lưu ở cấp cao hơn.

\par\vspace{0.3cm}
\hrule
\vspace{0.3cm}
\end{minipage}

\begin{minipage}{\textwidth}
\subsection*{Question 26 \quad \footnotesize{[0026\_q026\_Question\_26.png]}}
\noindent \textbf{What kind of state does a process execute until it needs I/O resources when I/O is complete?}

\begin{enumerate}[label=\Alph*.]
    \item Hold state
    \item Ready state
    \item Running state
    \item \textbf{Waiting state} \hfill \textcolor{blue}{\textbf{[Correct]}}
\end{enumerate}

\noindent \textbf{Correct Answer:} \textcolor{blue}{\textbf{D}} \\[2pt]
\noindent \textbf{Explanation:} Khe cắm mở rộng (expansion slot) trên mainboard cho phép cắm thêm card mở rộng như card đồ họa (GPU), card mạng (NIC), card âm thanh (sound card). Các chuẩn kết nối: PCI, PCI Express (PCIe), AGP (cho card đồ họa cũ).

\par\vspace{0.3cm}
\hrule
\vspace{0.3cm}
\end{minipage}

\begin{minipage}{\textwidth}
\subsection*{Question 27 \quad \footnotesize{[0027\_q027\_Question\_27.png]}}
\noindent \textbf{The construct tests a condition.}

\begin{enumerate}[label=\Alph*.]
    \item Sequence
    \item \textbf{Decision} \hfill \textcolor{blue}{\textbf{[Correct]}}
    \item Repetition
    \item Flow
\end{enumerate}

\noindent \textbf{Correct Answer:} \textcolor{blue}{\textbf{B}} \\[2pt]
\noindent \textbf{Explanation:} Port (cổng) I/O: cổng nối tiếp (serial port - RS-232) truyền từng bit một; cổng song song (parallel port) truyền nhiều bit đồng thời; USB (Universal Serial Bus) là chuẩn hiện đại, tốc độ cao, hỗ trợ cắm nóng (hot-swap).

\par\vspace{0.3cm}
\hrule
\vspace{0.3cm}
\end{minipage}

\begin{minipage}{\textwidth}
\subsection*{Question 28 \quad \footnotesize{[0028\_q028\_Question\_28.png]}}
\noindent \textbf{is an English-language-like representation of code.}

\begin{enumerate}[label=\Alph*.]
    \item A UML diagram
    \item A program
    \item \textbf{Pseudocode} \hfill \textcolor{blue}{\textbf{[Correct]}}
    \item An algorithm
\end{enumerate}

\noindent \textbf{Correct Answer:} \textcolor{blue}{\textbf{C}} \\[2pt]
\noindent \textbf{Explanation:} Thiết bị lưu trữ: ổ cứng HDD (Hard Disk Drive) dùng đĩa từ, dung lượng lớn, giá rẻ; SSD (Solid State Drive) dùng flash memory, nhanh hơn, bền hơn, đắt hơn. Tốc độ đọc/ghi của SSD nhanh hơn HDD nhiều lần.

\par\vspace{0.3cm}
\hrule
\vspace{0.3cm}
\end{minipage}

\begin{minipage}{\textwidth}
\subsection*{Question 29 \quad \footnotesize{[0029\_q029\_Question\_29.png]}}
\noindent \textbf{is a basic algorithm in which we want to find the location of a target in a list of items.}

\begin{enumerate}[label=\Alph*.]
    \item Sorting
    \item \textbf{Searching} \hfill \textcolor{blue}{\textbf{[Correct]}}
    \item Product
    \item Summation
\end{enumerate}

\noindent \textbf{Correct Answer:} \textcolor{blue}{\textbf{B}} \\[2pt]
\noindent \textbf{Explanation:} Máy tính thế hệ thứ nhất (1940-1956) dùng đèn chân không (vacuum tubes); thế hệ 2 (1956-1963) dùng transistor; thế hệ 3 (1964-1971) dùng IC (Integrated Circuit); thế hệ 4 (1971-nay) dùng VLSI (Very Large Scale Integration) - vi xử lý (microprocessor).

\par\vspace{0.3cm}
\hrule
\vspace{0.3cm}
\end{minipage}

\begin{minipage}{\textwidth}
\subsection*{Question 30 \quad \footnotesize{[0030\_q030\_Question\_30.png]}}
\noindent \textbf{In sort, the item that goes into the sorted list is always the first item in the unsorted list.}

\begin{enumerate}[label=\Alph*.]
    \item Selection
    \item Bubble
    \item \textbf{Insertion} \hfill \textcolor{blue}{\textbf{[Correct]}}
    \item Every
\end{enumerate}

\noindent \textbf{Correct Answer:} \textcolor{blue}{\textbf{C}} \\[2pt]
\noindent \textbf{Explanation:} Bộ xử lý đa nhân (multi-core processor) có nhiều CPU core trên một chip, cho phép thực thi nhiều tiến trình đồng thời (parallel processing). Mỗi core có cache L1/L2 riêng, chia sẻ cache L3. Đa nhân cải thiện hiệu năng đa nhiệm và xử lý song song.

\par\vspace{0.3cm}
\hrule
\vspace{0.3cm}
\end{minipage}

\begin{minipage}{\textwidth}
\subsection*{Question 31 \quad \footnotesize{[0031\_q031\_Question\_31.png]}}
\noindent \textbf{A \_\_\_analyzer reads the source code symbol by symbol and creates a list of tokens in the source language}

\begin{enumerate}[label=\Alph*.]
    \item \textbf{Lexical} \hfill \textcolor{blue}{\textbf{[Correct]}}
    \item Syntax
    \item Semantic
    \item Code generation
\end{enumerate}

\noindent \textbf{Correct Answer:} \textcolor{blue}{\textbf{A}} \\[2pt]
\noindent \textbf{Explanation:} Địa chỉ bộ nhớ (memory address) là số nhị phân xác định vị trí của một ô nhớ trong RAM. CPU dùng address bus để gửi địa chỉ đến bộ nhớ. Độ rộng address bus quyết định không gian địa chỉ tối đa: n-bit address bus có thể đánh địa chỉ 2\textasciicircum{}n ô nhớ.

\par\vspace{0.3cm}
\hrule
\vspace{0.3cm}
\end{minipage}

\begin{minipage}{\textwidth}
\subsection*{Question 32 \quad \footnotesize{[0032\_q032\_Question\_32.png]}}
\noindent \textbf{A notation used to express clearly an algorithm is known as}

\begin{enumerate}[label=\Alph*.]
    \item \textbf{Algorithmic language} \hfill \textcolor{blue}{\textbf{[Correct]}}
    \item Assembly language
    \item Machine language
    \item High level language
    \item None of the others
\end{enumerate}

\noindent \textbf{Correct Answer:} \textcolor{blue}{\textbf{A}} \\[2pt]
\noindent \textbf{Explanation:} Trong kiến trúc Von Neumann, bộ nhớ lưu trữ cả chương trình (lệnh) và dữ liệu trong cùng một không gian địa chỉ. Đây gọi là stored-program concept - ý tưởng cách mạng cho phép máy tính linh hoạt thay đổi chương trình mà không cần thay đổi phần cứng.

\par\vspace{0.3cm}
\hrule
\vspace{0.3cm}
\end{minipage}

\begin{minipage}{\textwidth}
\subsection*{Question 33 \quad \footnotesize{[0033\_q033\_Question\_33.png]}}
\noindent \textbf{In the we can think of a program as an active agent that manipulates passive objects.}

\begin{enumerate}[label=\Alph*.]
    \item \textbf{Procedural paradigm} \hfill \textcolor{blue}{\textbf{[Correct]}}
    \item Functional paradigm
    \item Declarative paradigm
    \item Object-oriented paradigm
\end{enumerate}

\noindent \textbf{Correct Answer:} \textcolor{blue}{\textbf{A}} \\[2pt]
\noindent \textbf{Explanation:} Bộ xử lý (processor) có thể là CPU (Central Processing Unit - xử lý đa năng), GPU (Graphics Processing Unit - xử lý đồ họa song song), DSP (Digital Signal Processor - xử lý tín hiệu số), hoặc vi điều khiển (microcontroller - CPU tích hợp với bộ nhớ và I/O).

\par\vspace{0.3cm}
\hrule
\vspace{0.3cm}
\end{minipage}

\begin{minipage}{\textwidth}
\subsection*{Question 34 \quad \footnotesize{[0034\_q034\_Question\_34.png]}}
\noindent \textbf{Ais a unit of code consisting of zero or more statements. It is also known as a block. Example: \{ x = 1; Y = 20; \} SÖ DÁP ÁN CÃN CHON}

\begin{enumerate}[label=\Alph*.]
    \item \textbf{Compound statement} \hfill \textcolor{blue}{\textbf{[Correct]}}
    \item Assignment statement
    \item Control statement
    \item None of others
\end{enumerate}

\noindent \textbf{Correct Answer:} \textcolor{blue}{\textbf{A}} \\[2pt]
\noindent \textbf{Explanation:} Hệ thống bus gồm: data bus (2 chiều, truyền dữ liệu), address bus (1 chiều từ CPU, truyền địa chỉ), control bus (2 chiều, truyền tín hiệu điều khiển). Bus còn được phân loại theo vị trí: internal bus (trong CPU) và external bus (kết nối các thành phần).

\par\vspace{0.3cm}
\hrule
\vspace{0.3cm}
\end{minipage}

\begin{minipage}{\textwidth}
\subsection*{Question 35 \quad \footnotesize{[0035\_q035\_Question\_35.png]}}
\noindent \textbf{In a the project team needs to choose a language or a set of languages from among the procedural languages}

\begin{enumerate}[label=\Alph*.]
    \item Analysis phase
    \item Design phase
    \item \textbf{Implementation phase} \hfill \textcolor{blue}{\textbf{[Correct]}}
    \item Testing phase
\end{enumerate}

\noindent \textbf{Correct Answer:} \textcolor{blue}{\textbf{C}} \\[2pt]
\noindent \textbf{Explanation:} Bộ vi xử lý (microprocessor) là CPU được tích hợp trên một chip IC duy nhất. Intel 4004 (1971) là vi xử lý đầu tiên. Các hãng sản xuất vi xử lý lớn: Intel, AMD, ARM, IBM. Vi xử lý hiện đại có hàng tỷ transistor và tốc độ GHz.

\par\vspace{0.3cm}
\hrule
\vspace{0.3cm}
\end{minipage}

\begin{minipage}{\textwidth}
\subsection*{Question 36 \quad \footnotesize{[0036\_q036\_Question\_36.png]}}
\noindent \textbf{Defining the users, requirements, and methods is part of the \_phase.}

\begin{enumerate}[label=\Alph*.]
    \item \textbf{Analysis} \hfill \textcolor{blue}{\textbf{[Correct]}}
    \item Design
    \item Implementation
    \item Testing
\end{enumerate}

\noindent \textbf{Correct Answer:} \textcolor{blue}{\textbf{A}} \\[2pt]
\noindent \textbf{Explanation:} RISC vs CISC: RISC (Reduced Instruction Set) có ít lệnh hơn nhưng mỗi lệnh đơn giản, thực thi nhanh, thường dùng trong mobile/embedded; CISC (Complex Instruction Set) có nhiều lệnh phức tạp, mỗi lệnh làm nhiều việc, thường dùng trong PC/server. Ví dụ: ARM là RISC, x86 là CISC.

\par\vspace{0.3cm}
\hrule
\vspace{0.3cm}
\end{minipage}

\begin{minipage}{\textwidth}
\subsection*{Question 37 \quad \footnotesize{[0037\_q037\_Question\_37.png]}}
\noindent \textbf{In the system development process, structure charts are tools used in the \_phase.}

\begin{enumerate}[label=\Alph*.]
    \item Analysis
    \item \textbf{Design} \hfill \textcolor{blue}{\textbf{[Correct]}}
    \item Implementation
    \item Testing
\end{enumerate}

\noindent \textbf{Correct Answer:} \textcolor{blue}{\textbf{B}} \\[2pt]
\noindent \textbf{Explanation:} Pipeline hazard (xung đột đường ống): (1) Structural hazard - xung đột tài nguyên, (2) Data hazard - dữ liệu phụ thuộc giữa các lệnh, (3) Control hazard - lệnh rẽ nhánh. Kỹ thuật khắc phục: forwarding (truyền trực tiếp), branch prediction (dự đoán rẽ nhánh), stalling (chèn lệnh NOP).

\par\vspace{0.3cm}
\hrule
\vspace{0.3cm}
\end{minipage}

\begin{minipage}{\textwidth}
\subsection*{Question 38 \quad \footnotesize{[0038\_q038\_Question\_38.png]}}
\noindent \textbf{What is quality referring to the ability to move data and/or a system from one platform to another and to reuse code? SÖ DÁP ÁN CÀN CHON}

\begin{enumerate}[label=\Alph*.]
    \item Flexibility
    \item \textbf{Interoperability} \hfill \textcolor{blue}{\textbf{[Correct]}}
    \item Resusability
    \item Changeability
\end{enumerate}

\noindent \textbf{Correct Answer:} \textcolor{blue}{\textbf{B}} \\[2pt]
\noindent \textbf{Explanation:} Virtual memory (bộ nhớ ảo) là kỹ thuật cho phép chạy chương trình lớn hơn RAM bằng cách dùng ổ cứng làm bộ nhớ mở rộng. MMU (Memory Management Unit) trong CPU quản lý ánh xạ địa chỉ ảo → địa chỉ vật lý. Page fault xảy ra khi dữ liệu cần ở ổ cứng, phải nạp vào RAM.

\par\vspace{0.3cm}
\hrule
\vspace{0.3cm}
\end{minipage}

\begin{minipage}{\textwidth}
\subsection*{Question 39 \quad \footnotesize{[0039\_q039\_Question\_39.png]}}
\noindent \textbf{What is the method in black-box testing that a subset of values in the input domain is selected for testing?.}

\begin{enumerate}[label=\Alph*.]
    \item Boundary-value testing
    \item \textbf{Random testing} \hfill \textcolor{blue}{\textbf{[Correct]}}
    \item Exhaustive testing
    \item Data flow testing
\end{enumerate}

\noindent \textbf{Correct Answer:} \textcolor{blue}{\textbf{B}} \\[2pt]
\noindent \textbf{Explanation:} Các loại máy in: laser (dùng bột mực, chất lượng cao, nhanh), inkjet (dùng mực nước, giá rẻ, in ảnh đẹp), dot matrix (dùng kim đập, in hóa đơn), 3D printer (in vật thể 3D). Máy in kết nối qua USB, Wi-Fi, Ethernet.

\par\vspace{0.3cm}
\hrule
\vspace{0.3cm}
\end{minipage}

\begin{minipage}{\textwidth}
\subsection*{Question 40 \quad \footnotesize{[0040\_q040\_Question\_40.png]}}
\noindent \textbf{What is more common technique of two-dimensional arrays that can be stored in memory?}

\begin{enumerate}[label=\Alph*.]
    \item Column-major storage
    \item DMA
    \item Isolated I/O
    \item \textbf{Row-major storage} \hfill \textcolor{blue}{\textbf{[Correct]}}
\end{enumerate}

\noindent \textbf{Correct Answer:} \textcolor{blue}{\textbf{D}} \\[2pt]
\noindent \textbf{Explanation:} Thiết bị hiển thị: CRT (ống phóng tia âm cực - cũ), LCD (màn hình tinh thể lỏng), LED (dùng đèn LED nền), OLED (tự phát sáng). Độ phân giải đo bằng pixel (VD: 1920×1080 = Full HD). Tốc độ làm tươi (refresh rate) đo bằng Hz (VD: 144Hz cho gaming).

\par\vspace{0.3cm}
\hrule
\vspace{0.3cm}
\end{minipage}


\newpage
\section*{2. Numbering Systems}
\addcontentsline{toc}{section}{2. Numbering Systems}
\hrule\vspace{0.5cm}

\begin{minipage}{\textwidth}
\subsection*{Question 41 \quad \footnotesize{[0041\_q041\_Question\_41.png]}}
\noindent \textbf{Each element in a record is called}

\begin{enumerate}[label=\Alph*.]
    \item A variable
    \item An index
    \item \textbf{A field} \hfill \textcolor{blue}{\textbf{[Correct]}}
    \item A node
\end{enumerate}

\noindent \textbf{Correct Answer:} \textcolor{blue}{\textbf{C}} \\[2pt]
\noindent \textbf{Explanation:} Hệ thập phân (decimal - base 10) dùng 10 ký số 0-9. Hệ nhị phân (binary - base 2) dùng 2 ký số 0-1. Máy tính dùng hệ nhị phân vì chỉ có hai trạng thái điện (bật/tắt). Mỗi ký số nhị phân gọi là bit. Chuyển đổi giữa các hệ đếm là kỹ năng cơ bản.

\par\vspace{0.3cm}
\hrule
\vspace{0.3cm}
\end{minipage}

\begin{minipage}{\textwidth}
\subsection*{Question 42 \quad \footnotesize{[0042\_q042\_Question\_42.png]}}
\noindent \textbf{Given a linked list called children, the pointer variable children identifies element of the linked list.}

\begin{enumerate}[label=\Alph*.]
    \item \textbf{The first} \hfill \textcolor{blue}{\textbf{[Correct]}}
    \item The second
    \item The last
    \item Any
\end{enumerate}

\noindent \textbf{Correct Answer:} \textcolor{blue}{\textbf{A}} \\[2pt]
\noindent \textbf{Explanation:} Hệ bát phân (octal - base 8) dùng 8 ký số 0-7. Hệ thập lục phân (hexadecimal - base 16) dùng 16 ký số 0-9 và A-F. Hex được dùng phổ biến trong lập trình và địa chỉ bộ nhớ vì mỗi ký số hex tương ứng 4 bit (nibble).

\par\vspace{0.3cm}
\hrule
\vspace{0.3cm}
\end{minipage}

\begin{minipage}{\textwidth}
\subsection*{Question 43 \quad \footnotesize{[0043\_q043\_Question\_43.png]}}
\noindent \textbf{The push operation of the stack.}

\begin{enumerate}[label=\Alph*.]
    \item Deletes an item from the top
    \item Deletes an item from the bottom
    \item \textbf{Inserts an item at the top} \hfill \textcolor{blue}{\textbf{[Correct]}}
    \item Inserts an item at the bottom
\end{enumerate}

\noindent \textbf{Correct Answer:} \textcolor{blue}{\textbf{C}} \\[2pt]
\noindent \textbf{Explanation:} Chuyển đổi từ nhị phân sang thập phân: nhân mỗi bit với 2\textasciicircum{}i (i là vị trí, đếm từ 0 từ phải). Ví dụ: 1011₂ = 1×2³ + 0×2² + 1×2¹ + 1×2⁰ = 8+0+2+1 = 11₁₀. Chuyển từ thập phân sang nhị phân: chia liên tiếp cho 2, lấy số dư từ dưới lên.

\par\vspace{0.3cm}
\hrule
\vspace{0.3cm}
\end{minipage}

\begin{minipage}{\textwidth}
\subsection*{Question 44 \quad \footnotesize{[0044\_q044\_Question\_44.png]}}
\noindent \textbf{In of a binary tree we process all the children of a node before proceeding with the next generation}

\begin{enumerate}[label=\Alph*.]
    \item Preorder traversal
    \item Inorder traversal
    \item Postorder traversal
    \item \textbf{Breadth-first traversal} \hfill \textcolor{blue}{\textbf{[Correct]}}
\end{enumerate}

\noindent \textbf{Correct Answer:} \textcolor{blue}{\textbf{D}} \\[2pt]
\noindent \textbf{Explanation:} Chuyển đổi thập phân sang nhị phân dùng phương pháp chia cho 2: (1) Chia số thập phân cho 2, (2) ghi số dư (0 hoặc 1), (3) lấy thương làm số mới, (4) lặp đến khi thương = 0. Đọc số dư từ dưới lên trên.

\par\vspace{0.3cm}
\hrule
\vspace{0.3cm}
\end{minipage}

\begin{minipage}{\textwidth}
\subsection*{Question 45 \quad \footnotesize{[0045\_q045\_Question\_45.png]}}
\noindent \textbf{The queue operation: "queue (queueName)" is used to}

\begin{enumerate}[label=\Alph*.]
    \item Delete queue
    \item Check status of queue
    \item \textbf{Create an empty queue} \hfill \textcolor{blue}{\textbf{[Correct]}}
    \item Insert queue
\end{enumerate}

\noindent \textbf{Correct Answer:} \textcolor{blue}{\textbf{C}} \\[2pt]
\noindent \textbf{Explanation:} Hệ hex (base 16) rất hữu ích vì mỗi chữ số hex tương ứng 4 bit (nibble). Ví dụ: AF₁₆ = 1010 1111₂. Để chuyển hex sang nhị phân: đổi từng chữ số hex thành 4 bit. Để chuyển nhị phân sang hex: gom 4 bit từ phải, đổi từng nhóm.

\par\vspace{0.3cm}
\hrule
\vspace{0.3cm}
\end{minipage}

\begin{minipage}{\textwidth}
\subsection*{Question 46 \quad \footnotesize{[0046\_q046\_Question\_46.png]}}
\noindent \textbf{In the sequential file we process the records one by one. After the operating system processes the last record the is detected and the loop is exited.}

\begin{enumerate}[label=\Alph*.]
    \item Hashed file
    \item Sequential file
    \item \textbf{EOF} \hfill \textcolor{blue}{\textbf{[Correct]}}
    \item None of others
\end{enumerate}

\noindent \textbf{Correct Answer:} \textcolor{blue}{\textbf{C}} \\[2pt]
\noindent \textbf{Explanation:} Phép cộng nhị phân: 0+0=0, 0+1=1, 1+0=1, 1+1=0 nhớ 1 (carry). Phép trừ nhị phân: 0-0=0, 1-0=1, 1-1=0, 0-1=1 mượn 1 (borrow). Phép nhân/chia nhị phân tương tự thập phân nhưng đơn giản hơn.

\par\vspace{0.3cm}
\hrule
\vspace{0.3cm}
\end{minipage}

\begin{minipage}{\textwidth}
\subsection*{Question 47 \quad \footnotesize{[0047\_q047\_Question\_47.png]}}
\noindent \textbf{In the hashing method, the key is divided by the file size, and the address is the remainder plus 1.}

\begin{enumerate}[label=\Alph*.]
    \item Direct
    \item Division remainder
    \item \textbf{Modulo division} \hfill \textcolor{blue}{\textbf{[Correct]}}
    \item Digit extraction
\end{enumerate}

\noindent \textbf{Correct Answer:} \textcolor{blue}{\textbf{C}} \\[2pt]
\noindent \textbf{Explanation:} Biểu diễn số âm trong máy tính: (1) Sign-magnitude: bit đầu là dấu (0=+, 1=-), (2) Ones' complement: đảo tất cả bit, (3) Two's complement: đảo bit + 1. Two's complement là phương pháp phổ biến nhất vì đơn giản hóa phép cộng/trừ.

\par\vspace{0.3cm}
\hrule
\vspace{0.3cm}
\end{minipage}

\begin{minipage}{\textwidth}
\subsection*{Question 48 \quad \footnotesize{[0048\_q048\_Question\_48.png]}}
\noindent \textbf{are provided by most operating systems for organizing files. This performs the same function as a folder in a filing cabinet.}

\begin{enumerate}[label=\Alph*.]
    \item Operating system
    \item Files
    \item Binary files
    \item \textbf{Directories} \hfill \textcolor{blue}{\textbf{[Correct]}}
\end{enumerate}

\noindent \textbf{Correct Answer:} \textcolor{blue}{\textbf{D}} \\[2pt]
\noindent \textbf{Explanation:} Two's complement: để biểu diễn số âm -N, lấy biểu diễn nhị phân của N, đảo tất cả bit (ones' complement) rồi cộng 1. Phạm vi biểu diễn với n bit: từ -2\textasciicircum{}(n-1) đến 2\textasciicircum{}(n-1)-1. Ví dụ 8 bit: -128 đến 127.

\par\vspace{0.3cm}
\hrule
\vspace{0.3cm}
\end{minipage}

\begin{minipage}{\textwidth}
\subsection*{Question 49 \quad \footnotesize{[0049\_q049\_Question\_49.png]}}
\noindent \textbf{In a three-level DBMS architecture, the level interacts directly with the users. SÖ DÁP ÁN CÃN CHON}

\begin{enumerate}[label=\Alph*.]
    \item \textbf{External} \hfill \textcolor{blue}{\textbf{[Correct]}}
    \item Conceptual
    \item Internal
    \item Physical
\end{enumerate}

\noindent \textbf{Correct Answer:} \textcolor{blue}{\textbf{A}} \\[2pt]
\noindent \textbf{Explanation:} Số thực dấu phẩy động (floating-point) biểu diễn theo chuẩn IEEE 754 gồm: 1 bit dấu (sign), m bit mũ (exponent), n bit phần lẻ (fraction/mantissa). Ví dụ với 32-bit (float): 1 bit dấu, 8 bit mũ, 23 bit phần lẻ. Công thức: (-1)\textasciicircum{}sign × 1.mantissa × 2\textasciicircum{}(exponent-127).

\par\vspace{0.3cm}
\hrule
\vspace{0.3cm}
\end{minipage}

\begin{minipage}{\textwidth}
\subsection*{Question 50 \quad \footnotesize{[0050\_q050\_Question\_50.png]}}
\noindent \textbf{In the network database model, the entities are organized in a in which some entities can be accessed through several paths PO}

\begin{enumerate}[label=\Alph*.]
    \item \textbf{Graph} \hfill \textcolor{blue}{\textbf{[Correct]}}
    \item Tree
    \item Array
    \item Record
\end{enumerate}

\noindent \textbf{Correct Answer:} \textcolor{blue}{\textbf{A}} \\[2pt]
\noindent \textbf{Explanation:} Overflow (tràn số) xảy ra khi kết quả vượt quá phạm vi biểu diễn. Ví dụ: cộng hai số 8-bit có dấu: 127 + 1 = 128 không biểu diễn được (vì max là 127). Underflow xảy ra khi số quá nhỏ so với khả năng biểu diễn.

\par\vspace{0.3cm}
\hrule
\vspace{0.3cm}
\end{minipage}

\begin{minipage}{\textwidth}
\subsection*{Question 51 \quad \footnotesize{[0051\_q051\_Question\_51.png]}}
\noindent \textbf{Each column in a relation is called}

\begin{enumerate}[label=\Alph*.]
    \item \textbf{An attribute} \hfill \textcolor{blue}{\textbf{[Correct]}}
    \item A tuple
    \item A union
    \item An attitude
\end{enumerate}

\noindent \textbf{Correct Answer:} \textcolor{blue}{\textbf{A}} \\[2pt]
\noindent \textbf{Explanation:} BCD (Binary-Coded Decimal): mỗi chữ số thập phân được biểu diễn bằng 4 bit nhị phân. Ví dụ: 273₁₀ = 0010 0111 0011 (BCD). BCD dùng trong tính toán tài chính vì tránh sai số làm tròn của floating-point.

\par\vspace{0.3cm}
\hrule
\vspace{0.3cm}
\end{minipage}

\begin{minipage}{\textwidth}
\subsection*{Question 52 \quad \footnotesize{[0052\_q052\_Question\_52.png]}}
\noindent \textbf{The unary operation always results in a relation that has exactly one more row than the original relation.}

\begin{enumerate}[label=\Alph*.]
    \item \textbf{Insert} \hfill \textcolor{blue}{\textbf{[Correct]}}
    \item Delete
    \item Update
    \item Select
\end{enumerate}

\noindent \textbf{Correct Answer:} \textcolor{blue}{\textbf{A}} \\[2pt]
\noindent \textbf{Explanation:} EBCDIC (Extended Binary Coded Decimal Interchange Code) là mã ký tự 8-bit do IBM phát triển. ASCII (American Standard Code for Information Interchange) là mã 7-bit (hoặc 8-bit mở rộng) phổ biến hơn. Unicode là chuẩn hiện đại, dùng 16-32 bit.

\par\vspace{0.3cm}
\hrule
\vspace{0.3cm}
\end{minipage}

\begin{minipage}{\textwidth}
\subsection*{Question 53 \quad \footnotesize{[0053\_q053\_Question\_53.png]}}
\noindent \textbf{The term users in a DBMS has a broad meaning. We can divide users into: SÖ DÁP ÁN CÀN CHON}

\begin{enumerate}[label=\Alph*.]
    \item End users and database asministrators
    \item End users and normal users
    \item \textbf{Application programs and end users} \hfill \textcolor{blue}{\textbf{[Correct]}}
    \item Application programs and database adminitrators
\end{enumerate}

\noindent \textbf{Correct Answer:} \textcolor{blue}{\textbf{C}} \\[2pt]
\noindent \textbf{Explanation:} Chuyển đổi giữa các hệ đếm: từ hệ b (ví dụ nhị phân) sang thập phân: nhân từng chữ số với b\textasciicircum{}i. Từ thập phân sang hệ b: chia liên tiếp cho b. Giữa nhị phân và hex: gom 4 bit. Giữa nhị phân và bát phân: gom 3 bit.

\par\vspace{0.3cm}
\hrule
\vspace{0.3cm}
\end{minipage}

\begin{minipage}{\textwidth}
\subsection*{Question 54 \quad \footnotesize{[0054\_q054\_Question\_54.png]}}
\noindent \textbf{they can obtain some other type of In the attacks threaten the confidentiality based on information by monitoring online traffic}

\begin{enumerate}[label=\Alph*.]
    \item Snooping
    \item Repudiation
    \item Denial of service
    \item \textbf{Traffic analysis} \hfill \textcolor{blue}{\textbf{[Correct]}}
\end{enumerate}

\noindent \textbf{Correct Answer:} \textcolor{blue}{\textbf{D}} \\[2pt]
\noindent \textbf{Explanation:} Biểu diễn số dùng n bit: số không dấu (unsigned): 0 đến 2\textasciicircum{}n-1. Số có dấu (signed) dùng MSB (Most Significant Bit) làm bit dấu. Two's complement là chuẩn cho số có dấu vì phép cộng/trừ không cần xử lý đặc biệt.

\par\vspace{0.3cm}
\hrule
\vspace{0.3cm}
\end{minipage}

\begin{minipage}{\textwidth}
\subsection*{Question 55 \quad \footnotesize{[0055\_q055\_Question\_55.png]}}
\noindent \textbf{Before transmission the sender transforms the plaintext into the ciphertext—this is}

\begin{enumerate}[label=\Alph*.]
    \item Plain text
    \item Cipher text
    \item \textbf{Encryption} \hfill \textcolor{blue}{\textbf{[Correct]}}
    \item Decryption
\end{enumerate}

\noindent \textbf{Correct Answer:} \textcolor{blue}{\textbf{C}} \\[2pt]
\noindent \textbf{Explanation:} Số lẻ thập phân sang nhị phân: nhân phần lẻ với 2 liên tiếp, lấy phần nguyên. Ví dụ: 0.375₁₀ = 0.011₂ vì 0.375×2=0.75 (bit 0), 0.75×2=1.5 (bit 1), 0.5×2=1.0 (bit 1). Đọc từ trên xuống.

\par\vspace{0.3cm}
\hrule
\vspace{0.3cm}
\end{minipage}

\begin{minipage}{\textwidth}
\subsection*{Question 56 \quad \footnotesize{[0056\_q056\_Question\_56.png]}}
\noindent \textbf{are unwanted programs that are hidden within other programs (host). When the user executes the host program}

\begin{enumerate}[label=\Alph*.]
    \item DoS
    \item Trojan horses
    \item Worms
    \item \textbf{Viruses} \hfill \textcolor{blue}{\textbf{[Correct]}}
\end{enumerate}

\noindent \textbf{Correct Answer:} \textcolor{blue}{\textbf{D}} \\[2pt]
\noindent \textbf{Explanation:} Chuẩn IEEE 754 cho số thực: single precision (32-bit): 1 bit dấu, 8 bit exponent (biased by 127), 23 bit mantissa. Double precision (64-bit): 1 bit dấu, 11 bit exponent (biased by 1023), 52 bit mantissa. Các giá trị đặc biệt: ±0, ±∞, NaN (Not a Number).

\par\vspace{0.3cm}
\hrule
\vspace{0.3cm}
\end{minipage}

\begin{minipage}{\textwidth}
\subsection*{Question 57 \quad \footnotesize{[0057\_q057\_Question\_57.png]}}
\noindent \textbf{attack. Sending a virus to a computer is called an}

\begin{enumerate}[label=\Alph*.]
    \item \textbf{Penetration} \hfill \textcolor{blue}{\textbf{[Correct]}}
    \item Denial of service
    \item Either a or b
    \item Neither a nor b
\end{enumerate}

\noindent \textbf{Correct Answer:} \textcolor{blue}{\textbf{A}} \\[2pt]
\noindent \textbf{Explanation:} Overflow trong số không dấu: kết quả cần nhiều bit hơn n bit cho phép. Overflow trong số có dấu: kết quả vượt phạm vi -2\textasciicircum{}(n-1) đến 2\textasciicircum{}(n-1)-1. Máy tính phát hiện overflow qua carry flag và overflow flag trong thanh ghi trạng thái.

\par\vspace{0.3cm}
\hrule
\vspace{0.3cm}
\end{minipage}

\begin{minipage}{\textwidth}
\subsection*{Question 58 \quad \footnotesize{[0058\_q058\_Question\_58.png]}}
\noindent \textbf{In a\_ the relationship between a character in the plaintext to a character in the ciphertext is one-to-many}

\begin{enumerate}[label=\Alph*.]
    \item Monoalphabetic cipher
    \item Substitution cipher
    \item \textbf{Polyalphabetic cipher} \hfill \textcolor{blue}{\textbf{[Correct]}}
    \item Transposition cipher
\end{enumerate}

\noindent \textbf{Correct Answer:} \textcolor{blue}{\textbf{C}} \\[2pt]
\noindent \textbf{Explanation:} Base (cơ số) của hệ đếm là số ký số dùng để biểu diễn. Binary (base 2): 0,1. Octal (base 8): 0-7. Decimal (base 10): 0-9. Hexadecimal (base 16): 0-9,A-F. Cơ số càng lớn, biểu diễn càng ngắn gọn.

\par\vspace{0.3cm}
\hrule
\vspace{0.3cm}
\end{minipage}

\begin{minipage}{\textwidth}
\subsection*{Question 59 \quad \footnotesize{[0059\_q059\_Question\_59.png]}}
\noindent \textbf{ciphers are sometimes called public-key ciphers}

\begin{enumerate}[label=\Alph*.]
    \item Symmetric-key
    \item \textbf{Asymmetric-key} \hfill \textcolor{blue}{\textbf{[Correct]}}
    \item Public-key
    \item Private-key
\end{enumerate}

\noindent \textbf{Correct Answer:} \textcolor{blue}{\textbf{B}} \\[2pt]
\noindent \textbf{Explanation:} Bit MSB (Most Significant Bit) là bit có trọng số lớn nhất (bên trái nhất). Bit LSB (Least Significant Bit) là bit có trọng số nhỏ nhất (bên phải nhất). Trong unsigned integer, MSB có giá trị 2\textasciicircum{}(n-1). Trong signed, MSB là bit dấu.

\par\vspace{0.3cm}
\hrule
\vspace{0.3cm}
\end{minipage}

\begin{minipage}{\textwidth}
\subsection*{Question 60 \quad \footnotesize{[0060\_q060\_Question\_60.png]}}
\noindent \textbf{In a computer, the subsystem stores data and programs}

\begin{enumerate}[label=\Alph*.]
    \item ALU
    \item Input/output
    \item \textbf{Memory} \hfill \textcolor{blue}{\textbf{[Correct]}}
    \item Control unit
\end{enumerate}

\noindent \textbf{Correct Answer:} \textcolor{blue}{\textbf{C}} \\[2pt]
\noindent \textbf{Explanation:} Kỹ thuật biểu diễn số: signed-magnitude dùng MSB làm dấu, phần còn lại là độ lớn. Ones' complement: đảo tất cả bit. Two's complement: đảo bit + 1. Excess system (biased): số được cộng thêm bias. Excess-127 dùng trong IEEE 754 exponent.

\par\vspace{0.3cm}
\hrule
\vspace{0.3cm}
\end{minipage}

\begin{minipage}{\textwidth}
\subsection*{Question 61 \quad \footnotesize{[0061\_q061\_Question\_61.png]}}
\noindent \textbf{A step-by-step solution to a problem is called}

\begin{enumerate}[label=\Alph*.]
    \item Hardware
    \item An operating system
    \item A computer language
    \item \textbf{An algorithm} \hfill \textcolor{blue}{\textbf{[Correct]}}
\end{enumerate}

\noindent \textbf{Correct Answer:} \textcolor{blue}{\textbf{D}} \\[2pt]
\noindent \textbf{Explanation:} Số thập lục phân (hex) thường dùng trong địa chỉ bộ nhớ, màu sắc (HTML color), và lập trình hệ thống. Ví dụ: màu trắng là \#FFFFFF. Địa chỉ IPv6 dùng hex. Mỗi byte (8 bit) được biểu diễn bằng 2 ký số hex.

\par\vspace{0.3cm}
\hrule
\vspace{0.3cm}
\end{minipage}

\begin{minipage}{\textwidth}
\subsection*{Question 62 \quad \footnotesize{[0062\_q062\_Question\_62.png]}}
\noindent \textbf{The first microprocessors produced by Intel Corpn. and Texas Instruments were used primarily to control small}

\begin{enumerate}[label=\Alph*.]
    \item Microwave ovens
    \item Washing machines
    \item \textbf{Calculators} \hfill \textcolor{blue}{\textbf{[Correct]}}
    \item Personal computers
    \item Robotics
\end{enumerate}

\noindent \textbf{Correct Answer:} \textcolor{blue}{\textbf{C}} \\[2pt]
\noindent \textbf{Explanation:} Với n bit, có thể biểu diễn 2\textasciicircum{}n giá trị khác nhau. Cho unsigned: 0 đến 2\textasciicircum{}n-1. Cho signed (two's complement): -2\textasciicircum{}(n-1) đến 2\textasciicircum{}(n-1)-1. Ví dụ 4 bit: unsigned 0-15, signed -8 đến 7.

\par\vspace{0.3cm}
\hrule
\vspace{0.3cm}
\end{minipage}

\begin{minipage}{\textwidth}
\subsection*{Question 63 \quad \footnotesize{[0063\_q063\_Question\_63.png]}}
\noindent \textbf{Off-line operation is the operation of devices without the control of}

\begin{enumerate}[label=\Alph*.]
    \item Memory
    \item \textbf{CPU} \hfill \textcolor{blue}{\textbf{[Correct]}}
    \item ALU
    \item Control unit
    \item None of the others
\end{enumerate}

\noindent \textbf{Correct Answer:} \textcolor{blue}{\textbf{B}} \\[2pt]
\noindent \textbf{Explanation:} BCD dùng 4 bit cho mỗi chữ số thập phân. Ví dụ: 123₁₀ = 0001 0010 0011 (BCD). So với nhị phân thuần: 123₁₀ = 1111011₂ (7 bit). BCD dùng nhiều bit hơn nhưng dễ chuyển đổi và tránh sai số.

\par\vspace{0.3cm}
\hrule
\vspace{0.3cm}
\end{minipage}

\begin{minipage}{\textwidth}
\subsection*{Question 64 \quad \footnotesize{[0064\_q064\_Question\_64.png]}}
\noindent \textbf{is a memory type with capacitors that need to be refreshed periodically.}

\begin{enumerate}[label=\Alph*.]
    \item SRAM
    \item \textbf{DRAM} \hfill \textcolor{blue}{\textbf{[Correct]}}
    \item ROM
    \item CROM
\end{enumerate}

\noindent \textbf{Correct Answer:} \textcolor{blue}{\textbf{B}} \\[2pt]
\noindent \textbf{Explanation:} Dùng bù 1 (ones' complement): đảo 0 thành 1 và 1 thành 0. Ví dụ: 0101 (5) → 1010 (-5 trong bù 1). Two's complement: bù 1 + 1: 1010 + 1 = 1011 (-5 trong bù 2). Hầu hết máy tính hiện đại dùng two's complement.

\par\vspace{0.3cm}
\hrule
\vspace{0.3cm}
\end{minipage}

\begin{minipage}{\textwidth}
\subsection*{Question 65 \quad \footnotesize{[0065\_q065\_Question\_65.png]}}
\noindent \textbf{In magnetic disks, the defines how fast the disk is spinning.}

\begin{enumerate}[label=\Alph*.]
    \item Seek time
    \item Tranfer time
    \item \textbf{Rotational speed} \hfill \textcolor{blue}{\textbf{[Correct]}}
    \item None of others
\end{enumerate}

\noindent \textbf{Correct Answer:} \textcolor{blue}{\textbf{C}} \\[2pt]
\noindent \textbf{Explanation:} Phép nhân trong nhị phân: tương tự thập phân nhưng chỉ có 0 và 1. Nhân với 0 cho kết quả 0, nhân với 1 cho số bị nhân. Phép nhân được thực hiện qua cộng và dịch (shift-and-add). Phần cứng nhân dùng ALU và shift register.

\par\vspace{0.3cm}
\hrule
\vspace{0.3cm}
\end{minipage}

\begin{minipage}{\textwidth}
\subsection*{Question 66 \quad \footnotesize{[0066\_q066\_Question\_66.png]}}
\noindent \textbf{What is the name given to the molecular-scale computer?}

\begin{enumerate}[label=\Alph*.]
    \item Femtocomputer
    \item \textbf{Nanocomputer} \hfill \textcolor{blue}{\textbf{[Correct]}}
    \item Supercomputer
    \item Microcomputer
    \item None of the others
\end{enumerate}

\noindent \textbf{Correct Answer:} \textcolor{blue}{\textbf{B}} \\[2pt]
\noindent \textbf{Explanation:} Excess (biased) representation: giá trị thực = giá trị lưu - bias. Ví dụ: excess-3: 0 lưu thành 3, 1 lưu thành 4. Trong IEEE 754, exponent dùng excess-127 (32-bit) hoặc excess-1023 (64-bit).

\par\vspace{0.3cm}
\hrule
\vspace{0.3cm}
\end{minipage}

\begin{minipage}{\textwidth}
\subsection*{Question 67 \quad \footnotesize{[0067\_q067\_Question\_67.png]}}
\noindent \textbf{Which is a unit representing the no bits of discrete.}

\begin{enumerate}[label=\Alph*.]
    \item \textbf{Baud} \hfill \textcolor{blue}{\textbf{[Correct]}}
    \item Byte
    \item Bit
    \item All of the others
    \item None of the others
\end{enumerate}

\noindent \textbf{Correct Answer:} \textcolor{blue}{\textbf{A}} \\[2pt]
\noindent \textbf{Explanation:} Số hữu tỷ có thể biểu diễn chính xác trong nhị phân nếu mẫu số là lũy thừa của 2 (VD: 1/2, 3/4, 7/8). Các số như 0.1₁₀ (1/10) không biểu diễn chính xác trong nhị phân, dẫn đến sai số floating-point.

\par\vspace{0.3cm}
\hrule
\vspace{0.3cm}
\end{minipage}

\begin{minipage}{\textwidth}
\subsection*{Question 68 \quad \footnotesize{[0068\_q068\_Question\_68.png]}}
\noindent \textbf{The base of the binary number system is}

\begin{enumerate}[label=\Alph*.]
    \item \textbf{2} \hfill \textcolor{blue}{\textbf{[Correct]}}
    \item 8
    \item 10
    \item 16
\end{enumerate}

\noindent \textbf{Correct Answer:} \textcolor{blue}{\textbf{A}} \\[2pt]
\noindent \textbf{Explanation:} Tính giá trị biểu diễn nhị phân: signed magnitude: (-1)\textasciicircum{}sign × giá\_trị\_phần\_còn\_lại. Ones' complement: số dương giống unsigned, số âm = đảo bit. Two's complement: số dương giống unsigned, số âm = đảo bit + 1.

\par\vspace{0.3cm}
\hrule
\vspace{0.3cm}
\end{minipage}

\begin{minipage}{\textwidth}
\subsection*{Question 69 \quad \footnotesize{[0069\_q069\_Question\_69.png]}}
\noindent \textbf{When converting a decimal fraction to base b, we repeatedly \_b.}

\begin{enumerate}[label=\Alph*.]
    \item \textbf{Divide by} \hfill \textcolor{blue}{\textbf{[Correct]}}
    \item Multiply by
    \item Add to
    \item Subtract from
\end{enumerate}

\noindent \textbf{Correct Answer:} \textcolor{blue}{\textbf{A}} \\[2pt]
\noindent \textbf{Explanation:} Biểu diễn ký tự: ASCII dùng 7 bit, mã hóa 128 ký tự (chữ cái, số, dấu câu). Extended ASCII dùng 8 bit, thêm 128 ký tự. Unicode dùng 16-32 bit, mã hóa hơn 143,000 ký tự từ mọi ngôn ngữ.

\par\vspace{0.3cm}
\hrule
\vspace{0.3cm}
\end{minipage}

\begin{minipage}{\textwidth}
\subsection*{Question 70 \quad \footnotesize{[0070\_q070\_Question\_70.png]}}
\noindent \textbf{The binary equivalent of the Octal number 13.54 is}

\begin{enumerate}[label=\Alph*.]
    \item \textbf{1011.1011} \hfill \textcolor{blue}{\textbf{[Correct]}}
    \item 1101.1110
    \item 1001.1110
    \item All of the others
    \item None of the others
\end{enumerate}

\noindent \textbf{Correct Answer:} \textcolor{blue}{\textbf{A}} \\[2pt]
\noindent \textbf{Explanation:} Range (phạm vi) của số không dấu n bit: [0, 2\textasciicircum{}n-1]. Range của số có dấu n bit (two's complement): [-2\textasciicircum{}(n-1), 2\textasciicircum{}(n-1)-1]. Ví dụ: 8-bit unsigned: 0-255, signed: -128 đến 127.

\par\vspace{0.3cm}
\hrule
\vspace{0.3cm}
\end{minipage}

\begin{minipage}{\textwidth}
\subsection*{Question 71 \quad \footnotesize{[0071\_q071\_Question\_71.png]}}
\noindent \textbf{In a positional number system with base b, we can always find the number of digits of an integer. So how many bits can we find in the decimal number 20 in hexadecimal system?}

\begin{enumerate}[label=\Alph*.]
    \item 1
    \item \textbf{2} \hfill \textcolor{blue}{\textbf{[Correct]}}
    \item 3
    \item 4
    \item None of the others
\end{enumerate}

\noindent \textbf{Correct Answer:} \textcolor{blue}{\textbf{B}} \\[2pt]
\noindent \textbf{Explanation:} Phân loại số trong máy tính: integer (số nguyên) dùng unsigned/signed; floating-point (số thực) dùng IEEE 754; BCD dùng cho tính toán chính xác; fixed-point dùng trong DSP và ứng dụng nhúng.

\par\vspace{0.3cm}
\hrule
\vspace{0.3cm}
\end{minipage}

\begin{minipage}{\textwidth}
\subsection*{Question 72 \quad \footnotesize{[0072\_q072\_Question\_72.png]}}
\noindent \textbf{CÅU HÒI In the sign-and-magnitude representation, the memory location can hold four bits, assume that we have stored integer -6 in a memory location and we then try to subtract 7 to the integer. The computer's response is\_?}

\begin{enumerate}[label=\Alph*.]
    \item 6
    \item -6
    \item \textbf{7} \hfill \textcolor{blue}{\textbf{[Correct]}}
    \item -7
\end{enumerate}

\noindent \textbf{Correct Answer:} \textcolor{blue}{\textbf{C}} \\[2pt]
\noindent \textbf{Explanation:} Bậc (order) của hệ đếm: cơ số b, mỗi chữ số ở vị trí i có giá trị × b\textasciicircum{}i. Vị trí đánh từ 0 bên phải dấu phẩy. Ví dụ: 123.45₁₀ = 1×10² + 2×10¹ + 3×10⁰ + 4×10⁻¹ + 5×10⁻².

\par\vspace{0.3cm}
\hrule
\vspace{0.3cm}
\end{minipage}

\begin{minipage}{\textwidth}
\subsection*{Question 73 \quad \footnotesize{[0073\_q073\_Question\_73.png]}}
\noindent \textbf{Different sets of bit patterns have been designed to represent text symbols. Each set is called a ?}

\begin{enumerate}[label=\Alph*.]
    \item bit patterns
    \item binary
    \item \textbf{code} \hfill \textcolor{blue}{\textbf{[Correct]}}
    \item binary text file
\end{enumerate}

\noindent \textbf{Correct Answer:} \textcolor{blue}{\textbf{C}} \\[2pt]
\noindent \textbf{Explanation:} Bit có trọng số (weighted): mỗi bit trong số nhị phân có trọng số là 2\textasciicircum{}i. Bit cuối (LSB) có trọng số 2⁰=1. Bit đầu (MSB) có trọng số 2\textasciicircum{}(n-1). Tổng các bit × trọng số cho giá trị thập phân.

\par\vspace{0.3cm}
\hrule
\vspace{0.3cm}
\end{minipage}

\begin{minipage}{\textwidth}
\subsection*{Question 74 \quad \footnotesize{[0074\_q074\_Question\_74.png]}}
\noindent \textbf{When we want to store music in a computer, the audio signal must be}

\begin{enumerate}[label=\Alph*.]
    \item Sampled only
    \item Quantized only only
    \item Coded only
    \item \textbf{Sampled, quantized, and coded} \hfill \textcolor{blue}{\textbf{[Correct]}}
\end{enumerate}

\noindent \textbf{Correct Answer:} \textcolor{blue}{\textbf{D}} \\[2pt]
\noindent \textbf{Explanation:} Tràn số (overflow) xảy ra khi carry vào bit dấu khác carry ra khỏi bit dấu trong two's complement. Phát hiện overflow: XOR carry-in và carry-out của MSB. Nếu khác nhau → overflow.

\par\vspace{0.3cm}
\hrule
\vspace{0.3cm}
\end{minipage}

\begin{minipage}{\textwidth}
\subsection*{Question 75 \quad \footnotesize{[0075\_q075\_Question\_75.png]}}
\noindent \textbf{What is NOT a standard of vector graphic?}

\begin{enumerate}[label=\Alph*.]
    \item CAD
    \item FLASH
    \item \textbf{TrueType} \hfill \textcolor{blue}{\textbf{[Correct]}}
    \item PostScript
    \item None of others
\end{enumerate}

\noindent \textbf{Correct Answer:} \textcolor{blue}{\textbf{C}} \\[2pt]
\noindent \textbf{Explanation:} Các phép toán số học trong hệ nhị phân: cộng (addition), trừ (subtraction), nhân (multiplication), chia (division). ALU thực hiện các phép này qua mạch logic tổ hợp: half adder, full adder, ALU slice.

\par\vspace{0.3cm}
\hrule
\vspace{0.3cm}
\end{minipage}

\begin{minipage}{\textwidth}
\subsection*{Question 76 \quad \footnotesize{[0076\_q076\_Question\_76.png]}}
\noindent \textbf{Use a simple left shift operation on the bit pattern 1001 1000}

\begin{enumerate}[label=\Alph*.]
    \item 0001 1001
    \item 0001 1010
    \item \textbf{0011 0000} \hfill \textcolor{blue}{\textbf{[Correct]}}
    \item 0011 0001
\end{enumerate}

\noindent \textbf{Correct Answer:} \textcolor{blue}{\textbf{C}} \\[2pt]
\noindent \textbf{Explanation:} Excess-127 (biased-127) dùng trong IEEE 754 single precision: exponent thực = exponent lưu - 127. Phạm vi exponent thực: -126 đến 127 (exponent lưu 1-254). Exponent 0 và 255 dành cho giá trị đặc biệt (0, ∞, NaN).

\par\vspace{0.3cm}
\hrule
\vspace{0.3cm}
\end{minipage}

\begin{minipage}{\textwidth}
\subsection*{Question 77 \quad \footnotesize{[0077\_q077\_Question\_77.png]}}
\noindent \textbf{WAN is a network that connects two communicating devices through a transmission A medium (cable or air) SÖ DÁP ÁN CÀN CHON}

\begin{enumerate}[label=\Alph*.]
    \item Peer-to-peer
    \item Node -to-node
    \item \textbf{Point-to-point} \hfill \textcolor{blue}{\textbf{[Correct]}}
    \item Data link
\end{enumerate}

\noindent \textbf{Correct Answer:} \textcolor{blue}{\textbf{C}} \\[2pt]
\noindent \textbf{Explanation:} Truncation (cắt bớt) và rounding (làm tròn) khi chuyển số thực: truncation bỏ phần dư; rounding chọn giá trị gần nhất. IEEE 754 có 4 chế độ rounding: round to nearest (default), round toward 0, round toward +∞, round toward -∞.

\par\vspace{0.3cm}
\hrule
\vspace{0.3cm}
\end{minipage}

\begin{minipage}{\textwidth}
\subsection*{Question 78 \quad \footnotesize{[0078\_q078\_Question\_78.png]}}
\noindent \textbf{In physical layer, an has infinitely many levels of intensity over a period of time}

\begin{enumerate}[label=\Alph*.]
    \item \textbf{Analog signals} \hfill \textcolor{blue}{\textbf{[Correct]}}
    \item Electromagnetic signals
    \item Digital signals
    \item Electronic signals
\end{enumerate}

\noindent \textbf{Correct Answer:} \textcolor{blue}{\textbf{A}} \\[2pt]
\noindent \textbf{Explanation:} Gray code là mã nhị phân trong đó hai số liên tiếp chỉ khác nhau 1 bit. Dùng trong bộ mã hóa vị trí (rotary encoder) để tránh lỗi đọc. Chuyển từ nhị phân sang Gray: XOR bit thứ i với bit thứ i+1.

\par\vspace{0.3cm}
\hrule
\vspace{0.3cm}
\end{minipage}


\newpage
\section*{3.1 Data Storage}
\addcontentsline{toc}{section}{3.1 Data Storage}
\hrule\vspace{0.5cm}

\begin{minipage}{\textwidth}
\subsection*{Question 79 \quad \footnotesize{[0079\_q079\_Question\_79.png]}}
\noindent \textbf{At the transport layer, addresses are called, and these define the application-layer programs at the source and destination.}

\begin{enumerate}[label=\Alph*.]
    \item TCP
    \item \textbf{Port numbers} \hfill \textcolor{blue}{\textbf{[Correct]}}
    \item UDP
    \item IP address
\end{enumerate}

\noindent \textbf{Correct Answer:} \textcolor{blue}{\textbf{B}} \\[2pt]
\noindent \textbf{Explanation:} Dữ liệu trong máy tính được lưu trữ dưới dạng nhị phân (0 và 1). Mỗi bit là đơn vị thông tin nhỏ nhất. 8 bit = 1 byte. Byte là đơn vị cơ bản để đo dung lượng lưu trữ và địa chỉ bộ nhớ.

\par\vspace{0.3cm}
\hrule
\vspace{0.3cm}
\end{minipage}

\begin{minipage}{\textwidth}
\subsection*{Question 80 \quad \footnotesize{[0080\_q080\_Question\_80.png]}}
\noindent \textbf{is the standard protocol provided by TCP/IP for copying a file from one host to another andtransferring files from one system to another seems simple and straightforward.}

\begin{enumerate}[label=\Alph*.]
    \item \textbf{FTP} \hfill \textcolor{blue}{\textbf{[Correct]}}
    \item SMTP
    \item URL
    \item HTTP
\end{enumerate}

\noindent \textbf{Correct Answer:} \textcolor{blue}{\textbf{A}} \\[2pt]
\noindent \textbf{Explanation:} Bộ nhớ máy tính gồm RAM (Random Access Memory) - bộ nhớ khả biến, và ROM (Read Only Memory) - bộ nhớ bất biến. RAM lưu dữ liệu tạm thời khi máy hoạt động. ROM lưu firmware và dữ liệu khởi động.

\par\vspace{0.3cm}
\hrule
\vspace{0.3cm}
\end{minipage}

\begin{minipage}{\textwidth}
\subsection*{Question 81 \quad \footnotesize{[0081\_q081\_Question\_81.png]}}
\noindent \textbf{The is a protocol that is used to define how the client-server programs can be written to retrieve web pages from the Web.}

\begin{enumerate}[label=\Alph*.]
    \item FTP
    \item SMTP
    \item TELNET
    \item \textbf{HTTP} \hfill \textcolor{blue}{\textbf{[Correct]}}
\end{enumerate}

\noindent \textbf{Correct Answer:} \textcolor{blue}{\textbf{D}} \\[2pt]
\noindent \textbf{Explanation:} Số nguyên (integer) được lưu dưới dạng nhị phân. Số nguyên không dấu (unsigned): tất cả bit dùng cho giá trị. Số nguyên có dấu (signed): bit MSB là dấu, phần còn lại là giá trị. Two's complement là chuẩn cho số có dấu.

\par\vspace{0.3cm}
\hrule
\vspace{0.3cm}
\end{minipage}

\begin{minipage}{\textwidth}
\subsection*{Question 82 \quad \footnotesize{[0082\_q082\_Question\_82.png]}}
\noindent \textbf{All the instructions and procedures needed to start up a computer, to the point it can load an operating system, are stored in ROM chips. In IBM compatible PCs this is called the}

\begin{enumerate}[label=\Alph*.]
    \item RAM BIOS
    \item \textbf{ROM BIOS} \hfill \textcolor{blue}{\textbf{[Correct]}}
    \item EPROM
    \item POST
    \item None of the others
\end{enumerate}

\noindent \textbf{Correct Answer:} \textcolor{blue}{\textbf{B}} \\[2pt]
\noindent \textbf{Explanation:} Số thực dấu phẩy động (floating-point) lưu theo IEEE 754: số được biểu diễn dưới dạng ±1.mantissa × 2\textasciicircum{}exponent. Phần mantissa là phần lẻ, exponent được bias. Đây là cách lưu số thực phổ biến nhất.

\par\vspace{0.3cm}
\hrule
\vspace{0.3cm}
\end{minipage}

\begin{minipage}{\textwidth}
\subsection*{Question 83 \quad \footnotesize{[0083\_q083\_Question\_83.png]}}
\noindent \textbf{Which is an electronic device that can store temporarily a single bit of data}

\begin{enumerate}[label=\Alph*.]
    \item Accumulator
    \item Buffer
    \item Memory
    \item \textbf{Latch} \hfill \textcolor{blue}{\textbf{[Correct]}}
    \item None of the others
\end{enumerate}

\noindent \textbf{Correct Answer:} \textcolor{blue}{\textbf{D}} \\[2pt]
\noindent \textbf{Explanation:} Dữ liệu văn bản: mỗi ký tự được gán một mã số. ASCII dùng 7 bit cho 128 ký tự tiếng Anh. Unicode (UTF-8, UTF-16) mở rộng cho mọi ngôn ngữ. UTF-8 phổ biến trên web, mỗi ký tự 1-4 byte.

\par\vspace{0.3cm}
\hrule
\vspace{0.3cm}
\end{minipage}

\begin{minipage}{\textwidth}
\subsection*{Question 84 \quad \footnotesize{[0084\_q084\_Question\_84.png]}}
\noindent \textbf{Which technique is used in multiprogramming that needs to be determined size beforehand by the memory manager and with each program occupying a contiguous space?}

\begin{enumerate}[label=\Alph*.]
    \item \textbf{Partitioning} \hfill \textcolor{blue}{\textbf{[Correct]}}
    \item Paging
    \item Demand paging
    \item Demand segmentation
\end{enumerate}

\noindent \textbf{Correct Answer:} \textcolor{blue}{\textbf{A}} \\[2pt]
\noindent \textbf{Explanation:} Dữ liệu âm thanh số (digital audio): lấy mẫu tín hiệu analog theo tần số (sample rate, Hz), mã hóa mỗi mẫu bằng số bit (bit depth). VD: CD audio dùng 44.1 kHz, 16-bit. Nyquist theorem: tần số lấy mẫu >= 2× tần số cao nhất.

\par\vspace{0.3cm}
\hrule
\vspace{0.3cm}
\end{minipage}

\begin{minipage}{\textwidth}
\subsection*{Question 85 \quad \footnotesize{[0085\_q085\_Question\_85.png]}}
\noindent \textbf{A is a program in execution. It is a program that has started but has not finished}

\begin{enumerate}[label=\Alph*.]
    \item Program
    \item Job
    \item \textbf{Process} \hfill \textcolor{blue}{\textbf{[Correct]}}
    \item Task
\end{enumerate}

\noindent \textbf{Correct Answer:} \textcolor{blue}{\textbf{C}} \\[2pt]
\noindent \textbf{Explanation:} Dữ liệu hình ảnh: bitmap (raster) lưu từng pixel (picture element). Mỗi pixel có màu sắc. Màu đen trắng: 1 bit/pixel. Grayscale: 8 bit/pixel (256 mức xám). True Color: 24 bit/pixel (16.7 triệu màu RGB).

\par\vspace{0.3cm}
\hrule
\vspace{0.3cm}
\end{minipage}

\begin{minipage}{\textwidth}
\subsection*{Question 86 \quad \footnotesize{[0086\_q086\_Question\_86.png]}}
\noindent \textbf{The device manager will be tracking every device constantly to ensure that the device is}

\begin{enumerate}[label=\Alph*.]
    \item Ready
    \item Finished
    \item \textbf{Functioning properly} \hfill \textcolor{blue}{\textbf{[Correct]}}
    \item Running
\end{enumerate}

\noindent \textbf{Correct Answer:} \textcolor{blue}{\textbf{C}} \\[2pt]
\noindent \textbf{Explanation:} Nén dữ liệu (data compression) giảm kích thước file. Lossless (nén không mất dữ liệu): giữ nguyên dữ liệu gốc (ZIP, PNG). Lossy (nén mất dữ liệu): bỏ bớt dữ liệu ít quan trọng (JPEG, MP3).

\par\vspace{0.3cm}
\hrule
\vspace{0.3cm}
\end{minipage}

\begin{minipage}{\textwidth}
\subsection*{Question 87 \quad \footnotesize{[0087\_q087\_Question\_87.png]}}
\noindent \textbf{The list of coded instructions is called}

\begin{enumerate}[label=\Alph*.]
    \item \textbf{Computer program} \hfill \textcolor{blue}{\textbf{[Correct]}}
    \item Algorithm
    \item Flowchart
    \item Utility programs
    \item None of the others
\end{enumerate}

\noindent \textbf{Correct Answer:} \textcolor{blue}{\textbf{A}} \\[2pt]
\noindent \textbf{Explanation:} Dữ liệu video: chuỗi các frame (khung hình) liên tiếp. Mỗi frame là một ảnh. Nén video dùng MPEG, H.264, H.265. Tần số frame (fps) và độ phân giải (resolution) ảnh hưởng đến chất lượng và kích thước.

\par\vspace{0.3cm}
\hrule
\vspace{0.3cm}
\end{minipage}

\begin{minipage}{\textwidth}
\subsection*{Question 88 \quad \footnotesize{[0088\_q088\_Question\_88.png]}}
\noindent \textbf{is an ordered set of unambiguous steps that produces a result and terminates in a finite time.}

\begin{enumerate}[label=\Alph*.]
    \item A construct
    \item A recursion
    \item An iteration
    \item \textbf{An algorithm} \hfill \textcolor{blue}{\textbf{[Correct]}}
\end{enumerate}

\noindent \textbf{Correct Answer:} \textcolor{blue}{\textbf{D}} \\[2pt]
\noindent \textbf{Explanation:} Địa chỉ bộ nhớ: mỗi byte trong RAM có địa chỉ duy nhất (đánh số từ 0). CPU dùng address bus để chỉ định byte cần đọc/ghi. Dung lượng bộ nhớ = số địa chỉ × số byte/địa chỉ.

\par\vspace{0.3cm}
\hrule
\vspace{0.3cm}
\end{minipage}

\begin{minipage}{\textwidth}
\subsection*{Question 89 \quad \footnotesize{[0089\_q089\_Question\_89.png]}}
\noindent \textbf{We use a search for an unordered list.}

\begin{enumerate}[label=\Alph*.]
    \item \textbf{Sequential} \hfill \textcolor{blue}{\textbf{[Correct]}}
    \item Binary
    \item Bubble
    \item Insertion
\end{enumerate}

\noindent \textbf{Correct Answer:} \textcolor{blue}{\textbf{A}} \\[2pt]
\noindent \textbf{Explanation:} Little-endian vs Big-endian: cách sắp xếp byte trong bộ nhớ. Little-endian: byte thấp (LSB) ở địa chỉ thấp. Big-endian: byte cao (MSB) ở địa chỉ thấp. Intel dùng little-endian, Motorola dùng big-endian. Network byte order là big-endian.

\par\vspace{0.3cm}
\hrule
\vspace{0.3cm}
\end{minipage}

\begin{minipage}{\textwidth}
\subsection*{Question 90 \quad \footnotesize{[0090\_q090\_Question\_90.png]}}
\noindent \textbf{In sort, the smallest item moves to the beginning of the unsorted list. There is no one-to- one swapping.}

\begin{enumerate}[label=\Alph*.]
    \item \textbf{Selection} \hfill \textcolor{blue}{\textbf{[Correct]}}
    \item Bubble
    \item Insertion
    \item Every
\end{enumerate}

\noindent \textbf{Correct Answer:} \textcolor{blue}{\textbf{A}} \\[2pt]
\noindent \textbf{Explanation:} Parity bit (bit chẵn lẻ): bit thêm vào để phát hiện lỗi. Parity chẵn (even parity): tổng số bit 1 (kể cả parity) là chẵn. Parity lẻ (odd parity): tổng số bit 1 là lẻ. Phát hiện được 1 bit lỗi nhưng không sửa được.

\par\vspace{0.3cm}
\hrule
\vspace{0.3cm}
\end{minipage}

\begin{minipage}{\textwidth}
\subsection*{Question 91 \quad \footnotesize{[0091\_q091\_Question\_91.png]}}
\noindent \textbf{The analyzer checks the sentences created by the syntax analyzer to be sure that they contain no ambiguity.}

\begin{enumerate}[label=\Alph*.]
    \item Lexical
    \item Syntax
    \item \textbf{Semantic} \hfill \textcolor{blue}{\textbf{[Correct]}}
    \item Code generation
\end{enumerate}

\noindent \textbf{Correct Answer:} \textcolor{blue}{\textbf{C}} \\[2pt]
\noindent \textbf{Explanation:} Bộ nhớ cache chia làm cache L1 (trong CPU core, nhỏ nhất, nhanh nhất), L2, L3 (lớn nhất, chậm nhất). Cache hit: dữ liệu có trong cache. Cache miss: dữ liệu không có, phải xuống RAM.

\par\vspace{0.3cm}
\hrule
\vspace{0.3cm}
\end{minipage}

\begin{minipage}{\textwidth}
\subsection*{Question 92 \quad \footnotesize{[0092\_q092\_Question\_92.png]}}
\noindent \textbf{An instruction that transfers program control to one or more possible paths is known as}

\begin{enumerate}[label=\Alph*.]
    \item Utility program
    \item System software
    \item \textbf{Broadband channel} \hfill \textcolor{blue}{\textbf{[Correct]}}
    \item Application program
    \item None of the others
\end{enumerate}

\noindent \textbf{Correct Answer:} \textcolor{blue}{\textbf{C}} \\[2pt]
\noindent \textbf{Explanation:} Hệ thống file (file system) là cách tổ chức lưu trữ file trên ổ đĩa. Các hệ thống file: FAT32 (tối đa 4GB/file), NTFS (Windows), ext4 (Linux), APFS (Mac). Mỗi file có tên, kích thước, thuộc tính, và dữ liệu.

\par\vspace{0.3cm}
\hrule
\vspace{0.3cm}
\end{minipage}

\begin{minipage}{\textwidth}
\subsection*{Question 93 \quad \footnotesize{[0093\_q093\_Question\_93.png]}}
\noindent \textbf{What is language of procedural paradigm that is high level language and also has some low-level instructions?}

\begin{enumerate}[label=\Alph*.]
    \item COBOL
    \item \textbf{C} \hfill \textcolor{blue}{\textbf{[Correct]}}
    \item Pascal
    \item Ada
\end{enumerate}

\noindent \textbf{Correct Answer:} \textcolor{blue}{\textbf{B}} \\[2pt]
\noindent \textbf{Explanation:} Nén dữ liệu lossless: Run-length encoding (RLE) nén chuỗi lặp, Huffman coding (dùng cây nhị phân, mã ngắn cho ký tự phổ biến), LZW (tìm mẫu lặp). Các thuật toán này giảm kích thước mà không mất thông tin.

\par\vspace{0.3cm}
\hrule
\vspace{0.3cm}
\end{minipage}

\begin{minipage}{\textwidth}
\subsection*{Question 94 \quad \footnotesize{[0094\_q094\_Question\_94.png]}}
\noindent \textbf{A causes an action to be performed by the program. It translates directly into one or more executable computer instructions.}

\begin{enumerate}[label=\Alph*.]
    \item Expression
    \item \textbf{Operator} \hfill \textcolor{blue}{\textbf{[Correct]}}
    \item Variables
    \item Statements
\end{enumerate}

\noindent \textbf{Correct Answer:} \textcolor{blue}{\textbf{B, E}} \\[2pt]
\noindent \textbf{Explanation:} Các định dạng file thông dụng: .txt (văn bản thuần), .docx (Word - XML nén), .pdf (Portable Document Format), .html (web), .xml (dữ liệu có cấu trúc), .json (dữ liệu JavaScript). Mỗi định dạng có cấu trúc riêng.

\par\vspace{0.3cm}
\hrule
\vspace{0.3cm}
\end{minipage}

\begin{minipage}{\textwidth}
\subsection*{Question 95 \quad \footnotesize{[0095\_q095\_Question\_95.png]}}
\noindent \textbf{What is the analysis process in the system implementation phase that will use a program as an active agent that manipulates passive object?}

\begin{enumerate}[label=\Alph*.]
    \item \textbf{Procedure-oriented analysis} \hfill \textcolor{blue}{\textbf{[Correct]}}
    \item Structured analysis
    \item Classical analysis
    \item All of others
\end{enumerate}

\noindent \textbf{Correct Answer:} \textcolor{blue}{\textbf{A}} \\[2pt]
\noindent \textbf{Explanation:} Bộ nhớ máy tính: SRAM (Static RAM) dùng flip-flop, nhanh, đắt, dung lượng nhỏ, dùng cho cache. DRAM (Dynamic RAM) dùng tụ điện, chậm hơn, rẻ hơn, dung lượng lớn, cần refresh định kỳ.

\par\vspace{0.3cm}
\hrule
\vspace{0.3cm}
\end{minipage}

\begin{minipage}{\textwidth}
\subsection*{Question 96 \quad \footnotesize{[0096\_q096\_Question\_96.png]}}
\noindent \textbf{One phase in system development is}

\begin{enumerate}[label=\Alph*.]
    \item \textbf{Analysis} \hfill \textcolor{blue}{\textbf{[Correct]}}
    \item Application
    \item Designing
    \item Collecting
\end{enumerate}

\noindent \textbf{Correct Answer:} \textcolor{blue}{\textbf{A}} \\[2pt]
\noindent \textbf{Explanation:} Nén ảnh lossless: PNG (Portable Network Graphics) dùng thuật toán DEFLATE, nén không mất dữ liệu, hỗ trợ trong suốt (alpha channel). GIF (Graphics Interchange Format) nén 8-bit, hỗ trợ hoạt ảnh.

\par\vspace{0.3cm}
\hrule
\vspace{0.3cm}
\end{minipage}

\begin{minipage}{\textwidth}
\subsection*{Question 97 \quad \footnotesize{[0097\_q097\_Question\_97.png]}}
\noindent \textbf{is the breaking up of a large project into smaller parts.}

\begin{enumerate}[label=\Alph*.]
    \item Coupling
    \item Incrementing
    \item Obsolescence
    \item \textbf{Modularization} \hfill \textcolor{blue}{\textbf{[Correct]}}
\end{enumerate}

\noindent \textbf{Correct Answer:} \textcolor{blue}{\textbf{D}} \\[2pt]
\noindent \textbf{Explanation:} Nén ảnh lossy: JPEG (Joint Photographic Experts Group) dùng DCT (Discrete Cosine Transform) và lượng tử hóa, giảm chất lượng để giảm kích thước. WebP của Google cho tỷ lệ nén tốt hơn JPEG.

\par\vspace{0.3cm}
\hrule
\vspace{0.3cm}
\end{minipage}

\begin{minipage}{\textwidth}
\subsection*{Question 98 \quad \footnotesize{[0098\_q098\_Question\_98.png]}}
\noindent \textbf{In the software quality, what factor measures can be mentioned for transferability??}

\begin{enumerate}[label=\Alph*.]
    \item Accuracy
    \item Timeliness
    \item \textbf{Resusability} \hfill \textcolor{blue}{\textbf{[Correct]}}
    \item Changeability
\end{enumerate}

\noindent \textbf{Correct Answer:} \textcolor{blue}{\textbf{C}} \\[2pt]
\noindent \textbf{Explanation:} Cache mapping: (1) Direct-mapped: mỗi khối RAM chỉ vào một vị trí cache. (2) Fully associative: bất kỳ vị trí nào. (3) Set-associative: kết hợp cả hai. LRU (Least Recently Used) là thuật toán thay thế cache phổ biến.

\par\vspace{0.3cm}
\hrule
\vspace{0.3cm}
\end{minipage}


\newpage
\section*{3.2 Operations on Data}
\addcontentsline{toc}{section}{3.2 Operations on Data}
\hrule\vspace{0.5cm}

\begin{minipage}{\textwidth}
\subsection*{Question 99 \quad \footnotesize{[0099\_q099\_Question\_99.png]}}
\noindent \textbf{What is method in black-box testing which tests are designed to include representatives of boundary values in a range?}

\begin{enumerate}[label=\Alph*.]
    \item \textbf{Boundary-value testing} \hfill \textcolor{blue}{\textbf{[Correct]}}
    \item Random testing
    \item Exhaustive testing
    \item Data flow testing
\end{enumerate}

\noindent \textbf{Correct Answer:} \textcolor{blue}{\textbf{A}} \\[2pt]
\noindent \textbf{Explanation:} Phép toán số học cơ bản trên dữ liệu: cộng (ADD), trừ (SUB), nhân (MUL), chia (DIV). ALU thực hiện các phép này. Phép cộng là cơ bản nhất - các phép khác đều có thể suy ra từ phép cộng và dịch bit.

\par\vspace{0.3cm}
\hrule
\vspace{0.3cm}
\end{minipage}

\begin{minipage}{\textwidth}
\subsection*{Question 100 \quad \footnotesize{[0100\_q100\_Question\_100.png]}}
\noindent \textbf{array. An array that consists of just rows and columns is a}

\begin{enumerate}[label=\Alph*.]
    \item One-dimensional
    \item \textbf{Two-dimensional} \hfill \textcolor{blue}{\textbf{[Correct]}}
    \item Three-dimensional
    \item Multidimensional
\end{enumerate}

\noindent \textbf{Correct Answer:} \textcolor{blue}{\textbf{B}} \\[2pt]
\noindent \textbf{Explanation:} Phép toán logic: AND (tích logic - cả 2 đều 1 thì kết quả 1), OR (tổng logic - ít nhất 1 là 1 thì kết quả 1), NOT (phủ định - đảo bit), XOR (cộng modulo 2 - khác nhau thì 1). Các phép này là nền tảng của mọi tính toán.

\par\vspace{0.3cm}
\hrule
\vspace{0.3cm}
\end{minipage}

\begin{minipage}{\textwidth}
\subsection*{Question 101 \quad \footnotesize{[0101\_q101\_Question\_101.png]}}
\noindent \textbf{is a set of characters that is treated like an array in some languages and as a type in others.}

\begin{enumerate}[label=\Alph*.]
    \item \textbf{A string} \hfill \textcolor{blue}{\textbf{[Correct]}}
    \item A record
    \item A symbol
    \item A link
\end{enumerate}

\noindent \textbf{Correct Answer:} \textcolor{blue}{\textbf{A}} \\[2pt]
\noindent \textbf{Explanation:} Phép AND: 0 AND 0 = 0, 0 AND 1 = 0, 1 AND 0 = 0, 1 AND 1 = 1. Dùng để xóa bit (masking): AND với 0 để xóa bit, AND với 1 để giữ nguyên. Ứng dụng: kiểm tra bit cờ, lọc dữ liệu.

\par\vspace{0.3cm}
\hrule
\vspace{0.3cm}
\end{minipage}

\begin{minipage}{\textwidth}
\subsection*{Question 102 \quad \footnotesize{[0102\_q102\_Question\_102.png]}}
\noindent \textbf{The name of a linked list is the name of the head pointer that points to the \_of the list.}

\begin{enumerate}[label=\Alph*.]
    \item \textbf{First node} \hfill \textcolor{blue}{\textbf{[Correct]}}
    \item Second node
    \item Null node
    \item Last node
\end{enumerate}

\noindent \textbf{Correct Answer:} \textcolor{blue}{\textbf{A}} \\[2pt]
\noindent \textbf{Explanation:} Phép OR: 0 OR 0 = 0, 0 OR 1 = 1, 1 OR 0 = 1, 1 OR 1 = 1. Dùng để bật bit (setting): OR với 1 để bật bit, OR với 0 để giữ nguyên. Ứng dụng: bật cờ, kết hợp bit field.

\par\vspace{0.3cm}
\hrule
\vspace{0.3cm}
\end{minipage}

\begin{minipage}{\textwidth}
\subsection*{Question 103 \quad \footnotesize{[0103\_q103\_Question\_103.png]}}
\noindent \textbf{The pop operation of the stack.}

\begin{enumerate}[label=\Alph*.]
    \item \textbf{Deletes an item from the top} \hfill \textcolor{blue}{\textbf{[Correct]}}
    \item Deletes an item from the bottom
    \item Inserts an item at the top
    \item Inserts an item at the bottom
\end{enumerate}

\noindent \textbf{Correct Answer:} \textcolor{blue}{\textbf{A}} \\[2pt]
\noindent \textbf{Explanation:} Phép XOR: 0 XOR 0 = 0, 0 XOR 1 = 1, 1 XOR 0 = 1, 1 XOR 1 = 0. Còn gọi là exclusive OR. Dùng để đảo bit (toggling): XOR với 1 để đảo bit. Ứng dụng: mã hóa đơn giản, kiểm tra parity.

\par\vspace{0.3cm}
\hrule
\vspace{0.3cm}
\end{minipage}

\begin{minipage}{\textwidth}
\subsection*{Question 104 \quad \footnotesize{[0104\_q104\_Question\_104.png]}}
\noindent \textbf{In preorder traversal of a binary tree, the}

\begin{enumerate}[label=\Alph*.]
    \item Left subtree is processed first
    \item Right subtree is processed first
    \item \textbf{Root is processed first} \hfill \textcolor{blue}{\textbf{[Correct]}}
    \item The root is never processed
\end{enumerate}

\noindent \textbf{Correct Answer:} \textcolor{blue}{\textbf{C}} \\[2pt]
\noindent \textbf{Explanation:} Phép NOT (complement): đảo tất cả bit. NOT 0 = 1, NOT 1 = 0. Trong two's complement, NOT là bước đầu tiên để tính số âm (NOT x + 1 = -x). Ứng dụng: tạo mặt nạ bit, đảo cờ.

\par\vspace{0.3cm}
\hrule
\vspace{0.3cm}
\end{minipage}

\begin{minipage}{\textwidth}
\subsection*{Question 105 \quad \footnotesize{[0105\_q105\_Question\_105.png]}}
\noindent \textbf{The enqueue operation \_an item at the \_ of the queue.}

\begin{enumerate}[label=\Alph*.]
    \item Delete/front
    \item Delete/end
    \item Inserts/front
    \item \textbf{Insert/rear} \hfill \textcolor{blue}{\textbf{[Correct]}}
\end{enumerate}

\noindent \textbf{Correct Answer:} \textcolor{blue}{\textbf{D}} \\[2pt]
\noindent \textbf{Explanation:} Shift (dịch chuyển) bit: logical shift dịch và thêm 0 vào bên trái/phải; arithmetic shift giữ dấu (MSB không đổi). Dịch trái 1 bit = nhân với 2. Dịch phải 1 bit = chia lấy nguyên cho 2.

\par\vspace{0.3cm}
\hrule
\vspace{0.3cm}
\end{minipage}

\begin{minipage}{\textwidth}
\subsection*{Question 106 \quad \footnotesize{[0106\_q106\_Question\_106.png]}}
\noindent \textbf{A is a collection of data stored in the internal format of the computer. This contains data that is meaningful only if it is properly interpreted by a program.}

\begin{enumerate}[label=\Alph*.]
    \item Text file
    \item \textbf{Bianry file} \hfill \textcolor{blue}{\textbf{[Correct]}}
    \item Sequential file
    \item EOF
\end{enumerate}

\noindent \textbf{Correct Answer:} \textcolor{blue}{\textbf{B}} \\[2pt]
\noindent \textbf{Explanation:} Arithmetic shift (dịch số học): shift right giữ nguyên bit dấu (bản sao MSB), shift left thêm 0. Logical shift: thêm 0 không điều kiện. Dịch trái k bit: nhân với 2\textasciicircum{}k. Dịch phải k bit: chia cho 2\textasciicircum{}k (với số không dấu).

\par\vspace{0.3cm}
\hrule
\vspace{0.3cm}
\end{minipage}

\begin{minipage}{\textwidth}
\subsection*{Question 107 \quad \footnotesize{[0107\_q107\_Question\_107.png]}}
\noindent \textbf{In the hashing method, there are no synonyms or collisions.}

\begin{enumerate}[label=\Alph*.]
    \item \textbf{Direct} \hfill \textcolor{blue}{\textbf{[Correct]}}
    \item Division remainder
    \item Modulo division
    \item Digit extraction
\end{enumerate}

\noindent \textbf{Correct Answer:} \textcolor{blue}{\textbf{A}} \\[2pt]
\noindent \textbf{Explanation:} Rotate (quay) bit: các bit dịch ra được đưa vào đầu bên kia. Rotate left: bit MSB → vị trí LSB. Rotate right: bit LSB → vị trí MSB. Khác với shift ở chỗ không mất bit. Dùng trong mã hóa và checksum.

\par\vspace{0.3cm}
\hrule
\vspace{0.3cm}
\end{minipage}

\begin{minipage}{\textwidth}
\subsection*{Question 108 \quad \footnotesize{[0108\_q108\_Question\_108.png]}}
\noindent \textbf{The address produced by a hashing algorithm is the address.}

\begin{enumerate}[label=\Alph*.]
    \item Probe
    \item Synonym
    \item Collision
    \item \textbf{Home} \hfill \textcolor{blue}{\textbf{[Correct]}}
\end{enumerate}

\noindent \textbf{Correct Answer:} \textcolor{blue}{\textbf{D}} \\[2pt]
\noindent \textbf{Explanation:} Phép cộng nhị phân có thể được thực hiện bằng mạch half adder (HA: cộng 2 bit, ra sum và carry) và full adder (FA: cộng 3 bit - 2 bit cộng + carry in). Ripple-carry adder: nhiều FA nối tiếp.

\par\vspace{0.3cm}
\hrule
\vspace{0.3cm}
\end{minipage}

\begin{minipage}{\textwidth}
\subsection*{Question 109 \quad \footnotesize{[0109\_q109\_Question\_109.png]}}
\noindent \textbf{A is a collection of related logically coherent data used by the application programs in an organization.}

\begin{enumerate}[label=\Alph*.]
    \item \textbf{Database} \hfill \textcolor{blue}{\textbf{[Correct]}}
    \item File
    \item Directory
    \item Flat-file
\end{enumerate}

\noindent \textbf{Correct Answer:} \textcolor{blue}{\textbf{A}} \\[2pt]
\noindent \textbf{Explanation:} Flag (cờ) trong CPU: carry flag (CF) - có nhớ từ MSB; zero flag (ZF) - kết quả = 0; sign flag (SF) - kết quả âm; overflow flag (OF) - tràn số có dấu. Các flag này ở thanh ghi trạng thái (status register/FLAGS).

\par\vspace{0.3cm}
\hrule
\vspace{0.3cm}
\end{minipage}

\begin{minipage}{\textwidth}
\subsection*{Question 110 \quad \footnotesize{[0110\_q110\_Question\_110.png]}}
\noindent \textbf{In the relational database model, data is organized in called relations.}

\begin{enumerate}[label=\Alph*.]
    \item Two-dimensional array
    \item \textbf{Two-dimensional tables} \hfill \textcolor{blue}{\textbf{[Correct]}}
    \item Two-dimensional record
    \item None of others
\end{enumerate}

\noindent \textbf{Correct Answer:} \textcolor{blue}{\textbf{B}} \\[2pt]
\noindent \textbf{Explanation:} Phép trừ nhị phân: dùng two's complement. a - b = a + two's complement(b). Two's complement của b = NOT b + 1. Cách này đơn giản hóa phần cứng vì chỉ cần mạch cộng và mạch đảo bit.

\par\vspace{0.3cm}
\hrule
\vspace{0.3cm}
\end{minipage}

\begin{minipage}{\textwidth}
\subsection*{Question 111 \quad \footnotesize{[0111\_q111\_Question\_111.png]}}
\noindent \textbf{Each row in a relation is called SÖ DÁP ÁN CÃN CHON}

\begin{enumerate}[label=\Alph*.]
    \item An attribute
    \item \textbf{A tuple} \hfill \textcolor{blue}{\textbf{[Correct]}}
    \item A union
    \item An attitude
\end{enumerate}

\noindent \textbf{Correct Answer:} \textcolor{blue}{\textbf{B}} \\[2pt]
\noindent \textbf{Explanation:} Phép nhân trong máy tính dùng mạch multiplier: thuật toán Booth và Wallace tree. Nhân số n-bit cho kết quả 2n-bit. Phần cứng nhân gồm ALU, shift register, và accumulator.

\par\vspace{0.3cm}
\hrule
\vspace{0.3cm}
\end{minipage}

\begin{minipage}{\textwidth}
\subsection*{Question 112 \quad \footnotesize{[0112\_q112\_Question\_112.png]}}
\noindent \textbf{operation. If you want to change the value of an attribute of a tuple, you use the}

\begin{enumerate}[label=\Alph*.]
    \item Project
    \item Join
    \item \textbf{Update} \hfill \textcolor{blue}{\textbf{[Correct]}}
    \item Select
\end{enumerate}

\noindent \textbf{Correct Answer:} \textcolor{blue}{\textbf{C}} \\[2pt]
\noindent \textbf{Explanation:} Phép chia trong máy tính: thuật toán chia phục hồi (restoring) và không phục hồi (non-restoring). Số bị chia (dividend) / số chia (divisor) = thương (quotient) + dư (remainder). Phần cứng dùng ALU và subtractor.

\par\vspace{0.3cm}
\hrule
\vspace{0.3cm}
\end{minipage}

\begin{minipage}{\textwidth}
\subsection*{Question 113 \quad \footnotesize{[0113\_q113\_Question\_113.png]}}
\noindent \textbf{is the process by which a given set of relations are transformed to a new set of relations with a more solid structure.}

\begin{enumerate}[label=\Alph*.]
    \item Relational database
    \item Database system
    \item Atomic operations
    \item \textbf{Normalization} \hfill \textcolor{blue}{\textbf{[Correct]}}
\end{enumerate}

\noindent \textbf{Correct Answer:} \textcolor{blue}{\textbf{D}} \\[2pt]
\noindent \textbf{Explanation:} Bit masking: dùng AND để xóa bit, OR để bật bit, XOR để đảo bit. Mặt nạ (mask) là mẫu bit dùng để thao tác trên các bit cụ thể. Ví dụ: 0x0F (00001111₂) mask để lấy 4 bit thấp.

\par\vspace{0.3cm}
\hrule
\vspace{0.3cm}
\end{minipage}

\begin{minipage}{\textwidth}
\subsection*{Question 114 \quad \footnotesize{[0114\_q114\_Question\_114.png]}}
\noindent \textbf{Which of the following attacks is threatening confidentiality? SÖ DÁP ÁN CÀN CHON}

\begin{enumerate}[label=\Alph*.]
    \item \textbf{Snooping} \hfill \textcolor{blue}{\textbf{[Correct]}}
    \item Repudiation
    \item Denial of service
    \item Modification
\end{enumerate}

\noindent \textbf{Correct Answer:} \textcolor{blue}{\textbf{A}} \\[2pt]
\noindent \textbf{Explanation:} Set (bật), clear (xóa), toggle (đảo) bit: set bit i: OR với (1<<i). Clear bit i: AND với NOT(1<<i). Toggle bit i: XOR với (1<<i). Kiểm tra bit i: AND với (1<<i) và kiểm tra kết quả khác 0.

\par\vspace{0.3cm}
\hrule
\vspace{0.3cm}
\end{minipage}

\begin{minipage}{\textwidth}
\subsection*{Question 115 \quad \footnotesize{[0115\_q115\_Question\_115.png]}}
\noindent \textbf{means concealing the message by covering it with something else.}

\begin{enumerate}[label=\Alph*.]
    \item Cryptography
    \item \textbf{Steganography} \hfill \textcolor{blue}{\textbf{[Correct]}}
    \item Compressing
    \item Authentication
\end{enumerate}

\noindent \textbf{Correct Answer:} \textcolor{blue}{\textbf{B}} \\[2pt]
\noindent \textbf{Explanation:} Các phép toán trên bit rất nhanh vì hoạt động trực tiếp ở mức phần cứng. Ứng dụng: cờ (flags), mã hóa (XOR cipher), kiểm tra parity, kiểm tra số chẵn/lẻ (x \& 1), nén dữ liệu, đồ họa máy tính.

\par\vspace{0.3cm}
\hrule
\vspace{0.3cm}
\end{minipage}


\newpage
\section*{4. Computer Networks and Internet}
\addcontentsline{toc}{section}{4. Computer Networks and Internet}
\hrule\vspace{0.5cm}

\begin{minipage}{\textwidth}
\subsection*{Question 116 \quad \footnotesize{[0116\_q116\_Question\_116.png]}}
\noindent \textbf{is an attack that reduce the capability of a computer system to function correctly or bring the system down altogether by exhausting its resources.}

\begin{enumerate}[label=\Alph*.]
    \item \textbf{DoS} \hfill \textcolor{blue}{\textbf{[Correct]}}
    \item Trojan horses
    \item Worms
    \item Viruses
\end{enumerate}

\noindent \textbf{Correct Answer:} \textcolor{blue}{\textbf{A}} \\[2pt]
\noindent \textbf{Explanation:} Tấn công từ chối dịch vụ (DoS - Denial of Service) là hình thức tấn công mạng nhằm làm cạn kiệt tài nguyên của hệ thống mục tiêu (như băng thông mạng, bộ nhớ hoặc năng lực xử lý của CPU), khiến hệ thống bị quá tải, suy giảm hiệu năng nghiêm trọng hoặc tê liệt hoàn toàn đối với người dùng hợp lệ (đáp án A). Trong khi đó, virus, sâu máy tính (worms) và trojan horse là các loại phần mềm độc hại (malware) tập trung vào việc tự nhân bản, phá hoại dữ liệu hoặc mở cửa hậu đánh cắp thông tin chứ không có định nghĩa chính là làm sập hệ thống bằng cách vắt kiệt tài nguyên.

\par\vspace{0.3cm}
\hrule
\vspace{0.3cm}
\end{minipage}

\begin{minipage}{\textwidth}
\subsection*{Question 117 \quad \footnotesize{[0117\_q117\_Question\_117.png]}}
\noindent \textbf{is the amateur. Usually their techniques are deployed out of ill motives such as revenge attacks SÖ DÁP ÁN CÀN CHON}

\begin{enumerate}[label=\Alph*.]
    \item Black Hat Hacker
    \item Grey Hat Hacker
    \item Red Hat Hacker
    \item \textbf{Blue Hat Hacker} \hfill \textcolor{blue}{\textbf{[Correct]}}
\end{enumerate}

\noindent \textbf{Correct Answer:} \textcolor{blue}{\textbf{D}} \\[2pt]
\noindent \textbf{Explanation:} Trong phân loại hacker theo động cơ và trình độ, Blue Hat Hacker (hacker mũ xanh) thường dùng để chỉ những người nghiệp dư (amateur) thực hiện các cuộc tấn công mạng không phải vì mục đích trục lợi tài chính chuyên nghiệp mà chủ yếu do động cơ cá nhân như muốn trả đũa hoặc gây chú ý (đáp án D). Ngược lại, Black Hat Hacker là những tin tặc chuyên nghiệp hoạt động phi pháp với kỹ thuật cao vì tư lợi, còn Grey Hat Hacker nằm ở ranh giới giữa bảo mật và xâm nhập bất hợp pháp nhưng thường không mang dã tâm phá hoại để trả thù.

\par\vspace{0.3cm}
\hrule
\vspace{0.3cm}
\end{minipage}

\begin{minipage}{\textwidth}
\subsection*{Question 118 \quad \footnotesize{[0118\_q118\_Question\_118.png]}}
\noindent \textbf{In a block cipher a is used to encrypt the whole block even if the key is made of multiple values}

\begin{enumerate}[label=\Alph*.]
    \item Key block
    \item \textbf{Single key} \hfill \textcolor{blue}{\textbf{[Correct]}}
    \item Key stream
    \item Auto key
\end{enumerate}

\noindent \textbf{Correct Answer:} \textcolor{blue}{\textbf{B}} \\[2pt]
\noindent \textbf{Explanation:} Trong mã hóa khối (block cipher), một khóa duy nhất (single key) được sử dụng để mã hóa toàn bộ một khối dữ liệu có kích thước cố định, ngay cả khi khóa chính này được thuật toán mở rộng thành nhiều khóa con (subkeys) ở các vòng mã hóa (đáp án B). Điều này phân biệt hoàn toàn với mã hóa dòng (stream cipher), vốn sử dụng một luồng khóa liên tục (key stream) để mã hóa tuần tự từng bit hoặc từng byte một. Các thuật ngữ như 'key block' hay 'auto key' không phải là cơ chế mã hóa nguyên khối chuẩn trong định nghĩa này.

\par\vspace{0.3cm}
\hrule
\vspace{0.3cm}
\end{minipage}

\begin{minipage}{\textwidth}
\subsection*{Question 119 \quad \footnotesize{[0119\_q119\_Question\_119.png]}}
\noindent \textbf{The RSA algorithm for confidentiality uses cryptography.}

\begin{enumerate}[label=\Alph*.]
    \item \textbf{Asymmetric-key} \hfill \textcolor{blue}{\textbf{[Correct]}}
    \item Symmetric-key
    \item Substitution
    \item Transposition
\end{enumerate}

\noindent \textbf{Correct Answer:} \textcolor{blue}{\textbf{A}} \\[2pt]
\noindent \textbf{Explanation:} Thuật toán RSA (do Rivest, Shamir và Adleman phát minh) là thuật toán mật mã hóa khóa bất đối xứng (asymmetric-key cryptography / public-key cryptography) kinh điển, dựa trên độ phức tạp của bài toán phân tích các số nguyên lớn ra thừa số nguyên tố (đáp án A). Hệ mật mã này sử dụng một cặp khóa riêng biệt gồm khóa công khai (public key) để mã hóa và khóa riêng tư (private key) để giải mã. Ngược lại, mật mã khóa đối xứng (symmetric-key) như AES hay DES sử dụng chung một khóa bí mật cho cả hai chiều mã hóa và giải mã.

\par\vspace{0.3cm}
\hrule
\vspace{0.3cm}
\end{minipage}

\begin{minipage}{\textwidth}
\subsection*{Question 120 \quad \footnotesize{[0120\_q120\_Question\_120.png]}}
\noindent \textbf{separated the programming task from computer operation tasks.}

\begin{enumerate}[label=\Alph*.]
    \item Algorithms
    \item Data processors
    \item \textbf{High-level programming languages} \hfill \textcolor{blue}{\textbf{[Correct]}}
    \item Operating systems
\end{enumerate}

\noindent \textbf{Correct Answer:} \textcolor{blue}{\textbf{C}} \\[2pt]
\noindent \textbf{Explanation:} Theo lịch sử phát triển máy tính trong giáo trình Forouzan, sự ra đời của các ngôn ngữ lập trình bậc cao (High-level programming languages như FORTRAN, COBOL) đã tách biệt nhiệm vụ lập trình ra khỏi các thao tác vận hành phần cứng máy tính trực tiếp (đáp án C). Trước đó, lập trình viên buộc phải hiểu rõ kiến trúc vật lý và viết mã máy/hợp ngữ gắn liền với từng cỗ máy cụ thể; ngôn ngữ bậc cao giúp trừu tượng hóa phần cứng để người viết mã chỉ cần tập trung vào giải thuật. Hệ điều hành ra đời để quản lý tài nguyên, nhưng việc giải phóng tư duy lập trình khỏi chi tiết vận hành máy tính thuộc về ngôn ngữ bậc cao.

\par\vspace{0.3cm}
\hrule
\vspace{0.3cm}
\end{minipage}

\begin{minipage}{\textwidth}
\subsection*{Question 121 \quad \footnotesize{[0121\_q121\_Question\_121.png]}}
\noindent \textbf{A 17th-century computing machine that could perform addition and subtraction was the}

\begin{enumerate}[label=\Alph*.]
    \item \textbf{Pascaline} \hfill \textcolor{blue}{\textbf{[Correct]}}
    \item Jacquard loom
    \item Analytical Engine
    \item Babbage machine
\end{enumerate}

\noindent \textbf{Correct Answer:} \textcolor{blue}{\textbf{A}} \\[2pt]
\noindent \textbf{Explanation:} Vào thế kỷ 17 (năm 1642), nhà toán học và triết học người Pháp Blaise Pascal đã chế tạo ra cỗ máy tính cơ học đầu tiên mang tên Pascaline, có khả năng thực hiện tự động các phép tính cộng và trừ trực tiếp thông qua hệ thống bánh răng xoay (đáp án A). Các cỗ máy nổi tiếng khác như máy dệt Jacquard loom hay cỗ máy phân tích Analytical Engine của Charles Babbage đều ra đời muộn hơn nhiều vào thế kỷ 19. Vì vậy, Pascaline là cỗ máy tính toán cơ học tiêu biểu của thế kỷ 17.

\par\vspace{0.3cm}
\hrule
\vspace{0.3cm}
\end{minipage}

\begin{minipage}{\textwidth}
\subsection*{Question 122 \quad \footnotesize{[0122\_q122\_Question\_122.png]}}
\noindent \textbf{After the instruction is decoded, the control unit sends the task order to a component in the CPU. What is this phase?.}

\begin{enumerate}[label=\Alph*.]
    \item Fetch phase
    \item Decode phase
    \item \textbf{Excute phase} \hfill \textcolor{blue}{\textbf{[Correct]}}
    \item None of others
\end{enumerate}

\noindent \textbf{Correct Answer:} \textcolor{blue}{\textbf{C}} \\[2pt]
\noindent \textbf{Explanation:} Chu kỳ thực thi lệnh (instruction cycle) của CPU bao gồm ba pha cơ bản: Nạp lệnh (Fetch), Giải mã lệnh (Decode) và Thực thi lệnh (Execute). Sau khi lệnh đã được giải mã xong trong pha Decode, Đơn vị điều khiển (Control Unit) sẽ phát tín hiệu điều khiển gửi lệnh thao tác tới các thành phần phù hợp trong CPU (như ALU hay các thanh ghi) để tiến hành xử lý; pha này được gọi là pha thực thi (Execute phase, đáp án C). Pha Fetch chỉ nạp mã lệnh từ bộ nhớ vào thanh ghi, còn pha Decode làm nhiệm vụ phân tích ý nghĩa của lệnh.

\par\vspace{0.3cm}
\hrule
\vspace{0.3cm}
\end{minipage}

\begin{minipage}{\textwidth}
\subsection*{Question 123 \quad \footnotesize{[0123\_q123\_Question\_123.png]}}
\noindent \textbf{RAM is used as a short memory because it is}

\begin{enumerate}[label=\Alph*.]
    \item \textbf{Volatile} \hfill \textcolor{blue}{\textbf{[Correct]}}
    \item Has small capacity
    \item Is very expensive
    \item Is programmable
    \item None of the others
\end{enumerate}

\noindent \textbf{Correct Answer:} \textcolor{blue}{\textbf{A}} \\[2pt]
\noindent \textbf{Explanation:} RAM (Bộ nhớ truy cập ngẫu nhiên) được sử dụng làm bộ nhớ tạm thời / ngắn hạn (short-term memory) vì nó có tính chất khả biến (volatile), nghĩa là mọi dữ liệu lưu trữ trên RAM sẽ lập tức biến mất khi ngắt nguồn điện cấp cho máy tính (đáp án A). Ngược lại, các thiết bị lưu trữ thứ cấp như ổ đĩa cứng (HDD) hay ổ SSD là bộ nhớ bất biến (non-volatile), cho phép lưu trữ thông tin vĩnh viễn ngay cả khi tắt máy. Đặc tính dễ bay hơi (khả biến) này là lý do dữ liệu trong RAM chỉ tồn tại tạm thời trong phiên làm việc.

\par\vspace{0.3cm}
\hrule
\vspace{0.3cm}
\end{minipage}

\begin{minipage}{\textwidth}
\subsection*{Question 124 \quad \footnotesize{[0124\_q124\_Question\_124.png]}}
\noindent \textbf{For surface organization of magnetic tape, the width of the tape is divided into nine tracks. Nine vertical locations can store 8 bits of information related to a byte plus a bit for}

\begin{enumerate}[label=\Alph*.]
    \item Next block
    \item DataUO
    \item Next tape information
    \item \textbf{Error detection} \hfill \textcolor{blue}{\textbf{[Correct]}}
\end{enumerate}

\noindent \textbf{Correct Answer:} \textcolor{blue}{\textbf{D}} \\[2pt]
\noindent \textbf{Explanation:} Trong tổ chức bề mặt của băng từ loại 9 rãnh (9-track magnetic tape), chiều rộng của băng được chia làm 9 rãnh song song. Tại mỗi lát cắt thẳng đứng, 8 rãnh được sử dụng để lưu 8 bit cấu thành một byte dữ liệu, và rãnh thứ 9 còn lại được sử dụng làm bit kiểm tra chẵn lẻ (parity bit) phục vụ mục đích phát hiện lỗi dữ liệu (Error detection, đáp án D). Kỹ thuật này giúp hệ thống kịp thời phát hiện các lỗi suy hao từ tính hoặc đọc sai lệch khi cuộn băng quay với tốc độ cao.

\par\vspace{0.3cm}
\hrule
\vspace{0.3cm}
\end{minipage}

\begin{minipage}{\textwidth}
\subsection*{Question 125 \quad \footnotesize{[0125\_q125\_Question\_125.png]}}
\noindent \textbf{IBM-PC is an example of}

\begin{enumerate}[label=\Alph*.]
    \item A main frame
    \item Machine-in-built
    \item \textbf{Micro computer} \hfill \textcolor{blue}{\textbf{[Correct]}}
    \item Special purpose
    \item None of the others
\end{enumerate}

\noindent \textbf{Correct Answer:} \textcolor{blue}{\textbf{C}} \\[2pt]
\noindent \textbf{Explanation:} IBM-PC (ra đời năm 1981) là biểu tượng kinh điển của dòng máy vi tính (microcomputer / personal computer - PC), loại máy tính được thiết kế dựa trên bộ vi xử lý (microprocessor) đơn chip phục vụ cho nhu cầu sử dụng của cá nhân (đáp án C). Ngược lại, máy tính lớn (mainframe) là các hệ thống máy tính khổng lồ phục vụ xử lý dữ liệu quy mô lớn cho doanh nghiệp, còn máy chuyên dụng (special-purpose) được chế tạo chỉ để thực thi một tác vụ cố định duy nhất. Do đó, IBM-PC chính xác là một microcomputer.

\par\vspace{0.3cm}
\hrule
\vspace{0.3cm}
\end{minipage}

\begin{minipage}{\textwidth}
\subsection*{Question 126 \quad \footnotesize{[0126\_q126\_Question\_126.png]}}
\noindent \textbf{Computers are extremely fast and have fantastic memories. However, the only thing they can remember is}

\begin{enumerate}[label=\Alph*.]
    \item Instruction
    \item \textbf{Series of 1s and 0s} \hfill \textcolor{blue}{\textbf{[Correct]}}
    \item Boolean algebra
    \item Logic theorems
    \item None of the others
\end{enumerate}

\noindent \textbf{Correct Answer:} \textcolor{blue}{\textbf{B}} \\[2pt]
\noindent \textbf{Explanation:} Về mặt bản chất phần cứng kỹ thuật số, bộ nhớ máy tính chỉ có khả năng lưu trữ thông tin dưới dạng các mức điện áp đại diện cho hai trạng thái nhị phân: bit 1 và bit 0 (Series of 1s and 0s, đáp án B). Mọi dạng dữ liệu phức tạp từ văn bản, âm thanh, hình ảnh, video cho đến các mã lệnh chương trình (instructions) đều phải được mã hóa thành các chuỗi số nhị phân 0 và 1 trước khi được máy tính lưu trữ và xử lý. Đại số Boole hay định lý logic chỉ là công cụ toán học mô hình hóa chứ không phải dạng lưu trữ vật lý của bộ nhớ.

\par\vspace{0.3cm}
\hrule
\vspace{0.3cm}
\end{minipage}

\begin{minipage}{\textwidth}
\subsection*{Question 127 \quad \footnotesize{[0127\_q127\_Question\_127.png]}}
\noindent \textbf{A group contains bits in the binary system are represented as one digit in the octal system. SÓ DÁP ÁN CÃN CHN}

\begin{enumerate}[label=\Alph*.]
    \item 1
    \item 2
    \item \textbf{3} \hfill \textcolor{blue}{\textbf{[Correct]}}
    \item 4
    \item None of the others
\end{enumerate}

\noindent \textbf{Correct Answer:} \textcolor{blue}{\textbf{C}} \\[2pt]
\noindent \textbf{Explanation:} Hệ bát phân (octal) có cơ số là 8, và vì 2\textasciicircum{}3 = 8 nên mỗi chữ số trong hệ bát phân (có giá trị từ 0 đến 7) biểu diễn tương đương chính xác một nhóm gồm 3 bit trong hệ nhị phân (đáp án C). Ví dụ, chuỗi nhị phân 101 tương ứng với chữ số 5 trong hệ bát phân. Trong khi đó, nhóm 4 bit nhị phân (nibble) được dùng để biểu diễn một chữ số trong hệ thập lục phân (hexadecimal - cơ số 16).

\par\vspace{0.3cm}
\hrule
\vspace{0.3cm}
\end{minipage}

\begin{minipage}{\textwidth}
\subsection*{Question 128 \quad \footnotesize{[0128\_q128\_Question\_128.png]}}
\noindent \textbf{What is the result of storing -25 in an 8-bit memory location using two's complement representation?}

\begin{enumerate}[label=\Alph*.]
    \item 00011001
    \item 00110010
    \item \textbf{11100111} \hfill \textcolor{blue}{\textbf{[Correct]}}
    \item 11110010
\end{enumerate}

\noindent \textbf{Correct Answer:} \textcolor{blue}{\textbf{C}} \\[2pt]
\noindent \textbf{Explanation:} Để tìm biểu diễn bù 2 (two's complement) 8 bit của số âm -25, trước hết ta biểu diễn số dương +25 ở dạng nhị phân 8 bit là 00011001 (16 + 8 + 1). Tiếp theo, ta lấy bù 1 bằng cách đảo tất cả các bit thành 11100110, sau đó cộng thêm 1 để được giá trị bù 2 là 11100111 (đáp án C). Phương án A (00011001) chính là biểu diễn của số +25, do đó không chính xác cho số âm.

\par\vspace{0.3cm}
\hrule
\vspace{0.3cm}
\end{minipage}

\begin{minipage}{\textwidth}
\subsection*{Question 129 \quad \footnotesize{[0129\_q129\_Question\_129.png]}}
\noindent \textbf{The ASCII code for the character J is:}

\begin{enumerate}[label=\Alph*.]
    \item 1001 0001
    \item 1001 1010
    \item \textbf{0100 1010} \hfill \textcolor{blue}{\textbf{[Correct]}}
    \item 1010 0001
    \item None of the others
\end{enumerate}

\noindent \textbf{Correct Answer:} \textcolor{blue}{\textbf{C}} \\[2pt]
\noindent \textbf{Explanation:} Trong bảng mã chuẩn ASCII 7 bit biểu diễn trên một byte 8 bit (với bit đầu tiên bằng 0), chữ cái in hoa 'A' có mã thập phân là 65 (01000001). Vì 'J' là chữ cái thứ 10 trong bảng chữ cái tiếng Anh, mã của 'J' bằng 65 + 9 = 74 trong hệ thập phân; đổi 74 sang nhị phân 8 bit ta được 0100 1010 (64 + 8 + 2, đáp án C). Các đáp án bắt đầu bằng bit 1 (A, B, D) đều không thuộc phạm vi mã ASCII chuẩn 7 bit.

\par\vspace{0.3cm}
\hrule
\vspace{0.3cm}
\end{minipage}

\begin{minipage}{\textwidth}
\subsection*{Question 130 \quad \footnotesize{[0130\_q130\_Question\_130.png]}}
\noindent \textbf{In the graphic method of representing an image in a computer, re-scaling of the image creates a ragged or grainy image.}

\begin{enumerate}[label=\Alph*.]
    \item \textbf{Bitmap} \hfill \textcolor{blue}{\textbf{[Correct]}}
    \item Vector
    \item Quantized
    \item Binary
\end{enumerate}

\noindent \textbf{Correct Answer:} \textcolor{blue}{\textbf{A}} \\[2pt]
\noindent \textbf{Explanation:} Trong phương pháp đồ họa raster/bitmap (đáp án A), hình ảnh được cấu tạo từ một ma trận các điểm ảnh (pixel) cố định; do đó khi phóng to hoặc thay đổi kích thước (re-scaling), các pixel bị phóng đại dẫn đến hiện tượng ảnh bị răng cưa, vỡ hạt (grainy/pixelated). Ngược lại, đồ họa vector (Vector graphic) mô tả hình ảnh bằng các công thức hình học toán học (điểm, đường, đường cong), cho phép co giãn kích thước tùy ý mà vẫn duy trì độ sắc nét hoàn hảo.

\par\vspace{0.3cm}
\hrule
\vspace{0.3cm}
\end{minipage}

\begin{minipage}{\textwidth}
\subsection*{Question 131 \quad \footnotesize{[0131\_q131\_Question\_131.png]}}
\noindent \textbf{Which physical topology uses a hub or switch?}

\begin{enumerate}[label=\Alph*.]
    \item Bus
    \item Ring
    \item \textbf{Star} \hfill \textcolor{blue}{\textbf{[Correct]}}
    \item Bus and ring
\end{enumerate}

\noindent \textbf{Correct Answer:} \textcolor{blue}{\textbf{C}} \\[2pt]
\noindent \textbf{Explanation:} Cấu trúc liên kết mạng hình sao (Star topology) là mô hình mạng trong đó tất cả các trạm hoặc máy trạm đều được kết nối trực tiếp vào một thiết bị trung tâm duy nhất như Hub hoặc Switch (đáp án C). Thiết bị trung tâm này đóng vai trò điều phối, lặp lại tín hiệu hoặc chuyển tiếp các khung dữ liệu giữa các máy tính. Ngược lại, cấu trúc Bus sử dụng một sợi cáp trục đơn duy nhất nối tất cả các máy tính, còn cấu trúc Ring nối các máy tuần tự thành một vòng lặp khép kín mà không có thiết bị tập trung.

\par\vspace{0.3cm}
\hrule
\vspace{0.3cm}
\end{minipage}

\begin{minipage}{\textwidth}
\subsection*{Question 132 \quad \footnotesize{[0132\_q132\_Question\_132.png]}}
\noindent \textbf{\_layers. The TCP/IP model has}

\begin{enumerate}[label=\Alph*.]
    \item \textbf{Five} \hfill \textcolor{blue}{\textbf{[Correct]}}
    \item Six
    \item Seven
    \item Eight
\end{enumerate}

\noindent \textbf{Correct Answer:} \textcolor{blue}{\textbf{A}} \\[2pt]
\noindent \textbf{Explanation:} Theo chuẩn trình bày trong giáo trình Foundations of Computer Science của tác giả Behrouz Forouzan, mô hình bộ giao thức TCP/IP hiện đại được định nghĩa gồm 5 tầng (Five layers, đáp án A): Tầng vật lý (Physical), Tầng liên kết dữ liệu (Data Link), Tầng mạng (Network/Internet), Tầng giao vận (Transport) và Tầng ứng dụng (Application). Việc phân chia 5 tầng này tương ứng trực tiếp với 5 tầng dưới của mô hình OSI (với tầng Application gộp cả Presentation và Session của OSI). Mô hình OSI 7 tầng (đáp án C) là mô hình tham chiếu lý thuyết, không phải là số tầng của bộ giao thức TCP/IP trong giáo trình Forouzan.

\par\vspace{0.3cm}
\hrule
\vspace{0.3cm}
\end{minipage}

\begin{minipage}{\textwidth}
\subsection*{Question 133 \quad \footnotesize{[0133\_q133\_Question\_133.png]}}
\noindent \textbf{The format of an Ipv6 address is referred to as}

\begin{enumerate}[label=\Alph*.]
    \item Dotted-hexadecimal notation
    \item Dotted-octal notation
    \item Dotted-decimal notation
    \item \textbf{8- Hexadecimal digits} \hfill \textcolor{blue}{\textbf{[Correct]}}
\end{enumerate}

\noindent \textbf{Correct Answer:} \textcolor{blue}{\textbf{D}} \\[2pt]
\noindent \textbf{Explanation:} Địa chỉ IPv6 có độ dài 128 bit, được cấu trúc và hiển thị dưới dạng 8 nhóm số thập lục phân (mỗi nhóm gồm 16 bit, tương ứng 4 ký số hex) phân tách nhau bởi dấu hai chấm (colon-hexadecimal notation, ứng với đáp án D: 8- Hexadecimal digits). Ngược lại, ký hiệu dạng chấm thập phân (dotted-decimal notation, đáp án C) là định dạng chuẩn dành riêng cho địa chỉ IPv4 (dài 32 bit gồm 4 octet số thập phân).

\par\vspace{0.3cm}
\hrule
\vspace{0.3cm}
\end{minipage}

\begin{minipage}{\textwidth}
\subsection*{Question 134 \quad \footnotesize{[0134\_q134\_Question\_134.png]}}
\noindent \textbf{explicitly defines connection establishment, data transfer, connection teardown phases to provide a connection-oriented service AND reliable protocol}

\begin{enumerate}[label=\Alph*.]
    \item \textbf{TCP} \hfill \textcolor{blue}{\textbf{[Correct]}}
    \item UDP
    \item IP
    \item HTTP
\end{enumerate}

\noindent \textbf{Correct Answer:} \textcolor{blue}{\textbf{A}} \\[2pt]
\noindent \textbf{Explanation:} Giao thức TCP (Transmission Control Protocol) là giao thức thuộc tầng giao vận (Transport layer) cung cấp dịch vụ hướng kết nối (connection-oriented) và đảm bảo truyền tin cậy (reliable) thông qua quy trình 3 pha chặt chẽ: thiết lập kết nối (bắt tay 3 bước), truyền tải dữ liệu có xác nhận/kiểm soát lỗi, và giải phóng kết nối (bắt tay 4 bước) (đáp án A). Ngược lại, UDP là giao thức phi kết nối (connectionless) và không tin cậy (unreliable), IP là giao thức tầng mạng không đảm bảo chuyển phát tin cậy, còn HTTP là giao thức tầng ứng dụng hoạt động dựa trên nền của TCP.

\par\vspace{0.3cm}
\hrule
\vspace{0.3cm}
\end{minipage}

\begin{minipage}{\textwidth}
\subsection*{Question 135 \quad \footnotesize{[0135\_q135\_Question\_135.png]}}
\noindent \textbf{A memory that does not change its contents without external causes is known as}

\begin{enumerate}[label=\Alph*.]
    \item Dynamic memory
    \item \textbf{Static memory} \hfill \textcolor{blue}{\textbf{[Correct]}}
    \item RAM
    \item EEPROM
    \item None of the others
\end{enumerate}

\noindent \textbf{Correct Answer:} \textcolor{blue}{\textbf{B}} \\[2pt]
\noindent \textbf{Explanation:} Bộ nhớ tĩnh (Static memory / SRAM) sử dụng các mạch lật logic (flip-flops) để lưu giữ dữ liệu, do đó nó duy trì nội dung ổn định mà không bị suy hao hay thay đổi nếu không có tác động từ bên ngoài hoặc khi nguồn điện còn duy trì (đáp án B). Ngược lại, bộ nhớ động (Dynamic memory / DRAM) lưu dữ liệu trên các tụ điện nhỏ li ti, các tụ điện này liên tục bị rò rỉ điện tích nên đòi hỏi mạch điều khiển phải thực hiện chu kỳ nạp/làm tươi (refresh) định kỳ liên tục để không bị mất dữ liệu.

\par\vspace{0.3cm}
\hrule
\vspace{0.3cm}
\end{minipage}

\begin{minipage}{\textwidth}
\subsection*{Question 136 \quad \footnotesize{[0136\_q136\_Question\_136.png]}}
\noindent \textbf{Which of the following is the user programmed semiconductor memory?}

\begin{enumerate}[label=\Alph*.]
    \item SRAM
    \item DRAM
    \item \textbf{EPROM} \hfill \textcolor{blue}{\textbf{[Correct]}}
    \item All of the others
    \item None of the others
\end{enumerate}

\noindent \textbf{Correct Answer:} \textcolor{blue}{\textbf{C}} \\[2pt]
\noindent \textbf{Explanation:} EPROM (Erasable Programmable Read-Only Memory) là loại bộ nhớ bán dẫn bất biến thuộc dòng ROM cho phép người dùng (kỹ sư/người phát triển) tự lập trình ghi dữ liệu/firmware vào chip bằng thiết bị ghi chuyên dụng (user-programmed) và có thể tái sử dụng bằng cách dùng tia cực tím để xóa nội dung cũ (đáp án C). Ngược lại, SRAM và DRAM là bộ nhớ RAM đọc/ghi dữ liệu tạm thời trong quá trình hệ thống vận hành chứ không phải là loại chip nhớ bất biến do người dùng nạp chương trình cố định.

\par\vspace{0.3cm}
\hrule
\vspace{0.3cm}
\end{minipage}

\begin{minipage}{\textwidth}
\subsection*{Question 137 \quad \footnotesize{[0137\_q137\_Question\_137.png]}}
\noindent \textbf{Which of the following holds data and processing instructions temporarily until the CPU needs it?}

\begin{enumerate}[label=\Alph*.]
    \item ROM
    \item Control unit
    \item \textbf{Main memory} \hfill \textcolor{blue}{\textbf{[Correct]}}
    \item Coprocessor chips
    \item None of the others
\end{enumerate}

\noindent \textbf{Correct Answer:} \textcolor{blue}{\textbf{C}} \\[2pt]
\noindent \textbf{Explanation:} Bộ nhớ chính (Main memory, hay RAM) là nơi lưu trữ tạm thời các dữ liệu và chỉ lệnh chương trình đang chờ được CPU xử lý trong quá trình máy tính hoạt động (đáp án C). Nhờ tốc độ truy cập nhanh hơn nhiều so với bộ nhớ ngoài, bộ nhớ chính đóng vai trò cầu nối cung cấp dữ liệu tức thì cho vi xử lý. Ngược lại, ROM lưu trữ vĩnh viễn các chỉ dẫn khởi động (BIOS) và không thể ghi dữ liệu tạm thời, còn Đơn vị điều khiển (Control Unit) chỉ chịu trách nhiệm giải mã và điều phối luồng lệnh chứ không dùng để lưu trữ dữ liệu.

\par\vspace{0.3cm}
\hrule
\vspace{0.3cm}
\end{minipage}

\begin{minipage}{\textwidth}
\subsection*{Question 138 \quad \footnotesize{[0138\_q138\_Question\_138.png]}}
\noindent \textbf{is a step-by-step method for solving a problem or doing a task.}

\begin{enumerate}[label=\Alph*.]
    \item A construct
    \item A recursion
    \item An iteration
    \item \textbf{An algorithm} \hfill \textcolor{blue}{\textbf{[Correct]}}
\end{enumerate}

\noindent \textbf{Correct Answer:} \textcolor{blue}{\textbf{D}} \\[2pt]
\noindent \textbf{Explanation:} Theo định nghĩa chuẩn trong khoa học máy tính và giáo trình Forouzan, thuật toán (Algorithm) là một tập hợp hữu hạn các chỉ dẫn rõ ràng, được thực hiện theo từng bước (step-by-step method) tuần tự để giải quyết một bài toán cụ thể hoặc hoàn thành một tác vụ đã định trước (đáp án D). Trong khi đó, đệ quy (recursion) và lặp (iteration) chỉ là các kỹ thuật lập trình hoặc cấu trúc điều khiển dùng để biểu diễn một phần logic bên trong thuật toán, chứ không phải khái niệm bao quát phương pháp giải quyết bài toán.

\par\vspace{0.3cm}
\hrule
\vspace{0.3cm}
\end{minipage}

\begin{minipage}{\textwidth}
\subsection*{Question 139 \quad \footnotesize{[0139\_q139\_Question\_139.png]}}
\noindent \textbf{hides all the details of an algorithm in an attempt to give the 'big picture' and to show how the algorithm flows from beginning to end.}

\begin{enumerate}[label=\Alph*.]
    \item Pseudocode
    \item \textbf{Unified Modeling Language (UML)} \hfill \textcolor{blue}{\textbf{[Correct]}}
    \item Flowchart
    \item State diagram
\end{enumerate}

\noindent \textbf{Correct Answer:} \textcolor{blue}{\textbf{B}} \\[2pt]
\noindent \textbf{Explanation:} Theo giáo trình Foundations of Computer Science của Behrouz Forouzan, biểu đồ UML (Unified Modeling Language - cụ thể là activity diagram) là công cụ biểu diễn thuật toán ở mức trừu tượng cao, giúp ẩn đi các chi tiết cài đặt vụn vặt nhằm mang lại bức tranh toàn cảnh trực quan ('big picture') về cách luồng giải thuật vận hành từ điểm bắt đầu đến điểm kết thúc (đáp án B). Ngược lại, mã giả (Pseudocode) và lưu đồ (Flowchart) truyền thống thường thể hiện chi tiết từng bước tính toán, rẽ nhánh và gán giá trị cụ thể của thuật toán.

\par\vspace{0.3cm}
\hrule
\vspace{0.3cm}
\end{minipage}

\begin{minipage}{\textwidth}
\subsection*{Question 140 \quad \footnotesize{[0140\_q140\_Question\_140.png]}}
\noindent \textbf{is a process in which an algorithm calls itself.}

\begin{enumerate}[label=\Alph*.]
    \item Insertion
    \item Searching
    \item \textbf{Recursion} \hfill \textcolor{blue}{\textbf{[Correct]}}
    \item Iteration
\end{enumerate}

\noindent \textbf{Correct Answer:} \textcolor{blue}{\textbf{C}} \\[2pt]
\noindent \textbf{Explanation:} Đệ quy (Recursion) là quá trình hoặc kỹ thuật lập trình trong đó một hàm/thuật toán tự gọi lại chính nó (calls itself) để giải quyết các trường hợp nhỏ hơn của cùng một bài toán cho đến khi đạt đến điều kiện dừng (base case) (đáp án C). Ngược lại, lặp (iteration) sử dụng các cấu trúc vòng lặp (như for, while) để lặp đi lặp lại một khối lệnh mà không tự gọi hàm, còn tìm kiếm (searching) và chèn (insertion) là các thao tác xử lý dữ liệu thông thường.

\par\vspace{0.3cm}
\hrule
\vspace{0.3cm}
\end{minipage}

\begin{minipage}{\textwidth}
\subsection*{Question 141 \quad \footnotesize{[0141\_q141\_Question\_141.png]}}
\noindent \textbf{In sort, the smallest item from the unsorted list is swapped with the item at the beginning of the unsorted list.}

\begin{enumerate}[label=\Alph*.]
    \item \textbf{Selection} \hfill \textcolor{blue}{\textbf{[Correct]}}
    \item Bubble
    \item Insertion
    \item Every
\end{enumerate}

\noindent \textbf{Correct Answer:} \textcolor{blue}{\textbf{A}} \\[2pt]
\noindent \textbf{Explanation:} Trong thuật toán sắp xếp chọn (Selection sort), ở mỗi lượt duyệt, thuật toán sẽ tìm kiếm phần tử có giá trị nhỏ nhất trong phần danh sách chưa được sắp xếp, sau đó hoán đổi (swap) phần tử này với phần tử nằm ở vị trí đầu tiên của danh sách chưa sắp xếp đó (đáp án A). Trái lại, sắp xếp nổi bọt (Bubble sort) chỉ so sánh và đổi chỗ từng cặp phần tử liền kề nhau, còn sắp xếp chèn (Insertion sort) lấy từng phần tử chưa sắp và chèn vào đúng vị trí tương đối trong mảng con đã có thứ tự.

\par\vspace{0.3cm}
\hrule
\vspace{0.3cm}
\end{minipage}

\begin{minipage}{\textwidth}
\subsection*{Question 142 \quad \footnotesize{[0142\_q142\_Question\_142.png]}}
\noindent \textbf{The program in a high-level language is called the}

\begin{enumerate}[label=\Alph*.]
    \item Programming language
    \item \textbf{Source program} \hfill \textcolor{blue}{\textbf{[Correct]}}
    \item Translation
    \item Interpretation
\end{enumerate}

\noindent \textbf{Correct Answer:} \textcolor{blue}{\textbf{B}} \\[2pt]
\noindent \textbf{Explanation:} Một chương trình được lập trình viên viết bằng ngôn ngữ lập trình bậc cao (như C++, Java, Python) được gọi là chương trình nguồn (Source program hoặc source code, đáp án B). Chương trình nguồn này cần phải trải qua quá trình dịch (translation/compilation) hoặc thông dịch (interpretation) để chuyển hóa thành chương trình đích (object/machine program) chứa các chỉ lệnh mã máy mà máy tính có thể trực tiếp thi hành. Do đó, 'Source program' là định nghĩa chính xác nhất.

\par\vspace{0.3cm}
\hrule
\vspace{0.3cm}
\end{minipage}

\begin{minipage}{\textwidth}
\subsection*{Question 143 \quad \footnotesize{[0143\_q143\_Question\_143.png]}}
\noindent \textbf{The code that should be adequate for present and anticipated data processing both for machine and human use is}

\begin{enumerate}[label=\Alph*.]
    \item Conciseness
    \item Uniqueness
    \item \textbf{Operability} \hfill \textcolor{blue}{\textbf{[Correct]}}
    \item All of the others
    \item None of the others
\end{enumerate}

\noindent \textbf{Correct Answer:} \textcolor{blue}{\textbf{C}} \\[2pt]
\noindent \textbf{Explanation:} Trong lý thuyết thiết kế và tiêu chuẩn hóa hệ thống mã dữ liệu, Tính khả năng vận hành (Operability) là đặc tính yêu cầu mã phải phù hợp, thuận tiện cho quy trình xử lý dữ liệu ở hiện tại lẫn các nhu cầu mở rộng trong tương lai, thỏa mãn khả năng sử dụng của cả con người và máy tính (đáp án C). Trong khi đó, Tính súc tích (Conciseness) yêu cầu mã càng ngắn gọn càng tốt, còn Tính duy nhất (Uniqueness) đảm bảo mỗi mã định danh chỉ tương ứng với đúng một thực thể duy nhất.

\par\vspace{0.3cm}
\hrule
\vspace{0.3cm}
\end{minipage}

\begin{minipage}{\textwidth}
\subsection*{Question 144 \quad \footnotesize{[0144\_q144\_Question\_144.png]}}
\noindent \textbf{are names for memory locations}

\begin{enumerate}[label=\Alph*.]
    \item \textbf{Variables} \hfill \textcolor{blue}{\textbf{[Correct]}}
    \item Imperative language
    \item Identifiers
    \item Data type
\end{enumerate}

\noindent \textbf{Correct Answer:} \textcolor{blue}{\textbf{A}} \\[2pt]
\noindent \textbf{Explanation:} Theo định nghĩa chuẩn trong giáo trình Forouzan, Biến (Variables) là những tên định danh được đặt cho các ô nhớ/vị trí bộ nhớ (memory locations) cụ thể trong máy tính, cho phép chương trình lưu trữ và thao tác dữ liệu thông qua tên gọi thay vì phải nhớ địa chỉ số phức tạp của bộ nhớ (đáp án A). Khi một biến được khai báo, hệ thống sẽ cấp phát một vùng nhớ tương ứng với kiểu dữ liệu của nó. Định danh (Identifiers) là khái niệm rộng hơn bao gồm cả tên hàm hay tên hằng, còn kiểu dữ liệu (data type) quy định kích thước và các phép toán hợp lệ trên vùng nhớ đó.

\par\vspace{0.3cm}
\hrule
\vspace{0.3cm}
\end{minipage}

\begin{minipage}{\textwidth}
\subsection*{Question 145 \quad \footnotesize{[0145\_q145\_Question\_145.png]}}
\noindent \textbf{gives the user's view of a system: it shows how the user communicates with the system}

\begin{enumerate}[label=\Alph*.]
    \item Class diagram
    \item State diagram
    \item \textbf{Use-case diagram} \hfill \textcolor{blue}{\textbf{[Correct]}}
    \item Sequence diagram
\end{enumerate}

\noindent \textbf{Correct Answer:} \textcolor{blue}{\textbf{C}} \\[2pt]
\noindent \textbf{Explanation:} Trong kỹ nghệ phần mềm và mô hình hóa UML, biểu đồ ca sử dụng (Use-case diagram) cung cấp góc nhìn của người dùng đối với hệ thống (user's view), thể hiện các tác nhân bên ngoài (actors) và cách họ tương tác, giao tiếp với hệ thống để đạt được một mục tiêu chức năng cụ thể (đáp án C). Ngược lại, biểu đồ lớp (Class diagram) mô tả cấu trúc tĩnh nội bộ, biểu đồ trạng thái (State diagram) mô tả các chu kỳ chuyển đổi trạng thái của đối tượng, còn biểu đồ tuần tự (Sequence diagram) biểu diễn chi tiết luồng thông điệp trao đổi theo thời gian giữa các thành phần.

\par\vspace{0.3cm}
\hrule
\vspace{0.3cm}
\end{minipage}

\begin{minipage}{\textwidth}
\subsection*{Question 146 \quad \footnotesize{[0146\_q146\_Question\_146.png]}}
\noindent \textbf{between modules in a software system must be maximized.}

\begin{enumerate}[label=\Alph*.]
    \item Coupling
    \item \textbf{Cohesion} \hfill \textcolor{blue}{\textbf{[Correct]}}
    \item Neither coupling nor cohesion
    \item Both coupling and cohesion
\end{enumerate}

\noindent \textbf{Correct Answer:} \textcolor{blue}{\textbf{B}} \\[2pt]
\noindent \textbf{Explanation:} Trong thiết kế phần mềm theo mô-đun, mục tiêu quan trọng hàng đầu là tối đa hóa độ kết dính (Cohesion) và tối thiểu hóa độ ghép nối (Coupling) giữa các mô-đun (đáp án B). Độ kết dính cao (high cohesion) đảm bảo các phần tử bên trong một mô-đun cùng thực hiện một nhiệm vụ thống nhất, rõ ràng, giúp mã nguồn dễ hiểu và bảo trì. Ngược lại, độ phụ thuộc (coupling) cần phải được giảm thiểu (minimized) để các mô-đun ít chịu ảnh hưởng từ sự thay đổi của nhau, giúp hệ thống linh hoạt hơn.

\par\vspace{0.3cm}
\hrule
\vspace{0.3cm}
\end{minipage}

\begin{minipage}{\textwidth}
\subsection*{Question 147 \quad \footnotesize{[0147\_q147\_Question\_147.png]}}
\noindent \textbf{What is quality referring to the ability to move software from one hardware platform to another}

\begin{enumerate}[label=\Alph*.]
    \item \textbf{Flexibility} \hfill \textcolor{blue}{\textbf{[Correct]}}
    \item Interoperability
    \item Resusability
    \item Changeability
\end{enumerate}

\noindent \textbf{Correct Answer:} \textcolor{blue}{\textbf{A}} \\[2pt]
\noindent \textbf{Explanation:} Trong đánh giá chất lượng phần mềm theo ngân hàng đề thi CSI106, Tính linh hoạt thích ứng (Flexibility, đáp án A) được định nghĩa là thước đo khả năng dễ dàng chuyển giao, cài đặt và vận hành phần mềm từ nền tảng phần cứng này sang nền tảng phần cứng khác. Trong khi đó, Tính tương tác (Interoperability) là khả năng kết nối và trao đổi dữ liệu giữa hai hay nhiều hệ thống phần mềm khác nhau, còn Tính tái sử dụng (Reusability) là khả năng sử dụng lại các mô-đun mã nguồn cho các dự án hoặc ứng dụng mới.

\par\vspace{0.3cm}
\hrule
\vspace{0.3cm}
\end{minipage}

\begin{minipage}{\textwidth}
\subsection*{Question 148 \quad \footnotesize{[0148\_q148\_Question\_148.png]}}
\noindent \textbf{Testing a software system can involve testing.}

\begin{enumerate}[label=\Alph*.]
    \item Black-box
    \item Glass-box
    \item Neither black-box nor glass-box
    \item \textbf{Both black-box and glass-box} \hfill \textcolor{blue}{\textbf{[Correct]}}
\end{enumerate}

\noindent \textbf{Correct Answer:} \textcolor{blue}{\textbf{D}} \\[2pt]
\noindent \textbf{Explanation:} Quy trình kiểm thử một hệ thống phần mềm toàn diện bao gồm cả phương pháp kiểm thử hộp đen (Black-box testing) và kiểm thử hộp trắng/hộp kính (Glass-box hay White-box testing) (đáp án D). Kiểm thử hộp đen tập trung đánh giá hành vi và chức năng bên ngoài của phần mềm dựa trên đầu vào và đầu ra mà không cần biết mã nguồn nội bộ. Ngược lại, kiểm thử hộp kính kiểm tra trực tiếp các đường rẽ nhánh, luồng thực thi và cấu trúc logic bên trong mã nguồn để phát hiện các lỗi thuật toán tiềm ẩn.

\par\vspace{0.3cm}
\hrule
\vspace{0.3cm}
\end{minipage}

\begin{minipage}{\textwidth}
\subsection*{Question 149 \quad \footnotesize{[0149\_q149\_Question\_149.png]}}
\noindent \textbf{A data structure can be}

\begin{enumerate}[label=\Alph*.]
    \item Only an array
    \item Only a record
    \item Only a linked list
    \item \textbf{An array, a record, or a linked list} \hfill \textcolor{blue}{\textbf{[Correct]}}
\end{enumerate}

\noindent \textbf{Correct Answer:} \textcolor{blue}{\textbf{D}} \\[2pt]
\noindent \textbf{Explanation:} Cấu trúc dữ liệu (Data structure) là phương thức có tổ chức dùng để lưu trữ và quản lý tập hợp các phần tử dữ liệu trong bộ nhớ máy tính nhằm giúp việc truy xuất và xử lý đạt hiệu quả tối ưu (đáp án D). Tùy thuộc vào yêu cầu bài toán, cấu trúc dữ liệu có thể được tổ chức dưới nhiều dạng hình thức khác nhau như mảng (array - tập hợp các phần tử cùng kiểu liên tiếp), bản ghi (record - tập hợp các trường có thể khác kiểu) hoặc danh sách liên kết (linked list - các nút trỏ đến nhau). Vì vậy, khẳng định nó chỉ giới hạn ở một cấu trúc đơn lẻ (A, B, C) là hoàn toàn sai.

\par\vspace{0.3cm}
\hrule
\vspace{0.3cm}
\end{minipage}

\begin{minipage}{\textwidth}
\subsection*{Question 150 \quad \footnotesize{[0150\_q150\_Question\_150.png]}}
\noindent \textbf{A is a collection of related elements possibly of different types having a single name}

\begin{enumerate}[label=\Alph*.]
    \item Array
    \item \textbf{Record} \hfill \textcolor{blue}{\textbf{[Correct]}}
    \item Linked list
    \item Tree
\end{enumerate}

\noindent \textbf{Correct Answer:} \textcolor{blue}{\textbf{B}} \\[2pt]
\noindent \textbf{Explanation:} Theo định nghĩa chuẩn trong khoa học máy tính và giáo trình Forouzan, Bản ghi (Record, hoặc struct trong một số ngôn ngữ) là một tập hợp các phần tử dữ liệu có liên quan với nhau, có thể mang các kiểu dữ liệu khác nhau (heterogeneous) và được quản lý chung dưới một tên duy nhất (đáp án B). Mỗi phần tử bên trong bản ghi được gọi là một trường (field). Ngược lại, Mảng (Array) chỉ lưu trữ tập hợp các phần tử bắt buộc phải có cùng một kiểu dữ liệu (homogeneous), còn cây (tree) hay danh sách liên kết (linked list) là các cấu trúc phân cấp hoặc liên kết tuyến tính phức tạp hơn.

\par\vspace{0.3cm}
\hrule
\vspace{0.3cm}
\end{minipage}

\begin{minipage}{\textwidth}
\subsection*{Question 151 \quad \footnotesize{[0151\_q151\_Question\_151.png]}}
\noindent \textbf{Inside the ADT are two different parts of the model that are \_and}

\begin{enumerate}[label=\Alph*.]
    \item \textbf{Data structure and operations} \hfill \textcolor{blue}{\textbf{[Correct]}}
    \item Application program and public operations
    \item Application program and private operations
    \item Data struture and application program
\end{enumerate}

\noindent \textbf{Correct Answer:} \textcolor{blue}{\textbf{A}} \\[2pt]
\noindent \textbf{Explanation:} Bên trong mô hình Kiểu dữ liệu trừu tượng (Abstract Data Type - ADT), hai thành phần cốt lõi được đóng gói kín đáo là cấu trúc dữ liệu (Data structure - dùng để lưu trữ các phần tử) và các thao tác/phép toán (Operations - các thuật toán xử lý dữ liệu đó) (đáp án A). Cơ chế ADT che giấu toàn bộ cấu trúc dữ liệu và chi tiết giải thuật khỏi chương trình ứng dụng bên ngoài (application program), chỉ cho phép truy cập thông qua một giao diện công khai được định nghĩa rõ ràng. Do đó, application program nằm ở bên ngoài ADT chứ không phải là thành phần bên trong.

\par\vspace{0.3cm}
\hrule
\vspace{0.3cm}
\end{minipage}

\begin{minipage}{\textwidth}
\subsection*{Question 152 \quad \footnotesize{[0152\_q152\_Question\_152.png]}}
\noindent \textbf{The of each node includes one of its children and all descendents of that child}

\begin{enumerate}[label=\Alph*.]
    \item Left tree
    \item Binary tree
    \item \textbf{Subtree} \hfill \textcolor{blue}{\textbf{[Correct]}}
    \item Right tree
\end{enumerate}

\noindent \textbf{Correct Answer:} \textcolor{blue}{\textbf{C}} \\[2pt]
\noindent \textbf{Explanation:} Trong cấu trúc dữ liệu cây (Tree), Cây con (Subtree) của một nút bất kỳ được định nghĩa bao gồm một trong những nút con (child node) trực tiếp của nó cùng toàn bộ tất cả các nút con cháu kế tiếp (descendants) thuộc về nút con đó (đáp án C). Cây con tự thân nó tạo thành một cấu trúc cây hoàn chỉnh có gốc chính là nút con được chọn. Các thuật ngữ như 'Left tree' hay 'Right tree' chỉ là các trường hợp nhánh con cụ thể trong cây nhị phân, trong khi 'Subtree' là định nghĩa tổng quát và chính xác nhất.

\par\vspace{0.3cm}
\hrule
\vspace{0.3cm}
\end{minipage}

\begin{minipage}{\textwidth}
\subsection*{Question 153 \quad \footnotesize{[0153\_q153\_Question\_153.png]}}
\noindent \textbf{When a sequential file is updated, the file contains a list of all errors occurring during the update process.}

\begin{enumerate}[label=\Alph*.]
    \item New master
    \item Old master
    \item Transaction
    \item \textbf{Error report} \hfill \textcolor{blue}{\textbf{[Correct]}}
\end{enumerate}

\noindent \textbf{Correct Answer:} \textcolor{blue}{\textbf{D}} \\[2pt]
\noindent \textbf{Explanation:} Trong quy trình cập nhật tệp tuần tự (Sequential file update) chuẩn, hệ thống tiếp nhận tệp chính cũ (Old master file) cùng tệp giao dịch (Transaction file) để xử lý tạo ra tệp chính mới (New master file) và một tệp báo cáo lỗi (Error report file) (đáp án D). Tệp báo cáo lỗi này lưu lại tất cả các giao dịch không hợp lệ phát sinh trong quá trình cập nhật, chẳng hạn như cố gắng thêm một bản ghi đã tồn tại hoặc cố sửa đổi/xóa một bản ghi không hề có trong cơ sở dữ liệu.

\par\vspace{0.3cm}
\hrule
\vspace{0.3cm}
\end{minipage}

\begin{minipage}{\textwidth}
\subsection*{Question 154 \quad \footnotesize{[0154\_q154\_Question\_154.png]}}
\noindent \textbf{When a hashing algorithm produces an address for an insertion key and that address is already occupied, it is called a}

\begin{enumerate}[label=\Alph*.]
    \item \textbf{Collision} \hfill \textcolor{blue}{\textbf{[Correct]}}
    \item Probe
    \item Synonym
    \item Linked list
\end{enumerate}

\noindent \textbf{Correct Answer:} \textcolor{blue}{\textbf{A}} \\[2pt]
\noindent \textbf{Explanation:} Trong cấu trúc dữ liệu bảng băm (hash table), hiện tượng Xung đột (Collision) xảy ra khi hàm băm ánh xạ một khóa mới cần chèn vào một địa chỉ ô nhớ mà tại đó đã có một phần tử khác đang chiếm giữ từ trước (đáp án A). Các khóa khác nhau nhưng được ánh xạ vào cùng một giá trị băm được gọi là các từ đồng nghĩa (synonyms), và hệ thống phải dùng các phương pháp giải quyết xung đột như thăm dò (probing) hoặc nối kết chuỗi (chaining/linked list) để xử lý. Vì vậy, định nghĩa chính xác cho hiện tượng này là Collision.

\par\vspace{0.3cm}
\hrule
\vspace{0.3cm}
\end{minipage}

\begin{minipage}{\textwidth}
\subsection*{Question 155 \quad \footnotesize{[0155\_q155\_Question\_155.png]}}
\noindent \textbf{The intersection of a column and row in a spreadsheet is called a bon or}

\begin{enumerate}[label=\Alph*.]
    \item Key
    \item Field
    \item \textbf{Cell} \hfill \textcolor{blue}{\textbf{[Correct]}}
    \item Menu
    \item None of the others
\end{enumerate}

\noindent \textbf{Correct Answer:} \textcolor{blue}{\textbf{C}} \\[2pt]
\noindent \textbf{Explanation:} Trong các phần mềm bảng tính (spreadsheet như Excel), giao điểm giữa một cột (column) và một hàng (row) được gọi là một Ô (Cell) hoặc hộp dữ liệu (box) (đáp án C). Mỗi ô được định danh duy nhất bằng sự kết hợp giữa chữ cái của cột và số thứ tự của hàng (ví dụ A1, C10), đóng vai trò là đơn vị cơ bản để người dùng nhập văn bản, giá trị số hoặc thiết lập công thức tính toán. Các khái niệm như Field (trường) thuộc về cơ sở dữ liệu quan hệ, còn Key thuộc về bảng băm hoặc khóa định danh.

\par\vspace{0.3cm}
\hrule
\vspace{0.3cm}
\end{minipage}


\newpage
\section*{5. Operating System}
\addcontentsline{toc}{section}{5. Operating System}
\hrule\vspace{0.5cm}

\begin{minipage}{\textwidth}
\subsection*{Question 156 \quad \footnotesize{[0156\_q156\_Question\_156.png]}}
\noindent \textbf{Which of the following attacks is threatening integrity?}

\begin{enumerate}[label=\Alph*.]
    \item \textbf{Masquerading} \hfill \textcolor{blue}{\textbf{[Correct]}}
    \item Traffic Analysis
    \item Denial of service
    \item Encoding
\end{enumerate}

\noindent \textbf{Correct Answer:} \textcolor{blue}{\textbf{A}} \\[2pt]
\noindent \textbf{Explanation:} Trong bảo mật thông tin, Tính toàn vẹn (Integrity) đảm bảo dữ liệu và tài nguyên hệ thống không bị sửa đổi, làm sai lệch hoặc tạo giả bởi các đối tượng không có thẩm quyền. Tấn công giả mạo (Masquerading) xảy ra khi kẻ tấn công mạo danh một người dùng hợp lệ để truy cập và có thể chèn hoặc sửa đổi dữ liệu trái phép, đe dọa trực tiếp đến tính toàn vẹn của hệ thống (đáp án A). Ngược lại, Phân tích lưu lượng (Traffic Analysis) đe dọa tính bảo mật (Confidentiality), còn Tấn công từ chối dịch vụ (Denial of service) đe dọa tính sẵn sàng (Availability).

\par\vspace{0.3cm}
\hrule
\vspace{0.3cm}
\end{minipage}

\begin{minipage}{\textwidth}
\subsection*{Question 157 \quad \footnotesize{[0157\_q157\_Question\_157.png]}}
\noindent \textbf{The ethical principle that dictates that an act is ethical when a majority of people in society agrees with it is referred to as}

\begin{enumerate}[label=\Alph*.]
    \item Moral rules
    \item Utilization
    \item \textbf{Social contract} \hfill \textcolor{blue}{\textbf{[Correct]}}
    \item None of the others
\end{enumerate}

\noindent \textbf{Correct Answer:} \textcolor{blue}{\textbf{C}} \\[2pt]
\noindent \textbf{Explanation:} Thuyết khế ước xã hội (Social contract theory) trong đạo đức học khẳng định rằng một hành vi được xem là hợp đạo đức khi nó nhận được sự đồng thuận của đa số mọi người trong xã hội dựa trên các quy ước chung nhằm duy trì trật tự và hòa bình (đáp án C). Trái lại, Thuyết vị lợi (Utilization / Utilitarianism) đánh giá đạo đức dựa trên kết quả mang lại lợi ích lớn nhất cho cộng đồng, còn Quy tắc đạo đức (Moral rules / Deontology) dựa vào việc tuân thủ các nghĩa vụ phổ quát bất kể số đông. Do đó, nguyên tắc đạo đức dựa trên sự đồng thuận của đa số xã hội là Social contract.

\par\vspace{0.3cm}
\hrule
\vspace{0.3cm}
\end{minipage}

\begin{minipage}{\textwidth}
\subsection*{Question 158 \quad \footnotesize{[0158\_q158\_Question\_158.png]}}
\noindent \textbf{who create algorithms to break existing internet networks so as to solve the loopholes in them}

\begin{enumerate}[label=\Alph*.]
    \item Black Hat Hacker
    \item Grey Hat Hacker
    \item \textbf{White Hat Hacker} \hfill \textcolor{blue}{\textbf{[Correct]}}
    \item Blue Hat Hacker
\end{enumerate}

\noindent \textbf{Correct Answer:} \textcolor{blue}{\textbf{C}} \\[2pt]
\noindent \textbf{Explanation:} Hacker mũ trắng (White Hat Hacker), hay hacker có đạo đức (ethical hacker), là những chuyên gia bảo mật xây dựng các thuật toán và phương pháp thâm nhập nhằm tìm kiếm và vá các lỗ hổng bảo mật trong hệ thống mạng trước khi chúng bị khai thác (đáp án C). Ngược lại, Hacker mũ đen (Black Hat Hacker) tấn công mạng với mục đích phá hoại hoặc trục lợi bất hợp pháp, còn Hacker mũ xám (Grey Hat Hacker) thâm nhập hệ thống mà không có sự cho phép trước nhưng không mang ý đồ phá hoại. Vì vậy, đối tượng tìm kiếm lỗ hổng để khắc phục và nâng cao bảo mật là White Hat Hacker.

\par\vspace{0.3cm}
\hrule
\vspace{0.3cm}
\end{minipage}

\begin{minipage}{\textwidth}
\subsection*{Question 159 \quad \footnotesize{[0159\_q159\_Question\_159.png]}}
\noindent \textbf{In a stream cipher encryption and decryption are done one symbol. We have a plaintext stream a ciphertext stream and a}

\begin{enumerate}[label=\Alph*.]
    \item Key block
    \item Single key
    \item \textbf{Key stream} \hfill \textcolor{blue}{\textbf{[Correct]}}
    \item Auto key
\end{enumerate}

\noindent \textbf{Correct Answer:} \textcolor{blue}{\textbf{C}} \\[2pt]
\noindent \textbf{Explanation:} Trong mật mã luồng (Stream cipher), quá trình mã hóa và giải mã được thực hiện tuần tự trên từng ký hiệu hoặc bit dữ liệu bằng cách kết hợp một luồng bản rõ (plaintext stream) với một luồng khóa (key stream) để tạo ra luồng bản mã (ciphertext stream) (đáp án C). Luồng khóa này thường được sinh ra từ một khóa bí mật thông qua bộ tạo số giả ngẫu nhiên. Ngược lại, việc chia dữ liệu và sử dụng khối khóa (key block) là đặc trưng của mật mã khối (Block cipher). Do đó, thành phần tương ứng trong mô hình là Key stream.

\par\vspace{0.3cm}
\hrule
\vspace{0.3cm}
\end{minipage}

\begin{minipage}{\textwidth}
\subsection*{Question 160 \quad \footnotesize{[0160\_q160\_Question\_160.png]}}
\noindent \textbf{According to the von Neumann model, are stored in memory.}

\begin{enumerate}[label=\Alph*.]
    \item Only data
    \item Only programs
    \item \textbf{Data and programs} \hfill \textcolor{blue}{\textbf{[Correct]}}
    \item Neither data nor programs
\end{enumerate}

\noindent \textbf{Correct Answer:} \textcolor{blue}{\textbf{C}} \\[2pt]
\noindent \textbf{Explanation:} Nguyên lý chương trình lưu trữ (stored-program concept) trong mô hình kiến trúc Von Neumann quy định rằng cả dữ liệu (data) và chương trình/chỉ thị (programs) đều được lưu trữ chung trong cùng một bộ nhớ chính của máy tính (đáp án C). Ý tưởng đột phá này cho phép máy tính có thể thay đổi tác vụ dễ dàng bằng cách nạp chương trình mới vào bộ nhớ mà không cần phải thiết kế hay đấu nối lại phần cứng. Vì vậy, phát biểu bộ nhớ lưu trữ cả dữ liệu và chương trình là hoàn toàn chính xác.

\par\vspace{0.3cm}
\hrule
\vspace{0.3cm}
\end{minipage}

\begin{minipage}{\textwidth}
\subsection*{Question 161 \quad \footnotesize{[0161\_q161\_Question\_161.png]}}
\noindent \textbf{What generation appeared high-level programming language?}

\begin{enumerate}[label=\Alph*.]
    \item First generation
    \item Second generation
    \item \textbf{Third generation} \hfill \textcolor{blue}{\textbf{[Correct]}}
    \item Fourth generation
\end{enumerate}

\noindent \textbf{Correct Answer:} \textcolor{blue}{\textbf{C}} \\[2pt]
\noindent \textbf{Explanation:} Trong lịch sử phát triển máy tính và ngôn ngữ lập trình, các ngôn ngữ lập trình bậc cao (High-level programming languages như FORTRAN, COBOL, ALGOL, BASIC) xuất hiện và phát triển mạnh mẽ ở thế hệ thứ ba (Third generation / 3GL) gắn liền với kỷ nguyên vi mạch tích hợp IC (đáp án C). Thế hệ thứ nhất (1GL) sử dụng ngôn ngữ máy nhị phân thuần túy, thế hệ thứ hai (2GL) sử dụng hợp ngữ (assembly) với các mã gợi nhớ, còn thế hệ thứ tư (4GL) là các ngôn ngữ phi thủ tục như SQL. Do đó, ngôn ngữ bậc cao xuất hiện ở thế hệ thứ ba.

\par\vspace{0.3cm}
\hrule
\vspace{0.3cm}
\end{minipage}

\begin{minipage}{\textwidth}
\subsection*{Question 162 \quad \footnotesize{[0162\_q162\_Question\_162.png]}}
\noindent \textbf{In the fetch phase, the control unit orders the system to copy the next instruction into the in the CPU.}

\begin{enumerate}[label=\Alph*.]
    \item Program counter
    \item \textbf{Instruction register} \hfill \textcolor{blue}{\textbf{[Correct]}}
    \item Memory
    \item Register
\end{enumerate}

\noindent \textbf{Correct Answer:} \textcolor{blue}{\textbf{B}} \\[2pt]
\noindent \textbf{Explanation:} Trong pha lấy lệnh (Fetch phase) của chu kỳ máy, Đơn vị điều khiển (Control Unit - CU) điều khiển việc sao chép lệnh tiếp theo từ bộ nhớ chính vào Thanh ghi chỉ thị (Instruction Register - IR) bên trong CPU để chuẩn bị cho pha giải mã (decode) (đáp án B). Trong khi đó, Bộ đếm chương trình (Program Counter - PC) chỉ chịu trách nhiệm lưu trữ địa chỉ bộ nhớ của lệnh tiếp theo chứ không chứa bản thân nội dung lệnh đó. Do đó, thanh ghi tiếp nhận chỉ thị được nạp vào CPU là Instruction register.

\par\vspace{0.3cm}
\hrule
\vspace{0.3cm}
\end{minipage}

\begin{minipage}{\textwidth}
\subsection*{Question 163 \quad \footnotesize{[0163\_q163\_Question\_163.png]}}
\noindent \textbf{A is a storage device that can undergo multiple writes and erasures.}

\begin{enumerate}[label=\Alph*.]
    \item CD-ROM
    \item CD-R
    \item \textbf{CD-RW} \hfill \textcolor{blue}{\textbf{[Correct]}}
    \item CD-RR
\end{enumerate}

\noindent \textbf{Correct Answer:} \textcolor{blue}{\textbf{C}} \\[2pt]
\noindent \textbf{Explanation:} Đĩa quang CD-RW (Compact Disc-ReWritable) là thiết bị lưu trữ quang học sử dụng công nghệ đổi pha (phase change), cho phép người dùng có thể ghi, xóa và ghi lại dữ liệu nhiều lần trên cùng một đĩa (đáp án C). Ngược lại, CD-ROM (Read-Only Memory) là đĩa chỉ đọc được dập dữ liệu cố định từ nhà sản xuất, còn CD-R (Recordable) là đĩa chỉ cho phép ghi dữ liệu đúng một lần duy nhất (WORM - Write Once, Read Many). Do đó, thiết bị lưu trữ cho phép ghi và xóa nhiều lần là CD-RW.

\par\vspace{0.3cm}
\hrule
\vspace{0.3cm}
\end{minipage}

\begin{minipage}{\textwidth}
\subsection*{Question 164 \quad \footnotesize{[0164\_q164\_Question\_164.png]}}
\noindent \textbf{What is the name of the display feature that highlights areas of the screen which require operator attention?}

\begin{enumerate}[label=\Alph*.]
    \item Pixel
    \item \textbf{Reverse video} \hfill \textcolor{blue}{\textbf{[Correct]}}
    \item Touch screen
    \item Cursor
    \item None of the others
\end{enumerate}

\noindent \textbf{Correct Answer:} \textcolor{blue}{\textbf{B}} \\[2pt]
\noindent \textbf{Explanation:} Chế độ hiển thị nghịch đảo (Reverse video) là kỹ thuật hiển thị hoán đổi màu sắc giữa nền và ký tự (ví dụ: chữ sáng trên nền tối chuyển thành chữ tối trên nền sáng hoặc ngược lại) nhằm làm nổi bật các khu vực quan trọng cần sự chú ý của người thao tác (đáp án B). Con trỏ (Cursor) chỉ vị trí nhập liệu hiện hành, Màn hình cảm ứng (Touch screen) là thiết bị ngoại vi tiếp nhận thao tác chạm, còn Pixel là điểm ảnh cơ bản cấu thành hình ảnh. Vì vậy, tính năng làm nổi bật vùng hiển thị trên màn hình chính là Reverse video.

\par\vspace{0.3cm}
\hrule
\vspace{0.3cm}
\end{minipage}

\begin{minipage}{\textwidth}
\subsection*{Question 165 \quad \footnotesize{[0165\_q165\_Question\_165.png]}}
\noindent \textbf{Binary numbers need more places for counting because:}

\begin{enumerate}[label=\Alph*.]
    \item 0's and 1's can be added in front of them
    \item 0's and 1's have to be properly placed
    \item They are always big numbers
    \item \textbf{Binary base is small} \hfill \textcolor{blue}{\textbf{[Correct]}}
    \item None of the others
\end{enumerate}

\noindent \textbf{Correct Answer:} \textcolor{blue}{\textbf{D}} \\[2pt]
\noindent \textbf{Explanation:} Hệ nhị phân (Binary) có cơ số rất nhỏ (cơ số bằng 2) và chỉ sử dụng hai chữ số là 0 và 1, do đó giá trị trọng số của mỗi vị trí tăng theo lũy thừa của 2 thay vì lũy thừa của 10 như hệ thập phân. Điều này khiến cho hệ nhị phân cần số lượng vị trí/chữ số nhiều hơn đáng kể để biểu diễn cùng một độ lớn giá trị so với các hệ đếm có cơ số lớn hơn (đáp án D). Bản chất của việc cần nhiều vị trí đếm không phải do số luôn lớn hay do vị trí đặt số, mà xuất phát từ việc cơ số nhị phân nhỏ (Binary base is small).

\par\vspace{0.3cm}
\hrule
\vspace{0.3cm}
\end{minipage}

\begin{minipage}{\textwidth}
\subsection*{Question 166 \quad \footnotesize{[0166\_q166\_Question\_166.png]}}
\noindent \textbf{Convert the FRACTION decimal number 0.625 to binary}

\begin{enumerate}[label=\Alph*.]
    \item 1.010
    \item \textbf{0.101} \hfill \textcolor{blue}{\textbf{[Correct]}}
    \item 1.110
    \item 0.110
\end{enumerate}

\noindent \textbf{Correct Answer:} \textcolor{blue}{\textbf{B}} \\[2pt]
\noindent \textbf{Explanation:} Để chuyển đổi phần thập phân 0.625 sang hệ nhị phân, ta nhân liên tiếp phần lẻ với 2 và ghi lại các phần nguyên thu được: 0.625 × 2 = 1.25 (lấy phần nguyên 1, còn 0.25); 0.25 × 2 = 0.5 (lấy phần nguyên 0, còn 0.5); 0.5 × 2 = 1.0 (lấy phần nguyên 1, phần lẻ bằng 0 thì dừng). Ghi các phần nguyên theo thứ tự từ trên xuống dưới ta được chuỗi nhị phân là 0.101 (đáp án B). Ta cũng có thể kiểm tra lại: 1×2\textasciicircum{}(-1) + 0×2\textasciicircum{}(-2) + 1×2\textasciicircum{}(-3) = 0.5 + 0.125 = 0.625.

\par\vspace{0.3cm}
\hrule
\vspace{0.3cm}
\end{minipage}

\begin{minipage}{\textwidth}
\subsection*{Question 167 \quad \footnotesize{[0167\_q167\_Question\_167.png]}}
\noindent \textbf{Assume a new Excess system uses 17 bits to represent the exponent section. What is he bias value in this system?}

\begin{enumerate}[label=\Alph*.]
    \item 17
    \item 16
    \item \textbf{65535} \hfill \textcolor{blue}{\textbf{[Correct]}}
    \item 65536
\end{enumerate}

\noindent \textbf{Correct Answer:} \textcolor{blue}{\textbf{C}} \\[2pt]
\noindent \textbf{Explanation:} Trong hệ thống biểu diễn số thừa số (Excess\_M hay Biased representation) với m bit được sử dụng cho phần số mũ (exponent), giá trị độ lệch (bias) chuẩn theo quy ước giáo trình và IEEE được tính theo công thức Bias = 2\textasciicircum{}(m - 1) - 1. Với m = 17 bit, ta tính được giá trị bias là 2\textasciicircum{}(17 - 1) - 1 = 2\textasciicircum{}16 - 1 = 65536 - 1 = 65535 (đáp án C). Giá trị 65536 là 2\textasciicircum{}(m-1) khi chưa trừ 1, còn 16 hay 17 chỉ là số lượng bit mũ chứ không phải giá trị bias.

\par\vspace{0.3cm}
\hrule
\vspace{0.3cm}
\end{minipage}

\begin{minipage}{\textwidth}
\subsection*{Question 168 \quad \footnotesize{[0168\_q168\_Question\_168.png]}}
\noindent \textbf{A 32-bit code called represents symbols in all languages.}

\begin{enumerate}[label=\Alph*.]
    \item ANSI
    \item \textbf{Unicode} \hfill \textcolor{blue}{\textbf{[Correct]}}
    \item EBCDIC
    \item Extended ASCII
\end{enumerate}

\noindent \textbf{Correct Answer:} \textcolor{blue}{\textbf{B}} \\[2pt]
\noindent \textbf{Explanation:} Unicode là chuẩn mã hóa ký tự toàn cầu sử dụng không gian mã rộng (với định dạng chuẩn UTF-32 sử dụng 32 bit cho mỗi ký tự), được thiết kế để biểu diễn đầy đủ tất cả các ký tự và biểu tượng của mọi ngôn ngữ viết trên thế giới (đáp án B). Ngược lại, bảng mã ASCII chỉ dùng 7 bit biểu diễn 128 ký tự tiếng Anh, Extended ASCII và ANSI dùng 8 bit (256 ký tự), còn EBCDIC của IBM cũng chỉ dùng 8 bit nên không thể chứa hết chữ viết của mọi ngôn ngữ. Do đó, bộ mã 32-bit bao quát mọi ngôn ngữ là Unicode.

\par\vspace{0.3cm}
\hrule
\vspace{0.3cm}
\end{minipage}

\begin{minipage}{\textwidth}
\subsection*{Question 169 \quad \footnotesize{[0169\_q169\_Question\_169.png]}}
\noindent \textbf{In storing audio, how many samples per second is good enough to reproduce an audio signal?}

\begin{enumerate}[label=\Alph*.]
    \item 10 000
    \item 20 000
    \item 30 000
    \item \textbf{40 000} \hfill \textcolor{blue}{\textbf{[Correct]}}
\end{enumerate}

\noindent \textbf{Correct Answer:} \textcolor{blue}{\textbf{D}} \\[2pt]
\noindent \textbf{Explanation:} Theo định lý lấy mẫu Nyquist-Shannon, để tái tạo chính xác và trung thực một tín hiệu âm thanh liên tục, tần số lấy mẫu phải đạt ít nhất gấp đôi tần số cao nhất có trong tín hiệu đó. Do tai người có khả năng cảm nhận âm thanh tối đa ở tần số khoảng 20.000 Hz (20 kHz), tần số lấy mẫu cần đạt tối thiểu 2 × 20.000 = 40.000 mẫu/giây (samples per second) để tái tạo trọn vẹn dải âm thanh nghe được (đáp án D). Các mức 10.000 hay 20.000 mẫu/giây là quá thấp và sẽ gây ra hiện tượng méo tiếng (aliasing).

\par\vspace{0.3cm}
\hrule
\vspace{0.3cm}
\end{minipage}

\begin{minipage}{\textwidth}
\subsection*{Question 170 \quad \footnotesize{[0170\_q170\_Question\_170.png]}}
\noindent \textbf{The scanning rate in image processing is called. If the resolution is sufficiently high, the human eye cannot recognize the discontinuity in reproduced images}

\begin{enumerate}[label=\Alph*.]
    \item Bit map
    \item Pixels
    \item \textbf{Resolution} \hfill \textcolor{blue}{\textbf{[Correct]}}
    \item Color depth
    \item None of others
\end{enumerate}

\noindent \textbf{Correct Answer:} \textcolor{blue}{\textbf{C}} \\[2pt]
\noindent \textbf{Explanation:} Trong xử lý ảnh kỹ thuật số, tốc độ quét hoặc mật độ điểm ảnh lấy mẫu trên một đơn vị chiều dài được gọi là độ phân giải (Resolution) (đáp án C). Khi độ phân giải của hình ảnh đạt mức đủ cao, khoảng cách giữa các điểm ảnh trở nên cực kỳ nhỏ khiến mắt người không thể nhận thấy sự gián đoạn hay hiện tượng vỡ hạt, tạo cảm giác hình ảnh liền mạch và sắc nét. Pixels là tên gọi của từng điểm ảnh riêng biệt, Color depth là số bit biểu diễn màu sắc cho mỗi pixel, còn Bitmap là định dạng lưu trữ ảnh điểm.

\par\vspace{0.3cm}
\hrule
\vspace{0.3cm}
\end{minipage}

\begin{minipage}{\textwidth}
\subsection*{Question 171 \quad \footnotesize{[0171\_q171\_Question\_171.png]}}
\noindent \textbf{Use a simple right shift operation on the bit pattern 1001 1000.}

\begin{enumerate}[label=\Alph*.]
    \item 1001 1001
    \item 1001 1000
    \item 0100 1101
    \item \textbf{0100 1100} \hfill \textcolor{blue}{\textbf{[Correct]}}
\end{enumerate}

\noindent \textbf{Correct Answer:} \textcolor{blue}{\textbf{D}} \\[2pt]
\noindent \textbf{Explanation:} Phép dịch phải đơn giản (simple / logical right shift) dịch chuyển toàn bộ các bit trong chuỗi sang phải một vị trí, trong đó bit ở rìa ngoài cùng bên phải bị đẩy ra ngoài (mất đi) và một bit 0 được điền vào vị trí trống ngoài cùng bên trái. Khi áp dụng lên mẫu bit 1001 1000, bit 0 ngoài cùng bên phải bị loại bỏ và bit 0 được chèn vào đầu, tạo thành mẫu bit kết quả là 0100 1100 (đáp án D). Các phương án khác có bit dấu hoặc bit cuối không chính xác với quy tắc dịch logic.

\par\vspace{0.3cm}
\hrule
\vspace{0.3cm}
\end{minipage}

\begin{minipage}{\textwidth}
\subsection*{Question 172 \quad \footnotesize{[0172\_q172\_Question\_172.png]}}
\noindent \textbf{A hybrid computer uses a to convert digital signals from a computer into analog signals.}

\begin{enumerate}[label=\Alph*.]
    \item \textbf{Modulator} \hfill \textcolor{blue}{\textbf{[Correct]}}
    \item Demodulator
    \item Modem
    \item Decoder
    \item None of the others
\end{enumerate}

\noindent \textbf{Correct Answer:} \textcolor{blue}{\textbf{A}} \\[2pt]
\noindent \textbf{Explanation:} Trong hệ thống máy tính lai (hybrid computer) và truyền thông dữ liệu, Bộ điều chế (Modulator) là thiết bị thực hiện việc chuyển đổi các tín hiệu số (digital signals) từ máy tính thành tín hiệu tương tự (analog signals) để truyền tải hoặc xử lý (đáp án A). Ngược lại, Bộ giải điều chế (Demodulator) thực hiện chức năng ngược lại là biến đổi tín hiệu tương tự thành tín hiệu số, còn Modem là thiết bị tích hợp cả Modulator và Demodulator. Do đó, thiết bị chuyển đổi digital sang analog là Modulator.

\par\vspace{0.3cm}
\hrule
\vspace{0.3cm}
\end{minipage}

\begin{minipage}{\textwidth}
\subsection*{Question 173 \quad \footnotesize{[0173\_q173\_Question\_173.png]}}
\noindent \textbf{In physical layer, a, can have only a limited number of defined values. Although each value can be any number, it is often as simple as 1 and 0. SÖ DÁP ÁN CÃN CHON}

\begin{enumerate}[label=\Alph*.]
    \item Analog signals
    \item Electromagnetic signals
    \item \textbf{Digital signals} \hfill \textcolor{blue}{\textbf{[Correct]}}
    \item Electronic signals
\end{enumerate}

\noindent \textbf{Correct Answer:} \textcolor{blue}{\textbf{C}} \\[2pt]
\noindent \textbf{Explanation:} Ở tầng vật lý (Physical layer), Tín hiệu số (Digital signals) là dạng tín hiệu rời rạc chỉ có thể nhận một số lượng hữu hạn các giá trị xác định, phổ biến nhất là hệ thống nhị phân chỉ gồm hai mức giá trị đơn giản tương ứng với 0 và 1 (đáp án C). Trái lại, Tín hiệu tương tự (Analog signals) là tín hiệu sóng liên tục có thể nhận vô số giá trị điện áp khác nhau trong một khoảng thời gian. Do đó, loại tín hiệu chỉ có số lượng giá trị xác định giới hạn là Digital signals.

\par\vspace{0.3cm}
\hrule
\vspace{0.3cm}
\end{minipage}

\begin{minipage}{\textwidth}
\subsection*{Question 174 \quad \footnotesize{[0174\_q174\_Question\_174.png]}}
\noindent \textbf{In paradigm, the service provider is an application program, it runs continuously, waiting for another application program, then to make a connection through the Internet and ask for service}

\begin{enumerate}[label=\Alph*.]
    \item P2P
    \item Tranditional
    \item \textbf{Client-server} \hfill \textcolor{blue}{\textbf{[Correct]}}
    \item Model
\end{enumerate}

\noindent \textbf{Correct Answer:} \textcolor{blue}{\textbf{C}} \\[2pt]
\noindent \textbf{Explanation:} Trong mô hình Khách - Chủ (Client-server paradigm), phía cung cấp dịch vụ là một chương trình ứng dụng máy chủ (server) hoạt động liên tục và luôn ở trạng thái lắng nghe trên mạng, chờ đợi các chương trình máy khách (clients) chủ động thiết lập kết nối Internet để gửi yêu cầu dịch vụ và nhận kết quả phản hồi (đáp án C). Ngược lại, trong mô hình mạng ngang hàng (P2P), không có máy chủ trung tâm cố định mà các nút mạng đóng vai trò bình đẳng vừa cung cấp vừa sử dụng dịch vụ. Vì vậy, mô hình được mô tả là Client-server.

\par\vspace{0.3cm}
\hrule
\vspace{0.3cm}
\end{minipage}

\begin{minipage}{\textwidth}
\subsection*{Question 175 \quad \footnotesize{[0175\_q175\_Question\_175.png]}}
\noindent \textbf{A program may need files located in different parts of the world. Networking and internetworking have created a new dimension in operating system that is}

\begin{enumerate}[label=\Alph*.]
    \item Personal systems
    \item \textbf{Distributed systems} \hfill \textcolor{blue}{\textbf{[Correct]}}
    \item Parallel systems
    \item Real-time systems
\end{enumerate}

\noindent \textbf{Correct Answer:} \textcolor{blue}{\textbf{B}} \\[2pt]
\noindent \textbf{Explanation:} Sự phát triển của mạng máy tính và liên mạng (internetworking) cho phép các chương trình truy xuất tệp tin và dữ liệu nằm rải rác ở khắp các vị trí địa lý khác nhau trên thế giới, dẫn đến sự ra đời của Hệ điều hành phân tán (Distributed systems) (đáp án B). Hệ thống phân tán kết nối nhiều máy tính độc lập qua mạng nhưng cung cấp giao diện thống nhất như thể toàn bộ tài nguyên đều nằm trên cùng một máy cục bộ. Trong khi đó, Personal systems dành cho máy tính cá nhân đơn lẻ, còn Parallel systems tập trung vào xử lý song song trên nhiều CPU dùng chung bộ nhớ.

\par\vspace{0.3cm}
\hrule
\vspace{0.3cm}
\end{minipage}

\begin{minipage}{\textwidth}
\subsection*{Question 176 \quad \footnotesize{[0176\_q176\_Question\_176.png]}}
\noindent \textbf{can occur if a process has too many resource restrictions.}

\begin{enumerate}[label=\Alph*.]
    \item \textbf{Starvation} \hfill \textcolor{blue}{\textbf{[Correct]}}
    \item Synchronization
    \item Paging
    \item Deadlock
\end{enumerate}

\noindent \textbf{Correct Answer:} \textcolor{blue}{\textbf{A}} \\[2pt]
\noindent \textbf{Explanation:} Hiện tượng 'đói tài nguyên' (Starvation) xảy ra trong quản lý tiến trình khi một tiến trình có thể sẵn sàng chạy nhưng liên tục bị từ chối cấp phát tài nguyên do có quá nhiều điều kiện hạn chế khắt khe hoặc do các tiến trình khác có độ ưu tiên cao hơn liên tục chiếm giữ tài nguyên (đáp án A). Ngược lại, Bế tắc (Deadlock) là tình trạng hai hoặc nhiều tiến trình nắm giữ tài nguyên và cùng chờ đợi tài nguyên của nhau tạo thành vòng lặp không thể tự giải quyết. Phân trang (Paging) là kỹ thuật quản lý bộ nhớ, còn Đồng bộ hóa (Synchronization) là cơ chế phối hợp hoạt động. Do đó, đáp án đúng là Starvation.

\par\vspace{0.3cm}
\hrule
\vspace{0.3cm}
\end{minipage}

\begin{minipage}{\textwidth}
\subsection*{Question 177 \quad \footnotesize{[0177\_q177\_Question\_177.png]}}
\noindent \textbf{The construct uses a set of actions one after another.}

\begin{enumerate}[label=\Alph*.]
    \item \textbf{Sequence} \hfill \textcolor{blue}{\textbf{[Correct]}}
    \item Repetition
    \item Flow
\end{enumerate}

\noindent \textbf{Correct Answer:} \textcolor{blue}{\textbf{A}} \\[2pt]
\noindent \textbf{Explanation:} Trong cấu trúc điều khiển thuật toán, Cấu trúc tuần tự (Sequence construct) là cấu trúc cơ bản thực hiện một tập hợp các hành động hoặc câu lệnh nối tiếp nhau lần lượt theo một trật tự cố định từ trên xuống dưới, hành động trước kết thúc thì hành động sau mới được bắt đầu (đáp án A). Trái lại, Cấu trúc lặp (Repetition) lặp lại một khối lệnh nhiều lần dựa trên điều kiện kiểm tra, còn Cấu trúc rẽ nhánh (Selection) chọn một trong các nhánh hành động dựa trên kết quả logic. Do đó, cấu trúc thực hiện các hành động kế tiếp nhau là Sequence.

\par\vspace{0.3cm}
\hrule
\vspace{0.3cm}
\end{minipage}

\begin{minipage}{\textwidth}
\subsection*{Question 178 \quad \footnotesize{[0178\_q178\_Question\_178.png]}}
\noindent \textbf{is a basic algorithm that arranges data according to its value.}

\begin{enumerate}[label=\Alph*.]
    \item Inquiry
    \item \textbf{Sorting} \hfill \textcolor{blue}{\textbf{[Correct]}}
    \item Searching
    \item Recursion
\end{enumerate}

\noindent \textbf{Correct Answer:} \textcolor{blue}{\textbf{B}} \\[2pt]
\noindent \textbf{Explanation:} Sắp xếp (Sorting) là thuật toán cơ bản trong khoa học máy tính có nhiệm vụ tổ chức và sắp đặt lại các phần tử trong một tập dữ liệu theo một trật tự xác định (tăng dần hoặc giảm dần) dựa trên giá trị khóa của từng phần tử (đáp án B). Việc sắp xếp dữ liệu giúp tối ưu hóa đáng kể hiệu suất của các thao tác tìm kiếm và truy xuất sau đó. Tìm kiếm (Searching) là thao tác xác định vị trí của một phần tử cho trước, còn Đệ quy (Recursion) là kỹ thuật giải thuật gọi lại chính nó. Vì vậy, thuật toán sắp đặt dữ liệu theo giá trị là Sorting.

\par\vspace{0.3cm}
\hrule
\vspace{0.3cm}
\end{minipage}

\begin{minipage}{\textwidth}
\subsection*{Question 179 \quad \footnotesize{[0179\_q179\_Question\_179.png]}}
\noindent \textbf{refers to the process of translating each line of the source program into the corresponding line of the object program and executing the line. SÖ DÁP ÁN CÀN CHON}

\begin{enumerate}[label=\Alph*.]
    \item Object program
    \item Source program
    \item Translation
    \item \textbf{Interpretation} \hfill \textcolor{blue}{\textbf{[Correct]}}
\end{enumerate}

\noindent \textbf{Correct Answer:} \textcolor{blue}{\textbf{D}} \\[2pt]
\noindent \textbf{Explanation:} Thông dịch (Interpretation) là quá trình phần mềm (trình thông dịch - interpreter) tiến hành đọc, dịch từng dòng lệnh đơn lẻ của chương trình nguồn sang ngôn ngữ máy tương ứng và thực thi ngay dòng lệnh đó trước khi chuyển sang dòng tiếp theo (đáp án D). Ngược lại, quá trình biên dịch (Compilation) dịch toàn bộ mã nguồn của chương trình thành một tệp mã đối tượng hoàn chỉnh một lần duy nhất trước khi chương trình được chạy độc lập. Do câu hỏi nhấn mạnh việc dịch và thực thi từng dòng một, đáp án chính xác là Interpretation.

\par\vspace{0.3cm}
\hrule
\vspace{0.3cm}
\end{minipage}

\begin{minipage}{\textwidth}
\subsection*{Question 180 \quad \footnotesize{[0180\_q180\_Question\_180.png]}}
\noindent \textbf{A uses the principle of logical reasoning to answer queries. It is based on formal logic defined by Greek mathematicians and later developed into firstorder predicate calculus SÖ DÁP ÁN CÃN CHON}

\begin{enumerate}[label=\Alph*.]
    \item Procedural paradigm
    \item Functional paradigm
    \item \textbf{Declarative paradigm} \hfill \textcolor{blue}{\textbf{[Correct]}}
    \item Object-oriented paradigm
\end{enumerate}

\noindent \textbf{Correct Answer:} \textcolor{blue}{\textbf{C}} \\[2pt]
\noindent \textbf{Explanation:} Mô hình lập trình khai báo / logic (Declarative paradigm, tiêu biểu là ngôn ngữ Prolog) hoạt động dựa trên các nguyên tắc suy luận logic hình thức và phép tính vị từ cấp một (predicate calculus) do các nhà toán học Hy Lạp cổ đại đặt nền móng và sau đó phát triển, nhằm tự động trả lời các câu hỏi truy vấn dựa trên các sự kiện (facts) và luật (rules) đã khai báo (đáp án C). Ngược lại, mô hình thủ tục (Procedural) và hướng đối tượng (OOP) tập trung vào việc mô tả các bước chỉ dẫn thuật toán chi tiết. Vì vậy, mô hình sử dụng suy luận logic để trả lời truy vấn là Declarative paradigm.

\par\vspace{0.3cm}
\hrule
\vspace{0.3cm}
\end{minipage}

\begin{minipage}{\textwidth}
\subsection*{Question 181 \quad \footnotesize{[0181\_q181\_Question\_181.png]}}
\noindent \textbf{What is statement that stores the value in the variable, which has already been created in the declaration section. We use the symbol ← in our algorithm to define it?}

\begin{enumerate}[label=\Alph*.]
    \item Compound statement
    \item \textbf{Assignment statement} \hfill \textcolor{blue}{\textbf{[Correct]}}
    \item Control statement
    \item None of above
\end{enumerate}

\noindent \textbf{Correct Answer:} \textcolor{blue}{\textbf{B}} \\[2pt]
\noindent \textbf{Explanation:} Câu lệnh gán (Assignment statement) là câu lệnh dùng để tính toán và lưu một giá trị cụ thể vào một biến đã được khai báo trước đó trong bộ nhớ, thường được biểu diễn bằng ký hiệu mũi tên trái `←` trong thuật toán/mã giả hoặc dấu `=` trong các ngôn ngữ lập trình (đáp án B). Câu lệnh ghép (Compound statement) là một khối gồm nhiều câu lệnh gộp lại, còn Câu lệnh điều khiển (Control statement) dùng để thay đổi luồng thực thi như rẽ nhánh hoặc lặp. Vì vậy, câu lệnh lưu giá trị vào biến được ký hiệu bằng `←` là Assignment statement.

\par\vspace{0.3cm}
\hrule
\vspace{0.3cm}
\end{minipage}

\begin{minipage}{\textwidth}
\subsection*{Question 182 \quad \footnotesize{[0182\_q182\_Question\_182.png]}}
\noindent \textbf{What is model a process of software development where requirements divided into multiple standalone modules of the software development cycle?}

\begin{enumerate}[label=\Alph*.]
    \item Waterfall model
    \item \textbf{Incremental model} \hfill \textcolor{blue}{\textbf{[Correct]}}
    \item Development process
    \item Software engineering
\end{enumerate}

\noindent \textbf{Correct Answer:} \textcolor{blue}{\textbf{B}} \\[2pt]
\noindent \textbf{Explanation:} Mô hình phát triển phần mềm tăng dần (Incremental model) là quy trình trong đó toàn bộ các yêu cầu của hệ thống được phân chia thành nhiều module/phân hệ độc lập, mỗi module sẽ lần lượt trải qua các bước của chu kỳ phát triển (phân tích, thiết kế, cài đặt, kiểm thử) và được bàn giao từng phần cho người dùng (đáp án B). Trái lại, Mô hình thác nước (Waterfall model) thực hiện toàn bộ hệ thống tuần tự theo từng giai đoạn và chỉ bàn giao sản phẩm cuối cùng khi toàn bộ dự án hoàn tất. Do đó, mô hình chia yêu cầu thành nhiều module độc lập theo chu kỳ phát triển là Incremental model.

\par\vspace{0.3cm}
\hrule
\vspace{0.3cm}
\end{minipage}

\begin{minipage}{\textwidth}
\subsection*{Question 183 \quad \footnotesize{[0183\_q183\_Question\_183.png]}}
\noindent \textbf{is a measure of how tightly two modules are bound to each other.}

\begin{enumerate}[label=\Alph*.]
    \item Modularity
    \item \textbf{Coupling} \hfill \textcolor{blue}{\textbf{[Correct]}}
    \item Interoperability
    \item Cohesion
\end{enumerate}

\noindent \textbf{Correct Answer:} \textcolor{blue}{\textbf{B}} \\[2pt]
\noindent \textbf{Explanation:} Độ ghép nối (Coupling) trong thiết kế phần mềm là thước đo mức độ gắn kết chặt chẽ hoặc mức độ phụ thuộc lẫn nhau giữa hai module độc lập trong một hệ thống (đáp án B). Một thiết kế phần mềm chất lượng cao luôn hướng tới việc giảm thiểu độ ghép nối (loose coupling) để giúp các module dễ bảo trì và tái sử dụng độc lập. Ngược lại, Độ gắn kết (Cohesion) đo lường mức độ liên quan giữa các thành phần bên trong nội bộ cùng một module, còn Modularity là tính module hóa tổng thể. Do đó, đáp án đúng là Coupling.

\par\vspace{0.3cm}
\hrule
\vspace{0.3cm}
\end{minipage}

\begin{minipage}{\textwidth}
\subsection*{Question 184 \quad \footnotesize{[0184\_q184\_Question\_184.png]}}
\noindent \textbf{refers to the ability to move data and/or a system from one platform to another and to reuse code.}

\begin{enumerate}[label=\Alph*.]
    \item Maintainability
    \item Operability
    \item Correctability
    \item \textbf{Transferability} \hfill \textcolor{blue}{\textbf{[Correct]}}
\end{enumerate}

\noindent \textbf{Correct Answer:} \textcolor{blue}{\textbf{D}} \\[2pt]
\noindent \textbf{Explanation:} Tính khả chuyển / tính chuyển giao (Transferability / Portability) trong công nghệ phần mềm chỉ khả năng di chuyển dữ liệu, ứng dụng hoặc toàn bộ hệ thống từ nền tảng phần cứng/hệ điều hành này sang nền tảng khác một cách thuận lợi cũng như khả năng tái sử dụng lại mã nguồn trên các môi trường khác nhau (đáp án D). Khả năng bảo trì (Maintainability) đề cập đến việc dễ dàng sửa lỗi và nâng cấp, Khả năng vận hành (Operability) liên quan đến sự thuận tiện khi người dùng sử dụng, và Khả năng sửa lỗi (Correctability) nói về việc khắc phục sai sót. Vì vậy, khả năng di chuyển hệ thống giữa các nền tảng là Transferability.

\par\vspace{0.3cm}
\hrule
\vspace{0.3cm}
\end{minipage}

\begin{minipage}{\textwidth}
\subsection*{Question 185 \quad \footnotesize{[0185\_q185\_Question\_185.png]}}
\noindent \textbf{To search for elements, we can use on sorted arrays that can help to decrease time complexity}

\begin{enumerate}[label=\Alph*.]
    \item Sequential search
    \item Insertion sort
    \item Selection sort
    \item \textbf{Binary search} \hfill \textcolor{blue}{\textbf{[Correct]}}
\end{enumerate}

\noindent \textbf{Correct Answer:} \textcolor{blue}{\textbf{D}} \\[2pt]
\noindent \textbf{Explanation:} Thuật toán tìm kiếm nhị phân (Binary search) được áp dụng trên các mảng đã được sắp xếp (sorted arrays), hoạt động bằng cách liên tục so sánh phần tử cần tìm với phần tử ở giữa mảng và loại bỏ một nửa số phần tử còn lại sau mỗi bước, giúp giảm độ phức tạp thời gian xuống mức O(log n) (đáp án D). Ngược lại, Tìm kiếm tuần tự (Sequential search) phải kiểm tra lần lượt từng phần tử với độ phức tạp O(n) và không tận dụng được lợi thế của mảng đã sắp xếp. Insertion sort và Selection sort là các thuật toán sắp xếp chứ không phải tìm kiếm. Vì vậy, đáp án đúng là Binary search.

\par\vspace{0.3cm}
\hrule
\vspace{0.3cm}
\end{minipage}

\begin{minipage}{\textwidth}
\subsection*{Question 186 \quad \footnotesize{[0186\_q186\_Question\_186.png]}}
\noindent \textbf{A stack is a structure. SÖ DÁP ÁN CÃN CHON}

\begin{enumerate}[label=\Alph*.]
    \item FIFO
    \item \textbf{LIFO} \hfill \textcolor{blue}{\textbf{[Correct]}}
    \item DIFO
    \item SIFO
\end{enumerate}

\noindent \textbf{Correct Answer:} \textcolor{blue}{\textbf{B}} \\[2pt]
\noindent \textbf{Explanation:} Ngăn xếp (Stack) là một cấu trúc dữ liệu tuyến tính hoạt động theo nguyên tắc Vào sau - Ra trước (LIFO - Last In, First Out), trong đó phần tử nào được đưa vào ngăn xếp sau cùng (thao tác push) sẽ là phần tử đầu tiên được lấy ra khỏi ngăn xếp (thao tác pop) tại cùng một đầu gọi là đỉnh (top) (đáp án B). Ngược lại, cấu trúc Vào trước - Ra trước (FIFO) là nguyên tắc hoạt động của hàng đợi (Queue). Vì vậy, cấu trúc của ngăn xếp là LIFO.

\par\vspace{0.3cm}
\hrule
\vspace{0.3cm}
\end{minipage}

\begin{minipage}{\textwidth}
\subsection*{Question 187 \quad \footnotesize{[0187\_q187\_Question\_187.png]}}
\noindent \textbf{A binary search tree (BST) is a binary tree with one extra property: the key value of root is the key values of all nodes in each left subtree}

\begin{enumerate}[label=\Alph*.]
    \item \textbf{Greater than} \hfill \textcolor{blue}{\textbf{[Correct]}}
    \item Less than
    \item Equal
    \item Smaller than
\end{enumerate}

\noindent \textbf{Correct Answer:} \textcolor{blue}{\textbf{A}} \\[2pt]
\noindent \textbf{Explanation:} Cây tìm kiếm nhị phân (Binary Search Tree - BST) là một cây nhị phân có tính chất sắp xếp đặc trưng: giá trị khóa tại nút gốc (root) luôn lớn hơn (Greater than) giá trị khóa của tất cả các nút nằm trong cây con bên trái (left subtree), và nhỏ hơn giá trị khóa của tất cả các nút nằm trong cây con bên phải (right subtree) (đáp án A). Tính chất này được áp dụng đệ quy cho mọi nút con trên cây, giúp cho việc tìm kiếm dữ liệu đạt hiệu suất tương tự tìm kiếm nhị phân. Do đó, nút gốc có giá trị lớn hơn các nút trong cây con trái.

\par\vspace{0.3cm}
\hrule
\vspace{0.3cm}
\end{minipage}

\begin{minipage}{\textwidth}
\subsection*{Question 188 \quad \footnotesize{[0188\_q188\_Question\_188.png]}}
\noindent \textbf{A(n) list is also known as a queue.}

\begin{enumerate}[label=\Alph*.]
    \item LIFO
    \item \textbf{FIFO} \hfill \textcolor{blue}{\textbf{[Correct]}}
    \item unordered
    \item ordered
\end{enumerate}

\noindent \textbf{Correct Answer:} \textcolor{blue}{\textbf{B}} \\[2pt]
\noindent \textbf{Explanation:} Hàng đợi (Queue) là một danh sách tuyến tính hoạt động theo nguyên tắc Vào trước - Ra trước (FIFO - First In, First Out), trong đó phần tử được thêm vào trước ở cuối hàng (rear) sẽ luôn là phần tử được lấy ra trước ở đầu hàng (front), mô phỏng cơ chế xếp hàng chờ phục vụ trong thực tế (đáp án B). Trái lại, danh sách Vào sau - Ra trước (LIFO) là tên gọi gắn liền với ngăn xếp (Stack). Do đó, hàng đợi còn được gọi là danh sách FIFO.

\par\vspace{0.3cm}
\hrule
\vspace{0.3cm}
\end{minipage}

\begin{minipage}{\textwidth}
\subsection*{Question 189 \quad \footnotesize{[0189\_q189\_Question\_189.png]}}
\noindent \textbf{If personal information about each employee in a company is stored in a file we can use access to retrieve each record at the end of the month to print the paychecks TELN}

\begin{enumerate}[label=\Alph*.]
    \item Hashed file
    \item \textbf{Sequential file} \hfill \textcolor{blue}{\textbf{[Correct]}}
    \item EOF
    \item Index file
\end{enumerate}

\noindent \textbf{Correct Answer:} \textcolor{blue}{\textbf{B}} \\[2pt]
\noindent \textbf{Explanation:} Khi lưu trữ thông tin nhân viên trong một tệp và cần truy xuất toàn bộ các bản ghi vào cuối tháng để xử lý và in phiếu lương cho tất cả nhân viên, phương pháp tệp tuần tự (Sequential file / Sequential access) là cách tiếp cận tự nhiên và hiệu quả nhất vì hệ thống sẽ duyệt lần lượt từng bản ghi từ đầu đến cuối tệp mà không cần chi phí quản lý bảng chỉ mục hay bảng băm (đáp án B). Tệp băm (Hashed file) và tệp chỉ mục (Index file) phù hợp cho việc tra cứu ngẫu nhiên một bản ghi riêng lẻ cụ thể. Vì vậy, lựa chọn phù hợp nhất cho việc duyệt toàn bộ định kỳ là Sequential file.

\par\vspace{0.3cm}
\hrule
\vspace{0.3cm}
\end{minipage}

\begin{minipage}{\textwidth}
\subsection*{Question 190 \quad \footnotesize{[0190\_q190\_Question\_190.png]}}
\noindent \textbf{If the transaction file key is 20 and the first master file key is 25, then we}

\begin{enumerate}[label=\Alph*.]
    \item \textbf{Add the new record to the new master file} \hfill \textcolor{blue}{\textbf{[Correct]}}
    \item Revise the contents of the old master file
    \item Delete the data
    \item Write the old master file record to the new master file
\end{enumerate}

\noindent \textbf{Correct Answer:} \textcolor{blue}{\textbf{A}} \\[2pt]
\noindent \textbf{Explanation:} Trong thuật toán cập nhật tệp tuần tự (Sequential file update): cả tệp giao dịch (transaction file) và tệp gốc (master file) đều được sắp xếp theo thứ tự khóa tăng dần. Khi khóa của bản ghi giao dịch (20) nhỏ hơn khóa của bản ghi đầu tiên trong tệp gốc (25), điều này chứng minh rằng bản ghi khóa 20 chưa từng có mặt trong tệp gốc cũ và đây là một yêu cầu thêm mới dữ liệu, do đó hệ thống sẽ thêm ngay bản ghi mới này vào tệp gốc mới (Add the new record to the new master file - đáp án A). Nếu khóa bằng nhau thì cập nhật/xóa, còn nếu khóa giao dịch lớn hơn thì sao chép bản ghi tệp gốc cũ sang tệp mới.

\par\vspace{0.3cm}
\hrule
\vspace{0.3cm}
\end{minipage}

\begin{minipage}{\textwidth}
\subsection*{Question 191 \quad \footnotesize{[0191\_q191\_Question\_191.png]}}
\noindent \textbf{Theis the highest level in the file system hierarchy.This does not have a parent directory}

\begin{enumerate}[label=\Alph*.]
    \item \textbf{Root directory} \hfill \textcolor{blue}{\textbf{[Correct]}}
    \item Home directory
    \item Working directory
    \item Parent directory
\end{enumerate}

\noindent \textbf{Correct Answer:} \textcolor{blue}{\textbf{A}} \\[2pt]
\noindent \textbf{Explanation:} Thư mục gốc (Root directory) là thư mục nằm ở cấp cao nhất trong cấu trúc phân cấp cây thư mục của hệ thống tệp và là thư mục duy nhất không có thư mục cha (parent directory) (đáp án A). Mọi thư mục con, thư mục nhánh và tệp tin trong toàn bộ hệ thống đều xuất phát từ thư mục gốc này (được ký hiệu là `/` trong các hệ điều hành Unix/Linux hoặc `\textbackslash\{\}` trên các ổ đĩa Windows). Ngược lại, Working directory là thư mục làm việc hiện thời, Home directory là thư mục cá nhân của người dùng, và Parent directory là thư mục chứa thư mục hiện tại.

\par\vspace{0.3cm}
\hrule
\vspace{0.3cm}
\end{minipage}

\begin{minipage}{\textwidth}
\subsection*{Question 192 \quad \footnotesize{[0192\_q192\_Question\_192.png]}}
\noindent \textbf{is needed to allow any relation in the database to be represented, to allow a language like SQL to use powerful retrieval operations composed of atomic operations, to remove anomalies in insertion, deletion, and updating, and reduce the need for restructuring the database as new data type are added.}

\begin{enumerate}[label=\Alph*.]
    \item Relational database
    \item Database system
    \item Atomic operations
    \item \textbf{Normalization} \hfill \textcolor{blue}{\textbf{[Correct]}}
\end{enumerate}

\noindent \textbf{Correct Answer:} \textcolor{blue}{\textbf{D}} \\[2pt]
\noindent \textbf{Explanation:} Chuẩn hóa dữ liệu (Normalization) là quá trình phân tích và tổ chức lại cấu trúc các bảng trong cơ sở dữ liệu quan hệ theo các dạng chuẩn (1NF, 2NF, 3NF, BCNF) nhằm loại bỏ sự dư thừa dữ liệu và triệt tiêu các bất thường (anomalies) trong quá trình thêm, xóa, sửa dữ liệu (đáp án D). Quá trình chuẩn hóa cũng giúp cấu trúc bảng trở nên linh hoạt hơn khi thêm các kiểu dữ liệu mới và tối ưu hóa hiệu quả thực thi của các thao tác truy vấn trong ngôn ngữ SQL. Do đó, quá trình giúp ngăn chặn các bất thường và cấu trúc lại cơ sở dữ liệu là Normalization.

\par\vspace{0.3cm}
\hrule
\vspace{0.3cm}
\end{minipage}

\begin{minipage}{\textwidth}
\subsection*{Question 193 \quad \footnotesize{[0193\_q193\_Question\_193.png]}}
\noindent \textbf{The of ethics says that an act is ethical if it brings about a useful for society}

\begin{enumerate}[label=\Alph*.]
    \item Moral rules
    \item Social contract
    \item \textbf{Ultilization} \hfill \textcolor{blue}{\textbf{[Correct]}}
    \item Nonmaleficence
\end{enumerate}

\noindent \textbf{Correct Answer:} \textcolor{blue}{\textbf{C}} \\[2pt]
\noindent \textbf{Explanation:} Thuyết vị lợi (Utilitarianism / trong đề ghi là Ultilization) là một học thuyết đạo đức mục đích luận khẳng định rằng một hành vi được coi là đúng đắn và hợp đạo đức nếu nó mang lại kết quả hữu ích, hạnh phúc và lợi ích cao nhất cho đại đa số thành viên trong xã hội (đáp án C). Ngược lại, Thuyết khế ước xã hội (Social contract) dựa trên sự đồng thuận của cộng đồng đối với các quy ước xã hội, còn Quy tắc đạo đức (Moral rules) đánh giá đạo đức dựa trên việc tuân thủ các nguyên tắc bổn phận tuyệt đối. Do đó, nguyên tắc đạo đức xem trọng sự hữu ích cho xã hội là Ultilization.

\par\vspace{0.3cm}
\hrule
\vspace{0.3cm}
\end{minipage}


\newpage
\section*{6. Algorithms}
\addcontentsline{toc}{section}{6. Algorithms}
\hrule\vspace{0.5cm}

\begin{minipage}{\textwidth}
\subsection*{Question 194 \quad \footnotesize{[0194\_q194\_Question\_194.png]}}
\noindent \textbf{these are the "bad guys". They are the types of hackers who break into computer networks with purely negative motives such as monetary gain or reputation.}

\begin{enumerate}[label=\Alph*.]
    \item \textbf{Black Hat Hacker} \hfill \textcolor{blue}{\textbf{[Correct]}}
    \item Grey Hat Hacker
    \item Red Hat Hacker
    \item Blue Hat Hacker
\end{enumerate}

\noindent \textbf{Correct Answer:} \textcolor{blue}{\textbf{A}} \\[2pt]
\noindent \textbf{Explanation:} Hacker mũ đen (Black Hat Hacker) là thuật ngữ chỉ những kẻ xâm nhập trái phép vào hệ thống hoặc mạng máy tính với động cơ xấu như trục lợi tài chính, đánh cắp dữ liệu hoặc phá hoại danh tiếng. Trái lại, White Hat Hacker (mũ trắng) là các chuyên gia bảo mật hợp pháp giúp kiểm tra và vá lỗ hổng hệ thống, còn Grey Hat Hacker (mũ xám) hoạt động ở ranh giới giữa hai nhóm nhưng thường không có ác ý phá hoại.

\par\vspace{0.3cm}
\hrule
\vspace{0.3cm}
\end{minipage}

\begin{minipage}{\textwidth}
\subsection*{Question 195 \quad \footnotesize{[0195\_q195\_Question\_195.png]}}
\noindent \textbf{A does not substitute one symbol for another; instead it changes the location of the symbols}

\begin{enumerate}[label=\Alph*.]
    \item Monoalphabetic cipher
    \item Substitution cipher
    \item Polyalphabetic cipher
    \item \textbf{Transposition cipher} \hfill \textcolor{blue}{\textbf{[Correct]}}
\end{enumerate}

\noindent \textbf{Correct Answer:} \textcolor{blue}{\textbf{D}} \\[2pt]
\noindent \textbf{Explanation:} Mật mã hoán vị (Transposition cipher) là phương pháp mã hóa không thay đổi bản thân các ký tự mà chỉ thay đổi vị trí hoặc thứ tự sắp xếp của chúng trong văn bản. Ngược lại, mật mã thay thế (Substitution cipher, bao gồm cả monoalphabetic và polyalphabetic) giữ nguyên vị trí nhưng thay thế ký tự này bằng một ký tự hoặc biểu tượng khác. Do đó, Transposition cipher là đáp án chính xác.

\par\vspace{0.3cm}
\hrule
\vspace{0.3cm}
\end{minipage}

\begin{minipage}{\textwidth}
\subsection*{Question 196 \quad \footnotesize{[0196\_q196\_Question\_196.png]}}
\noindent \textbf{In cryptography, everyone has access to everyone's public key.}

\begin{enumerate}[label=\Alph*.]
    \item Symmetric-key
    \item \textbf{Asymmetric-key} \hfill \textcolor{blue}{\textbf{[Correct]}}
    \item Secret-key
    \item Private-key
\end{enumerate}

\noindent \textbf{Correct Answer:} \textcolor{blue}{\textbf{B}} \\[2pt]
\noindent \textbf{Explanation:} Trong hệ mật mã khóa bất đối xứng (Asymmetric-key cryptography / Public-key cryptography), mỗi thực thể sở hữu một cặp khóa gồm khóa công khai (public key) và khóa bí mật (private key). Khóa công khai được phân phối rộng rãi để bất kỳ ai cũng có thể truy cập nhằm mã hóa thông điệp gửi cho chủ sở hữu hoặc xác thực chữ ký điện tử. Ngược lại, mật mã khóa đối xứng (Symmetric-key hay Secret-key) sử dụng chung một khóa bí mật duy nhất được giữ kín giữa hai bên.

\par\vspace{0.3cm}
\hrule
\vspace{0.3cm}
\end{minipage}

\begin{minipage}{\textwidth}
\subsection*{Question 197 \quad \footnotesize{[0197\_q197\_Question\_197.png]}}
\noindent \textbf{In the Turing model, the depends on the combination of two factors: the input data and the program.}

\begin{enumerate}[label=\Alph*.]
    \item \textbf{Output data} \hfill \textcolor{blue}{\textbf{[Correct]}}
    \item Instructions
    \item Algorithm
    \item Data processing
\end{enumerate}

\noindent \textbf{Correct Answer:} \textcolor{blue}{\textbf{A}} \\[2pt]
\noindent \textbf{Explanation:} Theo mô hình máy tính Turing (Turing model), dữ liệu đầu ra (Output data) được xác định bởi sự kết hợp đồng thời của hai yếu tố: dữ liệu đầu vào (input data) và chương trình/thuật toán (program). Máy tính là một bộ xử lý đa năng; khi nhận cùng một dữ liệu đầu vào nhưng áp dụng các tập lệnh chương trình khác nhau, kết quả đầu ra sẽ khác nhau. Vì vậy, dữ liệu đầu ra phụ thuộc trực tiếp vào cả dữ liệu vào và chương trình.

\par\vspace{0.3cm}
\hrule
\vspace{0.3cm}
\end{minipage}

\begin{minipage}{\textwidth}
\subsection*{Question 198 \quad \footnotesize{[0198\_q198\_Question\_198.png]}}
\noindent \textbf{Who is considered the 'father' of the minicomputer and one of the founder fathers of the modern computer industry world-wide?}

\begin{enumerate}[label=\Alph*.]
    \item George Tate
    \item \textbf{Kenneth H. Olsen} \hfill \textcolor{blue}{\textbf{[Correct]}}
    \item Seymour Cray
    \item Basic Pascal
    \item None of the others
\end{enumerate}

\noindent \textbf{Correct Answer:} \textcolor{blue}{\textbf{B}} \\[2pt]
\noindent \textbf{Explanation:} Kenneth H. Olsen là nhà đồng sáng lập công ty Digital Equipment Corporation (DEC) và được vinh danh là 'cha đẻ của máy tính mini' (minicomputer) nhờ chế tạo thành công các dòng máy tính nhỏ gọn, chi phí thấp tiêu biểu như PDP-8 vào thập niên 1960. Các dòng máy mini của ông đã góp phần quan trọng mở ra kỷ nguyên máy tính hiện đại và phá vỡ thế độc quyền của máy tính lớn (mainframe). Trong khi đó, Seymour Cray nổi tiếng là cha đẻ của siêu máy tính (supercomputer).

\par\vspace{0.3cm}
\hrule
\vspace{0.3cm}
\end{minipage}

\begin{minipage}{\textwidth}
\subsection*{Question 199 \quad \footnotesize{[0199\_q199\_Question\_199.png]}}
\noindent \textbf{A path by which communication is achieved between a central processor and other devices is called}

\begin{enumerate}[label=\Alph*.]
    \item Audit trail
    \item Network
    \item \textbf{Bus} \hfill \textcolor{blue}{\textbf{[Correct]}}
    \item Channel
    \item None of the others
\end{enumerate}

\noindent \textbf{Correct Answer:} \textcolor{blue}{\textbf{C}} \\[2pt]
\noindent \textbf{Explanation:} Trong kiến trúc máy tính, Bus là một tập hợp các đường dây dẫn vật lý dùng để truyền tải dữ liệu, địa chỉ và tín hiệu điều khiển giữa CPU với bộ nhớ chính và các thiết bị I/O ngoại vi. Bus đóng vai trò là tuyến đường giao tiếp nội bộ trung tâm giúp kết nối các thành phần bên trong hệ thống. Trong khi đó, Network kết nối các hệ thống máy tính độc lập từ xa, còn Channel hay Audit trail là các khái niệm quản lý và kênh truyền ở phạm vi khác.

\par\vspace{0.3cm}
\hrule
\vspace{0.3cm}
\end{minipage}

\begin{minipage}{\textwidth}
\subsection*{Question 200 \quad \footnotesize{[0200\_q200\_Question\_200.png]}}
\noindent \textbf{If the memory address space is 16 MB and the word size is 8 bits, thenbits are needed to access each word.}

\begin{enumerate}[label=\Alph*.]
    \item 8
    \item 16
    \item \textbf{24} \hfill \textcolor{blue}{\textbf{[Correct]}}
    \item 32
\end{enumerate}

\noindent \textbf{Correct Answer:} \textcolor{blue}{\textbf{C}} \\[2pt]
\noindent \textbf{Explanation:} Không gian địa chỉ bộ nhớ là 16 MB = 16 * 2\textasciicircum{}20 bytes = 2\textasciicircum{}24 bytes, và do mỗi từ nhớ (word) có kích thước 8 bits = 1 byte nên bộ nhớ chứa đúng 2\textasciicircum{}24 từ nhớ độc lập. Để định danh và truy xuất duy nhất từng từ nhớ trong số 2\textasciicircum{}24 vị trí này, thanh ghi địa chỉ cần có số bit là log2(2\textasciicircum{}24) = 24 bits. Do đó, cần đúng 24 bit địa chỉ.

\par\vspace{0.3cm}
\hrule
\vspace{0.3cm}
\end{minipage}

\begin{minipage}{\textwidth}
\subsection*{Question 201 \quad \footnotesize{[0201\_q201\_Question\_201.png]}}
\noindent \textbf{Which printer uses a combination of laser-beam \& electro photographic techniques.}

\begin{enumerate}[label=\Alph*.]
    \item \textbf{Laser printers} \hfill \textcolor{blue}{\textbf{[Correct]}}
    \item Dot-Matrix
    \item Line printer
    \item Daisy wheel
    \item None of the others
\end{enumerate}

\noindent \textbf{Correct Answer:} \textcolor{blue}{\textbf{A}} \\[2pt]
\noindent \textbf{Explanation:} Máy in laser (Laser printers) sử dụng công nghệ kết hợp giữa chùm tia laser và kỹ thuật chụp ảnh tĩnh điện (electrophotographic technique) tương tự như máy photocopy để tạo ảnh điện tích trên trống quang điện rồi hút mực bột (toner) và ép nhiệt lên giấy. Ngược lại, các loại máy in như Dot-Matrix (in kim), Line printer (in dòng) hay Daisy wheel (in hoa cúc) đều thuộc nhóm máy in tác động cơ học (impact printers) thế hệ cũ.

\par\vspace{0.3cm}
\hrule
\vspace{0.3cm}
\end{minipage}

\begin{minipage}{\textwidth}
\subsection*{Question 202 \quad \footnotesize{[0202\_q202\_Question\_202.png]}}
\noindent \textbf{RISC (reduced instruction set computer). The strategy behind RISC architecture is to have a small set of instructions that do a minimum number of}

\begin{enumerate}[label=\Alph*.]
    \item Complex instructions
    \item Complex operations
    \item \textbf{Simple operations} \hfill \textcolor{blue}{\textbf{[Correct]}}
    \item Simple instructions
\end{enumerate}

\noindent \textbf{Correct Answer:} \textcolor{blue}{\textbf{C}} \\[2pt]
\noindent \textbf{Explanation:} Kiến trúc RISC (Reduced Instruction Set Computer) được xây dựng dựa trên chiến lược cung cấp một tập lệnh thu gọn, trong đó mỗi lệnh chỉ thực hiện một số lượng tối thiểu các thao tác đơn giản (simple operations) trong đúng một chu kỳ xung nhịp. Triết lý này giúp đơn giản hóa mạch giải mã phần cứng và tối ưu hóa tốc độ thực thi so với kiến trúc CISC vốn sử dụng các lệnh phức tạp (complex instructions). Do đó, 'Simple operations' là đáp án chính xác theo giáo trình.

\par\vspace{0.3cm}
\hrule
\vspace{0.3cm}
\end{minipage}

\begin{minipage}{\textwidth}
\subsection*{Question 203 \quad \footnotesize{[0203\_q203\_Question\_203.png]}}
\noindent \textbf{The base of the octal number system is}

\begin{enumerate}[label=\Alph*.]
    \item 2
    \item \textbf{8} \hfill \textcolor{blue}{\textbf{[Correct]}}
    \item 10
    \item 16
\end{enumerate}

\noindent \textbf{Correct Answer:} \textcolor{blue}{\textbf{B}} \\[2pt]
\noindent \textbf{Explanation:} Hệ cơ số bát phân (Octal number system) có cơ số (base hoặc radix) là 8, sử dụng tập hợp gồm 8 chữ số từ 0 đến 7 để biểu diễn mọi giá trị số. Trong khoa học máy tính, hệ nhị phân (binary) có cơ số là 2, hệ thập phân (decimal) có cơ số 10, và hệ thập lục phân (hexadecimal) có cơ số 16.

\par\vspace{0.3cm}
\hrule
\vspace{0.3cm}
\end{minipage}

\begin{minipage}{\textwidth}
\subsection*{Question 204 \quad \footnotesize{[0204\_q204\_Question\_204.png]}}
\noindent \textbf{The maximum value of a binary integer with K digits. For example, if K = 5, then the maximum value is ?}

\begin{enumerate}[label=\Alph*.]
    \item 30
    \item \textbf{31} \hfill \textcolor{blue}{\textbf{[Correct]}}
    \item 32
    \item 33
\end{enumerate}

\noindent \textbf{Correct Answer:} \textcolor{blue}{\textbf{B}} \\[2pt]
\noindent \textbf{Explanation:} Với một số nguyên nhị phân không dấu có K chữ số (K bits), giá trị lớn nhất có thể biểu diễn được tính theo công thức 2\textasciicircum{}K - 1. Áp dụng với K = 5 bit, giá trị lớn nhất đạt được là 2\textasciicircum{}5 - 1 = 32 - 1 = 31 (tương ứng với mẫu bit 11111 trong hệ nhị phân). Vì vậy, 31 là đáp án chính xác.

\par\vspace{0.3cm}
\hrule
\vspace{0.3cm}
\end{minipage}

\begin{minipage}{\textwidth}
\subsection*{Question 205 \quad \footnotesize{[0205\_q205\_Question\_205.png]}}
\noindent \textbf{What is the number of bit patterns provided by a 7-bit code?}

\begin{enumerate}[label=\Alph*.]
    \item 256
    \item \textbf{128} \hfill \textcolor{blue}{\textbf{[Correct]}}
    \item 64
    \item 512
    \item None of the others
\end{enumerate}

\noindent \textbf{Correct Answer:} \textcolor{blue}{\textbf{B}} \\[2pt]
\noindent \textbf{Explanation:} Một bộ mã nhị phân có độ dài n bit có thể biểu diễn tối đa 2\textasciicircum{}n mẫu bit (bit patterns) riêng biệt. Với một mã gồm 7 bit (tiêu biểu như bộ mã hóa ký tự chuẩn ASCII 7-bit tiêu chuẩn), số lượng mẫu bit phân biệt được tạo ra là 2\textasciicircum{}7 = 128. Do đó, 128 là đáp án đúng.

\par\vspace{0.3cm}
\hrule
\vspace{0.3cm}
\end{minipage}

\begin{minipage}{\textwidth}
\subsection*{Question 206 \quad \footnotesize{[0206\_q206\_Question\_206.png]}}
\noindent \textbf{bits. A byte consists of}

\begin{enumerate}[label=\Alph*.]
    \item 2
    \item 4
    \item \textbf{8} \hfill \textcolor{blue}{\textbf{[Correct]}}
    \item 16
\end{enumerate}

\noindent \textbf{Correct Answer:} \textcolor{blue}{\textbf{C}} \\[2pt]
\noindent \textbf{Explanation:} Trong khoa học máy tính và kiến trúc phần cứng, một byte theo quy ước chuẩn quốc tế luôn được cấu tạo từ 8 bits (còn gọi là octet). Một nhóm gồm 4 bits được gọi là nibble, trong khi 16 bits thường tương ứng với một từ nhớ (word) trên các kiến trúc máy tính 16-bit.

\par\vspace{0.3cm}
\hrule
\vspace{0.3cm}
\end{minipage}

\begin{minipage}{\textwidth}
\subsection*{Question 207 \quad \footnotesize{[0207\_q207\_Question\_207.png]}}
\noindent \textbf{If we cannot record all the values of an audio signal over an interval, we can record some of them. means that we select only a finite number of points on the analog signal, measure their values, and record them}

\begin{enumerate}[label=\Alph*.]
    \item digital
    \item analog
    \item \textbf{sampling} \hfill \textcolor{blue}{\textbf{[Correct]}}
    \item binary
\end{enumerate}

\noindent \textbf{Correct Answer:} \textcolor{blue}{\textbf{C}} \\[2pt]
\noindent \textbf{Explanation:} Lấy mẫu (Sampling) là quá trình đo đạc biên độ của tín hiệu âm thanh analog liên tục tại các khoảng thời gian rời rạc định kỳ để thu được một chuỗi các giá trị hữu hạn đại diện. Đây là bước đầu tiên và cơ bản nhất trong quy trình số hóa tín hiệu âm thanh trước khi thực hiện lượng tử hóa (quantization) và mã hóa (encoding). Do đó, sampling là thuật ngữ chính xác mô tả quá trình này.

\par\vspace{0.3cm}
\hrule
\vspace{0.3cm}
\end{minipage}

\begin{minipage}{\textwidth}
\subsection*{Question 208 \quad \footnotesize{[0208\_q208\_Question\_208.png]}}
\noindent \textbf{Images are stored in computers using\_techniques?}

\begin{enumerate}[label=\Alph*.]
    \item Raster graphics and bit map
    \item \textbf{Vector graphics and bit map} \hfill \textcolor{blue}{\textbf{[Correct]}}
    \item True-color and raster graphics
    \item Index-color and vector graphics
    \item Bit map
\end{enumerate}

\noindent \textbf{Correct Answer:} \textcolor{blue}{\textbf{B}} \\[2pt]
\noindent \textbf{Explanation:} Hình ảnh kỹ thuật số trong máy tính được lưu trữ và biểu diễn chủ yếu bằng hai kỹ thuật chính: đồ họa raster/điểm ảnh (raster graphics hay bitmap) và đồ họa véc-tơ (vector graphics). Đồ họa bitmap chia hình ảnh thành một ma trận các điểm ảnh (pixels) có màu sắc cụ thể, trong khi đồ họa vector mô tả hình ảnh thông qua các công thức toán học gồm điểm, đường và hình học phẳng.

\par\vspace{0.3cm}
\hrule
\vspace{0.3cm}
\end{minipage}

\begin{minipage}{\textwidth}
\subsection*{Question 209 \quad \footnotesize{[0209\_q209\_Question\_209.png]}}
\noindent \textbf{Which network is a packet switching network?}

\begin{enumerate}[label=\Alph*.]
    \item Ring network
    \item LAN
    \item Star network
    \item \textbf{EuroNET} \hfill \textcolor{blue}{\textbf{[Correct]}}
    \item None of the others
\end{enumerate}

\noindent \textbf{Correct Answer:} \textcolor{blue}{\textbf{D}} \\[2pt]
\noindent \textbf{Explanation:} EuroNET là mạng viễn thông chuyển mạch gói (packet switching network) tiên phong tại châu Âu dựa trên chuẩn giao thức X.25, được triển khai từ thập niên 1970 để kết nối các trung tâm dữ liệu khoa học giữa các quốc gia thành viên. Trái lại, Ring network và Star network là các cấu trúc topo mạng (network topologies), còn LAN là mạng cục bộ nói chung chứ không phải tên riêng của một mạng chuyển mạch gói cụ thể.

\par\vspace{0.3cm}
\hrule
\vspace{0.3cm}
\end{minipage}

\begin{minipage}{\textwidth}
\subsection*{Question 210 \quad \footnotesize{[0210\_q210\_Question\_210.png]}}
\noindent \textbf{An IPv4 address is a address that uniquely and universally defines the connection of a host or a router to the Internet}

\begin{enumerate}[label=\Alph*.]
    \item \textbf{32-bit} \hfill \textcolor{blue}{\textbf{[Correct]}}
    \item 64-bit
    \item 128-bit
    \item None of others
\end{enumerate}

\noindent \textbf{Correct Answer:} \textcolor{blue}{\textbf{A}} \\[2pt]
\noindent \textbf{Explanation:} Địa chỉ IPv4 (Internet Protocol version 4) là một địa chỉ số có độ dài chính xác 32 bit, được sử dụng để định danh duy nhất và toàn cầu mỗi giao diện kết nối của máy tính (host) hoặc bộ định tuyến (router) với mạng Internet. Địa chỉ này thường được biểu diễn dưới dạng 4 số thập phân phân tách bởi dấu chấm (dotted-decimal notation, ví dụ: 192.168.1.1). Ngược lại, thế hệ tiếp theo là IPv6 sử dụng độ dài 128 bit.

\par\vspace{0.3cm}
\hrule
\vspace{0.3cm}
\end{minipage}

\begin{minipage}{\textwidth}
\subsection*{Question 211 \quad \footnotesize{[0211\_q211\_Question\_211.png]}}
\noindent \textbf{The acts as a liaison between a client program and a server program, a process-to- process connection.}

\begin{enumerate}[label=\Alph*.]
    \item Physical
    \item Data-link
    \item Network
    \item \textbf{Transport} \hfill \textcolor{blue}{\textbf{[Correct]}}
\end{enumerate}

\noindent \textbf{Correct Answer:} \textcolor{blue}{\textbf{D}} \\[2pt]
\noindent \textbf{Explanation:} Tầng giao vận (Transport layer) trong mô hình OSI/TCP-IP đóng vai trò trung gian chịu trách nhiệm cung cấp kết nối truyền thông tiến trình tới tiến trình (process-to-process connection) giữa tiến trình client và tiến trình server thông qua các số hiệu cổng (port numbers). Tầng mạng (Network layer) chỉ xử lý truyền thông từ máy tới máy (host-to-host), trong khi tầng liên kết dữ liệu (Data-link) chỉ chuyển giao chặng tới chặng (node-to-node).

\par\vspace{0.3cm}
\hrule
\vspace{0.3cm}
\end{minipage}

\begin{minipage}{\textwidth}
\subsection*{Question 212 \quad \footnotesize{[0212\_q212\_Question\_212.png]}}
\noindent \textbf{paradigm, there is no need for a server process to be running all the time and waiting for the client processes to connect.}

\begin{enumerate}[label=\Alph*.]
    \item \textbf{P2P} \hfill \textcolor{blue}{\textbf{[Correct]}}
    \item Tranditional
    \item Client-server
    \item Model
\end{enumerate}

\noindent \textbf{Correct Answer:} \textcolor{blue}{\textbf{A}} \\[2pt]
\noindent \textbf{Explanation:} Trong mô hình mạng ngang hàng (P2P - Peer-to-Peer), các máy tham gia (peers) có vai trò bình đẳng vừa là client vừa là server và chia sẻ dữ liệu trực tiếp với nhau, do đó không cần một tiến trình server chuyên dụng phải hoạt động liên tục 24/7 để chờ đợi kết nối. Ngược lại, mô hình Khách - Chủ (Client-Server truyền thống) luôn bắt buộc server process phải chạy thường trực để lắng nghe và xử lý yêu cầu từ các client.

\par\vspace{0.3cm}
\hrule
\vspace{0.3cm}
\end{minipage}

\begin{minipage}{\textwidth}
\subsection*{Question 213 \quad \footnotesize{[0213\_q213\_Question\_213.png]}}
\noindent \textbf{Every process is}

\begin{enumerate}[label=\Alph*.]
    \item Only a job
    \item Only a program
    \item Only a partition
    \item \textbf{A job and a program} \hfill \textcolor{blue}{\textbf{[Correct]}}
\end{enumerate}

\noindent \textbf{Correct Answer:} \textcolor{blue}{\textbf{D}} \\[2pt]
\noindent \textbf{Explanation:} Trong hệ điều hành, một tiến trình (process) được định nghĩa là một chương trình (program) đang được nạp vào bộ nhớ và thực thi, đồng thời nó cũng đóng vai trò là một tác vụ công việc (job) được hệ thống quản lý, lập lịch và cấp phát tài nguyên CPU. Vì vậy, mọi tiến trình vừa mang bản chất của một tác vụ công việc (job) vừa là một chương trình cụ thể (program).

\par\vspace{0.3cm}
\hrule
\vspace{0.3cm}
\end{minipage}

\begin{minipage}{\textwidth}
\subsection*{Question 214 \quad \footnotesize{[0214\_q214\_Question\_214.png]}}
\noindent \textbf{Why OS is responsible for the efficient use of input/output devices?}

\begin{enumerate}[label=\Alph*.]
    \item Because need to maintains a queue for each input/output device
    \item \textbf{Because these devices are slower in speed compared with the CPU and memory} \hfill \textcolor{blue}{\textbf{[Correct]}}
    \item It may use FIFO for one device
    \item The manager also needs to know when a device has finished serving one process
\end{enumerate}

\noindent \textbf{Correct Answer:} \textcolor{blue}{\textbf{B}} \\[2pt]
\noindent \textbf{Explanation:} Hệ điều hành chịu trách nhiệm quản lý và tối ưu hóa việc sử dụng các thiết bị vào/ra (I/O devices) vì tốc độ hoạt động cơ học và truyền tải của chúng chậm hơn rất nhiều so với tốc độ xử lý của CPU và bộ nhớ chính. Nếu không có cơ chế điều phối hiệu quả từ hệ điều hành (như hàng đợi I/O, bộ đệm buffering và spooling), CPU sẽ phải nhàn rỗi chờ đợi, gây suy giảm nghiêm trọng hiệu năng toàn hệ thống.

\par\vspace{0.3cm}
\hrule
\vspace{0.3cm}
\end{minipage}

\begin{minipage}{\textwidth}
\subsection*{Question 215 \quad \footnotesize{[0215\_q215\_Question\_215.png]}}
\noindent \textbf{is an English-language-like representation of an algorithm. There is no standard for it- some people use a lot of detail}

\begin{enumerate}[label=\Alph*.]
    \item \textbf{Pseudocode} \hfill \textcolor{blue}{\textbf{[Correct]}}
    \item Unified Modeling Language (UML)
    \item Flowchart
    \item State diagram
\end{enumerate}

\noindent \textbf{Correct Answer:} \textcolor{blue}{\textbf{A}} \\[2pt]
\noindent \textbf{Explanation:} Mã giả (Pseudocode) là phương thức mô tả thuật toán bằng ngôn ngữ tự nhiên tựa tiếng Anh kết hợp các từ khóa cấu trúc lập trình, không đòi hỏi tuân thủ cú pháp nghiêm ngặt của bất kỳ ngôn ngữ lập trình cụ thể nào. Lưu đồ (Flowchart) sử dụng các khối hình học trực quan, còn UML là ngôn ngữ mô hình hóa hướng đối tượng có tiêu chuẩn hóa quốc tế.

\par\vspace{0.3cm}
\hrule
\vspace{0.3cm}
\end{minipage}

\begin{minipage}{\textwidth}
\subsection*{Question 216 \quad \footnotesize{[0216\_q216\_Question\_216.png]}}
\noindent \textbf{We use the search for an ordered list.}

\begin{enumerate}[label=\Alph*.]
    \item Sequential
    \item \textbf{Binary} \hfill \textcolor{blue}{\textbf{[Correct]}}
    \item Bubble
    \item Insertion
\end{enumerate}

\noindent \textbf{Correct Answer:} \textcolor{blue}{\textbf{B}} \\[2pt]
\noindent \textbf{Explanation:} Tìm kiếm nhị phân (Binary search) là thuật toán tìm kiếm tối ưu được thiết kế chuyên biệt cho một danh sách đã được sắp xếp thứ tự (ordered list). Thuật toán hoạt động bằng cách liên tục so sánh phần tử cần tìm với phần tử ở vị trí giữa danh sách và loại bỏ một nửa số phần tử còn lại sau mỗi lần so sánh, đạt độ phức tạp thời gian O(log n). Trong khi đó, tìm kiếm tuần tự (Sequential search) áp dụng cho danh sách chưa có thứ tự, còn Bubble và Insertion là các thuật toán sắp xếp.

\par\vspace{0.3cm}
\hrule
\vspace{0.3cm}
\end{minipage}

\begin{minipage}{\textwidth}
\subsection*{Question 217 \quad \footnotesize{[0217\_q217\_Question\_217.png]}}
\noindent \textbf{The items are divided into two lists (sorted and unsorted) \_sort. SÖ DÁP ÁN CÀN CHON}

\begin{enumerate}[label=\Alph*.]
    \item Only in a selection
    \item Only in a bubble
    \item Only in an insertion
    \item \textbf{In selection, bubble, or insertion} \hfill \textcolor{blue}{\textbf{[Correct]}}
\end{enumerate}

\noindent \textbf{Correct Answer:} \textcolor{blue}{\textbf{D}} \\[2pt]
\noindent \textbf{Explanation:} Theo giáo trình Forouzan (Chương 6 - Thuật toán cơ bản), cả ba thuật toán sắp xếp cơ bản gồm Selection sort (sắp xếp chọn), Bubble sort (sắp xếp nổi bọt) và Insertion sort (sắp xếp chèn) đều hoạt động dựa trên nguyên lý logic chia mảng thành hai danh sách con: danh sách đã sắp xếp (sorted) và danh sách chưa sắp xếp (unsorted). Qua từng bước lặp, một phần tử từ phần chưa sắp xếp sẽ được chuyển vào vị trí thích hợp trong phần đã sắp xếp.

\par\vspace{0.3cm}
\hrule
\vspace{0.3cm}
\end{minipage}

\begin{minipage}{\textwidth}
\subsection*{Question 218 \quad \footnotesize{[0218\_q218\_Question\_218.png]}}
\noindent \textbf{First CAD system called AutoCAD for the PCs was announced in the year 1982. Can you tell who announced it?}

\begin{enumerate}[label=\Alph*.]
    \item Intel Corpn.
    \item Rolta India Ltd.
    \item \textbf{Autodesk Inc} \hfill \textcolor{blue}{\textbf{[Correct]}}
    \item AT \& T Corp.
    \item None of the others
\end{enumerate}

\noindent \textbf{Correct Answer:} \textcolor{blue}{\textbf{C}} \\[2pt]
\noindent \textbf{Explanation:} Hệ thống thiết kế hỗ trợ bởi máy tính đầu tiên chạy trên máy tính cá nhân (PC) là AutoCAD được công bố vào năm 1982 bởi công ty Autodesk Inc. (do John Walker cùng các cộng sự sáng lập). Autodesk đã tạo nên bước ngoặt lớn trong ngành kỹ thuật đồ họa khi đưa phần mềm CAD vốn chỉ dành cho các hệ thống máy trạm đắt tiền tiếp cận rộng rãi người dùng máy tính cá nhân.

\par\vspace{0.3cm}
\hrule
\vspace{0.3cm}
\end{minipage}

\begin{minipage}{\textwidth}
\subsection*{Question 219 \quad \footnotesize{[0219\_q219\_Question\_219.png]}}
\noindent \textbf{Programs designed to perform specific tasks is called known as}

\begin{enumerate}[label=\Alph*.]
    \item System software
    \item \textbf{Application software} \hfill \textcolor{blue}{\textbf{[Correct]}}
    \item Utility programs
    \item Operating system
    \item None of the others
\end{enumerate}

\noindent \textbf{Correct Answer:} \textcolor{blue}{\textbf{B}} \\[2pt]
\noindent \textbf{Explanation:} Phần mềm ứng dụng (Application software) là tập hợp các chương trình được thiết kế nhằm phục vụ người dùng thực hiện các công việc cụ thể như soạn thảo văn bản, bảng tính kế toán, đồ họa hay duyệt web. Ngược lại, phần mềm hệ thống (System software, bao gồm hệ điều hành) chịu trách nhiệm quản lý, điều khiển phần cứng và tạo môi trường nền tảng cho phần mềm ứng dụng hoạt động.

\par\vspace{0.3cm}
\hrule
\vspace{0.3cm}
\end{minipage}

\begin{minipage}{\textwidth}
\subsection*{Question 220 \quad \footnotesize{[0220\_q220\_Question\_220.png]}}
\noindent \textbf{The Stack operation returns a/an SÖ DÁP ÁN CÃN CHON}

\begin{enumerate}[label=\Alph*.]
    \item Full stack
    \item \textbf{Empty stack} \hfill \textcolor{blue}{\textbf{[Correct]}}
    \item Check status of stack FO
    \item All of the others
\end{enumerate}

\noindent \textbf{Correct Answer:} \textcolor{blue}{\textbf{B}} \\[2pt]
\noindent \textbf{Explanation:} Thao tác khởi tạo ngăn xếp (`Stack()` hoặc `createStack()`) trong cấu trúc dữ liệu trừu tượng (ADT) có chức năng tạo mới và trả về một ngăn xếp hoàn toàn rỗng (empty stack). Khi vừa được tạo, ngăn xếp chưa chứa bất kỳ phần tử nào (kích thước bằng 0) và sẵn sàng tiếp nhận các thao tác tiếp theo như `push` (thêm vào) hoặc `pop` (lấy ra).

\par\vspace{0.3cm}
\hrule
\vspace{0.3cm}
\end{minipage}

\begin{minipage}{\textwidth}
\subsection*{Question 221 \quad \footnotesize{[0221\_q221\_Question\_221.png]}}
\noindent \textbf{When a sequential file is updated, the file gets the actual update. SÖ DÁP ÁN CÃN CHON}

\begin{enumerate}[label=\Alph*.]
    \item \textbf{New master} \hfill \textcolor{blue}{\textbf{[Correct]}}
    \item Old master
    \item Transaction
    \item Error report
\end{enumerate}

\noindent \textbf{Correct Answer:} \textcolor{blue}{\textbf{A}} \\[2pt]
\noindent \textbf{Explanation:} Trong kỹ thuật cập nhật tệp tuần tự (Sequential file updating), dữ liệu không thể ghi đè trực tiếp lên tệp gốc. Quá trình cập nhật sẽ đọc đồng thời tệp chính cũ (old master file) cùng tệp giao dịch phát sinh (transaction file) để tổng hợp và ghi ra một tệp chính mới (new master file) chứa toàn bộ dữ liệu đã được cập nhật hoàn chỉnh. Tệp chính cũ sau đó được giữ lại làm dữ liệu dự phòng.

\par\vspace{0.3cm}
\hrule
\vspace{0.3cm}
\end{minipage}

\begin{minipage}{\textwidth}
\subsection*{Question 222 \quad \footnotesize{[0222\_q222\_Question\_222.png]}}
\noindent \textbf{The method is used when a collision occurs the prime area addresses are searched for an open or unoccupied record where the new data can be placed.}

\begin{enumerate}[label=\Alph*.]
    \item Linked list resolution
    \item \textbf{Open addressing resolution} \hfill \textcolor{blue}{\textbf{[Correct]}}
    \item Bucket hashing
    \item Digit extraction hashing
\end{enumerate}

\noindent \textbf{Correct Answer:} \textcolor{blue}{\textbf{B}} \\[2pt]
\noindent \textbf{Explanation:} Phương pháp định địa chỉ mở (Open addressing resolution) là kỹ thuật giải quyết đụng độ (collision) trong bảng băm bằng cách tìm kiếm một ô nhớ còn trống (open / unoccupied) ngay trong vùng nhớ chính (prime area) để ghi bản ghi mới vào. Ngược lại, phương pháp xâu chuỗi (linked list resolution / chaining) giải quyết đụng độ bằng cách liên kết các phần tử trùng khóa vào danh sách liên kết bên ngoài.

\par\vspace{0.3cm}
\hrule
\vspace{0.3cm}
\end{minipage}

\begin{minipage}{\textwidth}
\subsection*{Question 223 \quad \footnotesize{[0223\_q223\_Question\_223.png]}}
\noindent \textbf{Which of the following is a binary operator?}

\begin{enumerate}[label=\Alph*.]
    \item Select
    \item Update
    \item \textbf{Difference} \hfill \textcolor{blue}{\textbf{[Correct]}}
    \item All of the above
\end{enumerate}

\noindent \textbf{Correct Answer:} \textcolor{blue}{\textbf{C}} \\[2pt]
\noindent \textbf{Explanation:} Trong đại số quan hệ (Relational Algebra), phép hiệu / phép trừ quan hệ (Difference, ký hiệu R - S) là một toán tử hai ngôi (binary operator) vì nó đòi hỏi đúng hai quan hệ (bảng) tương thích làm đầu vào để xác định các bộ thuộc về quan hệ thứ nhất nhưng không thuộc quan hệ thứ hai. Trong khi đó, phép chọn (Select) và phép chiếu (Project) là các toán tử một ngôi (unary) chỉ thao tác trên một quan hệ đơn lẻ.

\par\vspace{0.3cm}
\hrule
\vspace{0.3cm}
\end{minipage}

\begin{minipage}{\textwidth}
\subsection*{Question 224 \quad \footnotesize{[0224\_q224\_Question\_224.png]}}
\noindent \textbf{Which of following that describe to protect data that information needs to be protected from unauthorized change?}

\begin{enumerate}[label=\Alph*.]
    \item Confidentiality
    \item \textbf{Integrity} \hfill \textcolor{blue}{\textbf{[Correct]}}
    \item Availability
    \item Masquerading
\end{enumerate}

\noindent \textbf{Correct Answer:} \textcolor{blue}{\textbf{B}} \\[2pt]
\noindent \textbf{Explanation:} Tính toàn vẹn (Integrity) là một trong ba trụ cột chính của an toàn thông tin (mô hình CIA), đảm bảo dữ liệu và thông tin được bảo vệ khỏi sự thay đổi, sửa đổi hoặc xóa bỏ trái phép bởi những đối tượng không có thẩm quyền. Trong khi đó, tính bảo mật (Confidentiality) ngăn chặn việc tiết lộ thông tin trái phép, còn tính khả dụng (Availability) đảm bảo dữ liệu luôn sẵn sàng cho người dùng hợp pháp truy cập.

\par\vspace{0.3cm}
\hrule
\vspace{0.3cm}
\end{minipage}

\begin{minipage}{\textwidth}
\subsection*{Question 225 \quad \footnotesize{[0225\_q225\_Question\_225.png]}}
\noindent \textbf{When the ciphertext is received the receiver must transform the ciphertext back into the plaintext-this is}

\begin{enumerate}[label=\Alph*.]
    \item Plain text
    \item Cipher text
    \item Encryption
    \item \textbf{Decryption} \hfill \textcolor{blue}{\textbf{[Correct]}}
\end{enumerate}

\noindent \textbf{Correct Answer:} \textcolor{blue}{\textbf{D}} \\[2pt]
\noindent \textbf{Explanation:} Giải mã (Decryption) là quá trình chuyển đổi bản mã (ciphertext) trở lại thành bản rõ ban đầu (plaintext) bằng cách sử dụng thuật toán giải mã cùng khóa bảo mật tương ứng. Quá trình ngược lại biến đổi bản rõ ban đầu thành bản mã để bảo vệ thông tin trước khi truyền đi được gọi là mã hóa (Encryption).

\par\vspace{0.3cm}
\hrule
\vspace{0.3cm}
\end{minipage}

\begin{minipage}{\textwidth}
\subsection*{Question 226 \quad \footnotesize{[0226\_q226\_Question\_226.png]}}
\noindent \textbf{is an independent program which can copy itself and which travels through the network. It tries to find weaknesses in the system to inflict harm}

\begin{enumerate}[label=\Alph*.]
    \item DoS
    \item Trojan horses
    \item \textbf{Worms} \hfill \textcolor{blue}{\textbf{[Correct]}}
    \item Viruses
\end{enumerate}

\noindent \textbf{Correct Answer:} \textcolor{blue}{\textbf{C}} \\[2pt]
\noindent \textbf{Explanation:} Sâu máy tính (Worms) là chương trình mã độc hoạt động độc lập (stand-alone program), có khả năng tự nhân bản và tự động lan truyền qua mạng máy tính bằng cách khai thác các lỗ hổng hệ thống mà không cần sự can thiệp của con người. Khác với sâu, virus bắt buộc phải đính kèm vào một tệp/chương trình chủ (host) để lây nhiễm, còn Trojan horse ngụy trang dưới dạng phần mềm hợp pháp và không tự sao chép.

\par\vspace{0.3cm}
\hrule
\vspace{0.3cm}
\end{minipage}

\begin{minipage}{\textwidth}
\subsection*{Question 227 \quad \footnotesize{[0227\_q227\_Question\_227.png]}}
\noindent \textbf{In RSA, if user A wants to send an encrypted message to user B, the plaintext is encrypted with the public key of}

\begin{enumerate}[label=\Alph*.]
    \item User A
    \item \textbf{User B} \hfill \textcolor{blue}{\textbf{[Correct]}}
    \item The network
    \item A third party.
\end{enumerate}

\noindent \textbf{Correct Answer:} \textcolor{blue}{\textbf{B}} \\[2pt]
\noindent \textbf{Explanation:} Trong hệ mật mã hóa công khai RSA, khi người dùng A muốn gửi thông điệp mã hóa bảo mật cho người dùng B, thông điệp đó phải được mã hóa bằng chính khóa công khai (public key) của người nhận B. Khi đó, chỉ duy nhất người nhận B sở hữu khóa bí mật (private key) tương ứng mới có thể giải mã và đọc được nội dung thông điệp.

\par\vspace{0.3cm}
\hrule
\vspace{0.3cm}
\end{minipage}

\begin{minipage}{\textwidth}
\subsection*{Question 228 \quad \footnotesize{[0228\_q228\_Question\_228.png]}}
\noindent \textbf{Which generation of computers is covered by the period 1964-77?}

\begin{enumerate}[label=\Alph*.]
    \item First
    \item Second
    \item \textbf{Third} \hfill \textcolor{blue}{\textbf{[Correct]}}
    \item Forth
    \item None of the others
\end{enumerate}

\noindent \textbf{Correct Answer:} \textcolor{blue}{\textbf{C}} \\[2pt]
\noindent \textbf{Explanation:} Thế hệ máy tính thứ ba (Third generation, từ năm 1964 đến khoảng năm 1977) được đặc trưng bởi việc sử dụng mạch tích hợp (Integrated Circuits - IC) thay thế cho các bóng bán dẫn riêng lẻ của thế hệ thứ hai. Sự phát triển này giúp máy tính có tốc độ xử lý nhanh hơn, kích thước nhỏ gọn hơn, tiết kiệm điện năng và độ tin cậy vượt trội so với các thế hệ trước.

\par\vspace{0.3cm}
\hrule
\vspace{0.3cm}
\end{minipage}

\begin{minipage}{\textwidth}
\subsection*{Question 229 \quad \footnotesize{[0229\_q229\_Question\_229.png]}}
\noindent \textbf{A control unit with five wires can define up to operations. PUO}

\begin{enumerate}[label=\Alph*.]
    \item 5
    \item 10
    \item 16
    \item \textbf{32} \hfill \textcolor{blue}{\textbf{[Correct]}}
\end{enumerate}

\noindent \textbf{Correct Answer:} \textcolor{blue}{\textbf{D}} \\[2pt]
\noindent \textbf{Explanation:} Trong khối điều khiển (Control Unit), mỗi đường dây tín hiệu mang một giá trị nhị phân (0 hoặc 1). Với n đường dây điều khiển, mạch giải mã có thể kết hợp để tạo ra tối đa 2\textasciicircum{}n trạng thái tín hiệu khác nhau, tương ứng với 2\textasciicircum{}n thao tác (operations). Vì vậy, khối điều khiển có 5 đường dây có thể định nghĩa tối đa 2\textasciicircum{}5 = 32 thao tác độc lập.

\par\vspace{0.3cm}
\hrule
\vspace{0.3cm}
\end{minipage}

\begin{minipage}{\textwidth}
\subsection*{Question 230 \quad \footnotesize{[0230\_q230\_Question\_230.png]}}
\noindent \textbf{Which 32-bit microprocessor is used in IBM's PS/2 model-80 computer?}

\begin{enumerate}[label=\Alph*.]
    \item 8088
    \item 80286
    \item \textbf{80386} \hfill \textcolor{blue}{\textbf{[Correct]}}
    \item 80486
    \item None of the others
\end{enumerate}

\noindent \textbf{Correct Answer:} \textcolor{blue}{\textbf{C}} \\[2pt]
\noindent \textbf{Explanation:} Dòng máy tính IBM Personal System/2 (PS/2) Model 80 ra mắt năm 1987 là mẫu máy tính cá nhân 32-bit tiên tiến đầu tiên của IBM được trang bị bộ vi xử lý Intel 80386 (386DX) cùng kiến trúc bus Micro Channel 32-bit. Các dòng chip 8088 và 80286 là vi xử lý 16-bit thế hệ trước, còn 80486 ra đời sau đó vào năm 1989.

\par\vspace{0.3cm}
\hrule
\vspace{0.3cm}
\end{minipage}

\begin{minipage}{\textwidth}
\subsection*{Question 231 \quad \footnotesize{[0231\_q231\_Question\_231.png]}}
\noindent \textbf{A billionth of a second is defined as a:}

\begin{enumerate}[label=\Alph*.]
    \item Millisecond
    \item Microsecond
    \item \textbf{Nanosecond} \hfill \textcolor{blue}{\textbf{[Correct]}}
    \item Picoseconds
    \item None of the others
\end{enumerate}

\noindent \textbf{Correct Answer:} \textcolor{blue}{\textbf{C}} \\[2pt]
\noindent \textbf{Explanation:} Một nano-giây (Nanosecond, viết tắt là ns) là đơn vị đo lường thời gian tương ứng với một phần tỷ của một giây (10\textasciicircum{}-9 giây). Trong các tiền tố đo lường thời gian, millisecond là một phần nghìn giây (10\textasciicircum{}-3 s), microsecond là một phần triệu giây (10\textasciicircum{}-6 s), còn picosecond là một phần nghìn tỷ giây (10\textasciicircum{}-12 s).

\par\vspace{0.3cm}
\hrule
\vspace{0.3cm}
\end{minipage}


\newpage
\section*{7. Programming}
\addcontentsline{toc}{section}{7. Programming}
\hrule\vspace{0.5cm}

\begin{minipage}{\textwidth}
\subsection*{Question 232 \quad \footnotesize{[0232\_q232\_Question\_232.png]}}
\noindent \textbf{Which of the following representations is erroneous?}

\begin{enumerate}[label=\Alph*.]
    \item (110)2
    \item (141)8
    \item (EF)16
    \item \textbf{22A} \hfill \textcolor{blue}{\textbf{[Correct]}}
\end{enumerate}

\noindent \textbf{Correct Answer:} \textcolor{blue}{\textbf{D}} \\[2pt]
\noindent \textbf{Explanation:} Trong hệ thống số theo vị trí (positional number system), một biểu diễn số hợp lệ phải tuân thủ tập ký số cho phép của cơ số đó và thường có chỉ số dưới chỉ rõ cơ số nếu không phải là hệ thập phân mặc định. Ký hiệu '22A' (đáp án D) là không hợp lệ vì chứa chữ cái 'A' (vốn chỉ thuộc hệ thập lục phân trở lên) nhưng lại không có chỉ số dưới xác định cơ số (16) hay tiền tố quy ước, khiến nó không thể xác định giá trị trong hệ thập phân thông thường. Trong khi đó, các phương án (110)₂, (141)₈ và (EF)₁₆ đều ghi rõ cơ số hợp lệ và sử dụng các chữ số nằm đúng trong phạm vi của từng hệ đếm.

\par\vspace{0.3cm}
\hrule
\vspace{0.3cm}
\end{minipage}

\begin{minipage}{\textwidth}
\subsection*{Question 233 \quad \footnotesize{[0233\_q233\_Question\_233.png]}}
\noindent \textbf{How many bits that we can find in the decimal number 20 in binary system?}

\begin{enumerate}[label=\Alph*.]
    \item 4
    \item \textbf{5} \hfill \textcolor{blue}{\textbf{[Correct]}}
    \item 6
    \item 7
\end{enumerate}

\noindent \textbf{Correct Answer:} \textcolor{blue}{\textbf{B}} \\[2pt]
\noindent \textbf{Explanation:} Để xác định số bit cần thiết biểu diễn một số nguyên dương trong hệ nhị phân, ta chuyển đổi số đó sang hệ nhị phân hoặc áp dụng công thức ⌊log₂(N)⌋ + 1. Số thập phân 20 được phân tích thành 20 = 16 + 4 = 1×2⁴ + 0×2³ + 1×2² + 0×2¹ + 0×2⁰, tức là 10100₂ trong hệ nhị phân, gồm đúng 5 chữ số nhị phân (bit). Do đó, số bit tìm thấy trong biểu diễn nhị phân của số 20 là 5 (đáp án B). Các phương án 4, 6 hay 7 không phản ánh đúng độ dài chuỗi bit tối thiểu của giá trị này.

\par\vspace{0.3cm}
\hrule
\vspace{0.3cm}
\end{minipage}

\begin{minipage}{\textwidth}
\subsection*{Question 234 \quad \footnotesize{[0234\_q234\_Question\_234.png]}}
\noindent \textbf{In a positional number system with base b, we can always find the number of digits of an integer. So how many digits can we find in the decimal number 20 in octal system?}

\begin{enumerate}[label=\Alph*.]
    \item 1
    \item \textbf{2} \hfill \textcolor{blue}{\textbf{[Correct]}}
    \item 3
    \item 4
    \item None of the others
\end{enumerate}

\noindent \textbf{Correct Answer:} \textcolor{blue}{\textbf{B}} \\[2pt]
\noindent \textbf{Explanation:} Số lượng chữ số k cần dùng để biểu diễn một số nguyên dương N trong hệ đếm cơ số b bất kỳ được xác định theo công thức k = ⌊log\_b(N)⌋ + 1. Khi chuyển đổi số thập phân 20 sang hệ bát phân (octal - cơ số 8), ta thực hiện phép chia liên tiếp cho 8: 20 chia 8 được thương là 2 dư 4, tức là 20₁₀ = 24₈. Biểu diễn này bao gồm đúng 2 chữ số (đáp án B). Vì vậy, câu trả lời chính xác là 2 chữ số.

\par\vspace{0.3cm}
\hrule
\vspace{0.3cm}
\end{minipage}

\begin{minipage}{\textwidth}
\subsection*{Question 235 \quad \footnotesize{[0235\_q235\_Question\_235.png]}}
\noindent \textbf{Assume memory location that can only hold four bits. What is the result, if we have stored the integer 11 in a memory location and then try to add 9 to the unsigned integer?}

\begin{enumerate}[label=\Alph*.]
    \item \textbf{4} \hfill \textcolor{blue}{\textbf{[Correct]}}
    \item 9
    \item 11
    \item 20
\end{enumerate}

\noindent \textbf{Correct Answer:} \textcolor{blue}{\textbf{A}} \\[2pt]
\noindent \textbf{Explanation:} Ô nhớ có kích thước 4 bit chỉ có thể lưu trữ các giá trị số nguyên không dấu trong phạm vi từ 0 đến 2⁴ - 1 = 15. Khi thực hiện phép cộng hai số không dấu 11 + 9 = 20, kết quả vượt quá khả năng lưu trữ tối đa của 4 bit dẫn đến hiện tượng tràn số (overflow). Trong số học 4 bit, giá trị được giữ lại là phần dư của phép chia cho 16: 20 mod 16 = 4 (tương ứng với 1011₂ + 1001₂ = 10100₂, khi cắt bỏ bit nhớ thứ 5 chỉ còn lại 0100₂ = 4, đáp án A). Các giá trị 9, 11 hay 20 đều không phải kết quả sau tràn số của phép tính này.

\par\vspace{0.3cm}
\hrule
\vspace{0.3cm}
\end{minipage}

\begin{minipage}{\textwidth}
\subsection*{Question 236 \quad \footnotesize{[0236\_q236\_Question\_236.png]}}
\noindent \textbf{The number of bits per sample is sometimes referred to as the}

\begin{enumerate}[label=\Alph*.]
    \item Quantization
    \item Sampling rate
    \item Encoding
    \item \textbf{Bit depth} \hfill \textcolor{blue}{\textbf{[Correct]}}
\end{enumerate}

\noindent \textbf{Correct Answer:} \textcolor{blue}{\textbf{D}} \\[2pt]
\noindent \textbf{Explanation:} Trong quá trình số hóa âm thanh (audio digitization), độ sâu bit (Bit depth - đáp án D) là thuật ngữ dùng để chỉ số lượng bit được phân bổ để lưu trữ cho mỗi mẫu âm thanh (bits per sample). Bit depth càng cao thì dải động và chất lượng âm thanh tái tạo càng chi tiết và trung thực (ví dụ 16 bit hoặc 24 bit). Ngược lại, Tốc độ lấy mẫu (Sampling rate) là số lượng mẫu thu được trong một giây (đo bằng Hz), còn Lượng tử hóa (Quantization) là quá trình làm tròn các giá trị biên độ liên tục thành các mức rời rạc.

\par\vspace{0.3cm}
\hrule
\vspace{0.3cm}
\end{minipage}

\begin{minipage}{\textwidth}
\subsection*{Question 237 \quad \footnotesize{[0237\_q237\_Question\_237.png]}}
\noindent \textbf{Use a circular left shift operation on the bit pattern 1001 1000}

\begin{enumerate}[label=\Alph*.]
    \item 0001 1001
    \item \textbf{0011 0001} \hfill \textcolor{blue}{\textbf{[Correct]}}
    \item 1001 1001
    \item 1001 1000
\end{enumerate}

\noindent \textbf{Correct Answer:} \textcolor{blue}{\textbf{B}} \\[2pt]
\noindent \textbf{Explanation:} Phép dịch vòng trái (circular left shift / rotate left) là thao tác dịch chuyển toàn bộ các bit sang trái một vị trí, trong đó bit ngoài cùng bên trái (MSB) bị đẩy ra ngoài sẽ xoay vòng quay trở lại lấp vào vị trí ngoài cùng bên phải (LSB). Với chuỗi bit ban đầu là 1001 1000, bit đầu tiên là 1 sẽ được đưa về cuối cùng, và 7 bit còn lại (001 1000) dịch sang trái tạo thành 0011 0001 (đáp án B). Các phương án khác đưa ra kết quả dịch sai vị trí hoặc nhầm lẫn với phép dịch vòng phải.

\par\vspace{0.3cm}
\hrule
\vspace{0.3cm}
\end{minipage}

\begin{minipage}{\textwidth}
\subsection*{Question 238 \quad \footnotesize{[0238\_q238\_Question\_238.png]}}
\noindent \textbf{In the internet today , Backbones and provider networks are also called}

\begin{enumerate}[label=\Alph*.]
    \item Customer networks
    \item Provider networks
    \item Peering points
    \item \textbf{Internet service providers (ISPs)} \hfill \textcolor{blue}{\textbf{[Correct]}}
\end{enumerate}

\noindent \textbf{Correct Answer:} \textcolor{blue}{\textbf{D}} \\[2pt]
\noindent \textbf{Explanation:} Trong kiến trúc mạng Internet hiện đại, các mạng đường trục (backbones) và mạng của các nhà cung cấp dịch vụ viễn thông chịu trách nhiệm duy trì cơ sở hạ tầng kết nối toàn cầu được gọi chung là Nhà cung cấp dịch vụ Internet (Internet Service Providers - ISPs, đáp án D). Mạng khách hàng (Customer networks) là các mạng ở biên sử dụng dịch vụ của ISP, còn Điểm kết nối ngang hàng (Peering points / IXPs) là các điểm trung chuyển nơi các ISP kết nối và trao đổi lưu lượng với nhau. Do đó, thuật ngữ chuẩn xác cho các mạng này là ISPs.

\par\vspace{0.3cm}
\hrule
\vspace{0.3cm}
\end{minipage}

\begin{minipage}{\textwidth}
\subsection*{Question 239 \quad \footnotesize{[0239\_q239\_Question\_239.png]}}
\noindent \textbf{The role of the physical layer is to transfer the bits received from the data-link layer and convert them to for transmission.}

\begin{enumerate}[label=\Alph*.]
    \item Analog signals
    \item \textbf{Electromagnetic signals} \hfill \textcolor{blue}{\textbf{[Correct]}}
    \item Digital signals
    \item Electronic signals
\end{enumerate}

\noindent \textbf{Correct Answer:} \textcolor{blue}{\textbf{B}} \\[2pt]
\noindent \textbf{Explanation:} Tầng vật lý (Physical layer) trong mô hình mạng chịu trách nhiệm tiếp nhận chuỗi bit dữ liệu từ tầng liên kết dữ liệu và chuyển đổi chúng thành các tín hiệu điện từ (Electromagnetic signals - đáp án B) thích hợp để truyền dẫn qua môi trường truyền thông vật lý (như cáp đồng, sợi quang hoặc không gian vô tuyến). Tín hiệu điện từ là khái niệm tổng quát bao gồm cả tín hiệu điện, tín hiệu quang và sóng vô tuyến. Tín hiệu tương tự (Analog) hay tín hiệu số (Digital) chỉ là các dạng biểu diễn hẹp hơn, không bao quát toàn diện vai trò chuyển đổi truyền dẫn vật lý này.

\par\vspace{0.3cm}
\hrule
\vspace{0.3cm}
\end{minipage}

\begin{minipage}{\textwidth}
\subsection*{Question 240 \quad \footnotesize{[0240\_q240\_Question\_240.png]}}
\noindent \textbf{What is the domain name in the e mail addresskayla@nasa.gov? SÖ DÁP ÁN CÀN CHON}

\begin{enumerate}[label=\Alph*.]
    \item kayla
    \item kayla@nasa.gov
    \item \textbf{nasa.gov} \hfill \textcolor{blue}{\textbf{[Correct]}}
    \item gov
\end{enumerate}

\noindent \textbf{Correct Answer:} \textcolor{blue}{\textbf{C}} \\[2pt]
\noindent \textbf{Explanation:} Một địa chỉ thư điện tử tiêu chuẩn luôn gồm hai thành phần chính ngăn cách bởi ký tự '@': tên hộp thư người dùng cục bộ (local-part) ở phía trước và tên miền của máy chủ thư (domain name) ở phía sau. Trong địa chỉ 'kayla@nasa.gov', 'kayla' là tên người dùng, còn 'nasa.gov' chính là tên miền xác định máy chủ thư của cơ quan NASA (đáp án C). 'gov' chỉ là tên miền cấp cao nhất (TLD - Top-Level Domain) chứ không phải tên miền đầy đủ của địa chỉ thư này.

\par\vspace{0.3cm}
\hrule
\vspace{0.3cm}
\end{minipage}

\begin{minipage}{\textwidth}
\subsection*{Question 241 \quad \footnotesize{[0241\_q241\_Question\_241.png]}}
\noindent \textbf{is a very simple protocol using a minimum of overhead. If a process wants to send a small message and does not care much about reliability.}

\begin{enumerate}[label=\Alph*.]
    \item TCP
    \item \textbf{UDP} \hfill \textcolor{blue}{\textbf{[Correct]}}
    \item IP
    \item HTTP
\end{enumerate}

\noindent \textbf{Correct Answer:} \textcolor{blue}{\textbf{B}} \\[2pt]
\noindent \textbf{Explanation:} UDP (User Datagram Protocol - đáp án B) là giao thức truyền vận phi kết nối (connectionless) cực kỳ đơn giản với chi phí mào đầu (overhead) tối thiểu và không cung cấp cơ chế đảm bảo độ tin cậy như bắt tay, báo nhận hay kiểm soát tắc nghẽn. Giao thức này được thiết kế tối ưu cho các tiến trình muốn gửi thông điệp nhỏ, nhanh chóng và không yêu cầu khắt khe về tính tin cậy tuyệt đối (chấp nhận mất gói tin nhỏ để đạt độ trễ thấp). Ngược lại, TCP là giao thức hướng kết nối có overhead lớn để đảm bảo độ tin cậy, còn IP thuộc tầng mạng và HTTP thuộc tầng ứng dụng.

\par\vspace{0.3cm}
\hrule
\vspace{0.3cm}
\end{minipage}

\begin{minipage}{\textwidth}
\subsection*{Question 242 \quad \footnotesize{[0242\_q242\_Question\_242.png]}}
\noindent \textbf{supervises the activity of each component in a computer system.}

\begin{enumerate}[label=\Alph*.]
    \item \textbf{An operating system} \hfill \textcolor{blue}{\textbf{[Correct]}}
    \item Hardware
    \item A queue
    \item An application program
\end{enumerate}

\noindent \textbf{Correct Answer:} \textcolor{blue}{\textbf{A}} \\[2pt]
\noindent \textbf{Explanation:} Hệ điều hành (An operating system - đáp án A) là phần mềm hệ thống cơ bản nhất đóng vai trò quản lý, điều phối và giám sát toàn bộ hoạt động của từng thành phần phần cứng cũng như phần mềm trong hệ thống máy tính. Hệ điều hành phân bổ tài nguyên CPU, bộ nhớ chính, thiết bị vào/ra một cách hiệu quả và an toàn cho các ứng dụng chạy đồng thời. Bản thân phần cứng chỉ cung cấp năng lực tính toán vật lý, còn chương trình ứng dụng chỉ thực hiện các tác vụ cụ thể của người dùng dưới sự giám sát của hệ điều hành.

\par\vspace{0.3cm}
\hrule
\vspace{0.3cm}
\end{minipage}

\begin{minipage}{\textwidth}
\subsection*{Question 243 \quad \footnotesize{[0243\_q243\_Question\_243.png]}}
\noindent \textbf{An operating system NO needs to manage what resources in a computer system.}

\begin{enumerate}[label=\Alph*.]
    \item Memory manager
    \item Process manager
    \item Device manager
    \item \textbf{Input/output data manager} \hfill \textcolor{blue}{\textbf{[Correct]}}
\end{enumerate}

\noindent \textbf{Correct Answer:} \textcolor{blue}{\textbf{D}} \\[2pt]
\noindent \textbf{Explanation:} Theo cấu trúc hệ điều hành được trình bày trong giáo trình Behrouz Forouzan, một hệ điều hành tiêu chuẩn quản lý 4 tài nguyên cốt lõi thông qua 4 thành phần: Bộ quản lý bộ nhớ (Memory manager), Bộ quản lý tiến trình (Process manager), Bộ quản lý thiết bị (Device manager), và Bộ quản lý tệp (File manager). Hệ thống không có một bộ phận riêng biệt mang tên 'Bộ quản lý dữ liệu vào/ra' (Input/output data manager - đáp án D), vì việc giao tiếp và quản lý I/O đã được tích hợp trong Device Manager và File Manager. Do đó, đây là lựa chọn hệ điều hành không cần quản lý dưới tên gọi này.

\par\vspace{0.3cm}
\hrule
\vspace{0.3cm}
\end{minipage}

\begin{minipage}{\textwidth}
\subsection*{Question 244 \quad \footnotesize{[0244\_q244\_Question\_244.png]}}
\noindent \textbf{is a multi-programming method in which multiple programs are entirely in memory with each program occupying a contiguous space.}

\begin{enumerate}[label=\Alph*.]
    \item \textbf{Partitioning} \hfill \textcolor{blue}{\textbf{[Correct]}}
    \item Paging
    \item Demand paging
    \item Demand segmentation
\end{enumerate}

\noindent \textbf{Correct Answer:} \textcolor{blue}{\textbf{A}} \\[2pt]
\noindent \textbf{Explanation:} Phân vùng (Partitioning - đáp án A) là phương pháp quản lý bộ nhớ đa chương cổ điển, trong đó bộ nhớ chính được chia thành các phân vùng (cố định hoặc thay đổi kích thước) và mỗi chương trình được nạp toàn bộ vào một vùng nhớ liên tục (contiguous space). Ngược lại, Phân trang (Paging) và Phân đoạn (Segmentation) cho phép chia chương trình thành các phần nhỏ (trang hoặc đoạn) có thể được nạp vào các khối bộ nhớ vật lý nằm rải rác không liên tục. Phân trang theo yêu cầu (Demand paging) thậm chí chỉ nạp các trang khi thực sự được CPU truy xuất đến.

\par\vspace{0.3cm}
\hrule
\vspace{0.3cm}
\end{minipage}

\begin{minipage}{\textwidth}
\subsection*{Question 245 \quad \footnotesize{[0245\_q245\_Question\_245.png]}}
\noindent \textbf{is a basic algorithm that adds a list of numbers.}

\begin{enumerate}[label=\Alph*.]
    \item \textbf{Summation} \hfill \textcolor{blue}{\textbf{[Correct]}}
    \item Product
    \item Smallest
    \item Largest
\end{enumerate}

\noindent \textbf{Correct Answer:} \textcolor{blue}{\textbf{A}} \\[2pt]
\noindent \textbf{Explanation:} Thuật toán tính tổng (Summation - đáp án A) là một trong những thuật toán cơ bản xử lý danh sách, có chức năng duyệt qua từng phần tử trong một danh sách các số và cộng dồn giá trị của chúng vào một biến tổng tích lũy. Trong khi đó, Thuật toán tính tích (Product) thực hiện việc nhân dồn các số lại với nhau, còn Thuật toán tìm giá trị nhỏ nhất (Smallest) và lớn nhất (Largest) dùng để tìm phần tử có giá trị cực trị. Do đó, thuật toán thực hiện việc cộng một danh sách số chính là Summation.

\par\vspace{0.3cm}
\hrule
\vspace{0.3cm}
\end{minipage}

\begin{minipage}{\textwidth}
\subsection*{Question 246 \quad \footnotesize{[0246\_q246\_Question\_246.png]}}
\noindent \textbf{A program that reads each of the instructions in mnemonic form and translates it into the machine-language equivalent}

\begin{enumerate}[label=\Alph*.]
    \item Machine language
    \item \textbf{Assembler} \hfill \textcolor{blue}{\textbf{[Correct]}}
    \item Interpreter
    \item C program
    \item None of the others
\end{enumerate}

\noindent \textbf{Correct Answer:} \textcolor{blue}{\textbf{B}} \\[2pt]
\noindent \textbf{Explanation:} Trình hợp dịch (Assembler - đáp án B) là chương trình dịch hệ thống có nhiệm vụ đọc các câu lệnh viết bằng ký hiệu gợi nhớ (mnemonic form như ADD, MOV, SUB) của hợp ngữ và chuyển đổi chúng thành các chỉ thị mã máy nhị phân tương đương để bộ vi xử lý có thể thực thi trực tiếp. Ngược lại, Trình thông dịch (Interpreter) dịch và chạy trực tiếp từng dòng lệnh của ngôn ngữ bậc cao, còn Chương trình C là mã nguồn bậc cao cần trình biên dịch (compiler) để tạo mã máy.

\par\vspace{0.3cm}
\hrule
\vspace{0.3cm}
\end{minipage}

\begin{minipage}{\textwidth}
\subsection*{Question 247 \quad \footnotesize{[0247\_q247\_Question\_247.webp]}}
\noindent \textbf{verflow4 (Choose 1 answer) [false] [true] A/an  is a program in a procedural language that is a set of statements. The statements are normally executed one after another. However, sometimes it is nesserary to change this sequential order. Statement 2 Statement 1 This figure illustrates for this statement}

\begin{enumerate}[label=\Alph*.]
    \item Compound statement DutgRun
    \item Assignment statement if (condition) statement1
    \item \textbf{Control statement else} \hfill \textcolor{blue}{\textbf{[Correct]}}
    \item None of others statement2; Code FUOVERFLOW.COM
\end{enumerate}

\noindent \textbf{Correct Answer:} \textcolor{blue}{\textbf{C}} \\[2pt]
\noindent \textbf{Explanation:} Trong ngôn ngữ lập trình thủ tục, các câu lệnh thông thường được CPU thực thi tuần tự từ trên xuống dưới; tuy nhiên, câu lệnh điều khiển (Control statement - đáp án C) được sử dụng để thay đổi trật tự thực thi tuần tự này dựa trên kết quả kiểm tra điều kiện (cấu trúc rẽ nhánh) hoặc thực hiện lặp lại một khối lệnh. Hình minh họa cấu trúc 'if (condition) statement1 else statement2' là ví dụ kinh điển của câu lệnh điều khiển rẽ nhánh. Câu lệnh gán (Assignment statement) chỉ gán giá trị cho biến, còn Câu lệnh ghép (Compound statement) là tập hợp nhiều câu lệnh được gom lại trong một khối mã.

\par\vspace{0.3cm}
\hrule
\vspace{0.3cm}
\end{minipage}

\begin{minipage}{\textwidth}
\subsection*{Question 248 \quad \footnotesize{[0248\_q248\_Question\_248.png]}}
\noindent \textbf{If modules are written so that they can be reused in other systems then they have high levels of}

\begin{enumerate}[label=\Alph*.]
    \item Flexibility
    \item Interoperability
    \item \textbf{Resusability} \hfill \textcolor{blue}{\textbf{[Correct]}}
    \item Changeability
\end{enumerate}

\noindent \textbf{Correct Answer:} \textcolor{blue}{\textbf{C}} \\[2pt]
\noindent \textbf{Explanation:} Trong kỹ nghệ phần mềm, Tính tái sử dụng (Reusability - đáp án C, trong đề viết là Resusability) là thuộc tính chất lượng của các mô-đun phần mềm được thiết kế và xây dựng độc lập, chuẩn hóa để có thể được sử dụng lại trong các hệ thống hoặc dự án khác mà không cần phải viết lại mã nguồn. Tính tương tác (Interoperability) thể hiện khả năng phối hợp hoạt động giữa các hệ thống phần mềm khác nhau, Tính linh hoạt (Flexibility) là khả năng thích ứng với các thay đổi yêu cầu, và Tính dễ sửa đổi (Changeability) là mức độ dễ dàng khi bảo trì mã nguồn.

\par\vspace{0.3cm}
\hrule
\vspace{0.3cm}
\end{minipage}

\begin{minipage}{\textwidth}
\subsection*{Question 249 \quad \footnotesize{[0249\_q249\_Question\_249.png]}}
\noindent \textbf{Which method in black-box testing is to test the software for all possible values in the input domain?}

\begin{enumerate}[label=\Alph*.]
    \item Boundary-value testing
    \item Random testing
    \item \textbf{Exhaustive testing} \hfill \textcolor{blue}{\textbf{[Correct]}}
    \item Data flow testing
\end{enumerate}

\noindent \textbf{Correct Answer:} \textcolor{blue}{\textbf{C}} \\[2pt]
\noindent \textbf{Explanation:} Kiểm thử vét cạn (Exhaustive testing - đáp án C) là phương pháp kiểm thử hộp đen trong đó phần mềm được thực thi và kiểm tra với mọi giá trị đầu vào có thể có trong toàn bộ miền giá trị đầu vào (input domain) nhằm phát hiện triệt để mọi lỗi tiềm ẩn. Mặc dù lý tưởng về lý thuyết, phương pháp này rất hiếm khi khả thi trong thực tế do số lượng trường hợp thử nghiệm có thể là vô hạn hoặc quá lớn. Kiểm thử giá trị biên (Boundary-value testing) chỉ kiểm tra các giá trị tại biên của miền đầu vào, Kiểm thử ngẫu nhiên (Random testing) chỉ chọn ngẫu nhiên một tập mẫu, còn Kiểm thử luồng dữ liệu (Data flow testing) là kỹ thuật kiểm thử hộp trắng.

\par\vspace{0.3cm}
\hrule
\vspace{0.3cm}
\end{minipage}

\begin{minipage}{\textwidth}
\subsection*{Question 250 \quad \footnotesize{[0250\_q250\_Question\_250.png]}}
\noindent \textbf{A(n) is a suitable structure when a small number of insertions and deletions are required but a lot of searching and retrieval is needed}

\begin{enumerate}[label=\Alph*.]
    \item \textbf{Array} \hfill \textcolor{blue}{\textbf{[Correct]}}
    \item Record
    \item Linked list
    \item Tree
\end{enumerate}

\noindent \textbf{Correct Answer:} \textcolor{blue}{\textbf{A}} \\[2pt]
\noindent \textbf{Explanation:} Mảng (Array - đáp án A) là cấu trúc dữ liệu lưu trữ các phần tử cùng kiểu tại các ô nhớ liền kề, cho phép truy xuất ngẫu nhiên cực nhanh tới bất kỳ phần tử nào dựa vào chỉ số với độ phức tạp O(1), rất phù hợp cho các bài toán đòi hỏi tìm kiếm và truy xuất nhiều. Tuy nhiên, thao tác chèn (insertion) và xóa (deletion) trong mảng đòi hỏi phải dịch chuyển toàn bộ các phần tử phía sau nên tốn kém chi phí (O(n)), vì thế mảng là lựa chọn lý tưởng nhất khi số lượng thao tác thêm/xóa là rất ít nhưng cần tìm kiếm và đọc dữ liệu thường xuyên. Danh sách liên kết (Linked list) tối ưu cho việc chèn/xóa nhưng truy xuất tuần tự lại chậm.

\par\vspace{0.3cm}
\hrule
\vspace{0.3cm}
\end{minipage}

\begin{minipage}{\textwidth}
\subsection*{Question 251 \quad \footnotesize{[0251\_q251\_Question\_251.png]}}
\noindent \textbf{An empty linked list consists of}

\begin{enumerate}[label=\Alph*.]
    \item A node
    \item Two nodes
    \item Data and a link
    \item \textbf{A null head pointer} \hfill \textcolor{blue}{\textbf{[Correct]}}
\end{enumerate}

\noindent \textbf{Correct Answer:} \textcolor{blue}{\textbf{D}} \\[2pt]
\noindent \textbf{Explanation:} Một danh sách liên kết (linked list) được quản lý và định vị thông qua một con trỏ đầu (head pointer) lưu trữ địa chỉ của nút đầu tiên trong danh sách. Khi danh sách liên kết rỗng (empty linked list), nghĩa là chưa có nút dữ liệu nào được khởi tạo trong bộ nhớ, con trỏ đầu này sẽ không trỏ đến bất kỳ ô nhớ nào và mang giá trị rỗng (A null head pointer - đáp án D). Các lựa chọn chứa một nút hay hai nút thể hiện danh sách đã có dữ liệu, còn dữ liệu kèm liên kết là cấu trúc thành phần của một nút đơn lẻ.

\par\vspace{0.3cm}
\hrule
\vspace{0.3cm}
\end{minipage}

\begin{minipage}{\textwidth}
\subsection*{Question 252 \quad \footnotesize{[0252\_q252\_Question\_252.png]}}
\noindent \textbf{In an inorder traversal of a binary tree, the root is processed}

\begin{enumerate}[label=\Alph*.]
    \item First
    \item \textbf{Second} \hfill \textcolor{blue}{\textbf{[Correct]}}
    \item Last
    \item Two times
\end{enumerate}

\noindent \textbf{Correct Answer:} \textcolor{blue}{\textbf{B}} \\[2pt]
\noindent \textbf{Explanation:} Trong phép duyệt cây nhị phân theo thứ tự giữa (Inorder traversal), quy tắc duyệt đệ quy tiêu chuẩn được tiến hành lần lượt theo thứ tự: Duyệt cây con bên trái (Left) → Xử lý nút gốc (Root) → Duyệt cây con bên phải (Right). Do đó, nút gốc luôn được xử lý ở vị trí thứ hai (Second - đáp án B), sau khi duyệt xong toàn bộ cây con bên trái và trước khi chuyển sang duyệt cây con bên phải. Ngược lại, trong phép duyệt tiền thứ tự (Preorder), nút gốc được duyệt đầu tiên (First); còn trong duyệt hậu thứ tự (Postorder), nút gốc được duyệt cuối cùng (Last).

\par\vspace{0.3cm}
\hrule
\vspace{0.3cm}
\end{minipage}

\begin{minipage}{\textwidth}
\subsection*{Question 253 \quad \footnotesize{[0253\_q253\_Question\_253.png]}}
\noindent \textbf{In the collision resolution method, we try to put data that cannot be placed in location 123 into location 124.}

\begin{enumerate}[label=\Alph*.]
    \item \textbf{Open addressing} \hfill \textcolor{blue}{\textbf{[Correct]}}
    \item Linked list
    \item Bucket hashing
    \item Random hashing
\end{enumerate}

\noindent \textbf{Correct Answer:} \textcolor{blue}{\textbf{A}} \\[2pt]
\noindent \textbf{Explanation:} Phương pháp địa chỉ mở (Open addressing - đáp án A), tiêu biểu là kỹ thuật dò tuyến tính (linear probing), là cơ chế xử lý xung đột trong bảng băm bằng cách tìm kiếm một ô trống tiếp theo ngay trong bảng khi vị trí được băm ban đầu đã bị phần tử khác chiếm đóng. Cụ thể, khi dữ liệu bị xung đột không thể chèn vào vị trí 123, thuật toán sẽ thử đưa dữ liệu vào ô nhớ kế tiếp là 124. Ngược lại, phương pháp danh sách liên kết (Linked list / separate chaining) lưu trữ các phần tử xung đột thành một danh sách móc nối tại cùng vị trí băm đó chứ không dời sang ô nhớ khác.

\par\vspace{0.3cm}
\hrule
\vspace{0.3cm}
\end{minipage}

\begin{minipage}{\textwidth}
\subsection*{Question 254 \quad \footnotesize{[0254\_q254\_Question\_254.png]}}
\noindent \textbf{The level of a three-level DBMS architecture defines the logical view of the data.}

\begin{enumerate}[label=\Alph*.]
    \item External
    \item \textbf{Conceptual} \hfill \textcolor{blue}{\textbf{[Correct]}}
    \item Internal
    \item Physical
\end{enumerate}

\noindent \textbf{Correct Answer:} \textcolor{blue}{\textbf{B}} \\[2pt]
\noindent \textbf{Explanation:} Trong kiến trúc hệ quản trị cơ sở dữ liệu 3 mức tiêu chuẩn ANSI/SPARC, Mức quan niệm (Conceptual level - đáp án B) là tầng mô tả toàn bộ cấu trúc logic tổng thể của cơ sở dữ liệu (bao gồm các thực thể, kiểu dữ liệu, các ràng buộc và mối quan hệ giữa chúng) một cách độc lập với phương thức lưu trữ vật lý. Mức ngoài (External level) mô tả các khung nhìn dữ liệu riêng biệt cho từng người dùng, còn Mức trong (Internal/Physical level) mô tả chi tiết cách thức dữ liệu được tổ chức và lưu trữ vật lý trên các thiết bị lưu trữ. Do đó, mức định nghĩa góc nhìn logic của dữ liệu chính là Conceptual level.

\par\vspace{0.3cm}
\hrule
\vspace{0.3cm}
\end{minipage}

\begin{minipage}{\textwidth}
\subsection*{Question 255 \quad \footnotesize{[0255\_q255\_Question\_255.png]}}
\noindent \textbf{relations (s) and creates an output of A binary operator is applied to relation(s).}

\begin{enumerate}[label=\Alph*.]
    \item One, one
    \item One, two
    \item \textbf{Two, one} \hfill \textcolor{blue}{\textbf{[Correct]}}
    \item Two, two
\end{enumerate}

\noindent \textbf{Correct Answer:} \textcolor{blue}{\textbf{C}} \\[2pt]
\noindent \textbf{Explanation:} Trong đại số quan hệ (relational algebra), một toán tử hai ngôi (binary operator, chẳng hạn như phép hợp, giao, trừ, tích Descartes hay phép kết nối - Join) là toán tử tác động lên đúng hai quan hệ đầu vào (two relations) và tạo ra kết quả đầu ra là một quan hệ mới duy nhất (one relation). Ngược lại, toán tử một ngôi (unary operator như Select hay Project) chỉ nhận đầu vào từ một quan hệ duy nhất. Do đó, toán tử hai ngôi được áp dụng cho 2 quan hệ và tạo ra 1 quan hệ đầu ra (Two, one - đáp án C).

\par\vspace{0.3cm}
\hrule
\vspace{0.3cm}
\end{minipage}

\begin{minipage}{\textwidth}
\subsection*{Question 256 \quad \footnotesize{[0256\_q256\_Question\_256.png]}}
\noindent \textbf{A database management system (DBMS) is a combination of components}

\begin{enumerate}[label=\Alph*.]
    \item 2
    \item 3
    \item 4
    \item \textbf{5} \hfill \textcolor{blue}{\textbf{[Correct]}}
\end{enumerate}

\noindent \textbf{Correct Answer:} \textcolor{blue}{\textbf{D}} \\[2pt]
\noindent \textbf{Explanation:} Theo giáo trình Foundations of Computer Science của Behrouz Forouzan (Chương 11 - Databases), một hệ quản trị cơ sở dữ liệu (DBMS) là sự kết hợp đồng bộ của đúng 5 thành phần chính gồm: Phần cứng (Hardware), Phần mềm (Software), Dữ liệu (Data), Người dùng (Users/People), và Quy trình/Thủ tục vận hành (Procedures). Năm thành phần này phối hợp chặt chẽ với nhau để thu thập, lưu trữ, xử lý và bảo vệ dữ liệu một cách an toàn và hiệu quả. Vì vậy, số lượng thành phần của DBMS là 5 (đáp án D).

\par\vspace{0.3cm}
\hrule
\vspace{0.3cm}
\end{minipage}

\begin{minipage}{\textwidth}
\subsection*{Question 257 \quad \footnotesize{[0257\_q257\_Question\_257.png]}}
\noindent \textbf{In the attacks threatent integrity with attack the attacker modifies the information to make it beneficial to them}

\begin{enumerate}[label=\Alph*.]
    \item Snooping
    \item \textbf{Repudiation} \hfill \textcolor{blue}{\textbf{[Correct]}}
    \item Denial of service
    \item Traffic analysis
\end{enumerate}

\noindent \textbf{Correct Answer:} \textcolor{blue}{\textbf{B}} \\[2pt]
\noindent \textbf{Explanation:} Trong phân loại các cuộc tấn công an ninh thông tin của Behrouz Forouzan (Chương 12), sự chối bỏ trách nhiệm (Repudiation - đáp án B) là một hình thức tấn công đe dọa tính toàn vẹn (integrity), trong đó kẻ tấn công hoặc một bên tham gia giao dịch chối bỏ hành động đã làm hoặc cố tình sửa đổi thông tin giao dịch để mang lại lợi ích cho bản thân (ví dụ phủ nhận đã gửi yêu cầu thanh toán hoặc thay đổi nội dung đơn hàng). Các hình thức như Snooping hay Traffic analysis đe dọa tính bảo mật (confidentiality), còn Denial of service đe dọa tính sẵn sàng (availability).

\par\vspace{0.3cm}
\hrule
\vspace{0.3cm}
\end{minipage}

\begin{minipage}{\textwidth}
\subsection*{Question 258 \quad \footnotesize{[0258\_q258\_Question\_258.png]}}
\noindent \textbf{cryptography is often used for short messages.}

\begin{enumerate}[label=\Alph*.]
    \item Symmetric-key
    \item \textbf{Asymmetric-key} \hfill \textcolor{blue}{\textbf{[Correct]}}
    \item Secret-key
    \item Open-key
\end{enumerate}

\noindent \textbf{Correct Answer:} \textcolor{blue}{\textbf{B}} \\[2pt]
\noindent \textbf{Explanation:} Mật mã học khóa bất đối xứng (Asymmetric-key cryptography - đáp án B), ví dụ như thuật toán RSA, sử dụng các phép tính số học mô-đun phức tạp trên các số nguyên rất lớn nên đòi hỏi năng lực xử lý cao và tốc độ chậm hơn hàng trăm lần so với mã hóa khóa đối xứng. Do đó, mật mã khóa bất đối xứng thường chỉ được ưu tiên dùng cho các thông điệp ngắn (short messages) như chữ ký số, chứng thực danh tính hoặc trao đổi khóa phiên (session key exchange). Ngược lại, mật mã khóa đối xứng (Symmetric-key) có tốc độ rất cao nên chuyên dùng để mã hóa khối lượng dữ liệu lớn.

\par\vspace{0.3cm}
\hrule
\vspace{0.3cm}
\end{minipage}

\begin{minipage}{\textwidth}
\subsection*{Question 259 \quad \footnotesize{[0259\_q259\_Question\_259.png]}}
\noindent \textbf{The ethical principle that dictates that a decision should be made according to the universally accepted principles of morality is referred to as}

\begin{enumerate}[label=\Alph*.]
    \item \textbf{Moral rules} \hfill \textcolor{blue}{\textbf{[Correct]}}
    \item Utilization
    \item Social contract
    \item None of the others
\end{enumerate}

\noindent \textbf{Correct Answer:} \textcolor{blue}{\textbf{A}} \\[2pt]
\noindent \textbf{Explanation:} Trong các lý thuyết đạo đức máy tính (Forouzan Chương 12), Nguyên tắc quy tắc đạo đức (Moral rules - đáp án A, dựa trên thuyết nghĩa vụ luận / deontology của Immanuel Kant) quy định rằng các quyết định và hành vi phải được đưa ra dựa trên các nguyên tắc luân lý được thừa nhận phổ quát (ví dụ luôn phải trung thực, không được gian lận), bất kể kết quả cuối cùng mang lại là gì. Ngược lại, Thuyết hữu dụng (Utilization / Utilitarianism) đánh giá hành động dựa trên việc tối đa hóa lợi ích cho số đông, còn Khế ước xã hội (Social contract) dựa trên các thỏa thuận và quy ước được cộng đồng đồng thuận.

\par\vspace{0.3cm}
\hrule
\vspace{0.3cm}
\end{minipage}

\begin{minipage}{\textwidth}
\subsection*{Question 260 \quad \footnotesize{[0260\_q260\_Question\_260.png]}}
\noindent \textbf{In the asymmetric-key method used for confidentiality, which key(s) is (are) publicly known?}

\begin{enumerate}[label=\Alph*.]
    \item \textbf{Encryption key only} \hfill \textcolor{blue}{\textbf{[Correct]}}
    \item Decryption key only
    \item Both encryption and decryption keys
    \item Neither encryption key nor decryption key
\end{enumerate}

\noindent \textbf{Correct Answer:} \textcolor{blue}{\textbf{A}} \\[2pt]
\noindent \textbf{Explanation:} Trong phương pháp mã hóa khóa bất đối xứng nhằm mục đích bảo vệ tính bí mật (confidentiality), người gửi sử dụng khóa công khai của người nhận (public key / encryption key) để mã hóa bức điện. Bản mã này chỉ có thể được giải mã bởi khóa riêng tư tương ứng của người nhận (private key / decryption key) được người nhận bí mật nắm giữ. Do đó, khóa duy nhất được công khai rộng rãi cho mọi người trong hệ thống là khóa mã hóa (Encryption key only - đáp án A). Nếu khóa giải mã bị công khai thì bất kỳ ai cũng có thể đọc trộm bức điện, làm mất tính bí mật.

\par\vspace{0.3cm}
\hrule
\vspace{0.3cm}
\end{minipage}

\begin{minipage}{\textwidth}
\subsection*{Question 261 \quad \footnotesize{[0261\_q261\_Question\_261.png]}}
\noindent \textbf{The result of the decode step is the for some operation that the system will perform.}

\begin{enumerate}[label=\Alph*.]
    \item Output
    \item Data processing
    \item \textbf{Binary code} \hfill \textcolor{blue}{\textbf{[Correct]}}
    \item Code
\end{enumerate}

\noindent \textbf{Correct Answer:} \textcolor{blue}{\textbf{C}} \\[2pt]
\noindent \textbf{Explanation:} Trong chu kỳ máy cơ bản của CPU (Fetch-Decode-Execute), bước giải mã (Decode step) do Khối điều khiển (Control Unit - CU) thực hiện có nhiệm vụ phân tích mã lệnh (opcode) trong thanh ghi lệnh (IR) và chuyển đổi thành mã nhị phân điều khiển (Binary code - đáp án C) tương ứng. Mã nhị phân này bao gồm các xung tín hiệu điều khiển trực tiếp các đường bus và mạch logic của ALU để chuẩn bị thực thi thao tác đó trong bước Execute tiếp theo. Kết quả dữ liệu đầu ra (Output) chỉ được tạo ra sau khi bước thực thi (Execute) đã hoàn tất.

\par\vspace{0.3cm}
\hrule
\vspace{0.3cm}
\end{minipage}

\begin{minipage}{\textwidth}
\subsection*{Question 262 \quad \footnotesize{[0262\_q262\_Question\_262.png]}}
\noindent \textbf{A number which is stored and processed but not in the standard exponential form is called}

\begin{enumerate}[label=\Alph*.]
    \item Floating point number
    \item \textbf{Fixed point number} \hfill \textcolor{blue}{\textbf{[Correct]}}
    \item Octal number
    \item Decimal number
    \item None of the others
\end{enumerate}

\noindent \textbf{Correct Answer:} \textcolor{blue}{\textbf{B}} \\[2pt]
\noindent \textbf{Explanation:} Số dấu phẩy tĩnh (Fixed point number - đáp án B) là phương pháp biểu diễn số trong máy tính mà vị trí của dấu phân cách thập phân (hoặc nhị phân) được cố định trước giữa các chữ số và không sử dụng dạng ký hiệu số mũ chuẩn. Ngược lại, Số dấu phẩy động (Floating point number) biểu diễn các số dưới dạng số mũ khoa học (mantissa × base\textasciicircum{}exponent), cho phép dấu phân cách có thể trôi nổi linh hoạt để bao quát dải giá trị rất lớn hoặc rất nhỏ. Do đó, số được lưu trữ và xử lý không theo dạng số mũ chuẩn được gọi là số dấu phẩy tĩnh.

\par\vspace{0.3cm}
\hrule
\vspace{0.3cm}
\end{minipage}

\begin{minipage}{\textwidth}
\subsection*{Question 263 \quad \footnotesize{[0263\_q263\_Question\_263.png]}}
\noindent \textbf{A group includes bits in the binary system are represented as one digit in the hexadecimal system. SÓ DÁP ÁN CÃN CHN}

\begin{enumerate}[label=\Alph*.]
    \item 1
    \item 2
    \item 3
    \item \textbf{4} \hfill \textcolor{blue}{\textbf{[Correct]}}
    \item None of the others
\end{enumerate}

\noindent \textbf{Correct Answer:} \textcolor{blue}{\textbf{D}} \\[2pt]
\noindent \textbf{Explanation:} Trong kiến trúc máy tính và biểu diễn số học, mỗi chữ số trong hệ thập lục phân (hexadecimal - cơ số 16) tương ứng chính xác với một nhóm gồm 4 bit trong hệ nhị phân (binary - cơ số 2), bởi vì 2⁴ = 16. Ví dụ, chữ số hex 'A' tương ứng với giá trị thập phân 10 và được biểu diễn bằng nhóm 4 bit là 1010₂. Việc gom nhóm 4 bit này giúp con người dễ dàng đọc, ghi nhớ và biểu diễn các chuỗi bit nhị phân dài một cách ngắn gọn và thuận tiện hơn. Vì vậy, số bit trong mỗi nhóm là 4 (đáp án D).

\par\vspace{0.3cm}
\hrule
\vspace{0.3cm}
\end{minipage}

\begin{minipage}{\textwidth}
\subsection*{Question 264 \quad \footnotesize{[0264\_q264\_Question\_264.png]}}
\noindent \textbf{In two's complement representation with a 4-bit allocation, we get when we add 5 to 5.}

\begin{enumerate}[label=\Alph*.]
    \item -5
    \item \textbf{-6} \hfill \textcolor{blue}{\textbf{[Correct]}}
    \item -7
    \item 10
\end{enumerate}

\noindent \textbf{Correct Answer:} \textcolor{blue}{\textbf{B}} \\[2pt]
\noindent \textbf{Explanation:} Trong hệ biểu diễn số nguyên có dấu bù 2 (two's complement) với độ dài 4 bit, dải giá trị biểu diễn hợp lệ là từ -2³ đến 2³ - 1 (tức từ -8 đến +7). Giá trị +5 được biểu diễn bằng 4 bit là 0101₂; khi cộng hai số +5 với nhau ta được: 0101₂ + 0101₂ = 1010₂. Vì bit ngoài cùng bên trái (MSB) là 1, kết quả này bị máy tính hiểu là một số âm có giá trị bù 2 là -8 + 2 = -6 (đáp án B). Đây là ví dụ kinh điển về hiện tượng tràn số (overflow) khi cộng hai số dương lại thu được một kết quả âm do vượt quá giá trị cực đại +7.

\par\vspace{0.3cm}
\hrule
\vspace{0.3cm}
\end{minipage}

\begin{minipage}{\textwidth}
\subsection*{Question 265 \quad \footnotesize{[0265\_q265\_Question\_265.png]}}
\noindent \textbf{EBCDIC can code up to how many different characters?}

\begin{enumerate}[label=\Alph*.]
    \item 8
    \item 16
    \item 32
    \item 64
    \item \textbf{256} \hfill \textcolor{blue}{\textbf{[Correct]}}
\end{enumerate}

\noindent \textbf{Correct Answer:} \textcolor{blue}{\textbf{E}} \\[2pt]
\noindent \textbf{Explanation:} Bảng mã EBCDIC (Extended Binary Coded Decimal Interchange Code - đáp án E) là chuẩn mã hóa ký tự sử dụng độ dài 8 bit do hãng IBM phát triển cho các dòng máy tính lớn (mainframe). Vì mỗi ký tự được mã hóa bằng một chuỗi 8 bit, tổng số tổ hợp ký tự phân biệt tối đa mà hệ thống này có thể biểu diễn là 2⁸ = 256 ký tự khác nhau (tương tự như bảng mã ASCII mở rộng 8 bit). Các giá trị 64 hay 128 tương ứng với các chuẩn mã hóa ký tự ít bit hơn như BCD 6 bit hay ASCII chuẩn 7 bit.

\par\vspace{0.3cm}
\hrule
\vspace{0.3cm}
\end{minipage}

\begin{minipage}{\textwidth}
\subsection*{Question 266 \quad \footnotesize{[0266\_q266\_Question\_266.png]}}
\noindent \textbf{In the storing audio, if the signal is changing rapidly, we need more samples. It has been shown that a of 40000 samples per second is good enough to reproduce an audio signal.}

\begin{enumerate}[label=\Alph*.]
    \item Quantization
    \item \textbf{Sampling rate} \hfill \textcolor{blue}{\textbf{[Correct]}}
    \item Encoding
    \item Bit depth
\end{enumerate}

\noindent \textbf{Correct Answer:} \textcolor{blue}{\textbf{B}} \\[2pt]
\noindent \textbf{Explanation:} Trong kỹ thuật số hóa âm thanh, Tốc độ lấy mẫu (Sampling rate - đáp án B) là số lượng mẫu biên độ tín hiệu âm thanh được ghi nhận trong mỗi giây (đo bằng Hertz - Hz hoặc samples/second). Khi tín hiệu âm thanh biến đổi nhanh, ta cần nhiều mẫu hơn để có thể tái tạo lại dạng sóng một cách chính xác mà không bị méo tiếng. Theo định lý lấy mẫu Nyquist, để tái tạo dải tần số âm thanh mà tai người nghe được (lên tới 20.000 Hz), tốc độ lấy mẫu cần đạt ít nhất gấp đôi là 40.000 mẫu/giây (ví dụ đĩa CD sử dụng 44.100 mẫu/giây). Do đó, thuật ngữ chính xác là Sampling rate.

\par\vspace{0.3cm}
\hrule
\vspace{0.3cm}
\end{minipage}

\begin{minipage}{\textwidth}
\subsection*{Question 267 \quad \footnotesize{[0267\_q267\_Question\_267.png]}}
\noindent \textbf{Use a circular right shift operation on the bit pattern 1001 1000}

\begin{enumerate}[label=\Alph*.]
    \item \textbf{0100 1100} \hfill \textcolor{blue}{\textbf{[Correct]}}
    \item 0011 0001
    \item 1001 1001
    \item 1001 1000
\end{enumerate}

\noindent \textbf{Correct Answer:} \textcolor{blue}{\textbf{A}} \\[2pt]
\noindent \textbf{Explanation:} Phép dịch vòng phải (circular right shift / rotate right) là thao tác dịch chuyển toàn bộ các bit sang bên phải một vị trí, trong đó bit ngoài cùng bên phải (LSB) bị đẩy ra ngoài sẽ quay vòng trở lại điền vào vị trí trống ngoài cùng bên trái (MSB). Với chuỗi bit ban đầu là 1001 1000, bit ngoài cùng bên phải là 0 sẽ được chuyển lên vị trí đầu tiên bên trái, và 7 bit đầu (1001 100) dịch sang phải một bước tạo thành chuỗi kết quả 0100 1100 (đáp án A). Các phương án còn lại đưa ra chuỗi bit không khớp với nguyên lý xoay vòng sang phải.

\par\vspace{0.3cm}
\hrule
\vspace{0.3cm}
\end{minipage}

\begin{minipage}{\textwidth}
\subsection*{Question 268 \quad \footnotesize{[0268\_q268\_Question\_268.png]}}
\noindent \textbf{The physic-layer addresses, sometimes called \_, are locally defined addresses, each of which defines a specific host or router in a network (LAN or WAN).}

\begin{enumerate}[label=\Alph*.]
    \item TCP
    \item Port numbers
    \item \textbf{MAC addresses} \hfill \textcolor{blue}{\textbf{[Correct]}}
    \item IP address
\end{enumerate}

\noindent \textbf{Correct Answer:} \textcolor{blue}{\textbf{C}} \\[2pt]
\noindent \textbf{Explanation:} Địa chỉ tầng vật lý, thường được gọi phổ biến là địa chỉ MAC (Media Access Control addresses - đáp án C) hay địa chỉ phần cứng, là định danh cục bộ duy nhất được nhà sản xuất ghi cố định vào bộ điều khiển giao tiếp mạng (NIC) để nhận diện cụ thể một máy tính hoặc router trong một mạng cục bộ (LAN hoặc WAN). Địa chỉ IP là địa chỉ logic thuộc tầng mạng (Network layer) dùng để định tuyến toàn cầu, còn Số hiệu cổng (Port numbers) thuộc tầng giao vận để định danh tiến trình ứng dụng. Do đó, địa chỉ tầng vật lý được gọi là địa chỉ MAC.

\par\vspace{0.3cm}
\hrule
\vspace{0.3cm}
\end{minipage}

\begin{minipage}{\textwidth}
\subsection*{Question 269 \quad \footnotesize{[0269\_q269\_Question\_269.png]}}
\noindent \textbf{Instead of using a numeric address, what unique address can we use to connect a host to the Internet in TCP/IP protocol?}

\begin{enumerate}[label=\Alph*.]
    \item \textbf{DNS} \hfill \textcolor{blue}{\textbf{[Correct]}}
    \item Port number
    \item MAC addresses
    \item IP address
\end{enumerate}

\noindent \textbf{Correct Answer:} \textcolor{blue}{\textbf{A}} \\[2pt]
\noindent \textbf{Explanation:} Trong bộ giao thức TCP/IP, thay vì phải ghi nhớ các dãy địa chỉ IP dạng số khô khan và phức tạp (ví dụ 142.250.190.46), người dùng có thể sử dụng các tên miền dạng chữ thân thiện và có ý nghĩa thông qua Hệ thống phân giải tên miền (DNS - Domain Name System, đáp án A) để kết nối máy tính vào Internet. Hệ thống DNS hoạt động như một danh bạ Internet, tự động chuyển đổi tên miền thân thiện thành địa chỉ IP tương ứng để định tuyến các gói tin. Địa chỉ MAC chỉ có phạm vi cục bộ, còn Port number dùng để phân biệt các ứng dụng trên cùng một máy.

\par\vspace{0.3cm}
\hrule
\vspace{0.3cm}
\end{minipage}


\newpage
\section*{8. Software Engineering}
\addcontentsline{toc}{section}{8. Software Engineering}
\hrule\vspace{0.5cm}

\begin{minipage}{\textwidth}
\subsection*{Question 270 \quad \footnotesize{[0270\_q270\_Question\_270.png]}}
\noindent \textbf{At the transport layer, TCP groups a number of bytes together into a packet called a}

\begin{enumerate}[label=\Alph*.]
    \item Frame
    \item Bit pattern
    \item \textbf{Segment} \hfill \textcolor{blue}{\textbf{[Correct]}}
    \item Packet
\end{enumerate}

\noindent \textbf{Correct Answer:} \textcolor{blue}{\textbf{C}} \\[2pt]
\noindent \textbf{Explanation:} Ở tầng giao vận (Transport layer), giao thức TCP phân chia luồng dữ liệu thành các gói tin có kích thước phù hợp gọi là Segment (phân đoạn). Mỗi TCP segment bao gồm phần tiêu đề điều khiển (TCP header) chứa thông tin cổng nguồn, cổng đích, số thứ tự (sequence number) và phần dữ liệu. Trong khi đó, Frame là đơn vị dữ liệu ở tầng liên kết dữ liệu (Data Link layer), còn Packet thường dùng để chỉ đơn vị dữ liệu ở tầng mạng (Network layer).

\par\vspace{0.3cm}
\hrule
\vspace{0.3cm}
\end{minipage}

\begin{minipage}{\textwidth}
\subsection*{Question 271 \quad \footnotesize{[0271\_q271\_Question\_271.png]}}
\noindent \textbf{A mechanism for arranging controlled access to a shared resource is}

\begin{enumerate}[label=\Alph*.]
    \item Retrieving
    \item Sorting
    \item Balleting
    \item \textbf{Lock-out} \hfill \textcolor{blue}{\textbf{[Correct]}}
    \item None of the others
\end{enumerate}

\noindent \textbf{Correct Answer:} \textcolor{blue}{\textbf{D}} \\[2pt]
\noindent \textbf{Explanation:} Trong quản lý tài nguyên và đồng bộ hóa tiến trình của hệ điều hành, cơ chế khóa loại trừ (Lock-out hay Locking mechanism) được sử dụng để kiểm soát và điều phối quyền truy cập độc quyền vào một tài nguyên dùng chung (shared resource). Cơ chế này đảm bảo tại một thời điểm chỉ có một tiến trình được phép thao tác trên tài nguyên, ngăn chặn hiện tượng tương tranh (race condition). Các phương án như Retrieving (truy xuất) hay Sorting (sắp xếp) là các thao tác xử lý dữ liệu thông thường chứ không phải cơ chế kiểm soát truy cập.

\par\vspace{0.3cm}
\hrule
\vspace{0.3cm}
\end{minipage}

\begin{minipage}{\textwidth}
\subsection*{Question 272 \quad \footnotesize{[0272\_q272\_Question\_272.png]}}
\noindent \textbf{In peripheral device management, what is the principle of hierarchical work on OS of third- generation computers?}

\begin{enumerate}[label=\Alph*.]
    \item Processor - Peripheral device - Control device
    \item \textbf{Processor - Control device - Peripheral device} \hfill \textcolor{blue}{\textbf{[Correct]}}
    \item Control device - Peripheral device - Processor
    \item None of the answers are correct
\end{enumerate}

\noindent \textbf{Correct Answer:} \textcolor{blue}{\textbf{B}} \\[2pt]
\noindent \textbf{Explanation:} Trong quản lý thiết bị ngoại vi của máy tính thế hệ thứ ba, hệ thống phân cấp điều khiển hoạt động theo thứ tự: Bộ xử lý (Processor) -> Thiết bị điều khiển (Control device) -> Thiết bị ngoại vi (Peripheral device). Bộ xử lý (CPU) không trực tiếp điều khiển các hoạt động cơ - điện tử chi tiết của ngoại vi mà truyền lệnh cho bộ điều khiển thiết bị (device controller); sau đó controller mới trực tiếp thao tác và điều khiển thiết bị ngoại vi tương ứng. Mô hình phân cấp này giúp giải phóng tải cho CPU và tăng hiệu quả xử lý đa nhiệm của hệ thống.

\par\vspace{0.3cm}
\hrule
\vspace{0.3cm}
\end{minipage}

\begin{minipage}{\textwidth}
\subsection*{Question 273 \quad \footnotesize{[0273\_q273\_Question\_273.png]}}
\noindent \textbf{The construct handles repeated actions.}

\begin{enumerate}[label=\Alph*.]
    \item Sequence
    \item Decision
    \item \textbf{Repetition} \hfill \textcolor{blue}{\textbf{[Correct]}}
    \item Flow
\end{enumerate}

\noindent \textbf{Correct Answer:} \textcolor{blue}{\textbf{C}} \\[2pt]
\noindent \textbf{Explanation:} Trong lập trình có cấu trúc và thiết kế thuật toán, cấu trúc lặp (Repetition construct) được sử dụng để thực hiện lặp đi lặp lại một hành động hoặc khối lệnh cho đến khi thỏa mãn một điều kiện dừng. Bên cạnh cấu trúc lặp, hai cấu trúc điều khiển cơ bản còn lại là cấu trúc tuần tự (Sequence) để thực hiện các lệnh nối tiếp nhau và cấu trúc lựa chọn (Decision) để rẽ nhánh luồng thực thi. Do đó, cấu trúc phụ trách xử lý các hành động lặp lại là Repetition.

\par\vspace{0.3cm}
\hrule
\vspace{0.3cm}
\end{minipage}

\begin{minipage}{\textwidth}
\subsection*{Question 274 \quad \footnotesize{[0274\_q274\_Question\_274.png]}}
\noindent \textbf{Time during which a job is processed by the computer is}

\begin{enumerate}[label=\Alph*.]
    \item Delay time
    \item Real time
    \item \textbf{Execution time} \hfill \textcolor{blue}{\textbf{[Correct]}}
    \item Down time
    \item None of the others
\end{enumerate}

\noindent \textbf{Correct Answer:} \textcolor{blue}{\textbf{C}} \\[2pt]
\noindent \textbf{Explanation:} Thời gian thực thi (Execution time hay Run time) là khoảng thời gian máy tính và bộ xử lý (CPU) đang trực tiếp thực hiện các câu lệnh của một công việc (job) hoặc tiến trình. Ngược lại, Delay time là khoảng thời gian công việc phải chờ đợi trong hàng đợi trước khi được xử lý, còn Down time là thời gian hệ thống ngừng hoạt động do lỗi hoặc bảo trì. Do đó, thời gian công việc thực sự được máy tính xử lý chính là Execution time.

\par\vspace{0.3cm}
\hrule
\vspace{0.3cm}
\end{minipage}

\begin{minipage}{\textwidth}
\subsection*{Question 275 \quad \footnotesize{[0275\_q275\_Question\_275.png]}}
\noindent \textbf{\_language. Prolog is an example of a(n)}

\begin{enumerate}[label=\Alph*.]
    \item Procedural
    \item Functional
    \item \textbf{Declarative} \hfill \textcolor{blue}{\textbf{[Correct]}}
    \item Object-oriented
\end{enumerate}

\noindent \textbf{Correct Answer:} \textcolor{blue}{\textbf{C}} \\[2pt]
\noindent \textbf{Explanation:} Prolog (Programming in Logic) là ngôn ngữ lập trình tiêu biểu đại diện cho mô hình lập trình khai báo (Declarative paradigm), cụ thể là lập trình logic. Thay vì mô tả từng bước thuật toán giải quyết bài toán như trong lập trình thủ tục (Procedural), lập trình viên trong Prolog chỉ cần khai báo các sự kiện (facts), quy tắc (rules) và mục tiêu để hệ thống tự suy diễn đưa ra kết luận. Các mô hình khác như Procedural bao gồm C/Pascal, hay Object-oriented bao gồm C++/Java.

\par\vspace{0.3cm}
\hrule
\vspace{0.3cm}
\end{minipage}

\begin{minipage}{\textwidth}
\subsection*{Question 276 \quad \footnotesize{[0276\_q276\_Question\_276.png]}}
\noindent \textbf{Pascal is a(n) \_language.}

\begin{enumerate}[label=\Alph*.]
    \item \textbf{Procedural} \hfill \textcolor{blue}{\textbf{[Correct]}}
    \item Functional
    \item Declarative
    \item Object-oriented
\end{enumerate}

\noindent \textbf{Correct Answer:} \textcolor{blue}{\textbf{A}} \\[2pt]
\noindent \textbf{Explanation:} Pascal là ngôn ngữ lập trình thuộc mô hình lập trình thủ tục (Procedural paradigm), được Niklaus Wirth thiết kế nhằm khuyến khích phong cách lập trình có cấu trúc. Trong ngôn ngữ thủ tục, chương trình được cấu thành từ một chuỗi các lệnh tuần tự, thủ tục (procedures) và hàm (functions) tác động làm thay đổi trạng thái của các biến trong bộ nhớ. Ngược lại, Prolog thuộc mô hình khai báo (Declarative), còn LISP/Haskell thuộc mô hình hàm (Functional).

\par\vspace{0.3cm}
\hrule
\vspace{0.3cm}
\end{minipage}

\begin{minipage}{\textwidth}
\subsection*{Question 277 \quad \footnotesize{[0277\_q277\_Question\_277.png]}}
\noindent \textbf{Which data of the following cannot change value with time?}

\begin{enumerate}[label=\Alph*.]
    \item \textbf{Constants} \hfill \textcolor{blue}{\textbf{[Correct]}}
    \item Literals
    \item Input/output
    \item Variables
\end{enumerate}

\noindent \textbf{Correct Answer:} \textcolor{blue}{\textbf{A}} \\[2pt]
\noindent \textbf{Explanation:} Hằng số (Constants) là các đại lượng dữ liệu được đặt tên trong chương trình mà giá trị của chúng được cố định ngay từ khi khởi tạo và không thể bị thay đổi trong suốt quá trình chương trình thực thi. Ngược lại, biến số (Variables) là vùng nhớ có giá trị có thể thay đổi linh hoạt theo thời gian thực thi, và dữ liệu vào/ra (Input/output) cũng biến đổi liên tục tùy thuộc vào tác vụ của người dùng hoặc hệ thống.

\par\vspace{0.3cm}
\hrule
\vspace{0.3cm}
\end{minipage}

\begin{minipage}{\textwidth}
\subsection*{Question 278 \quad \footnotesize{[0278\_q278\_Question\_278.png]}}
\noindent \textbf{between modules in a software system must be minimized.}

\begin{enumerate}[label=\Alph*.]
    \item \textbf{Coupling} \hfill \textcolor{blue}{\textbf{[Correct]}}
    \item Cohesion
    \item Neither coupling nor cohesion
    \item Both coupling and cohesion
\end{enumerate}

\noindent \textbf{Correct Answer:} \textcolor{blue}{\textbf{A}} \\[2pt]
\noindent \textbf{Explanation:} Trong thiết kế kiến trúc phần mềm, nguyên tắc cốt lõi là 'Coupling lỏng, Cohesion chặt' (Loose coupling, High cohesion), trong đó độ ghép nối (Coupling) giữa các module phải được cực tiểu hóa (minimized). Coupling đo lường mức độ phụ thuộc lẫn nhau giữa các module; độ ghép nối càng thấp thì các module càng độc lập, giúp việc sửa đổi, kiểm thử và tái sử dụng mã nguồn trở nên dễ dàng hơn. Ngược lại, độ gắn kết nội tại (Cohesion) bên trong mỗi module cần phải được tối đa hóa.

\par\vspace{0.3cm}
\hrule
\vspace{0.3cm}
\end{minipage}

\begin{minipage}{\textwidth}
\subsection*{Question 279 \quad \footnotesize{[0279\_q279\_Question\_279.png]}}
\noindent \textbf{In the software quality, what factor measures can be mentioned for maintainability?}

\begin{enumerate}[label=\Alph*.]
    \item Accuracy
    \item Timeliness
    \item Resusability
    \item \textbf{Changeability} \hfill \textcolor{blue}{\textbf{[Correct]}}
\end{enumerate}

\noindent \textbf{Correct Answer:} \textcolor{blue}{\textbf{D}} \\[2pt]
\noindent \textbf{Explanation:} Theo các tiêu chuẩn đo lường chất lượng phần mềm (như ISO 9126 và giáo trình Forouzan), khả năng bảo trì (Maintainability) đánh giá mức độ dễ dàng khi sửa đổi hệ thống, bao gồm thuộc tính đo lường trực tiếp là khả năng thay đổi (Changeability). Changeability đo lường nỗ lực cần thiết để điều chỉnh, cải tiến hoặc thích ứng phần mềm khi có sự thay đổi về yêu cầu kỹ thuật hay môi trường vận hành. Trong khi đó, Accuracy đo lường tính đúng đắn thuộc nhóm chức năng, còn Timeliness thuộc về hiệu năng xử lý.

\par\vspace{0.3cm}
\hrule
\vspace{0.3cm}
\end{minipage}

\begin{minipage}{\textwidth}
\subsection*{Question 280 \quad \footnotesize{[0280\_q280\_Question\_280.png]}}
\noindent \textbf{is a method in which each statement in the software is executed at least once.}

\begin{enumerate}[label=\Alph*.]
    \item \textbf{Basis path testing} \hfill \textcolor{blue}{\textbf{[Correct]}}
    \item Condition testing
    \item Loop testing
    \item Boundary testing
\end{enumerate}

\noindent \textbf{Correct Answer:} \textcolor{blue}{\textbf{A}} \\[2pt]
\noindent \textbf{Explanation:} Kiểm thử đường cơ sở (Basis path testing) là một kỹ thuật kiểm thử hộp trắng (White-box testing) dựa trên đồ thị dòng điều khiển (Control Flow Graph), đảm bảo rằng mọi câu lệnh trong chương trình đều được thực thi ít nhất một lần. Kỹ thuật này xác định tập hợp các đường đi độc lập cơ sở để xây dựng các ca kiểm thử đạt độ bao phủ câu lệnh (statement coverage) và bao phủ nhánh đầy đủ. Trong khi đó, Condition testing chỉ kiểm tra các biểu thức điều kiện logic, còn Loop testing chuyên biệt cho việc kiểm tra cấu trúc lặp.

\par\vspace{0.3cm}
\hrule
\vspace{0.3cm}
\end{minipage}

\begin{minipage}{\textwidth}
\subsection*{Question 281 \quad \footnotesize{[0281\_q281\_Question\_281.png]}}
\noindent \textbf{Before insertion into a linked list we first apply the searching algorithm because we DO NOT}

\begin{enumerate}[label=\Alph*.]
    \item Allow insertion into an empty list
    \item Allow insertion in the middle of the list
    \item Allow insertion at the beginning of the list
    \item \textbf{Allow data with duplicate values} \hfill \textcolor{blue}{\textbf{[Correct]}}
\end{enumerate}

\noindent \textbf{Correct Answer:} \textcolor{blue}{\textbf{D}} \\[2pt]
\noindent \textbf{Explanation:} Trong danh sách liên kết có thứ tự (ordered linked list), thao tác chèn phần tử luôn yêu cầu áp dụng thuật toán tìm kiếm trước tiên vì danh sách quy ước không cho phép chèn các phần tử có giá trị khóa trùng lặp (duplicate values). Việc tìm kiếm giúp định vị chính xác vị trí cần chèn thỏa mãn thứ tự sắp xếp và đồng thời phát hiện nếu khóa đã tồn tại để hủy bỏ thao tác chèn. Danh sách liên kết hoàn toàn cho phép chèn vào danh sách rỗng, đầu danh sách hay giữa danh sách một cách bình thường.

\par\vspace{0.3cm}
\hrule
\vspace{0.3cm}
\end{minipage}

\begin{minipage}{\textwidth}
\subsection*{Question 282 \quad \footnotesize{[0282\_q282\_Question\_282.png]}}
\noindent \textbf{In a three-level DBMS architecture, the layer that interacts directly with the hardware is the level.}

\begin{enumerate}[label=\Alph*.]
    \item External
    \item Conceptual
    \item \textbf{\_Internal} \hfill \textcolor{blue}{\textbf{[Correct]}}
    \item Physical
\end{enumerate}

\noindent \textbf{Correct Answer:} \textcolor{blue}{\textbf{C}} \\[2pt]
\noindent \textbf{Explanation:} Trong kiến trúc 3 mức của hệ quản trị cơ sở dữ liệu (ANSI/SPARC DBMS architecture), mức trong (Internal level) là tầng mô tả cấu trúc lưu trữ vật lý của dữ liệu và trực tiếp tương tác với hệ thống phần cứng cũng như hệ điều hành. Mức này quy định cách thức các bản ghi được lưu trên đĩa từ, tổ chức tệp tin, phân bổ dung lượng và xây dựng các chỉ mục (indexes). Trong khi đó, mức ngoài (External level) dành cho góc nhìn của người dùng, còn mức quan niệm (Conceptual level) mô tả toàn bộ cấu trúc logic tổng thể của CSDL.

\par\vspace{0.3cm}
\hrule
\vspace{0.3cm}
\end{minipage}

\begin{minipage}{\textwidth}
\subsection*{Question 283 \quad \footnotesize{[0283\_q283\_Question\_283.png]}}
\noindent \textbf{In the hierarchical database model, data is organized as an inverted}

\begin{enumerate}[label=\Alph*.]
    \item Graph
    \item \textbf{Tree} \hfill \textcolor{blue}{\textbf{[Correct]}}
    \item Array
    \item Record
\end{enumerate}

\noindent \textbf{Correct Answer:} \textcolor{blue}{\textbf{B}} \\[2pt]
\noindent \textbf{Explanation:} Trong mô hình cơ sở dữ liệu phân cấp (Hierarchical database model), dữ liệu được tổ chức dưới dạng một cây đảo ngược (inverted tree) với cấu trúc phân nhánh gồm các nút cha và nút con. Mỗi nút biểu diễn một thực thể/bản ghi, trong đó nút gốc (root) nằm ở vị trí cao nhất và mỗi nút con chỉ được phép có duy nhất một nút cha (quan hệ 1:N). Ngược lại, mô hình mạng (Network model) sử dụng cấu trúc đồ thị (graph) cho phép một nút con có thể có nhiều nút cha.

\par\vspace{0.3cm}
\hrule
\vspace{0.3cm}
\end{minipage}

\begin{minipage}{\textwidth}
\subsection*{Question 284 \quad \footnotesize{[0284\_q284\_Question\_284.png]}}
\noindent \textbf{What is the result of difference database operation from 2 relation tables?}

\begin{enumerate}[label=\Alph*.]
    \item Neither the first relation nor the second.
    \item Both of the first relation and the second.
    \item \textbf{The first relation but not the second.} \hfill \textcolor{blue}{\textbf{[Correct]}}
    \item The second but not the first
\end{enumerate}

\noindent \textbf{Correct Answer:} \textcolor{blue}{\textbf{C}} \\[2pt]
\noindent \textbf{Explanation:} Phép trừ trong đại số quan hệ (Difference operation: R - S) giữa hai bảng quan hệ tương thích hợp sẽ tạo ra một quan hệ mới chứa tất cả các bộ (tuples) thuộc quan hệ thứ nhất (R) nhưng không xuất hiện trong quan hệ thứ hai (S). Phép toán này tương tự như phép hiệu trong lý thuyết tập hợp, loại bỏ mọi phần tử trùng lặp với tập thứ hai. Nếu lấy các bộ thuộc cả hai bảng thì đó là phép giao (Intersection), còn lấy tất cả các bộ của cả hai bảng là phép hợp (Union).

\par\vspace{0.3cm}
\hrule
\vspace{0.3cm}
\end{minipage}

\begin{minipage}{\textwidth}
\subsection*{Question 285 \quad \footnotesize{[0285\_q285\_Question\_285.png]}}
\noindent \textbf{For each entity set in the E-R diagram we create a relation (table) in which there are n columns related to the attributes defined for that set.}

\begin{enumerate}[label=\Alph*.]
    \item n − 1
    \item \textbf{n} \hfill \textcolor{blue}{\textbf{[Correct]}}
    \item n + 1
    \item 1 − n
\end{enumerate}

\noindent \textbf{Correct Answer:} \textcolor{blue}{\textbf{B}} \\[2pt]
\noindent \textbf{Explanation:} Khi chuyển đổi một sơ đồ thực thể - liên kết (E-R diagram) sang mô hình cơ sở dữ liệu quan hệ, mỗi tập thực thể (entity set) sẽ được ánh xạ thành một bảng quan hệ riêng biệt. Nếu tập thực thể có n thuộc tính đơn thì bảng quan hệ tương ứng sẽ được tạo với đúng n cột để lưu trữ toàn bộ các thuộc tính đó, trong đó khóa chính của tập thực thể đóng vai trò là khóa chính của bảng. Do đó số lượng cột được tạo ra chính xác bằng n.

\par\vspace{0.3cm}
\hrule
\vspace{0.3cm}
\end{minipage}

\begin{minipage}{\textwidth}
\subsection*{Question 286 \quad \footnotesize{[0286\_q286\_Question\_286.png]}}
\noindent \textbf{The main principles of ethics are}

\begin{enumerate}[label=\Alph*.]
    \item \textbf{Moral rules, social contract, ultilization} \hfill \textcolor{blue}{\textbf{[Correct]}}
    \item Ethical decision, social justice, procedural justice
    \item Beneficence, procedural justice, fidelity
    \item Nonmaleficence, social justice, fidelity
\end{enumerate}

\noindent \textbf{Correct Answer:} \textcolor{blue}{\textbf{A}} \\[2pt]
\noindent \textbf{Explanation:} Trong nghiên cứu về đạo đức máy tính và trách nhiệm xã hội theo Forouzan, ba nguyên tắc đạo đức nền tảng kinh điển gồm: Quy tắc đạo đức (Moral rules - nghĩa vụ luận Kantianism), Khế ước xã hội (Social contract - tuân thủ các thỏa thuận chung vì lợi ích cộng đồng), và Thuyết vị lợi (Utilitarianism/utilization - hành vi mang lại điều tốt đẹp lớn nhất cho số đông). Đây là các khung lý thuyết chuẩn mực giúp các chuyên gia công nghệ thông tin đưa ra các phán đoán và quyết định đạo đức nghề nghiệp.

\par\vspace{0.3cm}
\hrule
\vspace{0.3cm}
\end{minipage}

\begin{minipage}{\textwidth}
\subsection*{Question 287 \quad \footnotesize{[0287\_q287\_Question\_287.png]}}
\noindent \textbf{who exploit the internet systems only to make public certain vast datasets of information that would be of benefit to everyone.}

\begin{enumerate}[label=\Alph*.]
    \item Black Hat Hacker
    \item \textbf{Grey Hat Hacker FO} \hfill \textcolor{blue}{\textbf{[Correct]}}
    \item Red Hat Hacker
    \item Blue Hat Hacker
\end{enumerate}

\noindent \textbf{Correct Answer:} \textcolor{blue}{\textbf{B}} \\[2pt]
\noindent \textbf{Explanation:} Hacker mũ xám (Grey Hat Hacker) là những cá nhân xâm nhập và khai thác hệ thống mạng mà không có sự cho phép chính thức, nhưng họ không nhằm mục đích phá hoại độc hại hay trục lợi tài chính bất chính như hacker mũ đen (Black Hat). Thay vào đó, họ thường công khai các tập dữ liệu lớn nhằm vạch trần sai phạm, lỗ hổng bảo mật hoặc phục vụ lợi ích thông tin của cộng đồng. Ngược lại, hacker mũ trắng (White Hat) hoạt động hoàn toàn hợp pháp và được tổ chức ủy quyền bảo mật.

\par\vspace{0.3cm}
\hrule
\vspace{0.3cm}
\end{minipage}

\begin{minipage}{\textwidth}
\subsection*{Question 288 \quad \footnotesize{[0288\_q288\_Question\_288.png]}}
\noindent \textbf{In , the identity of a party is verified once for the entire duration of system access.}

\begin{enumerate}[label=\Alph*.]
    \item \textbf{entity authentication} \hfill \textcolor{blue}{\textbf{[Correct]}}
    \item message integrity
    \item message authentication
    \item message encryption
\end{enumerate}

\noindent \textbf{Correct Answer:} \textcolor{blue}{\textbf{A}} \\[2pt]
\noindent \textbf{Explanation:} Xác thực thực thể (Entity authentication) là cơ chế bảo mật mà danh tính của một bên tham gia (người dùng, thiết bị hoặc tiến trình) được kiểm tra và xác minh một lần duy nhất tại thời điểm bắt đầu phiên làm việc để cấp quyền truy cập trong suốt phiên đó. Ngược lại, xác thực thông điệp (Message authentication) đòi hỏi phải kiểm tra nguồn gốc và tính toàn vẹn của từng thông điệp riêng lẻ trong suốt quá trình truyền tin. Mã hóa thông điệp (Message encryption) chỉ có nhiệm vụ che giấu nội dung để đảm bảo tính bí mật.

\par\vspace{0.3cm}
\hrule
\vspace{0.3cm}
\end{minipage}

\begin{minipage}{\textwidth}
\subsection*{Question 289 \quad \footnotesize{[0289\_q289\_Question\_289.png]}}
\noindent \textbf{The model is the basis for today's computers.}

\begin{enumerate}[label=\Alph*.]
    \item Leibnitz
    \item \textbf{Von Neumann} \hfill \textcolor{blue}{\textbf{[Correct]}}
    \item Pascal
    \item Charles Babbage
\end{enumerate}

\noindent \textbf{Correct Answer:} \textcolor{blue}{\textbf{B}} \\[2pt]
\noindent \textbf{Explanation:} Mô hình Von Neumann (Von Neumann model) là nền tảng kiến trúc phần cứng của hầu hết các máy tính hiện đại ngày nay, dựa trên khái niệm chương trình lưu trữ (stored-program concept) trong đó dữ liệu và chương trình được lưu chung trong một bộ nhớ. Kiến trúc này gồm 4 khối chức năng chính: Bộ nhớ (Memory), Bộ xử lý trung tâm (CPU gồm ALU và CU), Hệ thống vào/ra (I/O) và hệ thống đường truyền (Bus). Các nhà khoa học tiền bối như Pascal, Leibniz hay Babbage đã chế tạo máy tính cơ học nhưng không phải kiến trúc máy tính điện tử hiện đại.

\par\vspace{0.3cm}
\hrule
\vspace{0.3cm}
\end{minipage}

\begin{minipage}{\textwidth}
\subsection*{Question 290 \quad \footnotesize{[0290\_q290\_Question\_290.png]}}
\noindent \textbf{C++ and C\# language are a(n) language}

\begin{enumerate}[label=\Alph*.]
    \item Procedural
    \item Functional
    \item Declarative
    \item \textbf{Object-oriented} \hfill \textcolor{blue}{\textbf{[Correct]}}
\end{enumerate}

\noindent \textbf{Correct Answer:} \textcolor{blue}{\textbf{D}} \\[2pt]
\noindent \textbf{Explanation:} C++ và C\# là các ngôn ngữ lập trình tiêu biểu của mô hình lập trình hướng đối tượng (Object-oriented programming - OOP). Trong mô hình này, chương trình được xây dựng xoay quanh các đối tượng (objects) kết hợp giữa trạng thái dữ liệu (thuộc tính) và hành vi (phương thức), tuân thủ các nguyên lý cốt lõi như tính đóng gói, tính kế thừa và tính đa hình. Trong khi đó, mô hình thủ tục (Procedural) đại diện bởi C/Pascal, mô hình hàm (Functional) đại diện bởi Haskell/LISP, và mô hình khai báo (Declarative) đại diện bởi Prolog/SQL.

\par\vspace{0.3cm}
\hrule
\vspace{0.3cm}
\end{minipage}

\begin{minipage}{\textwidth}
\subsection*{Question 291 \quad \footnotesize{[0291\_q291\_Question\_291.png]}}
\noindent \textbf{Who is 1988 announced the latest version of CD-ROM called CD-ROMXA (extended architecture) which provides a documented means of integrating digital audio with computer data on a CD-ROM disk?}

\begin{enumerate}[label=\Alph*.]
    \item Sony
    \item Philips
    \item Microsoft
    \item \textbf{All of the others} \hfill \textcolor{blue}{\textbf{[Correct]}}
    \item None of the others
\end{enumerate}

\noindent \textbf{Correct Answer:} \textcolor{blue}{\textbf{D}} \\[2pt]
\noindent \textbf{Explanation:} Vào năm 1988, định dạng mở rộng CD-ROM XA (Extended Architecture) cho phép ghi xen kẽ (interleave) âm thanh số với dữ liệu máy tính trên cùng một đĩa quang đã được liên minh phát triển bởi cả ba tập đoàn công nghệ lớn: Sony, Philips và Microsoft. Chuẩn này đánh dấu bước tiến quan trọng giúp các hệ thống máy tính thời bấy giờ có thể đọc và phát nội dung đa phương tiện (multimedia) một cách mượt mà. Vì tất cả các bên trên đều tham gia công bố nên đáp án chính xác là All of the others.

\par\vspace{0.3cm}
\hrule
\vspace{0.3cm}
\end{minipage}

\begin{minipage}{\textwidth}
\subsection*{Question 292 \quad \footnotesize{[0292\_q292\_Question\_292.png]}}
\noindent \textbf{A device invented by Dr. Bobeck in 1966, for mass storage of data is}

\begin{enumerate}[label=\Alph*.]
    \item Punched paper tape
    \item \textbf{Magnetic bubble storage} \hfill \textcolor{blue}{\textbf{[Correct]}}
    \item Magnetic tape
    \item Magnetic disk
    \item None of the others
\end{enumerate}

\noindent \textbf{Correct Answer:} \textcolor{blue}{\textbf{B}} \\[2pt]
\noindent \textbf{Explanation:} Bộ nhớ bọt từ (Magnetic bubble storage) là thiết bị lưu trữ dữ liệu khối lượng lớn được phát minh vào năm 1966 bởi Tiến sĩ Andrew Bobeck tại phòng thí nghiệm Bell Labs. Công nghệ này sử dụng một màng vật liệu từ tính mỏng để lưu trữ các vùng từ hóa nhỏ hình trụ gọi là 'bọt từ' (magnetic bubbles), mỗi bọt tương ứng với một bit dữ liệu nhị phân không bay hơi. Các phương án khác như băng đục lỗ (Punched paper tape) hay băng từ (Magnetic tape) đã ra đời trước đó rất lâu.

\par\vspace{0.3cm}
\hrule
\vspace{0.3cm}
\end{minipage}

\begin{minipage}{\textwidth}
\subsection*{Question 293 \quad \footnotesize{[0293\_q293\_Question\_293.png]}}
\noindent \textbf{What is meant by a dedicated computer?}

\begin{enumerate}[label=\Alph*.]
    \item Which is used by one person only
    \item \textbf{Which is assigned one and only one task} \hfill \textcolor{blue}{\textbf{[Correct]}}
    \item Which uses on kind of software
    \item Which is meant for application software only
    \item None of the others
\end{enumerate}

\noindent \textbf{Correct Answer:} \textcolor{blue}{\textbf{B}} \\[2pt]
\noindent \textbf{Explanation:} Máy tính chuyên dụng (Dedicated computer) là loại máy tính được thiết kế, lập trình hoặc dành riêng để thực hiện một và chỉ một nhiệm vụ duy nhất (assigned one and only one task), ví dụ như máy tính trong hệ thống điều khiển phanh ABS xe hơi, máy bay hay thiết bị gia dụng thông minh. Điều này hoàn toàn khác biệt với máy tính đa năng (General-purpose computer) vốn có khả năng cài đặt nhiều phần mềm khác nhau để phục vụ đa dạng nhu cầu của người dùng. Máy tính cho một cá nhân sử dụng là máy tính cá nhân (Personal Computer).

\par\vspace{0.3cm}
\hrule
\vspace{0.3cm}
\end{minipage}

\begin{minipage}{\textwidth}
\subsection*{Question 294 \quad \footnotesize{[0294\_q294\_Question\_294.png]}}
\noindent \textbf{The base of the decimal number system is}

\begin{enumerate}[label=\Alph*.]
    \item 2
    \item 8
    \item \textbf{10} \hfill \textcolor{blue}{\textbf{[Correct]}}
    \item 16
\end{enumerate}

\noindent \textbf{Correct Answer:} \textcolor{blue}{\textbf{C}} \\[2pt]
\noindent \textbf{Explanation:} Cơ số (base hay radix) của một hệ thống đếm biểu thị số lượng ký số duy nhất được sử dụng để biểu diễn các giá trị trong hệ thống đó. Hệ thập phân (Decimal number system) là hệ đếm cơ số 10 vì nó sử dụng tập hợp gồm đúng 10 ký số từ 0 đến 9. Trong khi đó, cơ số 2 là của hệ nhị phân (Binary), cơ số 8 là của hệ bát phân (Octal), và cơ số 16 là của hệ thập lục phân (Hexadecimal).

\par\vspace{0.3cm}
\hrule
\vspace{0.3cm}
\end{minipage}

\begin{minipage}{\textwidth}
\subsection*{Question 295 \quad \footnotesize{[0295\_q295\_Question\_295.png]}}
\noindent \textbf{In two's complement addition, if there is a final carry after the left most column addition,}

\begin{enumerate}[label=\Alph*.]
    \item Add it to the right most column
    \item Add it to the left most column
    \item \textbf{Discard it} \hfill \textcolor{blue}{\textbf{[Correct]}}
    \item Increase the bit length
\end{enumerate}

\noindent \textbf{Correct Answer:} \textcolor{blue}{\textbf{C}} \\[2pt]
\noindent \textbf{Explanation:} Trong phép cộng số nguyên có dấu theo biểu diễn số bù 2 (Two's complement addition), nếu xuất hiện một bit nhớ cuối cùng (final carry) sinh ra từ phép cộng ở vị trí cột ngoài cùng bên trái (cột bit có trọng số lớn nhất - MSB), thì theo quy tắc chuẩn bit nhớ đó sẽ bị loại bỏ hoàn toàn (discard it). Kết quả thu được sau khi bỏ bit nhớ này chính là giá trị đúng của phép toán trong phạm vi độ dài bit quy định. Việc cộng dồn bit nhớ vào cột ngoài cùng bên phải (End-around carry) chỉ áp dụng cho biểu diễn số bù 1.

\par\vspace{0.3cm}
\hrule
\vspace{0.3cm}
\end{minipage}

\begin{minipage}{\textwidth}
\subsection*{Question 296 \quad \footnotesize{[0296\_q296\_Question\_296.png]}}
\noindent \textbf{A data item which is not broken down into smaller units is}

\begin{enumerate}[label=\Alph*.]
    \item Data element
    \item \textbf{Elementary data item} \hfill \textcolor{blue}{\textbf{[Correct]}}
    \item Data entry
    \item Database management
    \item None of the others
\end{enumerate}

\noindent \textbf{Correct Answer:} \textcolor{blue}{\textbf{B}} \\[2pt]
\noindent \textbf{Explanation:} Mục dữ liệu nguyên tố (Elementary data item) là một phần tử dữ liệu đơn giản nhất trong cấu trúc dữ liệu, không thể phân chia nhỏ hơn thành các đơn vị dữ liệu có ý nghĩa độc lập (ví dụ: mã định danh, số tuổi). Ngược lại, mục dữ liệu phức hợp (Composite data item) là mục dữ liệu có thể tách thành các phần tử nhỏ hơn (ví dụ: ngày tháng năm sinh hoặc họ và tên). Do đó, dữ liệu không thể phân rã thành đơn vị nhỏ hơn được gọi là Elementary data item.

\par\vspace{0.3cm}
\hrule
\vspace{0.3cm}
\end{minipage}

\begin{minipage}{\textwidth}
\subsection*{Question 297 \quad \footnotesize{[0297\_q297\_Question\_297.png]}}
\noindent \textbf{refers to a process that rounds the value of a sample to the closest integer value. For example, if the real value is 17.2, it can be rounded down to 17; if the value is 17.7, it can be rounded up to 18.}

\begin{enumerate}[label=\Alph*.]
    \item \textbf{Quantization} \hfill \textcolor{blue}{\textbf{[Correct]}}
    \item Sampling rate
    \item Encoding
    \item Bit depth
\end{enumerate}

\noindent \textbf{Correct Answer:} \textcolor{blue}{\textbf{A}} \\[2pt]
\noindent \textbf{Explanation:} Lượng tử hóa (Quantization) là quá trình làm tròn giá trị biên độ thực của một mẫu tín hiệu tương tự (analog) liên tục về mức giá trị số nguyên rời rạc gần nhất trong một dải mức cho phép. Đây là giai đoạn thứ hai trong quá trình số hóa âm thanh, diễn ra ngay sau bước lấy mẫu (Sampling) và trước bước mã hóa thành chuỗi bit (Encoding). Ví dụ, giá trị 17.2 được làm tròn xuống 17 và 17.7 được làm tròn lên 18 phản ánh chính xác bản chất của phép lượng tử hóa.

\par\vspace{0.3cm}
\hrule
\vspace{0.3cm}
\end{minipage}

\begin{minipage}{\textwidth}
\subsection*{Question 298 \quad \footnotesize{[0298\_q298\_Question\_298.png]}}
\noindent \textbf{A computer is a box full of electronic}

\begin{enumerate}[label=\Alph*.]
    \item \textbf{Switching devices} \hfill \textcolor{blue}{\textbf{[Correct]}}
    \item Chips
    \item Circuits
    \item Registers
    \item None of the others
\end{enumerate}

\noindent \textbf{Correct Answer:} \textcolor{blue}{\textbf{A}} \\[2pt]
\noindent \textbf{Explanation:} Theo góc nhìn kiến trúc máy tính phần cứng trong giáo trình Forouzan, về mặt bản chất điện tử cơ bản nhất, máy tính là một thiết bị chứa hàng triệu đến hàng tỷ phần tử chuyển mạch điện tử (electronic switching devices hay transistors). Các phần tử chuyển mạch này đóng vai trò bật hoặc tắt dòng điện để biểu diễn hai trạng thái nhị phân 0 và 1, tạo nên các cổng logic và mạch vi xử lý. Các linh kiện như chip hay thanh ghi đều được xây dựng từ tập hợp các thiết bị chuyển mạch này.

\par\vspace{0.3cm}
\hrule
\vspace{0.3cm}
\end{minipage}

\begin{minipage}{\textwidth}
\subsection*{Question 299 \quad \footnotesize{[0299\_q299\_Question\_299.png]}}
\noindent \textbf{An operating system intended for use on microprocessor based systems that support a single user is}

\begin{enumerate}[label=\Alph*.]
    \item PC-DOS
    \item MS-DOS
    \item UNIX
    \item \textbf{CP/M} \hfill \textcolor{blue}{\textbf{[Correct]}}
    \item None of the others
\end{enumerate}

\noindent \textbf{Correct Answer:} \textcolor{blue}{\textbf{D}} \\[2pt]
\noindent \textbf{Explanation:} CP/M (Control Program for Microcomputers) được Gary Kildall phát triển vào năm 1974, là hệ điều hành đơn người dùng (single user) tiêu chuẩn thương mại đầu tiên dành cho các hệ thống máy tính dựa trên bộ vi xử lý (như Intel 8080/Z80). CP/M đã đặt nền móng quản lý tệp và dòng lệnh cho các hệ điều hành cá nhân sau này, trước khi MS-DOS và PC-DOS ra đời trên máy tính IBM PC vào đầu thập niên 1980. Trong khi đó, UNIX ngay từ đầu đã được thiết kế là hệ điều hành đa người dùng (multi-user).

\par\vspace{0.3cm}
\hrule
\vspace{0.3cm}
\end{minipage}

\begin{minipage}{\textwidth}
\subsection*{Question 300 \quad \footnotesize{[0300\_q300\_Question\_300.png]}}
\noindent \textbf{What type of OS evolution that need for more speed and efficiency led to the design with multiple CPUs on the same machine? SÖ DÁP ÁN CÃN CHON}

\begin{enumerate}[label=\Alph*.]
    \item Personal systems
    \item Distributed systems
    \item \textbf{Parallel systems} \hfill \textcolor{blue}{\textbf{[Correct]}}
    \item Real-time systems
\end{enumerate}

\noindent \textbf{Correct Answer:} \textcolor{blue}{\textbf{C}} \\[2pt]
\noindent \textbf{Explanation:} Trong quá trình phát triển của hệ điều hành, nhu cầu gia tăng tốc độ tính toán và hiệu suất xử lý công việc đã thúc đẩy sự ra đời của các hệ thống song song (Parallel systems hay Multiprocessor systems), trong đó nhiều CPU cùng phối hợp hoạt động trên một cỗ máy phần cứng duy nhất. Các CPU này chia sẻ chung không gian bộ nhớ chính, bus hệ thống và các thiết bị ngoại vi để thực thi đồng thời nhiều tác vụ. Ngược lại, hệ thống phân tán (Distributed systems) kết nối nhiều máy tính độc lập với nhau thông qua mạng.

\par\vspace{0.3cm}
\hrule
\vspace{0.3cm}
\end{minipage}

\begin{minipage}{\textwidth}
\subsection*{Question 301 \quad \footnotesize{[0301\_q301\_Question\_301.png]}}
\noindent \textbf{A program becomes a when it is selected by the operating system and brought to the hold state.}

\begin{enumerate}[label=\Alph*.]
    \item \textbf{Job} \hfill \textcolor{blue}{\textbf{[Correct]}}
    \item Process
    \item Deadlock
    \item Partition
\end{enumerate}

\noindent \textbf{Correct Answer:} \textcolor{blue}{\textbf{A}} \\[2pt]
\noindent \textbf{Explanation:} Trong các hệ thống điều hành xử lý theo lô (Batch systems) và quản lý tiến trình theo Forouzan, một chương trình (program) lưu trên đĩa sẽ được gọi là công việc (Job) khi nó được hệ điều hành tiếp nhận và đưa vào hàng đợi tạm giữ (hold state / hold queue). Khi tài nguyên bộ nhớ được cấp phát, Job sẽ được chuyển vào bộ nhớ chính và chính thức trở thành một tiến trình (Process) ở trạng thái sẵn sàng (ready state) để CPU thực thi. Do đó, tại trạng thái hold, chương trình được định danh là một Job.

\par\vspace{0.3cm}
\hrule
\vspace{0.3cm}
\end{minipage}

\begin{minipage}{\textwidth}
\subsection*{Question 302 \quad \footnotesize{[0302\_q302\_Question\_302.png]}}
\noindent \textbf{is a pictorial representation of an algorithm.}

\begin{enumerate}[label=\Alph*.]
    \item \textbf{A UML diagram} \hfill \textcolor{blue}{\textbf{[Correct]}}
    \item A program
    \item Pseudocode
    \item An algorithm
\end{enumerate}

\noindent \textbf{Correct Answer:} \textcolor{blue}{\textbf{A}} \\[2pt]
\noindent \textbf{Explanation:} Sơ đồ UML (UML diagram, cụ thể là Activity diagram) là một biểu diễn trực quan bằng hình ảnh (pictorial representation) của một thuật toán hoặc quy trình xử lý, giúp con người dễ dàng quan sát luồng điều khiển và logic hoạt động. Ngược lại, mã giả (Pseudocode) là dạng văn bản bán cấu trúc sử dụng ngôn ngữ tự nhiên kết hợp cú pháp lập trình, còn chương trình (A program) là tập hợp các dòng lệnh mã nguồn cụ thể. Vì vậy, dạng biểu diễn bằng hình ảnh chính là sơ đồ UML.

\par\vspace{0.3cm}
\hrule
\vspace{0.3cm}
\end{minipage}

\begin{minipage}{\textwidth}
\subsection*{Question 303 \quad \footnotesize{[0303\_q303\_Question\_303.png]}}
\noindent \textbf{A computer language that is written in binary codes only is}

\begin{enumerate}[label=\Alph*.]
    \item \textbf{machine language FUO)} \hfill \textcolor{blue}{\textbf{[Correct]}}
    \item C
    \item C\#
    \item pascal
    \item None of the others
\end{enumerate}

\noindent \textbf{Correct Answer:} \textcolor{blue}{\textbf{A}} \\[2pt]
\noindent \textbf{Explanation:} Ngôn ngữ máy (Machine language) là ngôn ngữ lập trình cấp thấp nhất của máy tính, được viết hoàn toàn bằng các chuỗi mã nhị phân gồm các ký số 0 và 1 mà bộ xử lý trung tâm (CPU) có thể hiểu và thực thi trực tiếp trên phần cứng mà không cần qua bất kỳ bước dịch nào. Các ngôn ngữ như C, C\# hay Pascal là ngôn ngữ bậc cao sử dụng từ khóa tiếng Anh mang tính trừu tượng, đòi hỏi phải được biên dịch hoặc thông dịch sang mã máy trước khi chạy. Do đó đáp án chính xác là machine language.

\par\vspace{0.3cm}
\hrule
\vspace{0.3cm}
\end{minipage}

\begin{minipage}{\textwidth}
\subsection*{Question 304 \quad \footnotesize{[0304\_q304\_Question\_304.png]}}
\noindent \textbf{A(n) is a sequenced collection of elements of the same data type and use index to indicate the ordinal number of the element counting from the beginning}

\begin{enumerate}[label=\Alph*.]
    \item \textbf{Array} \hfill \textcolor{blue}{\textbf{[Correct]}}
    \item Record
    \item Linked list
    \item Tree
\end{enumerate}

\noindent \textbf{Correct Answer:} \textcolor{blue}{\textbf{A}} \\[2pt]
\noindent \textbf{Explanation:} Mảng (Array) là một cấu trúc dữ liệu tuyến tính lưu trữ một tập hợp tuần tự các phần tử có cùng kiểu dữ liệu (homogeneous), trong đó mỗi phần tử được truy xuất trực tiếp thông qua một chỉ số nguyên (index) biểu thị số thứ tự tính từ đầu mảng. Ngược lại, cấu trúc bản ghi (Record) cho phép lưu trữ các trường phần tử có kiểu dữ liệu khác nhau (heterogeneous), còn danh sách liên kết (Linked list) yêu cầu duyệt tuần tự qua các con trỏ chứ không hỗ trợ chỉ số trực tiếp với độ phức tạp O(1).

\par\vspace{0.3cm}
\hrule
\vspace{0.3cm}
\end{minipage}

\begin{minipage}{\textwidth}
\subsection*{Question 305 \quad \footnotesize{[0305\_q305\_Question\_305.png]}}
\noindent \textbf{\_pointer. To traverse a list, you need a}

\begin{enumerate}[label=\Alph*.]
    \item Null
    \item \textbf{Walking} \hfill \textcolor{blue}{\textbf{[Correct]}}
    \item Beginning
    \item Insertion
\end{enumerate}

\noindent \textbf{Correct Answer:} \textcolor{blue}{\textbf{B}} \\[2pt]
\noindent \textbf{Explanation:} Trong cấu trúc dữ liệu danh sách liên kết (Linked list) theo giáo trình Forouzan, để duyệt (traverse) qua tất cả các phần tử từ đầu đến cuối danh sách, ta cần sử dụng một con trỏ di chuyển gọi là con trỏ duyệt (Walking pointer hay Walker). Con trỏ này bắt đầu từ nút đầu tiên (head) và sau mỗi thao tác duyệt sẽ được cập nhật để trỏ đến nút kế tiếp (walker = walker->next) cho đến khi gặp con trỏ rỗng (Null pointer). Do đó, con trỏ cần dùng là Walking pointer.

\par\vspace{0.3cm}
\hrule
\vspace{0.3cm}
\end{minipage}

\begin{minipage}{\textwidth}
\subsection*{Question 306 \quad \footnotesize{[0306\_q306\_Question\_306.png]}}
\noindent \textbf{A is a node that can accommodate more than one record. The disadvantage of this method is that there may be a lot of wasted (unoccupied) locations.}

\begin{enumerate}[label=\Alph*.]
    \item Linked list resolution
    \item Open addressing resolution
    \item \textbf{Bucket hashing} \hfill \textcolor{blue}{\textbf{[Correct]}}
    \item Digit extraction hashing
\end{enumerate}

\noindent \textbf{Correct Answer:} \textcolor{blue}{\textbf{C}} \\[2pt]
\noindent \textbf{Explanation:} Băm theo khối/thùng (Bucket hashing) là một kỹ thuật giải quyết đụng độ (collision resolution) trong bảng băm bằng cách thiết kế mỗi địa chỉ băm như một thùng chứa (bucket) có khả năng lưu trữ được nhiều bản ghi cùng lúc. Tuy nhiên, nhược điểm lớn nhất của phương pháp này là gây lãng phí nhiều không gian bộ nhớ (wasted/unoccupied locations) khi dung lượng cố định của các thùng không được lấp đầy bởi các bản ghi thực tế. Ngược lại, phương pháp địa chỉ mở (Open addressing) hay danh sách liên kết có mức độ tận dụng bộ nhớ linh hoạt hơn.

\par\vspace{0.3cm}
\hrule
\vspace{0.3cm}
\end{minipage}

\begin{minipage}{\textwidth}
\subsection*{Question 307 \quad \footnotesize{[0307\_q307\_Question\_307.png]}}
\noindent \textbf{The union database operation is taking two relations and creating a new relation. What is a result of this operation?}

\begin{enumerate}[label=\Alph*.]
    \item Neither the first relation nor the second.
    \item \textbf{Both of the first relation and the second.} \hfill \textcolor{blue}{\textbf{[Correct]}}
    \item The first relation but not the second.
    \item The second but not the first.
\end{enumerate}

\noindent \textbf{Correct Answer:} \textcolor{blue}{\textbf{B}} \\[2pt]
\noindent \textbf{Explanation:} Phép hợp trong đại số quan hệ (Union operation: R ∪ S) lấy dữ liệu từ hai bảng quan hệ tương thích hợp để tạo ra một quan hệ mới chứa tất cả các bộ (tuples) xuất hiện ở quan hệ thứ nhất, quan hệ thứ hai hoặc ở cả hai quan hệ (với các bộ trùng lặp được loại bỏ). Vì kết quả thu được bao gồm các bản ghi đến từ cả bảng thứ nhất và bảng thứ hai (Both of the first relation and the second), nên đáp án B là hoàn toàn chính xác. Ngược lại, phép trừ (Difference) chỉ lấy các bộ thuộc bảng thứ nhất nhưng không thuộc bảng thứ hai.

\par\vspace{0.3cm}
\hrule
\vspace{0.3cm}
\end{minipage}


\newpage
\section*{9. Data Structures}
\addcontentsline{toc}{section}{9. Data Structures}
\hrule\vspace{0.5cm}

\begin{minipage}{\textwidth}
\subsection*{Question 308 \quad \footnotesize{[0308\_q308\_Question\_308.png]}}
\noindent \textbf{In the attacks threatent integrity is performed by one of the two parties in the communication: the sender or the receiver.}

\begin{enumerate}[label=\Alph*.]
    \item Snooping
    \item Modification
    \item \textbf{Repudiation} \hfill \textcolor{blue}{\textbf{[Correct]}}
    \item Masquerading
\end{enumerate}

\noindent \textbf{Correct Answer:} \textcolor{blue}{\textbf{C}} \\[2pt]
\noindent \textbf{Explanation:} Repudiation (sự chối bỏ) là hình thức tấn công đe dọa tính toàn vẹn (integrity), trong đó một trong hai bên tham gia truyền thông (người gửi hoặc người nhận) phủ nhận việc mình đã gửi hoặc nhận thông điệp. Điều này khác biệt so với các cuộc tấn công khác thường do bên thứ ba thực hiện. Ví dụ, Snooping đe dọa tính bảo mật (confidentiality), còn Modification và Masquerading là các hành vi sửa đổi hoặc giả mạo dữ liệu từ kẻ tấn công bên ngoài.

\par\vspace{0.3cm}
\hrule
\vspace{0.3cm}
\end{minipage}

\begin{minipage}{\textwidth}
\subsection*{Question 309 \quad \footnotesize{[0309\_q309\_Question\_309.png]}}
\noindent \textbf{In cryptography, the same key is used by the sender and the receiver.}

\begin{enumerate}[label=\Alph*.]
    \item \textbf{Symmetric-key} \hfill \textcolor{blue}{\textbf{[Correct]}}
    \item Asymmetric-key
    \item Public-key
    \item Open-key
\end{enumerate}

\noindent \textbf{Correct Answer:} \textcolor{blue}{\textbf{A}} \\[2pt]
\noindent \textbf{Explanation:} Trong mật mã khóa đối xứng (Symmetric-key cryptography), cả bên gửi và bên nhận đều sử dụng chung một khóa bí mật (shared secret key) để mã hóa và giải mã thông tin. Ngược lại, mật mã khóa bất đối xứng (Asymmetric-key hay Public-key) sử dụng một cặp khóa riêng biệt gồm khóa công khai và khóa riêng tư. Thuật ngữ Open-key không phải là một khái niệm chuẩn trong mật mã học.

\par\vspace{0.3cm}
\hrule
\vspace{0.3cm}
\end{minipage}

\begin{minipage}{\textwidth}
\subsection*{Question 310 \quad \footnotesize{[0310\_q310\_Question\_310.png]}}
\noindent \textbf{When Minicomputers appeared on the market?}

\begin{enumerate}[label=\Alph*.]
    \item First generation
    \item Second generation
    \item \textbf{Third generation} \hfill \textcolor{blue}{\textbf{[Correct]}}
    \item Fourth generation
\end{enumerate}

\noindent \textbf{Correct Answer:} \textcolor{blue}{\textbf{C}} \\[2pt]
\noindent \textbf{Explanation:} Máy tính mini (Minicomputer) xuất hiện trên thị trường vào thế hệ máy tính thứ ba (Third generation, thập niên 1960–1970) nhờ sự ra đời của công nghệ mạch tích hợp (Integrated Circuit - IC), giúp thu nhỏ đáng kể kích thước và giảm chi phí so với máy tính lớn (mainframe). Thế hệ thứ nhất sử dụng đèn chân không, thế hệ thứ hai sử dụng transistor, còn thế hệ thứ tư đặc trưng bởi bộ vi xử lý (VLSI) và máy tính cá nhân.

\par\vspace{0.3cm}
\hrule
\vspace{0.3cm}
\end{minipage}

\begin{minipage}{\textwidth}
\subsection*{Question 311 \quad \footnotesize{[0311\_q311\_Question\_311.png]}}
\noindent \textbf{A computer has 32 MB (megabytes) of memory. How many bits are needed to address any single byte in memory?}

\begin{enumerate}[label=\Alph*.]
    \item 5 bits
    \item 20 bits
    \item \textbf{25 bits} \hfill \textcolor{blue}{\textbf{[Correct]}}
    \item 32 bits
\end{enumerate}

\noindent \textbf{Correct Answer:} \textcolor{blue}{\textbf{C}} \\[2pt]
\noindent \textbf{Explanation:} Để đánh địa chỉ đến từng byte trong bộ nhớ dung lượng 32 MB, ta biểu diễn dung lượng theo lũy thừa của 2: 32 MB = 2\textasciicircum{}5 * 2\textasciicircum{}20 bytes = 2\textasciicircum{}25 bytes. Do đó, cần đúng 25 bit địa chỉ (k = 25) để định vị duy nhất từng byte ô nhớ vì 2\textasciicircum{}k = 2\textasciicircum{}25. Các lựa chọn khác như 20 bit chỉ đánh địa chỉ được 1 MB, còn 32 bit đánh địa chỉ được tới 4 GB.

\par\vspace{0.3cm}
\hrule
\vspace{0.3cm}
\end{minipage}

\begin{minipage}{\textwidth}
\subsection*{Question 312 \quad \footnotesize{[0312\_q312\_Question\_312.png]}}
\noindent \textbf{The device that can both feed data into and accept data from a computer is}

\begin{enumerate}[label=\Alph*.]
    \item ALU
    \item CPU
    \item \textbf{Input-Output device} \hfill \textcolor{blue}{\textbf{[Correct]}}
    \item All of the others
    \item None of the others
\end{enumerate}

\noindent \textbf{Correct Answer:} \textcolor{blue}{\textbf{C}} \\[2pt]
\noindent \textbf{Explanation:} Thiết bị Vào-Ra (Input-Output device) là thiết bị ngoại vi có khả năng thực hiện cả hai chức năng: đưa dữ liệu vào máy tính (input) và tiếp nhận dữ liệu xuất ra từ máy tính (output), ví dụ như màn hình cảm ứng hay ổ đĩa. Ngược lại, CPU và ALU là các đơn vị xử lý và tính toán trung tâm bên trong hệ thống máy tính, không phải là thiết bị ngoại vi vào-ra giao tiếp dữ liệu.

\par\vspace{0.3cm}
\hrule
\vspace{0.3cm}
\end{minipage}

\begin{minipage}{\textwidth}
\subsection*{Question 313 \quad \footnotesize{[0313\_q313\_Question\_313.png]}}
\noindent \textbf{How many connections (in the data bus) do we need to transfer 4 bytes of a word at the same time?}

\begin{enumerate}[label=\Alph*.]
    \item 4
    \item 8
    \item 16
    \item \textbf{32} \hfill \textcolor{blue}{\textbf{[Correct]}}
\end{enumerate}

\noindent \textbf{Correct Answer:} \textcolor{blue}{\textbf{D}} \\[2pt]
\noindent \textbf{Explanation:} Mỗi đường dây (connection) trong bus dữ liệu chịu trách nhiệm truyền tải 1 bit dữ liệu tại một thời điểm. Vì 1 byte bằng 8 bit nên một từ dữ liệu (word) gồm 4 byte tương đương với 4 * 8 = 32 bit. Do đó, để truyền đồng thời toàn bộ 4 byte này cùng lúc qua bus dữ liệu, hệ thống cần chính xác 32 đường kết nối vật lý.

\par\vspace{0.3cm}
\hrule
\vspace{0.3cm}
\end{minipage}

\begin{minipage}{\textwidth}
\subsection*{Question 314 \quad \footnotesize{[0314\_q314\_Question\_314.png]}}
\noindent \textbf{Which of the following is equivalent to 24 in decimal?}

\begin{enumerate}[label=\Alph*.]
    \item \textbf{(11000)2} \hfill \textcolor{blue}{\textbf{[Correct]}}
    \item (1A)16 UO
    \item (31)8
    \item None of the above
\end{enumerate}

\noindent \textbf{Correct Answer:} \textcolor{blue}{\textbf{A}} \\[2pt]
\noindent \textbf{Explanation:} Khi chuyển đổi số thập phân 24 sang hệ nhị phân (cơ số 2), ta phân tích 24 = 16 + 8 = 1*2\textasciicircum{}4 + 1*2\textasciicircum{}3 + 0*2\textasciicircum{}2 + 0*2\textasciicircum{}1 + 0*2\textasciicircum{}0 = (11000)2, do đó phương án A là chính xác. Trong khi đó, (1A)16 trong hệ thập lục phân tương đương với 1*16 + 10 = 26 trong hệ thập phân, và (31)8 trong hệ bát phân tương đương với 3*8 + 1 = 25 trong hệ thập phân.

\par\vspace{0.3cm}
\hrule
\vspace{0.3cm}
\end{minipage}

\begin{minipage}{\textwidth}
\subsection*{Question 315 \quad \footnotesize{[0315\_q315\_Question\_315.png]}}
\noindent \textbf{In an Excess conversion, we the number to be converted.}

\begin{enumerate}[label=\Alph*.]
    \item \textbf{Add the bias number to} \hfill \textcolor{blue}{\textbf{[Correct]}}
    \item Subtract the bias number from
    \item Multiply the bias number by
    \item Divide the bias number by
\end{enumerate}

\noindent \textbf{Correct Answer:} \textcolor{blue}{\textbf{A}} \\[2pt]
\noindent \textbf{Explanation:} Trong hệ biểu diễn số thừa (Excess system hay Biased representation), để biểu diễn một số nguyên, ta thực hiện cộng số độ lệch (bias number) vào số cần chuyển đổi theo công thức E = N + Bias. Phép cộng này giúp đưa tất cả các số âm và dương về phạm vi số không âm để lưu trữ dưới dạng nhị phân chuẩn. Ngược lại, thao tác trừ bias chỉ được thực hiện khi giải mã ngược từ biểu diễn Excess về lại số gốc.

\par\vspace{0.3cm}
\hrule
\vspace{0.3cm}
\end{minipage}

\begin{minipage}{\textwidth}
\subsection*{Question 316 \quad \footnotesize{[0316\_q316\_Question\_316.png]}}
\noindent \textbf{How many symbols can be represented by a bit pattern with ten bits?}

\begin{enumerate}[label=\Alph*.]
    \item 128
    \item 256
    \item 512
    \item \textbf{1024} \hfill \textcolor{blue}{\textbf{[Correct]}}
\end{enumerate}

\noindent \textbf{Correct Answer:} \textcolor{blue}{\textbf{D}} \\[2pt]
\noindent \textbf{Explanation:} Với một chuỗi bit có độ dài n bit, số lượng ký hiệu hoặc trạng thái khác nhau tối đa có thể biểu diễn được tính theo công thức 2\textasciicircum{}n. Với n = 10 bit, số ký hiệu có thể biểu diễn là 2\textasciicircum{}10 = 1024 ký hiệu. Các đáp án 128, 256 và 512 tương ứng với chuỗi có độ dài lần lượt là 7 bit, 8 bit và 9 bit.

\par\vspace{0.3cm}
\hrule
\vspace{0.3cm}
\end{minipage}

\begin{minipage}{\textwidth}
\subsection*{Question 317 \quad \footnotesize{[0317\_q317\_Question\_317.png]}}
\noindent \textbf{In the storing images, the number of bits used to represent a pixel, its, depends on how a pixel's color is handled by different encoding techniques.}

\begin{enumerate}[label=\Alph*.]
    \item Bit map
    \item Pixels
    \item Resolution
    \item \textbf{Color depth} \hfill \textcolor{blue}{\textbf{[Correct]}}
    \item None of others
\end{enumerate}

\noindent \textbf{Correct Answer:} \textcolor{blue}{\textbf{D}} \\[2pt]
\noindent \textbf{Explanation:} Độ sâu màu (Color depth, hay pixel depth) là số lượng bit được sử dụng để biểu diễn màu sắc của một điểm ảnh (pixel), quyết định số lượng màu sắc mà điểm ảnh đó có thể hiển thị. Ngược lại, độ phân giải (Resolution) là tổng số điểm ảnh hoặc mật độ điểm ảnh của bức ảnh, còn Bit map là định dạng biểu diễn ảnh theo lưới điểm ảnh.

\par\vspace{0.3cm}
\hrule
\vspace{0.3cm}
\end{minipage}

\begin{minipage}{\textwidth}
\subsection*{Question 318 \quad \footnotesize{[0318\_q318\_Question\_318.png]}}
\noindent \textbf{A defines the rules that both the sender and receiver and all intermediate devices need to follow to be able to communicate effectively.}

\begin{enumerate}[label=\Alph*.]
    \item \textbf{Protocol} \hfill \textcolor{blue}{\textbf{[Correct]}}
    \item Layers
    \item Provider networks
    \item Hardware and sofware
\end{enumerate}

\noindent \textbf{Correct Answer:} \textcolor{blue}{\textbf{A}} \\[2pt]
\noindent \textbf{Explanation:} Giao thức mạng (Protocol) là tập hợp các quy tắc chuẩn hóa mà bên gửi, bên nhận và mọi thiết bị trung gian đều phải tuân theo để đảm bảo quá trình truyền thông diễn ra chính xác và hiệu quả. Giao thức quy định cú pháp, ngữ nghĩa và thời gian truyền dữ liệu. Trong khi đó, tầng (Layers) chỉ là sự phân chia kiến trúc chức năng mạng, còn mạng nhà cung cấp và phần cứng/phần mềm không định nghĩa các quy tắc giao tiếp này.

\par\vspace{0.3cm}
\hrule
\vspace{0.3cm}
\end{minipage}

\begin{minipage}{\textwidth}
\subsection*{Question 319 \quad \footnotesize{[0319\_q319\_Question\_319.png]}}
\noindent \textbf{Communication by phone is indeed a activity; no party needs to be running forever waiting for the other party to cal}

\begin{enumerate}[label=\Alph*.]
    \item \textbf{P2P} \hfill \textcolor{blue}{\textbf{[Correct]}}
    \item Tranditional
    \item Client-server
    \item Model
\end{enumerate}

\noindent \textbf{Correct Answer:} \textcolor{blue}{\textbf{A}} \\[2pt]
\noindent \textbf{Explanation:} Giao tiếp qua điện thoại là hoạt động theo mô hình ngang hàng (P2P - Peer-to-Peer) vì hai bên tham gia có vai trò bình đẳng, kết nối trực tiếp khi có nhu cầu mà không bên nào phải chạy thường trực để chờ phục vụ. Ngược lại, trong mô hình Client-Server, máy chủ (server) bắt buộc phải luôn luôn hoạt động 24/7 để lắng nghe và xử lý các yêu cầu từ phía máy khách (client).

\par\vspace{0.3cm}
\hrule
\vspace{0.3cm}
\end{minipage}

\begin{minipage}{\textwidth}
\subsection*{Question 320 \quad \footnotesize{[0320\_q320\_Question\_320.png]}}
\noindent \textbf{The organization and interconnection of the various components of a computer system is}

\begin{enumerate}[label=\Alph*.]
    \item \textbf{Architecture} \hfill \textcolor{blue}{\textbf{[Correct]}}
    \item Networks
    \item Graphics
    \item Designing
    \item None of the others
\end{enumerate}

\noindent \textbf{Correct Answer:} \textcolor{blue}{\textbf{A}} \\[2pt]
\noindent \textbf{Explanation:} Kiến trúc máy tính (Architecture/Organization) là cấu trúc tổ chức và phương thức liên kết giữa các thành phần cơ bản của hệ thống máy tính như CPU, bộ nhớ, hệ thống bus và các thiết bị I/O. Khái niệm này xác định cách các khối phần cứng phối hợp hoạt động với nhau. Trong khi đó, mạng máy tính (Networks) kết nối nhiều hệ thống máy riêng biệt, còn đồ họa (Graphics) xử lý hình ảnh hiển thị.

\par\vspace{0.3cm}
\hrule
\vspace{0.3cm}
\end{minipage}

\begin{minipage}{\textwidth}
\subsection*{Question 321 \quad \footnotesize{[0321\_q321\_Question\_321.png]}}
\noindent \textbf{were designed in the 1950s to control mainframe computers. SÖ DÁP ÁN CÃN CHON}

\begin{enumerate}[label=\Alph*.]
    \item Time-sharing systems
    \item \textbf{Batch operating system} \hfill \textcolor{blue}{\textbf{[Correct]}}
    \item Personal systems
    \item Parallel systems
\end{enumerate}

\noindent \textbf{Correct Answer:} \textcolor{blue}{\textbf{B}} \\[2pt]
\noindent \textbf{Explanation:} Hệ điều hành theo lô (Batch operating system) ra đời vào những năm 1950 để quản lý và vận hành các máy tính lớn (mainframe) bằng cách gom nhóm các công việc tương tự nhau để xử lý tự động, liên tục mà không cần người dùng tương tác trực tiếp. Các hệ thống chia sẻ thời gian (Time-sharing) chỉ phát triển sau đó vào thập niên 1960, còn hệ thống máy tính cá nhân (Personal systems) xuất hiện từ cuối thập niên 1970.

\par\vspace{0.3cm}
\hrule
\vspace{0.3cm}
\end{minipage}

\begin{minipage}{\textwidth}
\subsection*{Question 322 \quad \footnotesize{[0322\_q322\_Question\_322.png]}}
\noindent \textbf{and storage is The difference between memory and storages is that memory is}

\begin{enumerate}[label=\Alph*.]
    \item \textbf{Temporary, permanent} \hfill \textcolor{blue}{\textbf{[Correct]}}
    \item Permanent, temporary
    \item Slow, fast
    \item All of the others
    \item None of the others
\end{enumerate}

\noindent \textbf{Correct Answer:} \textcolor{blue}{\textbf{A}} \\[2pt]
\noindent \textbf{Explanation:} Điểm khác biệt cơ bản giữa bộ nhớ chính (Memory - RAM) và thiết bị lưu trữ (Storage - HDD/SSD) là bộ nhớ mang tính tạm thời (temporary/volatile), nghĩa là dữ liệu sẽ mất đi khi ngắt nguồn điện, trong khi thiết bị lưu trữ mang tính lâu dài/vĩnh viễn (permanent/non-volatile) giúp giữ nguyên vẹn dữ liệu. Ngoài ra, bộ nhớ chính có tốc độ truy xuất nhanh hơn rất nhiều so với thiết bị lưu trữ.

\par\vspace{0.3cm}
\hrule
\vspace{0.3cm}
\end{minipage}

\begin{minipage}{\textwidth}
\subsection*{Question 323 \quad \footnotesize{[0323\_q323\_Question\_323.png]}}
\noindent \textbf{The scheduler moves a process from one process state to another.}

\begin{enumerate}[label=\Alph*.]
    \item Job
    \item \textbf{Process} \hfill \textcolor{blue}{\textbf{[Correct]}}
    \item Virtual
    \item Queue
\end{enumerate}

\noindent \textbf{Correct Answer:} \textcolor{blue}{\textbf{B}} \\[2pt]
\noindent \textbf{Explanation:} Bộ điều phối tiến trình (Process scheduler) trong hệ điều hành chịu trách nhiệm chuyển đổi một tiến trình từ trạng thái này sang trạng thái khác (giữa Ready, Running và Waiting). Trong khi đó, Job scheduler (bộ điều phối công việc) lựa chọn các chương trình từ bộ nhớ ngoài nạp vào hàng đợi công việc, chứ không trực tiếp quản lý việc chuyển trạng thái tiến trình thời gian thực.

\par\vspace{0.3cm}
\hrule
\vspace{0.3cm}
\end{minipage}

\begin{minipage}{\textwidth}
\subsection*{Question 324 \quad \footnotesize{[0324\_q324\_Question\_324.png]}}
\noindent \textbf{Codes consisting of light and dark marks which may be optically read is known as}

\begin{enumerate}[label=\Alph*.]
    \item Memonics
    \item \textbf{Bar code} \hfill \textcolor{blue}{\textbf{[Correct]}}
    \item Decoder
    \item All of the others
    \item None of the others
\end{enumerate}

\noindent \textbf{Correct Answer:} \textcolor{blue}{\textbf{B}} \\[2pt]
\noindent \textbf{Explanation:} Mã vạch (Bar code) là phương pháp biểu diễn thông tin bằng các vạch sáng và tối xen kẽ có độ rộng khác nhau để các thiết bị đọc quang học (optical reader) có thể nhận dạng và giải mã tự động. Ngược lại, Mnemonics là các ký hiệu gợi nhớ dùng trong hợp ngữ (assembly), còn bộ giải mã (Decoder) là mạch điện tử thực hiện chức năng chuyển đổi mã nhị phân.

\par\vspace{0.3cm}
\hrule
\vspace{0.3cm}
\end{minipage}

\begin{minipage}{\textwidth}
\subsection*{Question 325 \quad \footnotesize{[0325\_q325\_Question\_325.png]}}
\noindent \textbf{The language that is an input for statement translation is called}

\begin{enumerate}[label=\Alph*.]
    \item Assembly language
    \item \textbf{Source language} \hfill \textcolor{blue}{\textbf{[Correct]}}
    \item High-level language
    \item Object language
    \item None of the others
\end{enumerate}

\noindent \textbf{Correct Answer:} \textcolor{blue}{\textbf{B}} \\[2pt]
\noindent \textbf{Explanation:} Ngôn ngữ nguồn (Source language) là ngôn ngữ ban đầu được dùng làm đầu vào cho các trình biên dịch hoặc thông dịch để chuyển đổi câu lệnh sang ngôn ngữ khác. Kết quả đầu ra sau quá trình dịch được gọi là ngôn ngữ đối tượng hay ngôn ngữ đích (Object/Target language). Các khái niệm như High-level language hay Assembly language chỉ là các loại ngôn ngữ cụ thể chứ không khái quát vai trò đầu vào của bộ dịch.

\par\vspace{0.3cm}
\hrule
\vspace{0.3cm}
\end{minipage}

\begin{minipage}{\textwidth}
\subsection*{Question 326 \quad \footnotesize{[0326\_q326\_Question\_326.png]}}
\noindent \textbf{\_language. FORTRAN is a(n)}

\begin{enumerate}[label=\Alph*.]
    \item \textbf{Procedural} \hfill \textcolor{blue}{\textbf{[Correct]}}
    \item Functional
    \item Declarative
    \item Object-oriented
\end{enumerate}

\noindent \textbf{Correct Answer:} \textcolor{blue}{\textbf{A}} \\[2pt]
\noindent \textbf{Explanation:} FORTRAN (Formula Translation) là ngôn ngữ lập trình bậc cao thuộc mô hình thủ tục (Procedural paradigm), hoạt động dựa trên các hàm và quy trình thực thi lệnh tuần tự để giải quyết các bài toán khoa học kỹ thuật. Ngược lại, Lisp thuộc mô hình hàm (Functional), Prolog thuộc mô hình khai báo (Declarative), còn Java và C++ là các ngôn ngữ theo mô hình hướng đối tượng (Object-oriented).

\par\vspace{0.3cm}
\hrule
\vspace{0.3cm}
\end{minipage}

\begin{minipage}{\textwidth}
\subsection*{Question 327 \quad \footnotesize{[0327\_q327\_Question\_327.png]}}
\noindent \textbf{In the development process flows in only one direction. This means that a phase cannot be started until the previous phase is completed.}

\begin{enumerate}[label=\Alph*.]
    \item \textbf{Waterfall model} \hfill \textcolor{blue}{\textbf{[Correct]}}
    \item Incremental model
    \item Development proces
    \item Software engineering
\end{enumerate}

\noindent \textbf{Correct Answer:} \textcolor{blue}{\textbf{A}} \\[2pt]
\noindent \textbf{Explanation:} Mô hình Thác nước (Waterfall model) là mô hình phát triển phần mềm tuần tự tuyến tính, trong đó tiến trình phát triển chỉ chảy theo một chiều duy nhất và một giai đoạn chỉ được bắt đầu khi giai đoạn trước đó đã hoàn thành trọn vẹn. Trái lại, mô hình tăng trưởng (Incremental model) hay mô hình lặp (Iterative) cho phép phát triển từng phần và lặp lại các giai đoạn để tiếp nhận phản hồi.

\par\vspace{0.3cm}
\hrule
\vspace{0.3cm}
\end{minipage}

\begin{minipage}{\textwidth}
\subsection*{Question 328 \quad \footnotesize{[0328\_q328\_Question\_328.png]}}
\noindent \textbf{is the smallest element of named data that has meaning in a record.}

\begin{enumerate}[label=\Alph*.]
    \item Name of record
    \item Pointer
    \item \textbf{A field} \hfill \textcolor{blue}{\textbf{[Correct]}}
    \item Index
\end{enumerate}

\noindent \textbf{Correct Answer:} \textcolor{blue}{\textbf{C}} \\[2pt]
\noindent \textbf{Explanation:} Một trường dữ liệu (A field) là phần tử dữ liệu nhỏ nhất có tên gọi và mang ý nghĩa trọn vẹn bên trong một bản ghi (record), chẳng hạn như trường tên, tuổi hay mã số sinh viên. Tập hợp các trường liên quan sẽ cấu thành nên một bản ghi. Trong khi đó, con trỏ (Pointer) lưu địa chỉ vùng nhớ, còn chỉ mục (Index) là cấu trúc phụ trợ dùng để định vị bản ghi nhanh chóng.

\par\vspace{0.3cm}
\hrule
\vspace{0.3cm}
\end{minipage}

\begin{minipage}{\textwidth}
\subsection*{Question 329 \quad \footnotesize{[0329\_q329\_Question\_329.png]}}
\noindent \textbf{An array is normally a and so is more suitable when the number of data items is fixed}

\begin{enumerate}[label=\Alph*.]
    \item Dynamic data structure
    \item \textbf{Static data structure} \hfill \textcolor{blue}{\textbf{[Correct]}}
    \item Same data
    \item Related data
\end{enumerate}

\noindent \textbf{Correct Answer:} \textcolor{blue}{\textbf{B}} \\[2pt]
\noindent \textbf{Explanation:} Mảng (Array) thông thường là một cấu trúc dữ liệu tĩnh (Static data structure) vì số lượng phần tử và không gian bộ nhớ của nó được cấp phát cố định ngay từ đầu, nên rất phù hợp khi số lượng mục dữ liệu là cố định và đã biết trước. Ngược lại, danh sách liên kết (Linked list) là cấu trúc dữ liệu động (Dynamic data structure) cho phép thêm bớt phần tử linh hoạt trong thời gian chạy.

\par\vspace{0.3cm}
\hrule
\vspace{0.3cm}
\end{minipage}

\begin{minipage}{\textwidth}
\subsection*{Question 330 \quad \footnotesize{[0330\_q330\_Question\_330.png]}}
\noindent \textbf{\_is the first If A is the first data element input into a queue, followed by B, C, and D, then element to be removed.}

\begin{enumerate}[label=\Alph*.]
    \item \textbf{A} \hfill \textcolor{blue}{\textbf{[Correct]}}
    \item B
    \item C
    \item D
\end{enumerate}

\noindent \textbf{Correct Answer:} \textcolor{blue}{\textbf{A}} \\[2pt]
\noindent \textbf{Explanation:} Hàng đợi (Queue) hoạt động theo nguyên lý Vào trước - Ra trước (FIFO: First-In, First-Out), nghĩa là phần tử nào được đưa vào hàng đợi trước thì sẽ được lấy ra trước. Vì A là phần tử đầu tiên được đưa vào hàng đợi trước B, C và D nên A sẽ là phần tử đầu tiên được lấy ra khỏi hàng đợi. Nếu là cấu trúc Ngăn xếp (Stack - LIFO) thì phần tử D vào cuối cùng mới là phần tử được lấy ra đầu tiên.

\par\vspace{0.3cm}
\hrule
\vspace{0.3cm}
\end{minipage}

\begin{minipage}{\textwidth}
\subsection*{Question 331 \quad \footnotesize{[0331\_q331\_Question\_331.png]}}
\noindent \textbf{The dequeue operationthe item at theof the queue}

\begin{enumerate}[label=\Alph*.]
    \item \textbf{delete/front} \hfill \textcolor{blue}{\textbf{[Correct]}}
    \item delete/end
    \item inserts/front
    \item insert/rear
\end{enumerate}

\noindent \textbf{Correct Answer:} \textcolor{blue}{\textbf{A}} \\[2pt]
\noindent \textbf{Explanation:} Thao tác Dequeue trong cấu trúc hàng đợi (Queue) có chức năng xóa hoặc lấy phần tử ra khỏi vị trí đầu (front) của hàng đợi. Ngược lại, thao tác thêm mới một phần tử vào hàng đợi được gọi là Enqueue và luôn được thực hiện ở vị trí cuối (rear/end) của hàng đợi.

\par\vspace{0.3cm}
\hrule
\vspace{0.3cm}
\end{minipage}

\begin{minipage}{\textwidth}
\subsection*{Question 332 \quad \footnotesize{[0332\_q332\_Question\_332.png]}}
\noindent \textbf{The file structure allows random access is}

\begin{enumerate}[label=\Alph*.]
    \item \textbf{Hashed file} \hfill \textcolor{blue}{\textbf{[Correct]}}
    \item Sequential file
    \item EOF
    \item None of others
\end{enumerate}

\noindent \textbf{Correct Answer:} \textcolor{blue}{\textbf{A}} \\[2pt]
\noindent \textbf{Explanation:} Cấu trúc tệp băm (Hashed file) cho phép truy xuất ngẫu nhiên (random access hay direct access) đến từng bản ghi bằng cách áp dụng hàm băm (hash function) trực tiếp lên khóa dữ liệu để tìm ra địa chỉ lưu trữ. Ngược lại, tệp tuần tự (Sequential file) bắt buộc phải đọc lần lượt từ đầu tệp, còn EOF (End of File) chỉ là ký hiệu đánh dấu kết thúc tệp dữ liệu.

\par\vspace{0.3cm}
\hrule
\vspace{0.3cm}
\end{minipage}

\begin{minipage}{\textwidth}
\subsection*{Question 333 \quad \footnotesize{[0333\_q333\_Question\_333.png]}}
\noindent \textbf{file can be accessed sequentially.}

\begin{enumerate}[label=\Alph*.]
    \item \textbf{A sequential} \hfill \textcolor{blue}{\textbf{[Correct]}}
    \item An indexed
    \item A hashed
    \item No
\end{enumerate}

\noindent \textbf{Correct Answer:} \textcolor{blue}{\textbf{A}} \\[2pt]
\noindent \textbf{Explanation:} Tệp tuần tự (A sequential file) là loại tệp mà các bản ghi được lưu trữ nối tiếp nhau theo một trật tự cố định và phương thức truy xuất bản chất của nó là truy xuất tuần tự từ đầu đến cuối tệp. Trong khi đó, tệp băm (Hashed file) và tệp chỉ mục (Indexed file) được thiết kế chuyên biệt để hỗ trợ truy xuất ngẫu nhiên và trực tiếp tới bản ghi cần tìm.

\par\vspace{0.3cm}
\hrule
\vspace{0.3cm}
\end{minipage}

\begin{minipage}{\textwidth}
\subsection*{Question 334 \quad \footnotesize{[0334\_q334\_Question\_334.png]}}
\noindent \textbf{In a three-level DBMS architecture, the level determines where data is actually stored on the storage devices.}

\begin{enumerate}[label=\Alph*.]
    \item External
    \item Conceptual
    \item \textbf{Internal} \hfill \textcolor{blue}{\textbf{[Correct]}}
    \item Physical
\end{enumerate}

\noindent \textbf{Correct Answer:} \textcolor{blue}{\textbf{C}} \\[2pt]
\noindent \textbf{Explanation:} Trong kiến trúc 3 mức chuẩn của DBMS (ANSI/SPARC), mức trong (Internal level) xác định cấu trúc lưu trữ dữ liệu thực tế trên các thiết bị nhớ vật lý và cách tổ chức các đường dẫn truy xuất. Trong khi đó, mức ngoài (External level) định nghĩa các góc nhìn dữ liệu cho từng nhóm người dùng, còn mức quan niệm (Conceptual level) mô tả cấu trúc logic tổng thể của toàn bộ cơ sở dữ liệu.

\par\vspace{0.3cm}
\hrule
\vspace{0.3cm}
\end{minipage}

\begin{minipage}{\textwidth}
\subsection*{Question 335 \quad \footnotesize{[0335\_q335\_Question\_335.png]}}
\noindent \textbf{can be threatened by security attacks.}

\begin{enumerate}[label=\Alph*.]
    \item Confidentiality, cryptography, nonrepudiation
    \item Confidentiality, encryption, decryption
    \item \textbf{Confidentiality, integrity, availability} \hfill \textcolor{blue}{\textbf{[Correct]}}
    \item Confidentiality, denial of service, masquerading
\end{enumerate}

\noindent \textbf{Correct Answer:} \textcolor{blue}{\textbf{C}} \\[2pt]
\noindent \textbf{Explanation:} Ba mục tiêu bảo mật thông tin cốt lõi thường xuyên bị đe dọa bởi các cuộc tấn công mạng là Tam giác CIA: Tính bảo mật (Confidentiality), Tính toàn vẹn (Integrity) và Tính sẵn sàng (Availability). Các thuật ngữ như cryptography, encryption hay decryption là các công cụ và kỹ thuật dùng để phòng vệ an ninh, trong khi denial of service hay masquerading là tên các hình thức tấn công cụ thể.

\par\vspace{0.3cm}
\hrule
\vspace{0.3cm}
\end{minipage}

\begin{minipage}{\textwidth}
\subsection*{Question 336 \quad \footnotesize{[0336\_q336\_Question\_336.png]}}
\noindent \textbf{means that the data must arrive at the receiver exactly as they were sent.}

\begin{enumerate}[label=\Alph*.]
    \item Nonrepudiation
    \item \textbf{Message integrity} \hfill \textcolor{blue}{\textbf{[Correct]}}
    \item Authentication
    \item Secrecy
\end{enumerate}

\noindent \textbf{Correct Answer:} \textcolor{blue}{\textbf{B}} \\[2pt]
\noindent \textbf{Explanation:} Tính toàn vẹn thông điệp (Message integrity) đảm bảo rằng dữ liệu khi truyền đến bên nhận phải hoàn toàn nguyên vẹn, chính xác như khi bên gửi phát đi, không bị sửa đổi, xóa bỏ hay chèn thêm bất hợp pháp. Ngược lại, Secrecy (tính bảo mật) ngăn chặn kẻ xấu đọc trộm dữ liệu, Authentication giúp xác minh danh tính người gửi, và Nonrepudiation ngăn việc chối bỏ trách nhiệm gửi thông điệp.

\par\vspace{0.3cm}
\hrule
\vspace{0.3cm}
\end{minipage}

\begin{minipage}{\textwidth}
\subsection*{Question 337 \quad \footnotesize{[0337\_q337\_Question\_337.png]}}
\noindent \textbf{The attack that reduces the capability of a computer is called a \_attack. SÖ DÁP ÁN CÀN CHON}

\begin{enumerate}[label=\Alph*.]
    \item Penetration
    \item \textbf{Denial of service} \hfill \textcolor{blue}{\textbf{[Correct]}}
    \item Either a or b FUO)
    \item Neither a nor b
\end{enumerate}

\noindent \textbf{Correct Answer:} \textcolor{blue}{\textbf{B}} \\[2pt]
\noindent \textbf{Explanation:} Tấn công từ chối dịch vụ (Denial of service - DoS) là dạng tấn công làm suy giảm nghiêm trọng khả năng xử lý hoặc làm sập hệ thống bằng cách làm tràn ngập tài nguyên của máy tính, khiến người dùng hợp lệ không thể tiếp cận dịch vụ. Cuộc tấn công này trực tiếp đe dọa đến tính sẵn sàng (Availability) của hệ thống máy tính. Ngược lại, Penetration là hành vi đột nhập trái phép vào hệ thống để đánh cắp thông tin hoặc chiếm quyền điều khiển.

\par\vspace{0.3cm}
\hrule
\vspace{0.3cm}
\end{minipage}

\begin{minipage}{\textwidth}
\subsection*{Question 338 \quad \footnotesize{[0338\_q338\_Question\_338.png]}}
\noindent \textbf{In cryptography symbols are permuted or substituted; in asymmetric-key cryptography numbers are manipulated}

\begin{enumerate}[label=\Alph*.]
    \item \textbf{Symmetric-key} \hfill \textcolor{blue}{\textbf{[Correct]}}
    \item Asymmetric-key
    \item Private-key
\end{enumerate}

\noindent \textbf{Correct Answer:} \textcolor{blue}{\textbf{A}} \\[2pt]
\noindent \textbf{Explanation:} Trong mật mã khóa đối xứng (Symmetric-key cryptography), các thuật toán cổ điển sử dụng phép hoán vị (permutation) và phép thay thế (substitution) trên các ký tự hoặc khối bit để xáo trộn bản rõ. Ngược lại, mật mã khóa bất đối xứng (Asymmetric-key cryptography) hiện đại hoạt động dựa trên các phép toán số học phức tạp như số nguyên tố lớn và phép đồng dư lũy thừa. Vì vậy, Symmetric-key là đáp án chính xác.

\par\vspace{0.3cm}
\hrule
\vspace{0.3cm}
\end{minipage}

\begin{minipage}{\textwidth}
\subsection*{Question 339 \quad \footnotesize{[0339\_q339\_Question\_339.png]}}
\noindent \textbf{What is a major difference between procedural and object-oriented programming?}

\begin{enumerate}[label=\Alph*.]
    \item Procedural programming uses functions, while object-oriented programming uses loops.
    \item \textbf{Object-oriented programming organizes code into objects, while procedural programming uses a linear sequence of instructions.} \hfill \textcolor{blue}{\textbf{[Correct]}}
    \item Procedural programming supports inheritance, while object-oriented programming does not.
    \item Object-oriented programming does not allow code reuse, while procedural programming does.
\end{enumerate}

\noindent \textbf{Correct Answer:} \textcolor{blue}{\textbf{B}} \\[2pt]
\noindent \textbf{Explanation:} Khác biệt căn bản giữa hai phương pháp là lập trình hướng đối tượng (OOP) tổ chức chương trình thành các đối tượng chứa cả dữ liệu và hành vi liên quan, trong khi lập trình thủ tục (Procedural) tổ chức mã thành chuỗi các lệnh và hàm thực thi tuyến tính. Các phương án khác đều sai vì OOP hỗ trợ rất mạnh tính kế thừa (inheritance) và tái sử dụng mã (code reuse), đồng thời cả hai mô hình đều sử dụng cấu trúc vòng lặp và hàm.

\par\vspace{0.3cm}
\hrule
\vspace{0.3cm}
\end{minipage}

\begin{minipage}{\textwidth}
\subsection*{Question 340 \quad \footnotesize{[0340\_q340\_Question\_340.png]}}
\noindent \textbf{In which phase of the software lifecycle is the problem domain analyzed?}

\begin{enumerate}[label=\Alph*.]
    \item Design
    \item Testing
    \item \textbf{Analysis} \hfill \textcolor{blue}{\textbf{[Correct]}}
    \item Implementation
\end{enumerate}

\noindent \textbf{Correct Answer:} \textcolor{blue}{\textbf{C}} \\[2pt]
\noindent \textbf{Explanation:} Giai đoạn Phân tích (Analysis phase) trong vòng đời phát triển phần mềm (SDLC) là giai đoạn mà các chuyên viên tập trung tìm hiểu, khảo sát miền bài toán (problem domain) và xác định rõ các yêu cầu của hệ thống cần xây dựng. Sau đó, giai đoạn Thiết kế (Design) mới lên kiến trúc giải pháp kỹ thuật, giai đoạn Cài đặt (Implementation) thực hiện viết mã nguồn, và giai đoạn Kiểm thử (Testing) đánh giá chất lượng phần mềm.

\par\vspace{0.3cm}
\hrule
\vspace{0.3cm}
\end{minipage}

\begin{minipage}{\textwidth}
\subsection*{Question 341 \quad \footnotesize{[0341\_q341\_Question\_341.png]}}
\noindent \textbf{Which of the following software development methodologies emphasizes iterative development, collaboration among cross-functional teams, and frequent customer feedback?}

\begin{enumerate}[label=\Alph*.]
    \item Waterfall model
    \item \textbf{Agile methodology} \hfill \textcolor{blue}{\textbf{[Correct]}}
    \item Spiral model
    \item RAD (Rapid Application Development)
\end{enumerate}

\noindent \textbf{Correct Answer:} \textcolor{blue}{\textbf{B}} \\[2pt]
\noindent \textbf{Explanation:} Phương pháp luận Agile (Agile methodology) đặc trưng bởi quy trình phát triển lặp lại qua các chu kỳ ngắn, sự phối hợp chặt chẽ giữa các thành viên liên chức năng và liên tục tiếp nhận phản hồi từ khách hàng để điều chỉnh sản phẩm linh hoạt. Trái lại, mô hình Thác nước (Waterfall) phát triển theo trình tự cứng nhắc và khó thay đổi, còn mô hình Xoắn ốc (Spiral) tập trung chủ yếu vào việc phân tích và kiểm soát rủi ro.

\par\vspace{0.3cm}
\hrule
\vspace{0.3cm}
\end{minipage}

\begin{minipage}{\textwidth}
\subsection*{Question 342 \quad \footnotesize{[0342\_q342\_Question\_342.png]}}
\noindent \textbf{What is a key difference between glass-box testing and black-box testing?}

\begin{enumerate}[label=\Alph*.]
    \item \textbf{Glass-box testing focuses on the internal logic of the software, while black-box testing focuses on the external functionality} \hfill \textcolor{blue}{\textbf{[Correct]}}
    \item Glass-box testing does not require any knowledge of the internal code, while black-box testing does
    \item Black-box testing is performed by developers, while glass-box testing is performed by users.
    \item There is no significant difference between glass-box testing and black-box testing.
\end{enumerate}

\noindent \textbf{Correct Answer:} \textcolor{blue}{\textbf{A}} \\[2pt]
\noindent \textbf{Explanation:} Kiểm thử hộp kính (Glass-box testing hay White-box testing) tập trung kiểm tra cấu trúc mã nguồn và luồng logic bên trong của phần mềm, trong khi kiểm thử hộp đen (Black-box testing) chỉ đánh giá các chức năng bên ngoài dựa trên đầu vào và đầu ra mà không cần biết mã nguồn nội bộ. Phương án B nói ngược sự thật, còn phương án C không đúng vì người dùng thông thường không thực hiện kiểm thử hộp kính.

\par\vspace{0.3cm}
\hrule
\vspace{0.3cm}
\end{minipage}

\begin{minipage}{\textwidth}
\subsection*{Question 343 \quad \footnotesize{[0343\_q343\_Question\_343.png]}}
\noindent \textbf{Why is documentation considered important in software engineering?}

\begin{enumerate}[label=\Alph*.]
    \item It makes the software look more professional.
    \item \textbf{It helps ensure that the software can be maintained, understood, and used effectively by different stakeholders} \hfill \textcolor{blue}{\textbf{[Correct]}}
    \item It is required by law in most countries.
    \item It reduces the time needed for coding.
\end{enumerate}

\noindent \textbf{Correct Answer:} \textcolor{blue}{\textbf{B}} \\[2pt]
\noindent \textbf{Explanation:} Tài liệu phần mềm (Documentation) có vai trò sống còn trong kỹ thuật phần mềm nhằm giúp các bên liên quan (lập trình viên, người bảo trì và người dùng cuối) hiểu rõ cấu trúc hệ thống, từ đó đảm bảo phần mềm có thể được sử dụng, bảo trì và mở rộng dễ dàng trong tương lai. Mục đích cốt lõi của tài liệu là duy trì khả năng bảo trì và hiểu mã, chứ không phải vì lý do hình thức hay giúp giảm thời gian viết mã ban đầu.

\par\vspace{0.3cm}
\hrule
\vspace{0.3cm}
\end{minipage}

\begin{minipage}{\textwidth}
\subsection*{Question 344 \quad \footnotesize{[0344\_q344\_Question\_344.png]}}
\noindent \textbf{In a relational database, what operations can we define to create new relationships?}

\begin{enumerate}[label=\Alph*.]
    \item Insert, delete, update, select, join, union, intersection, and difference
    \item Insert, delete, update, select, project, union, intersection, and except
    \item \textbf{Insert, delete, update, select, project, join, union, intersection, and difference} \hfill \textcolor{blue}{\textbf{[Correct]}}
    \item Insert, delete, update, select, project, join, union, intersection, and except
    \item None of the others
\end{enumerate}

\noindent \textbf{Correct Answer:} \textcolor{blue}{\textbf{C}} \\[2pt]
\noindent \textbf{Explanation:} Theo giáo trình CSI106 của Behrouz Forouzan, các thao tác định nghĩa trên cơ sở dữ liệu quan hệ gồm các phép cập nhật dữ liệu (insert, delete, update) và các phép toán đại số quan hệ để tạo quan hệ mới: select (chọn), project (chiếu), join (kết nối), union (hợp), intersection (giao), và difference (hiệu). Phương án A bị thiếu phép chiếu (project), trong khi các phương án B và D sử dụng từ khóa except thay cho difference theo thuật ngữ chuẩn của giáo trình.

\par\vspace{0.3cm}
\hrule
\vspace{0.3cm}
\end{minipage}

\begin{minipage}{\textwidth}
\subsection*{Question 345 \quad \footnotesize{[0345\_q345\_Question\_345.png]}}
\noindent \textbf{Which of the following is NOT a component of a database management system (DBMS)?}

\begin{enumerate}[label=\Alph*.]
    \item Hardware
    \item Software
    \item \textbf{Network} \hfill \textcolor{blue}{\textbf{[Correct]}}
    \item Data
\end{enumerate}

\noindent \textbf{Correct Answer:} \textcolor{blue}{\textbf{C}} \\[2pt]
\noindent \textbf{Explanation:} Theo mô hình chuẩn của hệ quản trị cơ sở dữ liệu (DBMS), các thành phần cốt lõi cấu thành bao gồm phần cứng (Hardware), phần mềm (Software), dữ liệu (Data) và con người/quy trình. Mạng máy tính (Network) chỉ là hạ tầng liên lạc viễn thông hỗ trợ kết nối và truyền tải từ xa, không phải là một thành phần nội tại cấu tạo nên bản thân hệ DBMS.

\par\vspace{0.3cm}
\hrule
\vspace{0.3cm}
\end{minipage}


\newpage
\section*{10. File Structure}
\addcontentsline{toc}{section}{10. File Structure}
\hrule\vspace{0.5cm}

\begin{minipage}{\textwidth}
\subsection*{Question 346 \quad \footnotesize{[0346\_q346\_Question\_346.png]}}
\noindent \textbf{What is the primary key in a relational database table?}

\begin{enumerate}[label=\Alph*.]
    \item A column that can store null values
    \item \textbf{A column or combination of columns that uniquely identify each row in the table} \hfill \textcolor{blue}{\textbf{[Correct]}}
    \item A column used for indexing
    \item A foreign key reference in another table
\end{enumerate}

\noindent \textbf{Correct Answer:} \textcolor{blue}{\textbf{B}} \\[2pt]
\noindent \textbf{Explanation:} Trong mô hình cơ sở dữ liệu quan hệ, khóa chính (primary key) là một cột hoặc tập hợp các cột được dùng để định danh duy nhất (uniquely identify) mỗi hàng (bản ghi) trong bảng dữ liệu (đáp án B). Theo ràng buộc toàn vẹn thực thể, các giá trị của khóa chính bắt buộc phải là duy nhất và không được phép chứa giá trị rỗng (NULL), do đó đáp án A sai. Đáp án D mô tả khái niệm khóa ngoại (foreign key), còn đáp án C không chính xác vì chỉ số (index) có thể áp dụng cho bất kỳ cột nào chứ không chỉ riêng khóa chính.

\par\vspace{0.3cm}
\hrule
\vspace{0.3cm}
\end{minipage}

\begin{minipage}{\textwidth}
\subsection*{Question 347 \quad \footnotesize{[0347\_q347\_Question\_347.png]}}
\noindent \textbf{Which is true about organizing relational data storage? SÖ DÁP ÁN CÃN CHON}

\begin{enumerate}[label=\Alph*.]
    \item \textbf{The physical storage of the data is independent of the way in which the data is logically organized} \hfill \textcolor{blue}{\textbf{[Correct]}}
    \item The abstract storage of the data is independent of the way in which the data is logically organized
    \item The logical storage of the data is independent of the way in which the data is logically organized
    \item The virtual storage of the data is independent of the way in which the data is logically organized
\end{enumerate}

\noindent \textbf{Correct Answer:} \textcolor{blue}{\textbf{A}} \\[2pt]
\noindent \textbf{Explanation:} Trong hệ quản trị cơ sở dữ liệu quan hệ, nguyên lý độc lập dữ liệu vật lý (physical data independence) khẳng định rằng cấu trúc lưu trữ vật lý của dữ liệu trên bộ nhớ phụ hoàn toàn độc lập với cách dữ liệu được tổ chức dưới góc nhìn logic (bảng, hàng, cột) (đáp án A). Điều này cho phép người quản trị thay đổi thiết bị lưu trữ hoặc phương thức tổ chức tệp vật lý mà không làm ảnh hưởng đến các ứng dụng hay cách nhìn logic của người dùng. Các phương án B, C và D sử dụng các thuật ngữ sai lệch so với nguyên lý độc lập dữ liệu chuẩn.

\par\vspace{0.3cm}
\hrule
\vspace{0.3cm}
\end{minipage}

\begin{minipage}{\textwidth}
\subsection*{Question 348 \quad \footnotesize{[0348\_q348\_Question\_348.png]}}
\noindent \textbf{Which of the following examples demonstrates the "repetition" construct in an algorithm?}

\begin{enumerate}[label=\Alph*.]
    \item Assigning a value to a variable
    \item Evaluating a condition once
    \item \textbf{Using a for loop to iterate over a list} \hfill \textcolor{blue}{\textbf{[Correct]}}
    \item Breaking out of a loop
\end{enumerate}

\noindent \textbf{Correct Answer:} \textcolor{blue}{\textbf{C}} \\[2pt]
\noindent \textbf{Explanation:} Trong thiết kế thuật toán, cấu trúc lặp (repetition hay loop construct) là cấu trúc cho phép thực hiện lặp đi lặp lại một khối lệnh nhiều lần dựa trên một điều kiện hoặc duyệt qua các phần tử. Việc sử dụng vòng lặp for để duyệt tuần tự qua một danh sách là ví dụ minh họa trực tiếp cho cấu trúc lặp này (đáp án C). Trong khi đó, việc gán giá trị (A) thuộc cấu trúc tuần tự (sequence), đánh giá điều kiện một lần (B) thuộc cấu trúc rẽ nhánh (selection), còn thoát khỏi vòng lặp (D) chỉ là câu lệnh điều khiển luồng (break).

\par\vspace{0.3cm}
\hrule
\vspace{0.3cm}
\end{minipage}

\begin{minipage}{\textwidth}
\subsection*{Question 349 \quad \footnotesize{[0349\_q349\_Question\_349.png]}}
\noindent \textbf{Which of the following statements is true regarding search algorithms for a list of elements?}

\begin{enumerate}[label=\Alph*.]
    \item \textbf{Binary search is more efficient than sequential search for large lists, but it requires the list to be sorted.} \hfill \textcolor{blue}{\textbf{[Correct]}}
    \item Binary search can be used on an unsorted list to improve search efficiency compared to sequential search.
    \item Sequential search requires fewer comparisons than binary search in the worst case, regardless of whether the list is sorted or not.
    \item Sequential search and binary search require the same number of comparisons in the worst case if the list is sorted.
\end{enumerate}

\noindent \textbf{Correct Answer:} \textcolor{blue}{\textbf{A}} \\[2pt]
\noindent \textbf{Explanation:} Thuật toán tìm kiếm nhị phân (binary search) có độ phức tạp thời gian là O(log n), mang lại hiệu quả vượt trội so với tìm kiếm tuần tự (O(n)) trên các tập danh sách lớn, nhưng yêu cầu tiên quyết là danh sách dữ liệu bắt buộc phải được sắp xếp theo thứ tự (đáp án A). Đáp án B sai vì tìm kiếm nhị phân không thể áp dụng trên danh sách chưa sắp xếp; đáp án C và D sai vì trong trường hợp xấu nhất, tìm kiếm tuần tự đòi hỏi n phép so sánh, lớn hơn rất nhiều so với log2(n) phép so sánh của tìm kiếm nhị phân.

\par\vspace{0.3cm}
\hrule
\vspace{0.3cm}
\end{minipage}

\begin{minipage}{\textwidth}
\subsection*{Question 350 \quad \footnotesize{[0350\_q350\_Question\_350.png]}}
\noindent \textbf{What is an algorithm?}

\begin{enumerate}[label=\Alph*.]
    \item A set of programming commands or instructions that are executed in a specific order to perform a task
    \item \textbf{A step-by-step procedure to solve a problem or accomplish a task} \hfill \textcolor{blue}{\textbf{[Correct]}}
    \item A program that is loaded into the computer's memory and is ready for execution
    \item A method for compressing data to reduce its size for efficient storage or transmission
\end{enumerate}

\noindent \textbf{Correct Answer:} \textcolor{blue}{\textbf{B}} \\[2pt]
\noindent \textbf{Explanation:} Theo định nghĩa chuẩn trong khoa học máy tính (Forouzan), một thuật toán (algorithm) là một quy trình từng bước rõ ràng, có thứ tự logic nhằm giải quyết một bài toán cụ thể hoặc hoàn thành một nhiệm vụ định trước (đáp án B). Thuật toán tồn tại độc lập với bất kỳ ngôn ngữ lập trình cụ thể nào, nên nó không đồng nhất với một tập hợp mã lệnh chương trình cụ thể (đáp án A) hay một chương trình đã nạp vào bộ nhớ (đáp án C). Phương án D chỉ là định nghĩa của phương pháp nén dữ liệu.

\par\vspace{0.3cm}
\hrule
\vspace{0.3cm}
\end{minipage}

\begin{minipage}{\textwidth}
\subsection*{Question 351 \quad \footnotesize{[0351\_q351\_Question\_351.png]}}
\noindent \textbf{In which scenario is a binary search algorithm ineffective?}

\begin{enumerate}[label=\Alph*.]
    \item When the data is sorted
    \item \textbf{When the data is unsorted} \hfill \textcolor{blue}{\textbf{[Correct]}}
    \item When the data is in an array
    \item When the data is in a linked list
\end{enumerate}

\noindent \textbf{Correct Answer:} \textcolor{blue}{\textbf{B}} \\[2pt]
\noindent \textbf{Explanation:} Thuật toán tìm kiếm nhị phân dựa trên nguyên lý chia để trị, liên tục so sánh giá trị cần tìm với phần tử nằm ở giữa để loại bỏ một nửa không gian tìm kiếm, và nguyên lý này chỉ hoạt động chính xác khi dữ liệu đã được sắp xếp. Do đó, khi dữ liệu chưa được sắp xếp (unsorted), tìm kiếm nhị phân sẽ hoàn toàn vô hiệu và không thể tìm ra kết quả chính xác (đáp án B). Trong trường hợp dữ liệu chưa sắp xếp, phương pháp tìm kiếm duy nhất khả thi là tìm kiếm tuần tự (sequential search).

\par\vspace{0.3cm}
\hrule
\vspace{0.3cm}
\end{minipage}

\begin{minipage}{\textwidth}
\subsection*{Question 352 \quad \footnotesize{[0352\_q352\_Question\_352.png]}}
\noindent \textbf{Which of the following UML diagram types is most suitable for representing the flow of control within an algorithm, including loops and conditional statements? SÖ DÁP ÁN CÀN CHON}

\begin{enumerate}[label=\Alph*.]
    \item Class diagram
    \item Sequence diagram
    \item \textbf{Activity diagram} \hfill \textcolor{blue}{\textbf{[Correct]}}
    \item State diagram (FUO)
\end{enumerate}

\noindent \textbf{Correct Answer:} \textcolor{blue}{\textbf{C}} \\[2pt]
\noindent \textbf{Explanation:} Trong ngôn ngữ mô hình hóa thống nhất (UML), biểu đồ hoạt động (activity diagram) là dạng biểu đồ hành vi chuyên được sử dụng để trực quan hóa luồng điều khiển của một thuật toán hay quy trình nghiệp vụ, thể hiện rõ các bước tuần tự, điểm rẽ nhánh điều kiện và các vòng lặp (đáp án C). Trong khi đó, biểu đồ lớp (A) mô tả cấu trúc tĩnh của hệ thống, biểu đồ tuần tự (B) biểu diễn sự tương tác thông điệp giữa các đối tượng theo trục thời gian, còn biểu đồ trạng thái (D) mô tả vòng đời và các chuyển trạng thái của một đối tượng đơn lẻ.

\par\vspace{0.3cm}
\hrule
\vspace{0.3cm}
\end{minipage}

\begin{minipage}{\textwidth}
\subsection*{Question 353 \quad \footnotesize{[0353\_q353\_Question\_353.png]}}
\noindent \textbf{Each element in a record is called a(n)}

\begin{enumerate}[label=\Alph*.]
    \item Cell
    \item Word
    \item Item
    \item \textbf{Field} \hfill \textcolor{blue}{\textbf{[Correct]}}
\end{enumerate}

\noindent \textbf{Correct Answer:} \textcolor{blue}{\textbf{D}} \\[2pt]
\noindent \textbf{Explanation:} Trong cấu trúc dữ liệu và tổ chức tệp, một bản ghi (record) là một tập hợp các phần tử dữ liệu có quan hệ logic với nhau, trong đó mỗi phần tử dữ liệu thành phần lưu trữ một thuộc tính cụ thể được gọi là một trường (field - đáp án D). Ví dụ, trong bản ghi về một sinh viên, các mục thông tin như Mã số sinh viên, Họ tên, hay Điểm trung bình được gọi là các trường dữ liệu. Các thuật ngữ Cell (ô tính), Word (từ nhớ), hay Item không phải là thuật ngữ kỹ thuật chỉ thành phần của một record theo giáo trình của Forouzan.

\par\vspace{0.3cm}
\hrule
\vspace{0.3cm}
\end{minipage}

\begin{minipage}{\textwidth}
\subsection*{Question 354 \quad \footnotesize{[0354\_q354\_Question\_354.png]}}
\noindent \textbf{What are the characteristics of a binary search tree (BST) in terms of data order?}

\begin{enumerate}[label=\Alph*.]
    \item \textbf{The data in the parent node is always larger than the data in the left child node and smaller than the data in the right child node} \hfill \textcolor{blue}{\textbf{[Correct]}}
    \item The data in the parent node is always smaller than the data in the left child node and larger than the data in the right child node
    \item The data at the nodes does not have any order
    \item The data in the left child node and right child node are always equal
    \item None of the others
\end{enumerate}

\noindent \textbf{Correct Answer:} \textcolor{blue}{\textbf{A}} \\[2pt]
\noindent \textbf{Explanation:} Cây tìm kiếm nhị phân (Binary Search Tree - BST) là một cấu trúc dữ liệu cây có thứ tự nghiêm ngặt: với mỗi nút cha bất kỳ, dữ liệu tại nút cha luôn lớn hơn dữ liệu của nút con bên trái (và toàn bộ cây con bên trái), đồng thời nhỏ hơn dữ liệu của nút con bên phải (và toàn bộ cây con bên phải) (đáp án A). Đặc tính này giúp tối ưu hóa việc tìm kiếm với độ phức tạp trung bình là O(log n). Đáp án B đảo ngược hoàn toàn quy tắc thứ tự của BST, còn đáp án C và D mâu thuẫn trực tiếp với định nghĩa cây tìm kiếm nhị phân.

\par\vspace{0.3cm}
\hrule
\vspace{0.3cm}
\end{minipage}

\begin{minipage}{\textwidth}
\subsection*{Question 355 \quad \footnotesize{[0355\_q355\_Question\_355.png]}}
\noindent \textbf{What is Queue used for?}

\begin{enumerate}[label=\Alph*.]
    \item \textbf{Keep the same order} \hfill \textcolor{blue}{\textbf{[Correct]}}
    \item Reverse the order of elements
    \item Random elements
    \item Sort the elements
\end{enumerate}

\noindent \textbf{Correct Answer:} \textcolor{blue}{\textbf{A}} \\[2pt]
\noindent \textbf{Explanation:} Hàng đợi (Queue) là cấu trúc dữ liệu tuyến tính hoạt động theo nguyên lý FIFO (First-In, First-Out - Vào trước, Ra trước), nghĩa là các phần tử được lấy ra đúng theo thứ tự thời gian mà chúng được thêm vào, do đó Queue được sử dụng để giữ nguyên thứ tự ban đầu của các phần tử (Keep the same order - đáp án A). Trái lại, Ngăn xếp (Stack) với nguyên lý LIFO mới là cấu trúc giúp đảo ngược thứ tự các phần tử (đáp án B). Queue không có chức năng tạo ngẫu nhiên (C) hay tự động sắp xếp các phần tử (D).

\par\vspace{0.3cm}
\hrule
\vspace{0.3cm}
\end{minipage}

\begin{minipage}{\textwidth}
\subsection*{Question 356 \quad \footnotesize{[0356\_q356\_Question\_356.png]}}
\noindent \textbf{Inside the ADT (abstract data type) are two different parts of the model: SÖ DÁP ÁN CÃN CHON}

\begin{enumerate}[label=\Alph*.]
    \item Hardware and software
    \item Database and operations
    \item Data structure and algorithms
    \item \textbf{Data structure and operations} \hfill \textcolor{blue}{\textbf{[Correct]}}
\end{enumerate}

\noindent \textbf{Correct Answer:} \textcolor{blue}{\textbf{D}} \\[2pt]
\noindent \textbf{Explanation:} Theo giáo trình Forouzan, một kiểu dữ liệu trừu tượng (Abstract Data Type - ADT) bao gồm hai thành phần cốt lõi bên trong mô hình: cấu trúc dữ liệu (data structure) dùng để chứa dữ liệu và tập hợp các phép toán (operations) được phép thực hiện trên cấu trúc dữ liệu đó (đáp án D). Người dùng bên ngoài chỉ tương tác thông qua giao diện của các phép toán mà không cần biết chi tiết cài đặt bên trong của cấu trúc dữ liệu (tính bao đóng). Các đáp án A, B và C đều đưa vào các thành phần không thuộc định nghĩa cấu trúc của mô hình ADT.

\par\vspace{0.3cm}
\hrule
\vspace{0.3cm}
\end{minipage}

\begin{minipage}{\textwidth}
\subsection*{Question 357 \quad \footnotesize{[0357\_q357\_Question\_357.png]}}
\noindent \textbf{Which data structure is commonly used for First-In, First-Out (FIFO) data processing?}

\begin{enumerate}[label=\Alph*.]
    \item Stack
    \item \textbf{Queue} \hfill \textcolor{blue}{\textbf{[Correct]}}
    \item Tree
    \item Linked list
\end{enumerate}

\noindent \textbf{Correct Answer:} \textcolor{blue}{\textbf{B}} \\[2pt]
\noindent \textbf{Explanation:} Hàng đợi (Queue) là cấu trúc dữ liệu đặc trưng hoạt động theo cơ chế Vào trước - Ra trước (First-In, First-Out - FIFO), trong đó phần tử nào được đưa vào hàng đợi trước sẽ là phần tử đầu tiên được lấy ra xử lý (đáp án B). Ngược lại, Ngăn xếp (Stack - A) hoạt động theo nguyên lý LIFO (Last-In, First-Out); Cây (Tree - C) là cấu trúc phân cấp phi tuyến tính; còn Danh sách liên kết (Linked list - D) là một cấu trúc dữ liệu cơ bản có thể dùng để cài đặt cả Queue lẫn Stack nhưng bản thân nó không quy định thứ tự xử lý FIFO.

\par\vspace{0.3cm}
\hrule
\vspace{0.3cm}
\end{minipage}

\begin{minipage}{\textwidth}
\subsection*{Question 358 \quad \footnotesize{[0358\_q358\_Question\_358.png]}}
\noindent \textbf{is the mask to unset (clear) the four leftmost bits of the pattern 1010 0111?}

\begin{enumerate}[label=\Alph*.]
    \item \textbf{0000 1111} \hfill \textcolor{blue}{\textbf{[Correct]}}
    \item 1111 0000
    \item 0101 1000
    \item 1010 0000
\end{enumerate}

\noindent \textbf{Correct Answer:} \textcolor{blue}{\textbf{A}} \\[2pt]
\noindent \textbf{Explanation:} Để xóa (đưa về giá trị 0 / unset) các bit mong muốn trong một mẫu bit mà vẫn giữ nguyên các bit còn lại, ta thực hiện phép toán logic AND với một mặt nạ (mask) có bit 0 tại vị trí cần xóa và bit 1 tại vị trí cần giữ nguyên. Do đề bài yêu cầu xóa 4 bit ngoài cùng bên trái của mẫu bit 1010 0111, mặt nạ cần có 4 bit đầu bằng 0 và 4 bit sau bằng 1, tức là 0000 1111 (đáp án A). Khi thực hiện phép AND: 1010 0111 AND 0000 1111 = 0000 0111, 4 bit trái đã bị xóa thành công về 0 trong khi 4 bit phải được bảo toàn.

\par\vspace{0.3cm}
\hrule
\vspace{0.3cm}
\end{minipage}

\begin{minipage}{\textwidth}
\subsection*{Question 359 \quad \footnotesize{[0359\_q359\_Question\_359.png]}}
\noindent \textbf{Which IEEE standard is most commonly used for representing floating-point numbers in computers?}

\begin{enumerate}[label=\Alph*.]
    \item \textbf{IEEE 754} \hfill \textcolor{blue}{\textbf{[Correct]}}
    \item IEEE 802
    \item IEEE 1284
    \item IEEE 1394
\end{enumerate}

\noindent \textbf{Correct Answer:} \textcolor{blue}{\textbf{A}} \\[2pt]
\noindent \textbf{Explanation:} Chuẩn IEEE 754 là tiêu chuẩn công nghiệp phổ biến và quan trọng nhất được sử dụng trong các hệ thống máy tính hiện đại để biểu diễn và tính toán số thực dấu phẩy động (floating-point numbers), bao gồm định dạng độ chính xác đơn (32-bit) và kép (64-bit) (đáp án A). Trong khi đó, IEEE 802 là bộ tiêu chuẩn cho mạng máy tính cục bộ và đô thị (Ethernet, Wi-Fi), IEEE 1284 là tiêu chuẩn cổng song song của máy in, và IEEE 1394 là giao tiếp truyền dữ liệu tốc độ cao FireWire.

\par\vspace{0.3cm}
\hrule
\vspace{0.3cm}
\end{minipage}

\begin{minipage}{\textwidth}
\subsection*{Question 360 \quad \footnotesize{[0360\_q360\_Question\_360.png]}}
\noindent \textbf{To un-set (force to 0) all the bits of a bit pattern, we SÖ DÁP ÁN CÃN CHON}

\begin{enumerate}[label=\Alph*.]
    \item \textbf{make a mask of all 0s and then AND the bit pattern and the mask.} \hfill \textcolor{blue}{\textbf{[Correct]}}
    \item make a mask exactly the same as the bit pattern and then XOR the bit pattern and the mask.
    \item Both a and b are correct
    \item Both a and b are incorrect
\end{enumerate}

\noindent \textbf{Correct Answer:} \textcolor{blue}{\textbf{A}} \\[2pt]
\noindent \textbf{Explanation:} Theo giáo trình Forouzan, để un-set (ép về 0) các bit trong một mẫu bit, phương pháp chuẩn được định nghĩa là sử dụng phép toán logic AND với một mặt nạ (mask) chứa các bit 0 tại các vị trí tương ứng (đáp án A). Cụ thể, khi AND bất kỳ bit nào với 0, kết quả thu được luôn là 0 (x AND 0 = 0). Mặc dù phép XOR một mẫu bit với chính nó cũng cho kết quả toàn bit 0, nhưng trong ngữ cảnh kỹ thuật che bit (masking) theo giáo trình thì mặt nạ chứa toàn bit 0 kết hợp phép AND mới là định nghĩa chuẩn tắc.

\par\vspace{0.3cm}
\hrule
\vspace{0.3cm}
\end{minipage}

\begin{minipage}{\textwidth}
\subsection*{Question 361 \quad \footnotesize{[0361\_q361\_Question\_361.png]}}
\noindent \textbf{What does "sampling" refer to in the context of audio storage?}

\begin{enumerate}[label=\Alph*.]
    \item Converting audio into hexadecimal format
    \item \textbf{Capturing discrete values of the sound wave at regular intervals} \hfill \textcolor{blue}{\textbf{[Correct]}}
    \item Compressing the audio file
    \item Encoding text metadata for audio
\end{enumerate}

\noindent \textbf{Correct Answer:} \textcolor{blue}{\textbf{B}} \\[2pt]
\noindent \textbf{Explanation:} Trong quá trình số hóa và lưu trữ âm thanh, kỹ thuật lấy mẫu (sampling) là quá trình đo đạc và ghi lại biên độ điện áp của sóng âm thanh tương tự (analog) tại các khoảng thời gian rời rạc đều đặn (regular intervals) (đáp án B). Số lần lấy mẫu trong mỗi giây được gọi là tần số lấy mẫu (sampling rate), quyết định độ chi tiết của tín hiệu âm thanh kỹ thuật số thu được. Quá trình này hoàn toàn phân biệt với việc nén dung lượng tệp âm thanh (C) hay chuyển đổi sang hệ thập lục phân (A).

\par\vspace{0.3cm}
\hrule
\vspace{0.3cm}
\end{minipage}

\begin{minipage}{\textwidth}
\subsection*{Question 362 \quad \footnotesize{[0362\_q362\_Question\_362.png]}}
\noindent \textbf{Which shift operation would you use to divide a signed integer by 2? tup}

\begin{enumerate}[label=\Alph*.]
    \item Logical right shift
    \item \textbf{Arithmetic right shift} \hfill \textcolor{blue}{\textbf{[Correct]}}
    \item Logical left shift
    \item Arithmetic left shift
\end{enumerate}

\noindent \textbf{Correct Answer:} \textcolor{blue}{\textbf{B}} \\[2pt]
\noindent \textbf{Explanation:} Để thực hiện phép chia một số nguyên có dấu cho 2, ta sử dụng phép dịch phải số học (Arithmetic right shift - đáp án B). Phép toán này dịch chuyển toàn bộ các bit sang bên phải 1 vị trí, đồng thời sao chép giữ nguyên bit dấu ở vị trí ngoài cùng bên trái (sign-bit replication), qua đó bảo toàn tính chất âm hoặc dương của số nguyên. Trái lại, phép dịch phải logic (A) luôn chèn bit 0 vào vị trí ngoài cùng bên trái nên sẽ làm sai lệch dấu của các số âm, còn các phép dịch trái (C và D) tương đương với phép nhân với 2.

\par\vspace{0.3cm}
\hrule
\vspace{0.3cm}
\end{minipage}

\begin{minipage}{\textwidth}
\subsection*{Question 363 \quad \footnotesize{[0363\_q363\_Question\_363.png]}}
\noindent \textbf{Which file format is commonly used for storing vector images?}

\begin{enumerate}[label=\Alph*.]
    \item .jpg
    \item \textbf{.svg} \hfill \textcolor{blue}{\textbf{[Correct]}}
    \item .gif
    \item .mp4
\end{enumerate}

\noindent \textbf{Correct Answer:} \textcolor{blue}{\textbf{B}} \\[2pt]
\noindent \textbf{Explanation:} Định dạng SVG (Scalable Vector Graphics) là định dạng tệp chuẩn mở dựa trên XML chuyên dùng để lưu trữ đồ họa vector, biểu diễn hình ảnh bằng các đối tượng hình học và công thức toán học giúp ảnh không bị vỡ hạt hay giảm chất lượng khi co dãn kích thước (đáp án B). Ngược lại, .jpg và .gif (A, C) là các định dạng ảnh raster (bitmap) lưu trữ lưới các điểm ảnh (pixels), còn .mp4 (D) là định dạng tệp video đa phương tiện.

\par\vspace{0.3cm}
\hrule
\vspace{0.3cm}
\end{minipage}

\begin{minipage}{\textwidth}
\subsection*{Question 364 \quad \footnotesize{[0364\_q364\_Question\_364.png]}}
\noindent \textbf{To guarantee secure communication among 10 hosts using asymmetric-key encryption, how many keys are necessary?}

\begin{enumerate}[label=\Alph*.]
    \item \textbf{20} \hfill \textcolor{blue}{\textbf{[Correct]}}
    \item 21
    \item 100
    \item 90
\end{enumerate}

\noindent \textbf{Correct Answer:} \textcolor{blue}{\textbf{A}} \\[2pt]
\noindent \textbf{Explanation:} Trong hệ thống mã hóa khóa bất đối xứng (asymmetric-key / public-key cryptography), mỗi nút mạng (host) độc lập chỉ cần sở hữu đúng một cặp khóa gồm: một khóa công khai (public key) và một khóa riêng tư (private key). Do đó, đối với một hệ thống gồm 10 máy chủ, tổng số khóa cần thiết để đảm bảo liên lạc an toàn là 10 × 2 = 20 khóa (đáp án A). Nếu sử dụng phương pháp mã hóa khóa đối xứng (symmetric key), số lượng khóa bí mật cần thiết sẽ lớn hơn rất nhiều, cụ thể là n(n-1)/2 = (10 × 9)/2 = 45 khóa.

\par\vspace{0.3cm}
\hrule
\vspace{0.3cm}
\end{minipage}

\begin{minipage}{\textwidth}
\subsection*{Question 365 \quad \footnotesize{[0365\_q365\_Question\_365.png]}}
\noindent \textbf{Which statement describes red hat hackers?}

\begin{enumerate}[label=\Alph*.]
    \item Hackers who infiltrate computer networks with negative motives, such as monetary gain or reputation building
    \item \textbf{Hackers who track down black hat hackers and work to stop their activities} \hfill \textcolor{blue}{\textbf{[Correct]}}
    \item Hackers whose actions are often driven by revenge or personal agendas
    \item Hackers who expose large datasets online for the public good
\end{enumerate}

\noindent \textbf{Correct Answer:} \textcolor{blue}{\textbf{B}} \\[2pt]
\noindent \textbf{Explanation:} Hacker mũ đỏ (Red hat hackers) là thuật ngữ chỉ những hacker có vai trò như những hiệp sĩ săn lùng và vô hiệu hóa hoạt động của các tin tặc mũ đen (black hat hackers), thường chủ động sử dụng các kỹ thuật tấn công đáp trả trực diện vào chính hệ thống của kẻ tấn công (đáp án B). Ngược lại, hacker xâm nhập vì mục tiêu trục lợi cá nhân bất chính là hacker mũ đen (đáp án A), còn người hành động vì sự tò mò, khám phá hoặc công bố dữ liệu vì cộng đồng thuộc các nhóm hacker khác như mũ trắng hoặc hacktivist.

\par\vspace{0.3cm}
\hrule
\vspace{0.3cm}
\end{minipage}

\begin{minipage}{\textwidth}
\subsection*{Question 366 \quad \footnotesize{[0366\_q366\_Question\_366.png]}}
\noindent \textbf{How many types of cyber attacks exist?}

\begin{enumerate}[label=\Alph*.]
    \item 2
    \item 3
    \item \textbf{4} \hfill \textcolor{blue}{\textbf{[Correct]}}
    \item 5
\end{enumerate}

\noindent \textbf{Correct Answer:} \textcolor{blue}{\textbf{C}} \\[2pt]
\noindent \textbf{Explanation:} Trong lý thuyết nền tảng an ninh mạng và an toàn thông tin, các cuộc tấn công mạng được phân chia thành 4 dạng cơ bản dựa trên tác động tới luồng truyền thông: Chặn bắt (Interception - đe dọa tính bảo mật), Ngắt quãng (Interruption - đe dọa tính sẵn sàng), Sửa đổi (Modification - đe dọa tính toàn vẹn), và Giả mạo (Fabrication - tạo ra thông tin giả mạo) (đáp án C). Phân loại 4 nhóm này tạo thành cơ sở kinh điển để thiết kế các cơ chế bảo mật mạng tương ứng.

\par\vspace{0.3cm}
\hrule
\vspace{0.3cm}
\end{minipage}

\begin{minipage}{\textwidth}
\subsection*{Question 367 \quad \footnotesize{[0367\_q367\_Question\_367.png]}}
\noindent \textbf{How many main types of intellectual property are recognized?}

\begin{enumerate}[label=\Alph*.]
    \item 2
    \item 3
    \item \textbf{4} \hfill \textcolor{blue}{\textbf{[Correct]}}
    \item 5
\end{enumerate}

\noindent \textbf{Correct Answer:} \textcolor{blue}{\textbf{C}} \\[2pt]
\noindent \textbf{Explanation:} Theo giáo trình Forouzan (Chương các vấn đề xã hội và đạo đức), có 4 loại hình sở hữu trí tuệ chính được pháp luật và xã hội công nhận bao gồm: Nhãn hiệu thương mại (trademarks), Bí mật thương mại (trade secrets), Bằng sáng chế (patents), và Bản quyền tác giả (copyrights) (đáp án C). Bốn hình thức này bảo hộ các khía cạnh khác nhau của tài sản trí tuệ như biểu tượng nhận diện, công thức kinh doanh bí mật, phát minh công nghệ và các tác phẩm sáng tạo nguyên bản.

\par\vspace{0.3cm}
\hrule
\vspace{0.3cm}
\end{minipage}

\begin{minipage}{\textwidth}
\subsection*{Question 368 \quad \footnotesize{[0368\_q368\_Question\_368.png]}}
\noindent \textbf{Which of the following is NOT a computer language paradigm?}

\begin{enumerate}[label=\Alph*.]
    \item Procedural
    \item Object-oriented
    \item Functional
    \item Declarative
    \item \textbf{Mathematical} \hfill \textcolor{blue}{\textbf{[Correct]}}
\end{enumerate}

\noindent \textbf{Correct Answer:} \textcolor{blue}{\textbf{E}} \\[2pt]
\noindent \textbf{Explanation:} Trong khoa học máy tính, các mô hình ngôn ngữ lập trình (programming language paradigms) chính thống được phân loại thành: mô hình lập trình thủ tục (procedural), hướng đối tượng (object-oriented), hàm (functional), và khai báo/logic (declarative). Khái niệm 'Mathematical' (toán học) không phải là tên gọi của một mô hình ngôn ngữ lập trình chuẩn mực nào trong ngành (đáp án E). Mặc dù lập trình hàm dựa trên lý thuyết toán học (lambda calculus), nhưng mô hình chuẩn của nó vẫn là functional paradigm.

\par\vspace{0.3cm}
\hrule
\vspace{0.3cm}
\end{minipage}

\begin{minipage}{\textwidth}
\subsection*{Question 369 \quad \footnotesize{[0369\_q369\_Question\_369.png]}}
\noindent \textbf{How does a compiler function in the context of high-level languages?}

\begin{enumerate}[label=\Alph*.]
    \item It directly executes the high-level code
    \item \textbf{It translates high-level code into machine language all at once} \hfill \textcolor{blue}{\textbf{[Correct]}}
    \item It translates and executes the code line by line
    \item It optimizes the code for better performance
\end{enumerate}

\noindent \textbf{Correct Answer:} \textcolor{blue}{\textbf{B}} \\[2pt]
\noindent \textbf{Explanation:} Trình biên dịch (Compiler) là chương trình dịch có chức năng chuyển đổi toàn bộ mã nguồn của ngôn ngữ bậc cao sang mã máy hoặc mã đối tượng trong một lần duy nhất (all at once) trước khi chương trình được chạy (đáp án B). Quá trình này tạo ra tệp thực thi độc lập giúp chương trình chạy nhanh hơn. Trái lại, việc dịch và thực thi từng dòng lệnh một (line by line - đáp án C) là cơ chế làm việc của trình thông dịch (Interpreter), và bản thân compiler không trực tiếp thực thi mã (đáp án A).

\par\vspace{0.3cm}
\hrule
\vspace{0.3cm}
\end{minipage}

\begin{minipage}{\textwidth}
\subsection*{Question 370 \quad \footnotesize{[0370\_q370\_Question\_370.png]}}
\noindent \textbf{Choose the most suitable answer: Python is a/an language SÖ DÁP ÁN CÀN CHON}

\begin{enumerate}[label=\Alph*.]
    \item \textbf{object-oriented} \hfill \textcolor{blue}{\textbf{[Correct]}}
    \item procedural
    \item functional
    \item declarative
\end{enumerate}

\noindent \textbf{Correct Answer:} \textcolor{blue}{\textbf{A}} \\[2pt]
\noindent \textbf{Explanation:} Mặc dù Python là ngôn ngữ đa mô hình hỗ trợ cả lập trình thủ tục và hàm, nhưng trong phân loại kiến trúc cốt lõi của giáo trình khoa học máy tính, Python được xếp vào nhóm ngôn ngữ hướng đối tượng (object-oriented language - đáp án A). Trong Python, mọi cấu trúc dữ liệu từ các kiểu dữ liệu nguyên thủy (số, chuỗi) đến các hàm và module đều được thiết kế và xử lý dưới dạng các đối tượng (objects). Các mô hình khác chỉ là tính năng bổ trợ chứ không phải định danh đặc trưng nhất của Python.

\par\vspace{0.3cm}
\hrule
\vspace{0.3cm}
\end{minipage}

\begin{minipage}{\textwidth}
\subsection*{Question 371 \quad \footnotesize{[0371\_q371\_Question\_371.png]}}
\noindent \textbf{Who is known as the "Father of the Computer"?}

\begin{enumerate}[label=\Alph*.]
    \item Alan Turing
    \item \textbf{Charles Babbage} \hfill \textcolor{blue}{\textbf{[Correct]}}
    \item Steve Jobs
    \item Bill Gates
\end{enumerate}

\noindent \textbf{Correct Answer:} \textcolor{blue}{\textbf{B}} \\[2pt]
\noindent \textbf{Explanation:} Charles Babbage (1791–1871) được tôn vinh là 'Cha đẻ của máy tính' (Father of the Computer) vì ông là người đầu tiên thiết kế ra máy tính cơ học sai phân (Difference Engine) và máy phân tích (Analytical Engine) chứa đầy đủ các thành phần cơ bản của máy tính hiện đại như bộ nhớ và bộ xử lý (đáp án B). Trong khi đó, Alan Turing (A) được coi là cha đẻ của khoa học máy tính lý thuyết và AI; còn Steve Jobs (C) và Bill Gates (D) là những nhà lãnh đạo công nghệ tiên phong của kỷ nguyên máy tính cá nhân hiện đại.

\par\vspace{0.3cm}
\hrule
\vspace{0.3cm}
\end{minipage}

\begin{minipage}{\textwidth}
\subsection*{Question 372 \quad \footnotesize{[0372\_q372\_Question\_372.png]}}
\noindent \textbf{What model is the foundation of modern computers?}

\begin{enumerate}[label=\Alph*.]
    \item Leibnitz
    \item \textbf{von Neumann} \hfill \textcolor{blue}{\textbf{[Correct]}}
    \item Pascal
    \item Charles Babbage
\end{enumerate}

\noindent \textbf{Correct Answer:} \textcolor{blue}{\textbf{B}} \\[2pt]
\noindent \textbf{Explanation:} Mô hình kiến trúc Von Neumann (do nhà toán học John von Neumann đề xuất năm 1945) là nền tảng kiến trúc phần cứng của hầu hết mọi máy tính kỹ thuật số hiện đại ngày nay (đáp án B). Nguyên lý cơ bản của mô hình này là khái niệm chương trình được lưu trữ (stored-program concept), trong đó cả dữ liệu và tập lệnh thực thi đều được đặt chung trong bộ nhớ chính. Các nhà khoa học cơ học cổ điển như Pascal hay Leibniz chỉ chế tạo các máy tính cộng trừ cơ học đơn giản.

\par\vspace{0.3cm}
\hrule
\vspace{0.3cm}
\end{minipage}

\begin{minipage}{\textwidth}
\subsection*{Question 373 \quad \footnotesize{[0373\_q373\_Question\_373.png]}}
\noindent \textbf{According to the Von Neumann model, which one is NOT a computer hardware subsystem ?}

\begin{enumerate}[label=\Alph*.]
    \item Arithmetic logic unit (ALU)
    \item Input/Output (I/O)
    \item Control unit (CU)
    \item \textbf{Power supply unit (PSU)} \hfill \textcolor{blue}{\textbf{[Correct]}}
\end{enumerate}

\noindent \textbf{Correct Answer:} \textcolor{blue}{\textbf{D}} \\[2pt]
\noindent \textbf{Explanation:} Theo mô hình Von Neumann chuẩn mực, phần cứng máy tính gồm 4 phân hệ chính: Bộ nhớ (Memory), Đơn vị số học và logic (ALU), Đơn vị điều khiển (Control Unit - CU), và Phân hệ Nhập/Xuất (Input/Output - I/O). Bộ nguồn cấp điện (Power supply unit - PSU) là linh kiện phần cứng phụ trợ có nhiệm vụ cung cấp năng lượng điện cho máy hoạt động, không thuộc vào 4 phân hệ kiến trúc logic nền tảng của mô hình Von Neumann (đáp án D).

\par\vspace{0.3cm}
\hrule
\vspace{0.3cm}
\end{minipage}

\begin{minipage}{\textwidth}
\subsection*{Question 374 \quad \footnotesize{[0374\_q374\_Question\_374.png]}}
\noindent \textbf{A memory with 128 kilobytes and a word size of 1 byte has an address space that ranges from to}

\begin{enumerate}[label=\Alph*.]
    \item \textbf{0; 131 071} \hfill \textcolor{blue}{\textbf{[Correct]}}
    \item 1;131 071
    \item 0;131 072
    \item 1; 131 072
\end{enumerate}

\noindent \textbf{Correct Answer:} \textcolor{blue}{\textbf{A}} \\[2pt]
\noindent \textbf{Explanation:} Bộ nhớ có dung lượng 128 kilobytes (KB) tương đương với 128 × 1024 = 131,072 bytes. Với kích thước từ nhớ (word size) là 1 byte, bộ nhớ có tổng cộng đúng 131,072 vị trí nhớ cần được đánh số. Trong kiến trúc máy tính, không gian địa chỉ luôn bắt đầu từ 0 và kết thúc tại N - 1, do đó dải địa chỉ hợp lệ sẽ trải dài từ 0 đến 131,072 - 1 = 131,071 (đáp án A). Các đáp án bắt đầu từ địa chỉ 1 (B, D) hoặc kết thúc ở 131,072 (C) đều không tuân thủ quy tắc đánh chỉ số từ 0.

\par\vspace{0.3cm}
\hrule
\vspace{0.3cm}
\end{minipage}

\begin{minipage}{\textwidth}
\subsection*{Question 375 \quad \footnotesize{[0375\_q375\_Question\_375.png]}}
\noindent \textbf{Which of the following is closest to the CPU in the memory hierarchy?}

\begin{enumerate}[label=\Alph*.]
    \item \textbf{Cache memory} \hfill \textcolor{blue}{\textbf{[Correct]}}
    \item Secondary storage
    \item Main memory
    \item Virtual memory
\end{enumerate}

\noindent \textbf{Correct Answer:} \textcolor{blue}{\textbf{A}} \\[2pt]
\noindent \textbf{Explanation:} Trong hệ thống phân cấp bộ nhớ (Memory Hierarchy), bộ nhớ đệm (Cache memory) là tầng bộ nhớ nằm gần CPU nhất (thường được tích hợp ngay trên cùng phiến bán dẫn với CPU), có tốc độ truy xuất cực cao nhằm đồng bộ tốc độ giữa các thanh ghi của CPU và bộ nhớ chính (đáp án A). Bộ nhớ chính (RAM - C) nằm ở cấp tiếp theo với khoảng cách xa hơn và tốc độ chậm hơn; bộ nhớ phụ (Secondary storage như HDD, SSD - B) nằm ở tầng thấp nhất; còn bộ nhớ ảo (D) là cơ chế mở rộng không gian nhớ do hệ điều hành quản lý trên ổ đĩa.

\par\vspace{0.3cm}
\hrule
\vspace{0.3cm}
\end{minipage}

\begin{minipage}{\textwidth}
\subsection*{Question 376 \quad \footnotesize{[0376\_q376\_Question\_376.png]}}
\noindent \textbf{Which type of file structure uses hash functions for indexing?}

\begin{enumerate}[label=\Alph*.]
    \item Sequential files
    \item Indexed files
    \item \textbf{Hashed files} \hfill \textcolor{blue}{\textbf{[Correct]}}
    \item Binary files
\end{enumerate}

\noindent \textbf{Correct Answer:} \textcolor{blue}{\textbf{C}} \\[2pt]
\noindent \textbf{Explanation:} Trong các phương thức tổ chức tệp (File Structure - Chapter 10), cấu trúc tệp băm (Hashed files) là loại cấu trúc sử dụng một hàm băm (hash function) để tính toán và ánh xạ trực tiếp giá trị của trường khóa thành địa chỉ của khối dữ liệu hoặc bản ghi trên thiết bị lưu trữ (đáp án C). Trong khi đó, tệp tuần tự (A) chỉ lưu trữ các bản ghi liên tiếp theo thứ tự thời gian hoặc khóa; tệp có chỉ mục (Indexed files - B) sử dụng các bảng chỉ mục độc lập chứ không biến đổi khóa bằng hàm băm; còn tệp nhị phân (D) chỉ là định dạng lưu trữ mức bit/byte.

\par\vspace{0.3cm}
\hrule
\vspace{0.3cm}
\end{minipage}

\begin{minipage}{\textwidth}
\subsection*{Question 377 \quad \footnotesize{[0377\_q377\_Question\_377.png]}}
\noindent \textbf{What does the bit string '0110000101100011' mean?}

\begin{enumerate}[label=\Alph*.]
    \item It represents the two characters 'A' and 'C'
    \item It represents the two characters 'A' and 'b'
    \item \textbf{It represents the two characters 'a' and 'c'} \hfill \textcolor{blue}{\textbf{[Correct]}}
    \item It represents the decimal number 24903
\end{enumerate}

\noindent \textbf{Correct Answer:} \textcolor{blue}{\textbf{C}} \\[2pt]
\noindent \textbf{Explanation:} Chuỗi 16 bit 0110000101100011 gồm hai byte 8-bit riêng biệt mã hóa theo bảng mã ký tự chuẩn ASCII. Byte đầu tiên 01100001 tương ứng với số thập phân 97, là mã ASCII của ký tự chữ thường 'a'; byte thứ hai 01100011 tương ứng với số thập phân 99, là mã ASCII của ký tự chữ thường 'c', do đó chuỗi bit biểu diễn hai ký tự 'a' và 'c' (đáp án C). Chữ in hoa 'A' có mã ASCII là 65 (01000001) và 'C' là 67 (01000011) nên đáp án A không chính xác.

\par\vspace{0.3cm}
\hrule
\vspace{0.3cm}
\end{minipage}

\begin{minipage}{\textwidth}
\subsection*{Question 378 \quad \footnotesize{[0378\_q378\_Question\_378.png]}}
\noindent \textbf{If we need to access one specific record without having to retrieve all records before it, we use a}

\begin{enumerate}[label=\Alph*.]
    \item sequential file structure
    \item \textbf{random file structure} \hfill \textcolor{blue}{\textbf{[Correct]}}
    \item Both a and b are correct
    \item Both a and b are incorrect
\end{enumerate}

\noindent \textbf{Correct Answer:} \textcolor{blue}{\textbf{B}} \\[2pt]
\noindent \textbf{Explanation:} Trong cấu trúc tệp truy cập ngẫu nhiên hay trực tiếp (random / direct file structure), hệ thống cho phép định vị và truy xuất thẳng tới một bản ghi mong muốn ở bất kỳ vị trí nào mà không cần phải đọc qua toàn bộ các bản ghi đứng trước nó (đáp án B). Điều này mang lại tốc độ truy xuất nhanh chóng cho các ứng dụng tương tác thời gian thực. Ngược lại, cấu trúc tệp tuần tự (sequential file structure - đáp án A) bắt buộc phải đọc lần lượt từ đầu tệp qua tất cả các bản ghi phía trước để đến được bản ghi cần tìm.

\par\vspace{0.3cm}
\hrule
\vspace{0.3cm}
\end{minipage}

\begin{minipage}{\textwidth}
\subsection*{Question 379 \quad \footnotesize{[0379\_q379\_Question\_379.png]}}
\noindent \textbf{Which access method typically requires additional space for maintaining an index structure but provides faster access to records compared to sequential access?}

\begin{enumerate}[label=\Alph*.]
    \item Sequential access
    \item Direct access
    \item \textbf{Indexed sequential access} \hfill \textcolor{blue}{\textbf{[Correct]}}
    \item Hashing
\end{enumerate}

\noindent \textbf{Correct Answer:} \textcolor{blue}{\textbf{C}} \\[2pt]
\noindent \textbf{Explanation:} Phương pháp truy cập tuần tự có chỉ mục (Indexed sequential access) sử dụng thêm một tệp chỉ mục (index file) lưu trữ các cặp khóa và địa chỉ nhằm tăng tốc độ định vị bản ghi so với truy cập tuần tự thông thường, đổi lại hệ thống phải tốn thêm không gian bộ nhớ phụ để duy trì cấu trúc chỉ mục này (đáp án C). Truy cập tuần tự (Sequential access - A) không sử dụng chỉ mục nên không tốn thêm không gian chỉ mục; còn Direct access (B) hoặc Hashing (D) tính toán địa chỉ trực tiếp mà không lưu cấu trúc chỉ mục tuần tự.

\par\vspace{0.3cm}
\hrule
\vspace{0.3cm}
\end{minipage}

\begin{minipage}{\textwidth}
\subsection*{Question 380 \quad \footnotesize{[0380\_q380\_Question\_380.png]}}
\noindent \textbf{What type of network device operates at the Data Link layer of the OSI model?}

\begin{enumerate}[label=\Alph*.]
    \item Router
    \item \textbf{Switch} \hfill \textcolor{blue}{\textbf{[Correct]}}
    \item Firewall
    \item Hub
\end{enumerate}

\noindent \textbf{Correct Answer:} \textcolor{blue}{\textbf{B}} \\[2pt]
\noindent \textbf{Explanation:} Bộ chuyển mạch (Switch) là thiết bị mạng phần cứng hoạt động chủ yếu tại tầng Liên kết dữ liệu (Data Link layer - Tầng 2 của mô hình OSI), sử dụng bảng địa chỉ vật lý MAC để chuyển tiếp các khung dữ liệu (frames) một cách chính xác đến cổng của thiết bị đích (đáp án B). Trong khi đó, Bộ định tuyến (Router - A) hoạt động ở tầng Mạng (Network layer - Tầng 3) dựa trên địa chỉ IP; Hub (D) hoạt động ở tầng Vật lý (Physical layer - Tầng 1); còn Firewall (C) thường hoạt động tại các tầng cao hơn như Mạng, Giao vận hoặc Ứng dụng.

\par\vspace{0.3cm}
\hrule
\vspace{0.3cm}
\end{minipage}

\begin{minipage}{\textwidth}
\subsection*{Question 381 \quad \footnotesize{[0381\_q381\_Question\_381.png]}}
\noindent \textbf{What is the meaning of the first part of a URL? (for example: rtsp://192.168.1.1:8080/cam1)}

\begin{enumerate}[label=\Alph*.]
    \item \textbf{Define a protocol} \hfill \textcolor{blue}{\textbf{[Correct]}}
    \item Define a host name
    \item Define a port
    \item Define a device
\end{enumerate}

\noindent \textbf{Correct Answer:} \textcolor{blue}{\textbf{A}} \\[2pt]
\noindent \textbf{Explanation:} Trong cấu trúc của một địa chỉ định danh tài nguyên đồng nhất URL (ví dụ: rtsp://192.168.1.1:8080/cam1), thành phần đầu tiên đứng trước chuỗi phân cách :// (ở đây là rtsp) biểu thị giao thức mạng (protocol) được sử dụng để thiết lập kết nối và trao đổi dữ liệu (đáp án A). Các thành phần kế tiếp sau :// lần lượt là địa chỉ máy chủ (host IP: 192.168.1.1 - đáp án B), cổng dịch vụ (port: 8080 - đáp án C), và đường dẫn tài nguyên cụ thể (/cam1).

\par\vspace{0.3cm}
\hrule
\vspace{0.3cm}
\end{minipage}


\newpage
\section*{11. Databases}
\addcontentsline{toc}{section}{11. Databases}
\hrule\vspace{0.5cm}

\begin{minipage}{\textwidth}
\subsection*{Question 382 \quad \footnotesize{[0382\_q382\_Question\_382.png]}}
\noindent \textbf{Which cryptographic protocol forms the foundation of HTTPS for secure communication over the internet?}

\begin{enumerate}[label=\Alph*.]
    \item RSA
    \item \textbf{SSL/TLS} \hfill \textcolor{blue}{\textbf{[Correct]}}
    \item AES
    \item SHA-256
\end{enumerate}

\noindent \textbf{Correct Answer:} \textcolor{blue}{\textbf{B}} \\[2pt]
\noindent \textbf{Explanation:} Giao thức HTTPS (Hypertext Transfer Protocol Secure) là phiên bản bảo mật của HTTP được xây dựng trên nền tảng của giao thức mã hóa SSL (Secure Sockets Layer) hoặc phiên bản kế thừa nâng cấp của nó là TLS (Transport Layer Security). SSL/TLS cung cấp khả năng mã hóa dữ liệu truyền tải trên mạng Internet, xác thực danh tính máy chủ và đảm bảo tính toàn vẹn của thông điệp (đáp án B). Các thuật toán như RSA (bất đối xứng), AES (đối xứng) và SHA-256 (hàm băm) chỉ là các thuật toán mật mã học thành phần được tích hợp bên trong bộ giao thức SSL/TLS chứ không phải là giao thức mạng nền tảng của HTTPS.

\par\vspace{0.3cm}
\hrule
\vspace{0.3cm}
\end{minipage}

\begin{minipage}{\textwidth}
\subsection*{Question 383 \quad \footnotesize{[0383\_q383\_Question\_383.png]}}
\noindent \textbf{Which layer of the TCP/IP model ensures reliable data transfer?}

\begin{enumerate}[label=\Alph*.]
    \item Physical layer
    \item \textbf{Transport layer} \hfill \textcolor{blue}{\textbf{[Correct]}}
    \item Application layer
    \item Data link layer
\end{enumerate}

\noindent \textbf{Correct Answer:} \textcolor{blue}{\textbf{B}} \\[2pt]
\noindent \textbf{Explanation:} Trong mô hình TCP/IP, Tầng vận chuyển (Transport layer) đóng vai trò điều phối việc truyền dữ liệu giữa các tiến trình (process-to-process delivery) và đảm bảo tính tin cậy của quá trình truyền tải thông qua giao thức TCP với các cơ chế kiểm soát lỗi, kiểm soát luồng và đánh số thứ tự gói tin (đáp án B). Tầng vật lý (Physical) và Tầng liên kết dữ liệu (Data link) chỉ xử lý tín hiệu và khung dữ liệu giữa hai nút mạng kế cận, trong khi Tầng ứng dụng (Application) cung cấp giao diện trực tiếp cho các dịch vụ mạng. Do đó, tầng chịu trách nhiệm đảm bảo truyền dữ liệu tin cậy là Transport layer.

\par\vspace{0.3cm}
\hrule
\vspace{0.3cm}
\end{minipage}

\begin{minipage}{\textwidth}
\subsection*{Question 384 \quad \footnotesize{[0384\_q384\_Question\_384.png]}}
\noindent \textbf{A device is sending out data at the rate of 1000 bps. How long does it take to send out 10 bits?}

\begin{enumerate}[label=\Alph*.]
    \item \textbf{10 ms} \hfill \textcolor{blue}{\textbf{[Correct]}}
    \item 1 ms
    \item 10 000 ms
    \item 100 ms
\end{enumerate}

\noindent \textbf{Correct Answer:} \textcolor{blue}{\textbf{A}} \\[2pt]
\noindent \textbf{Explanation:} Tốc độ truyền 1000 bps (bits per second) biểu thị rằng thiết bị có khả năng gửi đi 1000 bit trong khoảng thời gian 1 giây (tương đương 1000 mili-giây). Do đó, thời gian cần thiết để gửi đi 10 bit được tính theo công thức: Thời gian = 10 bit / 1000 bps = 0.01 giây = 10 ms (đáp án A). Các phương án 1 ms, 100 ms hay 10 000 ms đều không chính xác về mặt tính toán toán học.

\par\vspace{0.3cm}
\hrule
\vspace{0.3cm}
\end{minipage}

\begin{minipage}{\textwidth}
\subsection*{Question 385 \quad \footnotesize{[0385\_q385\_Question\_385.png]}}
\noindent \textbf{What is the hexadecimal equivalent of the binary number (11011011)2?}

\begin{enumerate}[label=\Alph*.]
    \item 6D
    \item \textbf{DB} \hfill \textcolor{blue}{\textbf{[Correct]}}
    \item D6
    \item B6
\end{enumerate}

\noindent \textbf{Correct Answer:} \textcolor{blue}{\textbf{B}} \\[2pt]
\noindent \textbf{Explanation:} Để chuyển đổi một số nhị phân sang hệ thập lục phân (hexadecimal - cơ số 16), ta nhóm các bit thành từng cụm 4 bit từ phải sang trái. Chuỗi nhị phân 11011011 được phân tách thành hai cụm 4 bit là 1101 và 1011: cụm 1101 tương ứng với giá trị 13 trong hệ thập phân tức ký tự 'D', và cụm 1011 tương ứng với giá trị 11 tức ký tự 'B'. Kết hợp hai ký tự theo đúng thứ tự ta được giá trị thập lục phân là DB (đáp án B).

\par\vspace{0.3cm}
\hrule
\vspace{0.3cm}
\end{minipage}

\begin{minipage}{\textwidth}
\subsection*{Question 386 \quad \footnotesize{[0386\_q386\_Question\_386.png]}}
\noindent \textbf{What is the octal representation of the binary number 11010?}

\begin{enumerate}[label=\Alph*.]
    \item 25
    \item \textbf{32} \hfill \textcolor{blue}{\textbf{[Correct]}}
    \item 32
    \item 26
\end{enumerate}

\noindent \textbf{Correct Answer:} \textcolor{blue}{\textbf{B}} \\[2pt]
\noindent \textbf{Explanation:} Để chuyển đổi một số nhị phân sang hệ bát phân (octal - cơ số 8), ta chia chuỗi bit thành các nhóm 3 bit từ phải sang trái, bổ sung thêm các bit 0 vào đầu bên trái nếu thiếu. Chuỗi nhị phân 11010 được nhóm thành (011) và (010), trong đó nhóm (011) tương ứng với giá trị 3 và nhóm (010) tương ứng với giá trị 2 trong hệ bát phân. Kết quả ghép hai nhóm lại cho ra giá trị bát phân là 32 (đáp án B).

\par\vspace{0.3cm}
\hrule
\vspace{0.3cm}
\end{minipage}

\begin{minipage}{\textwidth}
\subsection*{Question 387 \quad \footnotesize{[0387\_q387\_Question\_387.png]}}
\noindent \textbf{Convert the binary number 1011011011 to decimal.}

\begin{enumerate}[label=\Alph*.]
    \item 721
    \item \textbf{731} \hfill \textcolor{blue}{\textbf{[Correct]}}
    \item 739
    \item 743
\end{enumerate}

\noindent \textbf{Correct Answer:} \textcolor{blue}{\textbf{B}} \\[2pt]
\noindent \textbf{Explanation:} Để chuyển đổi số nhị phân 1011011011 sang hệ thập phân, ta tính tổng giá trị của các lũy thừa cơ số 2 tại các vị trí có bit bằng 1 (tính từ vị trí 0 ở ngoài cùng bên phải): 2\textasciicircum{}9 + 2\textasciicircum{}7 + 2\textasciicircum{}6 + 2\textasciicircum{}4 + 2\textasciicircum{}3 + 2\textasciicircum{}1 + 2\textasciicircum{}0 = 512 + 128 + 64 + 16 + 8 + 2 + 1 = 731. Vì vậy, giá trị thập phân tương đương của số nhị phân đã cho là 731 (đáp án B).

\par\vspace{0.3cm}
\hrule
\vspace{0.3cm}
\end{minipage}

\begin{minipage}{\textwidth}
\subsection*{Question 388 \quad \footnotesize{[0388\_q388\_Question\_388.png]}}
\noindent \textbf{Convert the decimal number 213.75 to octal}

\begin{enumerate}[label=\Alph*.]
    \item \textbf{325.6} \hfill \textcolor{blue}{\textbf{[Correct]}}
    \item 225.7
    \item 235.6
    \item 322.5
\end{enumerate}

\noindent \textbf{Correct Answer:} \textcolor{blue}{\textbf{A}} \\[2pt]
\noindent \textbf{Explanation:} Để chuyển số thập phân 213.75 sang hệ bát phân, ta chuyển đổi riêng biệt phần nguyên và phần phân số: Phần nguyên 213 chia liên tiếp cho 8 được 213 = 26 × 8 + 5, 26 = 3 × 8 + 2, 3 = 0 × 8 + 3, đọc ngược số dư ta thu được 325; phần thập phân 0.75 nhân liên tiếp với 8: 0.75 × 8 = 6.0 (phần nguyên là 6 và phần dư bằng 0, dừng lại). Ghép hai phần lại ta được biểu diễn bát phân hoàn chỉnh là 325.6 (đáp án A).

\par\vspace{0.3cm}
\hrule
\vspace{0.3cm}
\end{minipage}

\begin{minipage}{\textwidth}
\subsection*{Question 389 \quad \footnotesize{[0389\_q389\_Question\_389.png]}}
\noindent \textbf{Which component of an operating system is responsible for accepting and interpreting user requests, often through a command-line interface?}

\begin{enumerate}[label=\Alph*.]
    \item \textbf{Shell} \hfill \textcolor{blue}{\textbf{[Correct]}}
    \item Kernel
    \item Graphical interface
    \item File system
\end{enumerate}

\noindent \textbf{Correct Answer:} \textcolor{blue}{\textbf{A}} \\[2pt]
\noindent \textbf{Explanation:} Trong hệ điều hành, Shell (trình bao) là thành phần đóng vai trò làm giao diện thông dịch lệnh (command interpreter), có nhiệm vụ tiếp nhận, phân tích cú pháp và thực thi các yêu cầu của người dùng, phổ biến nhất thông qua giao diện dòng lệnh CLI (đáp án A). Ngược lại, Kernel (nhân) là lõi trung tâm của hệ điều hành quản lý trực tiếp tài nguyên phần cứng, còn File system chỉ quản lý cấu trúc lưu trữ tập tin. Do đó, thành phần tiếp nhận và thông dịch yêu cầu từ người dùng chính là Shell.

\par\vspace{0.3cm}
\hrule
\vspace{0.3cm}
\end{minipage}

\begin{minipage}{\textwidth}
\subsection*{Question 390 \quad \footnotesize{[0390\_q390\_Question\_390.png]}}
\noindent \textbf{Which of the following is NOT a necessary condition for deadlock to occur?}

\begin{enumerate}[label=\Alph*.]
    \item Mutual exclusion: Only one process can hold a resource at a time
    \item Hold and wait: A process can hold a resource while waiting for another resource
    \item \textbf{Preemption: A resource can be forcibly taken away from a process} \hfill \textcolor{blue}{\textbf{[Correct]}}
    \item Circular wait: A circular chain of processes exists, where each process is waiting for a resource held by the next process in the chain
\end{enumerate}

\noindent \textbf{Correct Answer:} \textcolor{blue}{\textbf{C}} \\[2pt]
\noindent \textbf{Explanation:} Theo lý thuyết hệ điều hành của Coffman, bế tắc (deadlock) chỉ có thể xảy ra khi hội tụ đủ 4 điều kiện cần: Loại trừ tương hỗ (Mutual exclusion), Giữ và chờ (Hold and wait), Không trưng dụng (No preemption - tài nguyên không thể bị thu hồi cưỡng bức), và Chờ đợi vòng tròn (Circular wait). Đặc tính 'Preemption' (cho phép tước đoạt tài nguyên cưỡng bức) là một giải pháp nhằm phá vỡ hoặc ngăn chặn deadlock chứ không phải là điều kiện để deadlock phát sinh (đáp án C). Vì vậy, Preemption không phải là điều kiện cần của deadlock.

\par\vspace{0.3cm}
\hrule
\vspace{0.3cm}
\end{minipage}

\begin{minipage}{\textwidth}
\subsection*{Question 391 \quad \footnotesize{[0391\_q391\_Question\_391.png]}}
\noindent \textbf{CÅU HòI In the context of demand paging and demand segmentation, a memory size of 100 MB can execute ten programs, each of size 30 MB, for a total of 300 MB. At any moment, 100 MB of the ten programs are in memory and 200 MB are on disk. There is therefore a virtual memory size of MB.}

\begin{enumerate}[label=\Alph*.]
    \item \textbf{300} \hfill \textcolor{blue}{\textbf{[Correct]}}
    \item 100
    \item 200
    \item 400
\end{enumerate}

\noindent \textbf{Correct Answer:} \textcolor{blue}{\textbf{A}} \\[2pt]
\noindent \textbf{Explanation:} Bộ nhớ ảo (Virtual Memory) là kỹ thuật mở rộng cho phép hệ thống thực thi các chương trình có tổng kích thước lớn hơn dung lượng bộ nhớ vật lý thực tế bằng cách kết hợp RAM và ổ đĩa. Khi 10 chương trình (mỗi chương trình 30 MB) cùng chạy đồng thời, tổng không gian địa chỉ nhớ mà hệ điều hành cung cấp cho các tiến trình này là 10 × 30 MB = 300 MB, trong đó 100 MB nằm trên RAM và 200 MB được lưu trữ trên đĩa (đáp án A). Do đó, dung lượng bộ nhớ ảo của hệ thống ở đây là 300 MB.

\par\vspace{0.3cm}
\hrule
\vspace{0.3cm}
\end{minipage}

\begin{minipage}{\textwidth}
\subsection*{Question 392 \quad \footnotesize{[0392\_q392\_Question\_392.png]}}
\noindent \textbf{In UNIX, which of the following commands is used to manage processes?}

\begin{enumerate}[label=\Alph*.]
    \item Is
    \item \textbf{ps} \hfill \textcolor{blue}{\textbf{[Correct]}}
    \item mv
    \item cat
\end{enumerate}

\noindent \textbf{Correct Answer:} \textcolor{blue}{\textbf{B}} \\[2pt]
\noindent \textbf{Explanation:} Trong hệ điều hành UNIX/Linux, lệnh `ps` (viết tắt của Process Status) là lệnh tiêu chuẩn dùng để kiểm tra, liệt kê và giám sát trạng thái của các tiến trình đang hoạt động trong hệ thống (đáp án B). Trong khi đó, lệnh `ls` dùng để liệt kê danh sách tệp/thư mục, `mv` dùng để di chuyển hoặc đổi tên tệp, và `cat` dùng để hiển thị hoặc nối nội dung tệp tin. Do đó, lệnh quản lý tiến trình là `ps`.

\par\vspace{0.3cm}
\hrule
\vspace{0.3cm}
\end{minipage}

\begin{minipage}{\textwidth}
\subsection*{Question 393 \quad \footnotesize{[0393\_q393\_Question\_393.png]}}
\noindent \textbf{Which file system is most commonly associated with Linux?}

\begin{enumerate}[label=\Alph*.]
    \item NTFS
    \item FAT32
    \item \textbf{ext4} \hfill \textcolor{blue}{\textbf{[Correct]}}
    \item HFS+
\end{enumerate}

\noindent \textbf{Correct Answer:} \textcolor{blue}{\textbf{C}} \\[2pt]
\noindent \textbf{Explanation:} Hệ thống tệp ext4 (Fourth Extended Filesystem) là hệ thống tệp nhật ký (journaling file system) tiêu chuẩn và phổ biến nhất được sử dụng mặc định trong hầu hết các bản phân phối Linux hiện đại (đáp án C). Trong khi đó, NTFS và FAT32 là các hệ thống tệp gắn liền với hệ điều hành Windows của Microsoft, còn HFS+ là hệ thống tệp truyền thống của macOS do Apple phát triển. Vì vậy, hệ thống tệp gắn liền nhất với Linux là ext4.

\par\vspace{0.3cm}
\hrule
\vspace{0.3cm}
\end{minipage}

\begin{minipage}{\textwidth}
\subsection*{Question 394 \quad \footnotesize{[0394\_q394\_Question\_394.png]}}
\noindent \textbf{What is the role of a data scientist in an Al project?}

\begin{enumerate}[label=\Alph*.]
    \item To design and build the Al algorithms.
    \item \textbf{To collect, clean, and prepare the data for analysis.} \hfill \textcolor{blue}{\textbf{[Correct]}}
    \item To deploy the Al model into production.
    \item To interpret the results of the Al model and communicate them to stakeholders.
\end{enumerate}

\noindent \textbf{Correct Answer:} \textcolor{blue}{\textbf{B}} \\[2pt]
\noindent \textbf{Explanation:} Trong một dự án trí tuệ nhân tạo (AI), nhiệm vụ trọng tâm và tốn nhiều thời gian nhất của một nhà khoa học dữ liệu (Data Scientist) là thu thập, làm sạch và chuẩn bị dữ liệu (collect, clean, and prepare data) nhằm đảm bảo dữ liệu đầu vào đạt chuẩn cho việc huấn luyện mô hình (đáp án B). Việc thiết kế các thuật toán học máy cốt lõi thường thuộc về các nhà nghiên cứu AI (AI Researcher), còn việc đóng gói và triển khai mô hình lên hệ thống vận hành thực tế (production) là trách nhiệm của kỹ sư Machine Learning/MLOps. Do đó, vai trò của data scientist được mô tả chính xác nhất ở phương án B.

\par\vspace{0.3cm}
\hrule
\vspace{0.3cm}
\end{minipage}

\begin{minipage}{\textwidth}
\subsection*{Question 395 \quad \footnotesize{[0395\_q395\_Question\_395.png]}}
\noindent \textbf{What is an Al program?}

\begin{enumerate}[label=\Alph*.]
    \item \textbf{ChatGPT} \hfill \textcolor{blue}{\textbf{[Correct]}}
    \item Microsoft Word
    \item React
    \item MongoDB
\end{enumerate}

\noindent \textbf{Correct Answer:} \textcolor{blue}{\textbf{A}} \\[2pt]
\noindent \textbf{Explanation:} ChatGPT (do OpenAI phát triển) là một ứng dụng trí tuệ nhân tạo (AI program) dựa trên các mô hình ngôn ngữ lớn (LLM), có khả năng xử lý ngôn ngữ tự nhiên, hiểu ngữ cảnh và hội thoại tương tác thông minh với con người (đáp án A). Ngược lại, Microsoft Word là phần mềm soạn thảo văn bản thông thường, React là thư viện JavaScript dùng để xây dựng giao diện web, và MongoDB là một hệ quản trị cơ sở dữ liệu phi quan hệ (NoSQL). Do đó, ChatGPT là chương trình AI duy nhất trong các lựa chọn.

\par\vspace{0.3cm}
\hrule
\vspace{0.3cm}
\end{minipage}

\begin{minipage}{\textwidth}
\subsection*{Question 396 \quad \footnotesize{[0396\_q396\_Question\_396.png]}}
\noindent \textbf{What is the most common language to write Al models?}

\begin{enumerate}[label=\Alph*.]
    \item C++
    \item PROLOG
    \item LISP
    \item \textbf{Python} \hfill \textcolor{blue}{\textbf{[Correct]}}
\end{enumerate}

\noindent \textbf{Correct Answer:} \textcolor{blue}{\textbf{D}} \\[2pt]
\noindent \textbf{Explanation:} Python là ngôn ngữ lập trình phổ biến nhất hiện nay để xây dựng và huấn luyện các mô hình trí tuệ nhân tạo (AI) nhờ cú pháp trực quan, dễ học cùng hệ sinh thái các thư viện máy học chuyên sâu vô cùng mạnh mẽ như TensorFlow, PyTorch và Scikit-learn (đáp án D). Mặc dù LISP và PROLOG là các ngôn ngữ AI kinh điển trong lịch sử và C++ thường được dùng để tối ưu hóa hiệu năng tầng thấp, Python vẫn là ngôn ngữ thống trị trong phát triển mô hình AI hiện đại.

\par\vspace{0.3cm}
\hrule
\vspace{0.3cm}
\end{minipage}

\begin{minipage}{\textwidth}
\subsection*{Question 397 \quad \footnotesize{[0397\_q397\_Question\_397.png]}}
\noindent \textbf{Which of the following is not an advantage of AI?}

\begin{enumerate}[label=\Alph*.]
    \item Reduces the time taken to perform a task
    \item Decision-making by making the process faster and smarter
    \item Increased Automation
    \item \textbf{The cost hight to develop, implement, and maintain Al systems} \hfill \textcolor{blue}{\textbf{[Correct]}}
\end{enumerate}

\noindent \textbf{Correct Answer:} \textcolor{blue}{\textbf{D}} \\[2pt]
\noindent \textbf{Explanation:} Chi phí nghiên cứu, phát triển, trang bị hạ tầng phần cứng chuyên dụng và bảo trì hệ thống AI rất đắt đỏ (The cost high to develop, implement, and maintain AI systems), đây là một trong những rào cản và nhược điểm lớn nhất của công nghệ này chứ không phải là một ưu điểm (đáp án D). Ngược lại, việc rút ngắn thời gian xử lý công việc (A), hỗ trợ ra quyết định nhanh hơn và chính xác hơn (B), cùng khả năng đẩy mạnh tự động hóa quy trình (C) đều là những lợi ích then chốt mà AI mang lại.

\par\vspace{0.3cm}
\hrule
\vspace{0.3cm}
\end{minipage}

\begin{minipage}{\textwidth}
\subsection*{Question 398 \quad \footnotesize{[0398\_q398\_Question\_398.png]}}
\noindent \textbf{How do vector graphics differ from raster graphics?}

\begin{enumerate}[label=\Alph*.]
    \item \textbf{Vector graphics store images as mathematical equations, while raster graphics store images as grids of pixels.} \hfill \textcolor{blue}{\textbf{[Correct]}}
    \item Raster graphics are faster to render than vector graphics.
    \item Vector graphics are limited to black and white, while raster supports color.
    \item Raster graphics are more scalable than vector graphics.
\end{enumerate}

\noindent \textbf{Correct Answer:} \textcolor{blue}{\textbf{A}} \\[2pt]
\noindent \textbf{Explanation:} Đồ họa vector (vector graphics) định nghĩa và lưu trữ hình ảnh dưới dạng các công thức và phương trình toán học biểu diễn điểm, đường và hình khối học, cho phép phóng to thu nhỏ tự do mà không bao giờ bị giảm chất lượng ảnh; trong khi đồ họa raster (raster graphics) lưu trữ ảnh dưới dạng một lưới các điểm ảnh (pixels) có độ phân giải cố định và sẽ bị mờ vỡ hạt khi phóng to (đáp án A). Cả hai định dạng đều hỗ trợ hiển thị màu sắc đầy đủ, và đồ họa vector có khả năng thay đổi kích thước (scalability) vượt trội hơn hẳn so với raster.

\par\vspace{0.3cm}
\hrule
\vspace{0.3cm}
\end{minipage}

\begin{minipage}{\textwidth}
\subsection*{Question 399 \quad \footnotesize{[0399\_q399\_Question\_399.png]}}
\noindent \textbf{Which of the following is the main challenge during the testing phase of the software lifecycle?}

\begin{enumerate}[label=\Alph*.]
    \item \textbf{Ensuring complete coverage of all possible test cases} \hfill \textcolor{blue}{\textbf{[Correct]}}
    \item Writing code that adheres to design specifications
    \item Gathering requirements from stakeholders
    \item Selecting the appropriate programming language
\end{enumerate}

\noindent \textbf{Correct Answer:} \textcolor{blue}{\textbf{A}} \\[2pt]
\noindent \textbf{Explanation:} Trong vòng đời phát triển phần mềm (SDLC), thách thức lớn nhất trong giai đoạn kiểm thử (testing phase) là việc bảo đảm bao phủ toàn bộ các trường hợp kiểm thử khả dĩ (Ensuring complete coverage of all possible test cases), do số lượng kịch bản đầu vào và đường dẫn thực thi tiềm năng là gần như vô hạn nên không thể kiểm thử cạn kiệt (exhaustive testing) (đáp án A). Các công việc như viết mã theo thiết kế thuộc pha cài đặt, thu thập yêu cầu thuộc pha phân tích, và lựa chọn ngôn ngữ thuộc pha thiết kế kiến trúc.

\par\vspace{0.3cm}
\hrule
\vspace{0.3cm}
\end{minipage}

\begin{minipage}{\textwidth}
\subsection*{Question 400 \quad \footnotesize{[0400\_q400\_Question\_400.png]}}
\noindent \textbf{Which software engineering principle advocates breaking down complex systems into smaller, manageable components that are easier to understand, develop, and maintain?}

\begin{enumerate}[label=\Alph*.]
    \item Abstraction
    \item Encapsulation
    \item \textbf{Modularity} \hfill \textcolor{blue}{\textbf{[Correct]}}
    \item Inheritance
\end{enumerate}

\noindent \textbf{Correct Answer:} \textcolor{blue}{\textbf{C}} \\[2pt]
\noindent \textbf{Explanation:} Tính mô-đun (Modularity) là nguyên lý kỹ thuật phần mềm cốt lõi ủng hộ việc phân chia một hệ thống phức tạp thành các thành phần (mô-đun) nhỏ hơn, có tính độc lập cao và dễ quản lý, giúp quá trình phát triển, kiểm thử, bảo trì và tái sử dụng mã nguồn trở nên hiệu quả hơn (đáp án C). Trong khi đó, Tính trừu tượng (Abstraction) tập trung vào việc ẩn giấu chi tiết cài đặt, Tính đóng gói (Encapsulation) bảo vệ dữ liệu nội bộ bên trong đối tượng, và Kế thừa (Inheritance) dùng để mở rộng lớp trong OOP.

\par\vspace{0.3cm}
\hrule
\vspace{0.3cm}
\end{minipage}

\begin{minipage}{\textwidth}
\subsection*{Question 401 \quad \footnotesize{[0401\_q401\_Question\_401.png]}}
\noindent \textbf{In the context of software design, what is the significance of the "Software Architecture"?}

\begin{enumerate}[label=\Alph*.]
    \item It focuses on coding individual functions and procedures.
    \item \textbf{It defines the overall structure and organization of the software system.} \hfill \textcolor{blue}{\textbf{[Correct]}}
    \item It is a tool for project management.
    \item It involves testing the software against user requirements.
\end{enumerate}

\noindent \textbf{Correct Answer:} \textcolor{blue}{\textbf{B}} \\[2pt]
\noindent \textbf{Explanation:} Trong thiết kế phần mềm, Kiến trúc phần mềm (Software Architecture) đóng vai trò là bản thiết kế tổng thể định hình cấu trúc, tổ chức các hệ thống con và cách thức giao tiếp, tương tác giữa các thành phần chính trong toàn bộ hệ thống (đáp án B). Nó không đi sâu vào việc viết mã cụ thể cho từng hàm/thủ tục (thuộc mức thiết kế chi tiết và lập trình), cũng không phải là một công cụ quản lý dự án hay giai đoạn kiểm thử phần mềm.

\par\vspace{0.3cm}
\hrule
\vspace{0.3cm}
\end{minipage}

\begin{minipage}{\textwidth}
\subsection*{Question 402 \quad \footnotesize{[0402\_q402\_Question\_402.png]}}
\noindent \textbf{Which of the following statements about records' data types is correct?}

\begin{enumerate}[label=\Alph*.]
    \item Records cannot contain other records as fields.
    \item Records are immutable data types.
    \item Records have a fixed size that cannot be changed at runtime.
    \item \textbf{Records can have fields of different data types.} \hfill \textcolor{blue}{\textbf{[Correct]}}
\end{enumerate}

\noindent \textbf{Correct Answer:} \textcolor{blue}{\textbf{D}} \\[2pt]
\noindent \textbf{Explanation:} Trong khoa học máy tính (theo giáo trình Foundations of Computer Science - Chương 11), bản ghi (record / struct) là một kiểu dữ liệu có cấu trúc bao gồm tập hợp các trường (fields), trong đó các trường có thể lưu trữ các kiểu dữ liệu khác nhau (heterogeneous data types), khác biệt với mảng (array) vốn chỉ chứa các phần tử có cùng kiểu dữ liệu (đáp án D). Bản ghi hoàn toàn có thể chứa các bản ghi con lồng nhau và các trường trong bản ghi thường có thể thay đổi giá trị (mutable) trong suốt quá trình chạy chương trình.

\par\vspace{0.3cm}
\hrule
\vspace{0.3cm}
\end{minipage}

\begin{minipage}{\textwidth}
\subsection*{Question 403 \quad \footnotesize{[0403\_q403\_Question\_403.png]}}
\noindent \textbf{What is the order of steps to design a database?}

\begin{enumerate}[label=\Alph*.]
    \item \textbf{Collect the information needed to be stored --> Build an entity-relation model --> Normalization} \hfill \textcolor{blue}{\textbf{[Correct]}}
    \item Collect the information needed to be stored --> Normalization --> Build an entity-relation model
    \item Build an entity-relation model --> Normalization --> Collect the information needed to be stored
    \item Not need be done through a step-by-step process
\end{enumerate}

\noindent \textbf{Correct Answer:} \textcolor{blue}{\textbf{A}} \\[2pt]
\noindent \textbf{Explanation:} Theo giáo trình Foundations of Computer Science của Behrouz Forouzan (Chương 11 - Databases), quy trình thiết kế cơ sở dữ liệu chuẩn gồm 3 bước tuần tự: (1) Thu thập thông tin cần lưu trữ (Collect the information), (2) Xây dựng mô hình thực thể - mối quan hệ (Build an entity-relationship model / ER model), và (3) Chuẩn hóa cơ sở dữ liệu (Normalization) để loại bỏ hiện tượng dư thừa và bất thường dữ liệu (đáp án A). Quá trình chuẩn hóa chỉ có thể thực hiện sau khi mô hình dữ liệu quan hệ đã được xác lập.

\par\vspace{0.3cm}
\hrule
\vspace{0.3cm}
\end{minipage}

\begin{minipage}{\textwidth}
\subsection*{Question 404 \quad \footnotesize{[0404\_q404\_Question\_404.png]}}
\noindent \textbf{What describes the hierarchical model?}

\begin{enumerate}[label=\Alph*.]
    \item Data is organized as a binary tree and each entity can have several children
    \item Data is organized as a binary tree and each entity have only one children
    \item \textbf{Data is organized as an inverted tree and each entity can have several children} \hfill \textcolor{blue}{\textbf{[Correct]}}
    \item Data is organized as an inverted tree and each entity have only one child
\end{enumerate}

\noindent \textbf{Correct Answer:} \textcolor{blue}{\textbf{C}} \\[2pt]
\noindent \textbf{Explanation:} Trong mô hình cơ sở dữ liệu phân cấp (hierarchical model), dữ liệu được tổ chức dưới dạng một cây đảo ngược (inverted tree), bắt đầu từ một nút gốc (root) ở trên cùng tỏa xuống các nút bên dưới. Trong cấu trúc này, mỗi nút cha có thể liên kết và quản lý nhiều nút con (several children), nhưng mỗi nút con chỉ có duy nhất một nút cha đại diện cho mối quan hệ một - nhiều (1-to-N) (đáp án C). Đây là cấu trúc cây tổng quát chứ không bị giới hạn chỉ có 2 con như cây nhị phân (binary tree).

\par\vspace{0.3cm}
\hrule
\vspace{0.3cm}
\end{minipage}

\begin{minipage}{\textwidth}
\subsection*{Question 405 \quad \footnotesize{[0405\_q405\_Question\_405.png]}}
\noindent \textbf{Which of the following is NOT one of the three traditional database models?}

\begin{enumerate}[label=\Alph*.]
    \item Hierarchical
    \item Networking
    \item Relational
    \item \textbf{Object-oriented} \hfill \textcolor{blue}{\textbf{[Correct]}}
\end{enumerate}

\noindent \textbf{Correct Answer:} \textcolor{blue}{\textbf{D}} \\[2pt]
\noindent \textbf{Explanation:} Theo giáo trình của Behrouz Forouzan (Chương 11), ba mô hình cơ sở dữ liệu truyền thống (traditional database models) bao gồm: Mô hình phân cấp (Hierarchical model), Mô hình mạng (Network model), và Mô hình quan hệ (Relational model). Mô hình hướng đối tượng (Object-oriented model) ra đời muộn hơn nhằm đáp ứng sự phát triển của lập trình hướng đối tượng, thuộc nhóm mô hình thế hệ mới chứ không nằm trong ba mô hình truyền thống (đáp án D).

\par\vspace{0.3cm}
\hrule
\vspace{0.3cm}
\end{minipage}

\begin{minipage}{\textwidth}
\subsection*{Question 406 \quad \footnotesize{[0406\_q406\_Question\_406.png]}}
\noindent \textbf{In a relational database, which of the following is a unary operation ?}

\begin{enumerate}[label=\Alph*.]
    \item union
    \item join
    \item \textbf{project} \hfill \textcolor{blue}{\textbf{[Correct]}}
    \item all of these above
\end{enumerate}

\noindent \textbf{Correct Answer:} \textcolor{blue}{\textbf{C}} \\[2pt]
\noindent \textbf{Explanation:} Trong đại số quan hệ (relational algebra) của cơ sở dữ liệu, phép toán một ngôi (unary operation) là phép toán chỉ tác động lên duy nhất một bảng quan hệ đầu vào để tạo ra bảng kết quả mới; phép chiếu (project) là phép toán một ngôi điển hình dùng để trích chọn một số cột thuộc tính từ bảng quan hệ (đáp án C). Ngược lại, phép hợp (union) và phép kết nối (join) là các phép toán hai ngôi (binary operations) vì bắt buộc phải có hai bảng quan hệ đầu vào để thực hiện.

\par\vspace{0.3cm}
\hrule
\vspace{0.3cm}
\end{minipage}

\begin{minipage}{\textwidth}
\subsection*{Question 407 \quad \footnotesize{[0407\_q407\_Question\_407.png]}}
\noindent \textbf{Unified Modeling Language (UML) is}

\begin{enumerate}[label=\Alph*.]
    \item A high level language
    \item A machine language
    \item An English-language-like representation of code
    \item \textbf{A pictorial representation of an algorithm} \hfill \textcolor{blue}{\textbf{[Correct]}}
\end{enumerate}

\noindent \textbf{Correct Answer:} \textcolor{blue}{\textbf{D}} \\[2pt]
\noindent \textbf{Explanation:} Ngôn ngữ mô hình hóa thống nhất UML (Unified Modeling Language) là ngôn ngữ chuẩn hóa sử dụng các sơ đồ và biểu đồ trực quan (pictorial representation) để mô tả cấu trúc, thiết kế các thành phần và biểu diễn các luồng thuật toán của một hệ thống phần mềm (đáp án D). UML là công cụ thiết kế và mô hình hóa hình ảnh, không phải là một ngôn ngữ lập trình bậc cao hay ngôn ngữ máy dùng để thực thi trực tiếp trên máy tính.

\par\vspace{0.3cm}
\hrule
\vspace{0.3cm}
\end{minipage}

\begin{minipage}{\textwidth}
\subsection*{Question 408 \quad \footnotesize{[0408\_q408\_Question\_408.png]}}
\noindent \textbf{In which case would an iterative algorithm be preferred over a recursive algorithm?}

\begin{enumerate}[label=\Alph*.]
    \item When solving problems that require backtracking
    \item \textbf{When stack overflow is a concern} \hfill \textcolor{blue}{\textbf{[Correct]}}
    \item When the problem involves tree traversal
    \item When solving problems with divide-and-conquer techniques
\end{enumerate}

\noindent \textbf{Correct Answer:} \textcolor{blue}{\textbf{B}} \\[2pt]
\noindent \textbf{Explanation:} Thuật toán đệ quy (recursion) sử dụng ngăn xếp lời gọi hàm (call stack) của hệ thống để lưu trữ trạng thái qua mỗi lần gọi hàm, nên khi số tầng đệ quy quá sâu sẽ dễ dẫn đến hiện tượng tràn bộ nhớ ngăn xếp (stack overflow). Do đó, thuật toán lặp (iterative algorithm) được ưu tiên sử dụng khi nguy cơ tràn ngăn xếp là mối lo ngại lớn (When stack overflow is a concern), vì cấu trúc lặp sử dụng bộ nhớ ngăn xếp cố định O(1) và kiểm soát bộ nhớ hiệu quả hơn (đáp án B). Các kỹ thuật như quay lui hay duyệt cây thường được biểu diễn tự nhiên hơn bằng đệ quy.

\par\vspace{0.3cm}
\hrule
\vspace{0.3cm}
\end{minipage}

\begin{minipage}{\textwidth}
\subsection*{Question 409 \quad \footnotesize{[0409\_q409\_Question\_409.png]}}
\noindent \textbf{What distinguishes the Internet from an internet?}

\begin{enumerate}[label=\Alph*.]
    \item \textbf{The Internet is a global network, whereas an internet is any interconnected network.} \hfill \textcolor{blue}{\textbf{[Correct]}}
    \item The Internet uses proprietary protocols, while an internet uses TCP/IP.
    \item The Internet is limited to government organizations, and an internet is public.
    \item An internet is faster than the Internet.
\end{enumerate}

\noindent \textbf{Correct Answer:} \textcolor{blue}{\textbf{A}} \\[2pt]
\noindent \textbf{Explanation:} Trong thuật ngữ mạng máy tính (theo giáo trình Forouzan), 'the Internet' (viết hoa chữ 'I') là một danh từ riêng chỉ mạng liên mạng toàn cầu duy nhất kết nối hàng triệu mạng công cộng và tư nhân trên thế giới thông qua bộ giao thức TCP/IP; trong khi 'an internet' (viết chữ 'i' thường) là danh từ chung chỉ bất kỳ hệ thống mạng nào được hình thành do việc kết nối hai hay nhiều mạng riêng biệt lại với nhau (đáp án A). Mạng Internet mở rộng cho toàn xã hội chứ không chỉ giới hạn cho các cơ quan chính phủ, và cả hai đều có thể sử dụng các giao thức tiêu chuẩn.

\par\vspace{0.3cm}
\hrule
\vspace{0.3cm}
\end{minipage}

\begin{minipage}{\textwidth}
\subsection*{Question 410 \quad \footnotesize{[0410\_q410\_Question\_410.png]}}
\noindent \textbf{Which of the following statements about binary search is true?}

\begin{enumerate}[label=\Alph*.]
    \item Binary search can be used on unsorted data
    \item Binary search divides the search space into three parts
    \item \textbf{Binary search is faster than linear search for sorted data} \hfill \textcolor{blue}{\textbf{[Correct]}}
    \item Binary search has a time complexity of O(n)
\end{enumerate}

\noindent \textbf{Correct Answer:} \textcolor{blue}{\textbf{C}} \\[2pt]
\noindent \textbf{Explanation:} Thuật toán tìm kiếm nhị phân (Binary search) trên mảng đã sắp xếp có độ phức tạp thời gian là O(log n), nhanh hơn vượt trội so với tìm kiếm tuyến tính (Linear search) vốn có độ phức tạp là O(n) trên các tập dữ liệu lớn (đáp án C). Điều kiện tiên quyết bắt buộc của tìm kiếm nhị phân là dữ liệu phải được sắp xếp trước (không dùng được trên dữ liệu lộn xộn), và ở mỗi bước thuật toán chia đôi không gian tìm kiếm thành hai nửa chứ không phải ba phần.

\par\vspace{0.3cm}
\hrule
\vspace{0.3cm}
\end{minipage}

\begin{minipage}{\textwidth}
\subsection*{Question 411 \quad \footnotesize{[0411\_q411\_Question\_411.png]}}
\noindent \textbf{What is the primary difference between sequential search and binary search algorithms?}

\begin{enumerate}[label=\Alph*.]
    \item Sequential search requires the data to be sorted in ascending order, while binary search can work on unsorted data as well
    \item \textbf{Binary search divides the search space into halves, while sequential search checks elements one by one} \hfill \textcolor{blue}{\textbf{[Correct]}}
    \item Sequential search exhibits greater efficiency and speed compared to binary search when applied to large datasets containing a significant number of elements
    \item Binary search is a versatile algorithm that can be applied to various types of data structures, including arrays, linked lists, and trees, while sequential search is specifically designed to work only with arrays
\end{enumerate}

\noindent \textbf{Correct Answer:} \textcolor{blue}{\textbf{B}} \\[2pt]
\noindent \textbf{Explanation:} Điểm khác biệt cốt lõi nhất giữa hai thuật toán tìm kiếm là: Tìm kiếm nhị phân (Binary search) so sánh phần tử tìm kiếm với phần tử nằm ở giữa và chia đôi không gian tìm kiếm thành hai nửa sau mỗi bước kiểm tra; trong khi tìm kiếm tuần tự (Sequential search) tiến hành kiểm tra lần lượt từng phần tử một từ đầu đến cuối danh sách (đáp án B). Tìm kiếm tuần tự có thể áp dụng cho dữ liệu chưa sắp xếp và có độ phức tạp O(n), kém hiệu quả hơn nhiều so với O(log n) của tìm kiếm nhị phân trên các tập dữ liệu lớn.

\par\vspace{0.3cm}
\hrule
\vspace{0.3cm}
\end{minipage}

\begin{minipage}{\textwidth}
\subsection*{Question 412 \quad \footnotesize{[0412\_q412\_Question\_412.png]}}
\noindent \textbf{How is a linked list typically implemented?}

\begin{enumerate}[label=\Alph*.]
    \item Using arrays
    \item Using stacks
    \item Using queues
    \item \textbf{Using pointers} \hfill \textcolor{blue}{\textbf{[Correct]}}
\end{enumerate}

\noindent \textbf{Correct Answer:} \textcolor{blue}{\textbf{D}} \\[2pt]
\noindent \textbf{Explanation:} Danh sách liên kết (Linked list) là cấu trúc dữ liệu tuyến tính động được tạo thành từ tập hợp các nút (nodes), trong đó mỗi nút chứa dữ liệu và một hoặc nhiều trường liên kết sử dụng con trỏ (pointers / references) để trỏ đến địa chỉ bộ nhớ của nút kế tiếp (đáp án D). Cơ chế sử dụng con trỏ giúp danh sách liên kết có thể cấp phát bộ nhớ động một cách linh hoạt mà không đòi hỏi các ô nhớ phải nằm liên tiếp nhau như mảng thông thường.

\par\vspace{0.3cm}
\hrule
\vspace{0.3cm}
\end{minipage}

\begin{minipage}{\textwidth}
\subsection*{Question 413 \quad \footnotesize{[0413\_q413\_Question\_413.png]}}
\noindent \textbf{What is the time complexity of finding the height of a binary tree with n nodes?}

\begin{enumerate}[label=\Alph*.]
    \item \textbf{O(n)} \hfill \textcolor{blue}{\textbf{[Correct]}}
    \item O(log n)
    \item O(n²)
    \item O(n log n)
\end{enumerate}

\noindent \textbf{Correct Answer:} \textcolor{blue}{\textbf{A}} \\[2pt]
\noindent \textbf{Explanation:} Để xác định chiều cao của một cây nhị phân bất kỳ gồm n nút, thuật toán (thường được cài đặt bằng đệ quy duyệt cây) bắt buộc phải thăm qua toàn bộ tất cả n nút của cây để tìm ra đường đi dài nhất từ nút gốc đến nút lá. Do mỗi nút trong cây đều được duyệt qua một số lần cố định, độ phức tạp thời gian tổng quát của thuật toán là O(n) (đáp án A). Giá trị O(log n) chỉ là chiều cao của một cây nhị phân cân bằng hoàn hảo, không phải là thời gian tính toán của thuật toán duyệt tìm chiều cao.

\par\vspace{0.3cm}
\hrule
\vspace{0.3cm}
\end{minipage}

\begin{minipage}{\textwidth}
\subsection*{Question 414 \quad \footnotesize{[0414\_q414\_Question\_414.png]}}
\noindent \textbf{Which method is used to add an element to a stack?}

\begin{enumerate}[label=\Alph*.]
    \item \textbf{push()} \hfill \textcolor{blue}{\textbf{[Correct]}}
    \item pop()
    \item delete()
    \item top()
\end{enumerate}

\noindent \textbf{Correct Answer:} \textcolor{blue}{\textbf{A}} \\[2pt]
\noindent \textbf{Explanation:} Trong cấu trúc dữ liệu ngăn xếp (Stack), phương thức `push()` là thao tác tiêu chuẩn được sử dụng để thêm một phần tử mới vào vị trí đỉnh của ngăn xếp (đáp án A). Ngược lại, phương thức `pop()` dùng để loại bỏ và lấy phần tử ở đỉnh ra khỏi ngăn xếp, còn `top()` (hoặc `peek()`) dùng để xem giá trị của phần tử trên đỉnh mà không làm thay đổi cấu trúc của ngăn xếp.

\par\vspace{0.3cm}
\hrule
\vspace{0.3cm}
\end{minipage}

\begin{minipage}{\textwidth}
\subsection*{Question 415 \quad \footnotesize{[0415\_q415\_Question\_415.png]}}
\noindent \textbf{What is the maximum number of children a node in a binary tree can have?}

\begin{enumerate}[label=\Alph*.]
    \item 1
    \item \textbf{2} \hfill \textcolor{blue}{\textbf{[Correct]}}
    \item 3
    \item Any number of children
\end{enumerate}

\noindent \textbf{Correct Answer:} \textcolor{blue}{\textbf{B}} \\[2pt]
\noindent \textbf{Explanation:} Theo định nghĩa chuẩn trong khoa học máy tính, cây nhị phân (binary tree - tiền tố 'bi' biểu thị số hai) là một cấu trúc dữ liệu cây trong đó mỗi nút chỉ có thể có tối đa 2 nút con, thường được gọi là nút con bên trái (left child) và nút con bên phải (right child) (đáp án B). Cây mà mỗi nút có thể chứa số lượng con tùy ý được gọi là cây tổng quát (general tree), còn cây có tối đa 3 nút con được gọi là cây tam phân (ternary tree).

\par\vspace{0.3cm}
\hrule
\vspace{0.3cm}
\end{minipage}

\begin{minipage}{\textwidth}
\subsection*{Question 416 \quad \footnotesize{[0416\_q416\_Question\_416.png]}}
\noindent \textbf{CÀU HÒI What is the difference between a stack and a queue?}

\begin{enumerate}[label=\Alph*.]
    \item A stack follows FIFO, and a queue follows LIFO
    \item A stack is a hierarchical data structure, and a queue is not
    \item \textbf{A stack is LIFO, and a queue is FIFO} \hfill \textcolor{blue}{\textbf{[Correct]}}
    \item There is no difference between a stack and a queue
\end{enumerate}

\noindent \textbf{Correct Answer:} \textcolor{blue}{\textbf{C}} \\[2pt]
\noindent \textbf{Explanation:} Sự khác biệt căn bản nhất giữa ngăn xếp (stack) và hàng đợi (queue) là cơ chế thêm và trích xuất dữ liệu: Ngăn xếp hoạt động theo nguyên lý LIFO (Last In, First Out - vào sau ra trước), tức phần tử đưa vào sau cùng sẽ được lấy ra đầu tiên; trong khi hàng đợi hoạt động theo nguyên lý FIFO (First In, First Out - vào trước ra trước), tức phần tử đưa vào trước sẽ được phục vụ và lấy ra trước (đáp án C). Cả hai đều là các cấu trúc dữ liệu tuyến tính chứ không phải cấu trúc phân cấp.

\par\vspace{0.3cm}
\hrule
\vspace{0.3cm}
\end{minipage}

\begin{minipage}{\textwidth}
\subsection*{Question 417 \quad \footnotesize{[0417\_q417\_Question\_417.png]}}
\noindent \textbf{How does a queue differ from a stack?}

\begin{enumerate}[label=\Alph*.]
    \item A queue operates on a LIFO principle, while a stack operates on a FIFO principle.
    \item \textbf{A queue operates on a FIFO principle, while a stack operates on a LIFO principle.} \hfill \textcolor{blue}{\textbf{[Correct]}}
    \item A queue can only store characters, while a stack can store any data type.
    \item A queue is used for recursion, while a stack is used for sorting.
\end{enumerate}

\noindent \textbf{Correct Answer:} \textcolor{blue}{\textbf{B}} \\[2pt]
\noindent \textbf{Explanation:} Hàng đợi (Queue) tuân theo nguyên lý FIFO (First-In, First-Out - vào trước ra trước), trong đó các phần tử được thêm vào ở phần đuôi (rear) và được lấy ra ở phần đầu (front), tương tự như hành vi xếp hàng trong thực tế. Ngược lại, ngăn xếp (Stack) tuân theo nguyên lý LIFO (Last-In, First-Out - vào sau ra trước), trong đó mọi thao tác đẩy vào và rút ra đều chỉ thực hiện tại đỉnh (top) (đáp án B). Cả hàng đợi và ngăn xếp đều là cấu trúc dữ liệu tổng quát có thể lưu trữ bất kỳ kiểu dữ liệu nào.

\par\vspace{0.3cm}
\hrule
\vspace{0.3cm}
\end{minipage}

\begin{minipage}{\textwidth}
\subsection*{Question 418 \quad \footnotesize{[0418\_q418\_Question\_418.png]}}
\noindent \textbf{What is the main benefit of using the .jpg (JPEG) format for images?}

\begin{enumerate}[label=\Alph*.]
    \item Lossless compression
    \item High-quality video storage
    \item \textbf{Lossy compression, which reduces file size} \hfill \textcolor{blue}{\textbf{[Correct]}}
    \item Support for animation
\end{enumerate}

\noindent \textbf{Correct Answer:} \textcolor{blue}{\textbf{C}} \\[2pt]
\noindent \textbf{Explanation:} Định dạng hình ảnh .jpg (JPEG) sử dụng phương pháp nén có tổn hao dữ liệu (lossy compression) để loại bỏ các thông tin màu sắc và chi tiết dư thừa mà mắt người khó nhận biết, từ đó giảm thiểu đáng kể dung lượng tệp tin giúp lưu trữ và truyền tải nhanh chóng trên mạng (đáp án C). Nén không tổn hao (lossless) là đặc trưng của định dạng PNG, khả năng hiển thị ảnh động (animation) thuộc về định dạng GIF, và JPEG chỉ áp dụng cho ảnh tĩnh chứ không dùng để lưu trữ video.

\par\vspace{0.3cm}
\hrule
\vspace{0.3cm}
\end{minipage}


\newpage
\section*{12. Security and Ethical Issues}
\addcontentsline{toc}{section}{12. Security and Ethical Issues}
\hrule\vspace{0.5cm}

\begin{minipage}{\textwidth}
\subsection*{Question 419 \quad \footnotesize{[0419\_q419\_Question\_419.png]}}
\noindent \textbf{What is the result of the AND logic operation when applied to two bits, one with the value 1 and the other with the value 0?}

\begin{enumerate}[label=\Alph*.]
    \item \textbf{0} \hfill \textcolor{blue}{\textbf{[Correct]}}
    \item 1
    \item -1
    \item 2
\end{enumerate}

\noindent \textbf{Correct Answer:} \textcolor{blue}{\textbf{A}} \\[2pt]
\noindent \textbf{Explanation:} Phép toán logic AND chỉ trả về giá trị 1 khi và chỉ khi cả hai bit đầu vào đều bằng 1; nếu có ít nhất một bit bằng 0 thì kết quả luôn bằng 0. Do đó, khi áp dụng phép AND cho hai bit có giá trị 1 và 0, kết quả thu được là 0 (đáp án A). Các giá trị -1 và 2 (C và D) không phải là các giá trị bit nhị phân hợp lệ, còn đáp án 1 (B) chỉ xảy ra khi cả hai toán hạng đều là 1.

\par\vspace{0.3cm}
\hrule
\vspace{0.3cm}
\end{minipage}

\begin{minipage}{\textwidth}
\subsection*{Question 420 \quad \footnotesize{[0420\_q420\_Question\_420.png]}}
\noindent \textbf{Which operation may change the sign (leftmost) bit? SÖ DÁP ÁN CÃN CHON}

\begin{enumerate}[label=\Alph*.]
    \item Arithmetic right shift
    \item \textbf{Arithmetic left shift} \hfill \textcolor{blue}{\textbf{[Correct]}}
    \item Logical left shift
    \item None of these above
\end{enumerate}

\noindent \textbf{Correct Answer:} \textcolor{blue}{\textbf{B}} \\[2pt]
\noindent \textbf{Explanation:} Trong biểu diễn số nguyên có dấu (Two's complement), bit ngoài cùng bên trái (leftmost bit) đóng vai trò là bit dấu (0 là dương, 1 là âm). Phép dịch trái số học (Arithmetic left shift) dịch chuyển toàn bộ các bit sang trái và điền bit 0 vào vị trí ngoài cùng bên phải, điều này có thể đẩy bit dữ liệu vào vị trí bit dấu và làm thay đổi dấu của số (gây ra hiện tượng tràn số - overflow). Ngược lại, phép dịch phải số học (Arithmetic right shift) luôn nhân bản bit dấu ban đầu để bảo toàn dấu của số.

\par\vspace{0.3cm}
\hrule
\vspace{0.3cm}
\end{minipage}

\begin{minipage}{\textwidth}
\subsection*{Question 421 \quad \footnotesize{[0421\_q421\_Question\_421.png]}}
\noindent \textbf{Which is a unary operator?}

\begin{enumerate}[label=\Alph*.]
    \item AND
    \item OR
    \item XOR
    \item \textbf{NOT} \hfill \textcolor{blue}{\textbf{[Correct]}}
\end{enumerate}

\noindent \textbf{Correct Answer:} \textcolor{blue}{\textbf{D}} \\[2pt]
\noindent \textbf{Explanation:} Toán tử một ngôi (unary operator) là toán tử chỉ tác động lên một toán hạng duy nhất để tạo ra kết quả. Trong các phép toán logic cơ bản, NOT là toán tử một ngôi vì nó chỉ cần một bit đầu vào để thực hiện phép nghịch đảo giá trị (0 thành 1 và 1 thành 0). Ngược lại, các toán tử AND, OR và XOR đều là toán tử hai ngôi (binary operators) vì bắt buộc phải có hai toán hạng đầu vào.

\par\vspace{0.3cm}
\hrule
\vspace{0.3cm}
\end{minipage}

\begin{minipage}{\textwidth}
\subsection*{Question 422 \quad \footnotesize{[0422\_q422\_Question\_422.png]}}
\noindent \textbf{How does a computer system interpret a sequence of bits to represent a symbol in a character encoding scheme? SÓ DÁP ÁN CÀN CHN}

\begin{enumerate}[label=\Alph*.]
    \item \textbf{By following a predefined code table to determine the symbol} \hfill \textcolor{blue}{\textbf{[Correct]}}
    \item By mapping each bit to a corresponding word in a dictionary
    \item By interpreting bits as colors in a display
    \item By using random bit sequences to represent symbols
    \item None of the others
\end{enumerate}

\noindent \textbf{Correct Answer:} \textcolor{blue}{\textbf{A}} \\[2pt]
\noindent \textbf{Explanation:} Trong các hệ thống máy tính, việc mã hóa ký tự (character encoding scheme như ASCII hay Unicode) được thực hiện bằng cách ánh xạ chuỗi bit sang ký tự tương ứng dựa trên một bảng mã được định nghĩa trước (predefined code table). Ví dụ, theo bảng mã ASCII chuẩn, chuỗi 7 bit 1000001 luôn được máy tính tra cứu và hiểu là ký tự 'A'. Máy tính không dùng chuỗi bit ngẫu nhiên hay ánh xạ từng bit thành một từ trong từ điển để biểu diễn ký tự.

\par\vspace{0.3cm}
\hrule
\vspace{0.3cm}
\end{minipage}

\begin{minipage}{\textwidth}
\subsection*{Question 423 \quad \footnotesize{[0423\_q423\_Question\_423.png]}}
\noindent \textbf{In a steganographic algorithm, if each comma (,) in the encrypted message equals digit 0 and each dot (.) in the message equals digit 1. What is the hidden information in following message? "I am a human, a student. I studied hard, so I got a good score."}

\begin{enumerate}[label=\Alph*.]
    \item \textbf{0101} \hfill \textcolor{blue}{\textbf{[Correct]}}
    \item 1010
    \item 0011
    \item 1100
\end{enumerate}

\noindent \textbf{Correct Answer:} \textcolor{blue}{\textbf{A}} \\[2pt]
\noindent \textbf{Explanation:} Kỹ thuật giấu tin (steganography) là phương pháp che giấu thông điệp bí mật bên trong một đối tượng dữ liệu bình thường. Trong câu văn đã cho: 'I am a human, a student. I studied hard, so I got a good score.', thứ tự xuất hiện của các dấu câu lần lượt là: dấu phẩy (,) sau 'human' ứng với bit 0, dấu chấm (.) sau 'student' ứng với bit 1, dấu phẩy (,) sau 'hard' ứng với bit 0, và dấu chấm (.) sau 'score' ứng với bit 1. Do đó, chuỗi thông tin ẩn được trích xuất là 0101 (đáp án A).

\par\vspace{0.3cm}
\hrule
\vspace{0.3cm}
\end{minipage}

\begin{minipage}{\textwidth}
\subsection*{Question 424 \quad \footnotesize{[0424\_q424\_Question\_424.png]}}
\noindent \textbf{When employing symmetric-key encryption for data communication between 10 hosts, how many keys are necessary to guarantee the security of all transactions?}

\begin{enumerate}[label=\Alph*.]
    \item \textbf{45} \hfill \textcolor{blue}{\textbf{[Correct]}}
    \item 46
    \item 47
    \item 48
\end{enumerate}

\noindent \textbf{Correct Answer:} \textcolor{blue}{\textbf{A}} \\[2pt]
\noindent \textbf{Explanation:} Trong hệ mật mã hóa khóa đối xứng (symmetric-key encryption), mỗi cặp máy chủ giao tiếp riêng biệt cần chia sẻ một khóa bí mật chung. Với hệ thống gồm n máy chủ kết nối với nhau, tổng số khóa bí mật cần thiết để đảm bảo an toàn cho mọi giao dịch được tính theo công thức tổ hợp N = n(n - 1) / 2. Áp dụng với n = 10 máy chủ, số khóa cần thiết là 10 × (10 - 1) / 2 = 45 khóa (đáp án A).

\par\vspace{0.3cm}
\hrule
\vspace{0.3cm}
\end{minipage}

\begin{minipage}{\textwidth}
\subsection*{Question 425 \quad \footnotesize{[0425\_q425\_Question\_425.png]}}
\noindent \textbf{Which security goal focuses on protecting data from unauthorized access or disclosure?}

\begin{enumerate}[label=\Alph*.]
    \item \textbf{Confidentiality} \hfill \textcolor{blue}{\textbf{[Correct]}}
    \item Integrity
    \item Availability
    \item Non-repudiation
\end{enumerate}

\noindent \textbf{Correct Answer:} \textcolor{blue}{\textbf{A}} \\[2pt]
\noindent \textbf{Explanation:} Tính bảo mật / bí mật (Confidentiality) là mục tiêu an ninh mạng tập trung vào việc ngăn chặn sự truy cập, tiết lộ dữ liệu trái phép và đảm bảo chỉ những đối tượng được cấp quyền mới có thể đọc thông tin. Trong khi đó, Tính toàn vẹn (Integrity) bảo vệ dữ liệu không bị sửa đổi trái phép, còn Tính sẵn sàng (Availability) đảm bảo dữ liệu luôn có thể truy cập khi người dùng hợp lệ yêu cầu. Vì vậy, mục tiêu bảo vệ thông tin trước truy cập hoặc tiết lộ trái phép chính là Confidentiality (đáp án A).

\par\vspace{0.3cm}
\hrule
\vspace{0.3cm}
\end{minipage}

\begin{minipage}{\textwidth}
\subsection*{Question 426 \quad \footnotesize{[0426\_q426\_Question\_426.png]}}
\noindent \textbf{How many ethical principles are there?}

\begin{enumerate}[label=\Alph*.]
    \item 2
    \item \textbf{3} \hfill \textcolor{blue}{\textbf{[Correct]}}
    \item 4
    \item 5
    \item 1
\end{enumerate}

\noindent \textbf{Correct Answer:} \textcolor{blue}{\textbf{B}} \\[2pt]
\noindent \textbf{Explanation:} Trong giáo trình Foundations of Computer Science của Behrouz Forouzan (Chương 12 về Đạo đức và Bảo mật), các nguyên tắc đạo đức lý thuyết cơ bản trong đạo đức máy tính được quy về 3 nguyên tắc chính (virtue ethics, deontological ethics, và teleological/utilitarian ethics). Các nguyên tắc này cung cấp nền tảng để đánh giá hành vi và đưa ra các quyết định chuẩn mực trong lĩnh vực công nghệ thông tin. Do đó, số lượng nguyên tắc đạo đức được trình bày là 3 (đáp án B).

\par\vspace{0.3cm}
\hrule
\vspace{0.3cm}
\end{minipage}

\begin{minipage}{\textwidth}
\subsection*{Question 427 \quad \footnotesize{[0427\_q427\_Question\_427.png]}}
\noindent \textbf{Which principle of OOP encourages bundling data and methods that operate on the data into a single unit, known as a class?}

\begin{enumerate}[label=\Alph*.]
    \item Inheritance
    \item \textbf{Encapsulation} \hfill \textcolor{blue}{\textbf{[Correct]}}
    \item Polymorphism
    \item Abstraction
\end{enumerate}

\noindent \textbf{Correct Answer:} \textcolor{blue}{\textbf{B}} \\[2pt]
\noindent \textbf{Explanation:} Tính đóng gói (Encapsulation) là một trong những nguyên lý nền tảng của lập trình hướng đối tượng (OOP), cho phép gom nhóm dữ liệu (thuộc tính) và các phương thức xử lý trên dữ liệu đó vào chung một đơn vị cấu trúc gọi là lớp (class). Cơ chế này giúp bảo vệ trạng thái nội bộ của đối tượng khỏi sự can thiệp trực tiếp từ bên ngoài bằng cách ẩn chi tiết cài đặt (data hiding). Ngược lại, Tính kế thừa (Inheritance) dùng để tái sử dụng mã nguồn, còn Tính đa hình (Polymorphism) cho phép đối tượng thực thi cùng một hành vi theo nhiều cách khác nhau.

\par\vspace{0.3cm}
\hrule
\vspace{0.3cm}
\end{minipage}

\begin{minipage}{\textwidth}
\subsection*{Question 428 \quad \footnotesize{[0428\_q428\_Question\_428.png]}}
\noindent \textbf{What is an example of a declarative programming language?}

\begin{enumerate}[label=\Alph*.]
    \item Python
    \item \textbf{Prolog} \hfill \textcolor{blue}{\textbf{[Correct]}}
    \item SQL
    \item Java
\end{enumerate}

\noindent \textbf{Correct Answer:} \textcolor{blue}{\textbf{B}} \\[2pt]
\noindent \textbf{Explanation:} Lập trình khai báo (declarative programming) là mô hình lập trình tập trung vào việc mô tả mục tiêu cần đạt được (kết quả mong muốn) thay vì chỉ rõ từng bước thuật toán chi tiết để đạt được kết quả đó. Prolog là ngôn ngữ lập trình tiêu biểu nhất đại diện cho mô hình khai báo/logic, nơi lập trình viên chỉ định các sự kiện (facts) và quy tắc suy luận (rules) để hệ thống tự động suy ra kết quả. Trong khi đó, Python và Java chủ yếu thuộc mô hình mệnh lệnh và hướng đối tượng (imperative/OOP).

\par\vspace{0.3cm}
\hrule
\vspace{0.3cm}
\end{minipage}

\begin{minipage}{\textwidth}
\subsection*{Question 429 \quad \footnotesize{[0429\_q429\_Question\_429.png]}}
\noindent \textbf{In the procedural paradigm, how are data and program units (functions) typically related?}

\begin{enumerate}[label=\Alph*.]
    \item Functions and data are separate entities with no direct relationship.
    \item \textbf{Functions operate directly on data passed to them as parameters.} \hfill \textcolor{blue}{\textbf{[Correct]}}
    \item Functions encapsulate both data and behavior.
    \item Data and functions are globally shared without restrictions.
\end{enumerate}

\noindent \textbf{Correct Answer:} \textcolor{blue}{\textbf{B}} \\[2pt]
\noindent \textbf{Explanation:} Trong mô hình lập trình thủ tục (procedural paradigm), chương trình được chia thành các đơn vị hàm/thủ tục và dữ liệu tách biệt nhau, trong đó các hàm nhận dữ liệu được truyền vào dưới dạng tham số (parameters) để xử lý. Dữ liệu không được đóng gói kèm theo hành vi như trong mô hình hướng đối tượng, mà được luân chuyển giữa các hàm thông qua các đối số. Vì vậy, đặc điểm các hàm thao tác trực tiếp trên dữ liệu truyền qua tham số là hoàn toàn chính xác (đáp án B).

\par\vspace{0.3cm}
\hrule
\vspace{0.3cm}
\end{minipage}

\begin{minipage}{\textwidth}
\subsection*{Question 430 \quad \footnotesize{[0430\_q430\_Question\_430.png]}}
\noindent \textbf{The Arithmetic Logic Unit (ALU) in the von Neumann model is primarily responsible for:}

\begin{enumerate}[label=\Alph*.]
    \item Controlling the execution of a program
    \item \textbf{Performing arithmetic and logic operations} \hfill \textcolor{blue}{\textbf{[Correct]}}
    \item Managing input and output operations
    \item Storing data and programs
\end{enumerate}

\noindent \textbf{Correct Answer:} \textcolor{blue}{\textbf{B}} \\[2pt]
\noindent \textbf{Explanation:} Trong mô hình kiến trúc Von Neumann, Đơn vị số học và logic (Arithmetic Logic Unit - ALU) là hệ thống con của CPU có chức năng chính là thực hiện tất cả các phép tính số học (cộng, trừ, nhân, chia) và các phép toán logic (AND, OR, NOT, so sánh). Việc điều khiển luồng thực thi chương trình thuộc về Đơn vị điều khiển (Control Unit - CU), việc lưu trữ thuộc về Bộ nhớ (Memory), và nhập xuất do khối I/O quản lý. Do đó, vai trò chính của ALU là thực hiện các phép toán số học và logic (đáp án B).

\par\vspace{0.3cm}
\hrule
\vspace{0.3cm}
\end{minipage}

\begin{minipage}{\textwidth}
\subsection*{Question 431 \quad \footnotesize{[0431\_q431\_Question\_431.png]}}
\noindent \textbf{Which component of the CPU is responsible for performing arithmetic and logical operations?}

\begin{enumerate}[label=\Alph*.]
    \item Control Unit
    \item Memory Unit
    \item \textbf{Arithmetic Logic Unit (ALU)} \hfill \textcolor{blue}{\textbf{[Correct]}}
    \item Input/Output Unit
\end{enumerate}

\noindent \textbf{Correct Answer:} \textcolor{blue}{\textbf{C}} \\[2pt]
\noindent \textbf{Explanation:} Khối tính toán số học và logic (ALU - Arithmetic Logic Unit) là bộ phận cốt lõi bên trong CPU chịu trách nhiệm xử lý các phép toán số học cơ bản và các phép so sánh logic theo lệnh của chương trình. Khối điều khiển (Control Unit) đóng vai trò giải mã lệnh và điều phối các bộ phận, còn Khối bộ nhớ (Memory Unit/Registers) chỉ dùng để lưu trữ dữ liệu tạm thời. Vì vậy, thành phần thực hiện các phép toán số học và logic chính là ALU (đáp án C).

\par\vspace{0.3cm}
\hrule
\vspace{0.3cm}
\end{minipage}

\begin{minipage}{\textwidth}
\subsection*{Question 432 \quad \footnotesize{[0432\_q432\_Question\_432.png]}}
\noindent \textbf{What does the term 'bus' refer to in a computer system? SÓ DÁP ÁN CÃN CHN}

\begin{enumerate}[label=\Alph*.]
    \item \textbf{A physical path used for data transmission between components} \hfill \textcolor{blue}{\textbf{[Correct]}}
    \item A type of memory
    \item A backup storage system
    \item A processor's internal clock
\end{enumerate}

\noindent \textbf{Correct Answer:} \textcolor{blue}{\textbf{A}} \\[2pt]
\noindent \textbf{Explanation:} Trong kiến trúc máy tính, bus hệ thống là tập hợp các đường dây dẫn vật lý dùng để truyền tải dữ liệu, địa chỉ bộ nhớ và tín hiệu điều khiển giữa các thành phần phần cứng khác nhau (như CPU, RAM và các cổng I/O). Bus không phải là một dạng bộ nhớ, không phải thiết bị lưu trữ dự phòng và cũng không phải là xung nhịp đồng hồ của CPU. Do đó, định nghĩa chính xác của bus là đường dẫn vật lý dùng để truyền dữ liệu giữa các bộ phận (đáp án A).

\par\vspace{0.3cm}
\hrule
\vspace{0.3cm}
\end{minipage}

\begin{minipage}{\textwidth}
\subsection*{Question 433 \quad \footnotesize{[0433\_q433\_Question\_433.png]}}
\noindent \textbf{What is the name of the first electronic computer designed for a special-purpose?}

\begin{enumerate}[label=\Alph*.]
    \item Pascal
    \item Pascaline
    \item \textbf{ABC} \hfill \textcolor{blue}{\textbf{[Correct]}}
    \item ENIAC
\end{enumerate}

\noindent \textbf{Correct Answer:} \textcolor{blue}{\textbf{C}} \\[2pt]
\noindent \textbf{Explanation:} Máy tính Atanasoff-Berry Computer (ABC), được chế tạo bởi John Vincent Atanasoff và Clifford Berry giai đoạn 1937–1942, là máy tính điện tử kỹ thuật số đầu tiên được thiết kế cho mục đích chuyên dụng (special-purpose) nhằm giải hệ phương trình đại số tuyến tính. ENIAC được hoàn thành sau đó là máy tính điện tử đa năng đầu tiên, còn Pascaline là máy cộng cơ học bằng bánh răng của thế kỷ 17. Vì vậy, ABC là máy tính điện tử chuyên dụng đầu tiên trong lịch sử (đáp án C).

\par\vspace{0.3cm}
\hrule
\vspace{0.3cm}
\end{minipage}

\begin{minipage}{\textwidth}
\subsection*{Question 434 \quad \footnotesize{[0434\_q434\_Question\_434.png]}}
\noindent \textbf{Computers built on the Von Neumann model divide the computer hardware into subsystems}

\begin{enumerate}[label=\Alph*.]
    \item 3
    \item \textbf{4} \hfill \textcolor{blue}{\textbf{[Correct]}}
    \item 5
    \item 6
\end{enumerate}

\noindent \textbf{Correct Answer:} \textcolor{blue}{\textbf{B}} \\[2pt]
\noindent \textbf{Explanation:} Kiến trúc máy tính Von Neumann kinh điển phân chia phần cứng của một hệ thống máy tính thành 4 hệ thống con (subsystems) chính gồm: Bộ nhớ (Memory), Đơn vị số học logic (ALU), Đơn vị điều khiển (Control Unit), và Hệ thống nhập/xuất (Input/Output). ALU và Đơn vị điều khiển kết hợp với nhau tạo thành Bộ xử lý trung tâm (CPU). Do đó, máy tính xây dựng theo mô hình Von Neumann chia phần cứng thành 4 hệ thống con (đáp án B).

\par\vspace{0.3cm}
\hrule
\vspace{0.3cm}
\end{minipage}

\begin{minipage}{\textwidth}
\subsection*{Question 435 \quad \footnotesize{[0435\_q435\_Question\_435.png]}}
\noindent \textbf{Which subsystem is responsible for performing arithmetic and logic operations?}

\begin{enumerate}[label=\Alph*.]
    \item \textbf{Arithmetic Logic Unit} \hfill \textcolor{blue}{\textbf{[Correct]}}
    \item Input subsystem
    \item Output subsystem
    \item Storage subsystem
    \item None of the others
\end{enumerate}

\noindent \textbf{Correct Answer:} \textcolor{blue}{\textbf{A}} \\[2pt]
\noindent \textbf{Explanation:} Đơn vị số học và logic (ALU - Arithmetic Logic Unit) là hệ thống con đảm nhận việc thực hiện tất cả các phép toán số học và các phép tính logic trên dữ liệu nhị phân. Các hệ thống con khác như Input chịu trách nhiệm tiếp nhận dữ liệu từ người dùng/thiết bị ngoại vi, Output gửi kết quả ra ngoài, và Storage/Memory lưu trữ dữ liệu và chương trình. Do đó, hệ thống con thực hiện các phép toán số học và logic là ALU (đáp án A).

\par\vspace{0.3cm}
\hrule
\vspace{0.3cm}
\end{minipage}

\begin{minipage}{\textwidth}
\subsection*{Question 436 \quad \footnotesize{[0436\_q436\_Question\_436.png]}}
\noindent \textbf{In the TCP/IP model, which layer corresponds to both Data Link and Physical layers of the OSI model?}

\begin{enumerate}[label=\Alph*.]
    \item \textbf{Network Access Layer} \hfill \textcolor{blue}{\textbf{[Correct]}}
    \item Internet Layer
    \item Transport Layer
    \item Application Layer
\end{enumerate}

\noindent \textbf{Correct Answer:} \textcolor{blue}{\textbf{A}} \\[2pt]
\noindent \textbf{Explanation:} Trong bộ giao thức TCP/IP chuẩn 4 tầng, Tầng truy cập mạng (Network Access layer, còn gọi là Host-to-Network hay Network Interface) gộp chung chức năng của cả Tầng liên kết dữ liệu (Data Link) và Tầng vật lý (Physical) của mô hình tham chiếu OSI 7 tầng. Tầng này chịu trách nhiệm định dạng khung dữ liệu, quản lý địa chỉ MAC và chuyển đổi tín hiệu vật lý trên đường truyền. Tầng Internet trong TCP/IP tương ứng với tầng Network OSI, và Tầng Transport tương ứng với tầng Transport OSI. Vì vậy, Network Access Layer là đáp án đúng (A).

\par\vspace{0.3cm}
\hrule
\vspace{0.3cm}
\end{minipage}

\begin{minipage}{\textwidth}
\subsection*{Question 437 \quad \footnotesize{[0437\_q437\_Question\_437.png]}}
\noindent \textbf{What is the most correct statement about LAN topologies?}

\begin{enumerate}[label=\Alph*.]
    \item Bus topology is a sophisticated structure and robust with high traffic.
    \item Ring topology is a high risk for collision between 2 nodes.
    \item \textbf{Star topology is a structure that has a central switch.} \hfill \textcolor{blue}{\textbf{[Correct]}}
    \item Mesh topology is the most complicated structure.
\end{enumerate}

\noindent \textbf{Correct Answer:} \textcolor{blue}{\textbf{C}} \\[2pt]
\noindent \textbf{Explanation:} Trong mạng cục bộ (LAN), cấu trúc liên kết hình sao (Star topology) là mô hình trong đó mỗi nút mạng đều được kết nối trực tiếp đến một thiết bị điều khiển trung tâm như bộ chuyển mạch (switch) hoặc hub, giúp mạng dễ quản lý và một sự cố đứt cáp không làm ảnh hưởng đến các máy còn lại (đáp án C). Phát biểu A sai vì Bus topology có cấu trúc đơn giản chứ không tinh vi (sophisticated) và dễ bị suy giảm hiệu năng khi lưu lượng cao; phát biểu B sai vì cấu trúc Ring truyền dữ liệu tuần tự có kiểm soát nên hạn chế xung đột; còn Mesh topology chủ yếu phức tạp ở mạng diện rộng (WAN). Do đó, phát biểu C là câu trả lời chính xác nhất.

\par\vspace{0.3cm}
\hrule
\vspace{0.3cm}
\end{minipage}

\begin{minipage}{\textwidth}
\subsection*{Question 438 \quad \footnotesize{[0438\_q438\_Question\_438.png]}}
\noindent \textbf{Which of the following accurately describes the differences between TCP and UDP in the Transport layer of the TCP/IP model?}

\begin{enumerate}[label=\Alph*.]
    \item \textbf{TCP provides error recovery, congestion control, and flow control, while UDP provides none of these features} \hfill \textcolor{blue}{\textbf{[Correct]}}
    \item UDP is connection-oriented and ensures data reliability, while TCP is connectionless and faster
    \item Both TCP and UDP are connection-oriented protocols but differ in flow control mechanisms
    \item TCP is primarily used for real-time communication, while UDP is used for file transfers
\end{enumerate}

\noindent \textbf{Correct Answer:} \textcolor{blue}{\textbf{A}} \\[2pt]
\noindent \textbf{Explanation:} TCP (Transmission Control Protocol) là giao thức hướng kết nối, đảm bảo truyền dữ liệu tin cậy thông qua các cơ chế phát hiện và phục hồi lỗi (error recovery), kiểm soát lưu lượng (flow control), và kiểm soát tắc nghẽn (congestion control); trong khi đó UDP (User Datagram Protocol) là giao thức phi kết nối đơn giản, nhanh và không cung cấp bất kỳ cơ chế kiểm soát nào kể trên (đáp án A). Phát biểu B đảo ngược bản chất của hai giao thức, C sai vì UDP không phải giao thức hướng kết nối, và D sai vì UDP thường được ưu tiên cho các ứng dụng thời gian thực nhờ độ trễ thấp.

\par\vspace{0.3cm}
\hrule
\vspace{0.3cm}
\end{minipage}

\begin{minipage}{\textwidth}
\subsection*{Question 439 \quad \footnotesize{[0439\_q439\_Question\_439.png]}}
\noindent \textbf{Router is operating at layer in the TCP/IP protocol suite}

\begin{enumerate}[label=\Alph*.]
    \item Transport
    \item \textbf{Network} \hfill \textcolor{blue}{\textbf{[Correct]}}
    \item Data link
    \item Physical
\end{enumerate}

\noindent \textbf{Correct Answer:} \textcolor{blue}{\textbf{B}} \\[2pt]
\noindent \textbf{Explanation:} Router (bộ định tuyến) là thiết bị hoạt động tại Tầng mạng (Network layer / Internet layer trong mô hình TCP/IP). Chức năng chính của router là đọc địa chỉ IP đích của gói tin (packet), tra cứu bảng định tuyến và chọn đường dẫn tối ưu để chuyển tiếp gói tin giữa các mạng khác nhau. Các thiết bị như Switch thông thường hoạt động ở tầng Data link (tầng 2), còn Hub và Repeater hoạt động ở tầng Physical (tầng 1). Do đó, router hoạt động ở Network layer (đáp án B).

\par\vspace{0.3cm}
\hrule
\vspace{0.3cm}
\end{minipage}

\begin{minipage}{\textwidth}
\subsection*{Question 440 \quad \footnotesize{[0440\_q440\_Question\_440.png]}}
\noindent \textbf{What is the minimum number of octal digits required to represent the binary number 1111101010?}

\begin{enumerate}[label=\Alph*.]
    \item 3
    \item \textbf{4} \hfill \textcolor{blue}{\textbf{[Correct]}}
    \item 5
    \item 6
\end{enumerate}

\noindent \textbf{Correct Answer:} \textcolor{blue}{\textbf{B}} \\[2pt]
\noindent \textbf{Explanation:} Trong hệ bát phân (Octal - cơ số 8), mỗi chữ số bát phân biểu diễn tương ứng đúng 3 bit nhị phân vì 2\textasciicircum{}3 = 8. Số nhị phân 1111101010 có tổng cộng 10 bit, khi chia thành các nhóm 3 bit từ phải sang trái ta được: (001)(111)(101)(010), gồm 4 nhóm tương ứng với 4 chữ số bát phân (giá trị là 1752 trong hệ bát phân). Do đó, cần tối thiểu 4 chữ số bát phân để biểu diễn số nhị phân này (đáp án B).

\par\vspace{0.3cm}
\hrule
\vspace{0.3cm}
\end{minipage}

\begin{minipage}{\textwidth}
\subsection*{Question 441 \quad \footnotesize{[0441\_q441\_Question\_441.png]}}
\noindent \textbf{What distinguishes positional number systems from nonpositional systems?}

\begin{enumerate}[label=\Alph*.]
    \item They represent values using symbols only.
    \item \textbf{The position of a digit determines its value.} \hfill \textcolor{blue}{\textbf{[Correct]}}
    \item They use Roman numerals.
    \item They are limited to base 10.
\end{enumerate}

\noindent \textbf{Correct Answer:} \textcolor{blue}{\textbf{B}} \\[2pt]
\noindent \textbf{Explanation:} Đặc điểm cốt lõi phân biệt hệ thống số có trọng số vị trí (positional number system) với hệ thống số không theo vị trí là giá trị thực tế của một chữ số phụ thuộc hoàn toàn vào vị trí của nó trong số đó (với mỗi vị trí nhân với lũy thừa tương ứng của cơ số). Ví dụ, số 1 trong '10' có giá trị khác hoàn toàn số 1 trong '1'. Ngược lại, trong hệ thống không theo vị trí như chữ số La Mã (Roman numerals), mỗi ký hiệu biểu thị một giá trị cố định bất kể đứng ở đâu. Do đó, đáp án đúng là B.

\par\vspace{0.3cm}
\hrule
\vspace{0.3cm}
\end{minipage}

\begin{minipage}{\textwidth}
\subsection*{Question 442 \quad \footnotesize{[0442\_q442\_Question\_442.png]}}
\noindent \textbf{Convert the hexadecimal number C9 to binary.}

\begin{enumerate}[label=\Alph*.]
    \item 11001101
    \item \textbf{11001001} \hfill \textcolor{blue}{\textbf{[Correct]}}
    \item 10101101
    \item 11101001
\end{enumerate}

\noindent \textbf{Correct Answer:} \textcolor{blue}{\textbf{B}} \\[2pt]
\noindent \textbf{Explanation:} Để chuyển đổi một số từ hệ thập lục phân (hexadecimal - cơ số 16) sang hệ nhị phân (binary - cơ số 2), ta thay thế từng ký số hex bằng nhóm 4 bit nhị phân tương đương. Ký số 'C' tương ứng với giá trị 12 trong hệ thập phân, có biểu diễn 4 bit là 1100; còn ký số '9' có biểu diễn 4 bit là 1001. Ghép hai nhóm lại, ta thu được chuỗi nhị phân hoàn chỉnh là 11001001 (đáp án B).

\par\vspace{0.3cm}
\hrule
\vspace{0.3cm}
\end{minipage}

\begin{minipage}{\textwidth}
\subsection*{Question 443 \quad \footnotesize{[0443\_q443\_Question\_443.png]}}
\noindent \textbf{In the octal system, the base is:}

\begin{enumerate}[label=\Alph*.]
    \item 2
    \item \textbf{8} \hfill \textcolor{blue}{\textbf{[Correct]}}
    \item 10
    \item 16
\end{enumerate}

\noindent \textbf{Correct Answer:} \textcolor{blue}{\textbf{B}} \\[2pt]
\noindent \textbf{Explanation:} Hệ bát phân (octal system) là hệ thống đếm có cơ số (base hoặc radix) bằng 8, sử dụng tập hợp 8 chữ số từ 0 đến 7 để biểu diễn giá trị. Trong khoa học máy tính, hệ nhị phân có cơ số 2, hệ thập phân có cơ số 10, và hệ thập lục phân có cơ số 16. Do đó, cơ số của hệ bát phân chính là 8 (đáp án B).

\par\vspace{0.3cm}
\hrule
\vspace{0.3cm}
\end{minipage}

\begin{minipage}{\textwidth}
\subsection*{Question 444 \quad \footnotesize{[0444\_q444\_Question\_444.png]}}
\noindent \textbf{If a byte represents a binary number, how many different values can it represent?}

\begin{enumerate}[label=\Alph*.]
    \item 64
    \item 128
    \item \textbf{256} \hfill \textcolor{blue}{\textbf{[Correct]}}
    \item 512
\end{enumerate}

\noindent \textbf{Correct Answer:} \textcolor{blue}{\textbf{C}} \\[2pt]
\noindent \textbf{Explanation:} Một byte là đơn vị thông tin tiêu chuẩn gồm đúng 8 bit. Mỗi bit chỉ có thể nhận 1 trong 2 giá trị nhị phân là 0 hoặc 1, do đó một chuỗi 8 bit có thể tạo ra tổng cộng 2\textasciicircum{}8 = 256 tổ hợp giá trị khác nhau (từ 0 đến 255 đối với số nguyên không dấu, hoặc từ -128 đến 127 đối với số nguyên có dấu biểu diễn bù 2). Vì vậy, một byte có thể biểu diễn được 256 giá trị khác nhau (đáp án C).

\par\vspace{0.3cm}
\hrule
\vspace{0.3cm}
\end{minipage}

\begin{minipage}{\textwidth}
\subsection*{Question 445 \quad \footnotesize{[0445\_q445\_Question\_445.png]}}
\noindent \textbf{In a binary positional number system, how much do consecutive digits typically differ?}

\begin{enumerate}[label=\Alph*.]
    \item 1 unit
    \item \textbf{2 units} \hfill \textcolor{blue}{\textbf{[Correct]}}
    \item 5 units
    \item 10 units
\end{enumerate}

\noindent \textbf{Correct Answer:} \textcolor{blue}{\textbf{B}} \\[2pt]
\noindent \textbf{Explanation:} Trong một hệ thống số có trọng số theo vị trí (positional system), trọng số của mỗi vị trí kế tiếp nhau luôn gấp cơ số (base) lần so với vị trí liền kề bên phải nó. Vì hệ nhị phân có cơ số bằng 2, nên giá trị trọng số của hai vị trí liên tiếp (ví dụ 2\textasciicircum{}0, 2\textasciicircum{}1, 2\textasciicircum{}2, ...) luôn chênh lệch nhau theo cấp số nhân là 2 đơn vị (2 units). Do đó, đáp án đúng là 2 units (B).

\par\vspace{0.3cm}
\hrule
\vspace{0.3cm}
\end{minipage}

\begin{minipage}{\textwidth}
\subsection*{Question 446 \quad \footnotesize{[0446\_q446\_Question\_446.png]}}
\noindent \textbf{Which of the following is a common definition of an operating system?}

\begin{enumerate}[label=\Alph*.]
    \item \textbf{An interface between hardware and users} \hfill \textcolor{blue}{\textbf{[Correct]}}
    \item A specific application program
    \item A document outlining system policies
    \item A type of computer hardware
    \item None of the others
\end{enumerate}

\noindent \textbf{Correct Answer:} \textcolor{blue}{\textbf{A}} \\[2pt]
\noindent \textbf{Explanation:} Hệ điều hành (Operating System - OS) là một chương trình phần mềm hệ thống cơ bản đóng vai trò làm giao diện/cầu nối (interface) trung gian giữa phần cứng máy tính và người dùng hoặc các ứng dụng. Nhiệm vụ của hệ điều hành là trừu tượng hóa các phần cứng phức tạp và phân bổ tài nguyên hệ thống một cách hiệu quả và an toàn. Hệ điều hành không phải là một phần cứng máy tính và cũng không phải là một ứng dụng người dùng đơn thuần. Do đó, định nghĩa đúng nhất là giao diện giữa phần cứng và người dùng (đáp án A).

\par\vspace{0.3cm}
\hrule
\vspace{0.3cm}
\end{minipage}

\begin{minipage}{\textwidth}
\subsection*{Question 447 \quad \footnotesize{[0447\_q447\_Question\_447.png]}}
\noindent \textbf{Which memory management technique divides the memory into variable-sized segments?}

\begin{enumerate}[label=\Alph*.]
    \item Paging
    \item \textbf{Segmentation} \hfill \textcolor{blue}{\textbf{[Correct]}}
    \item Virtual memory
    \item Swapping
\end{enumerate}

\noindent \textbf{Correct Answer:} \textcolor{blue}{\textbf{B}} \\[2pt]
\noindent \textbf{Explanation:} Kỹ thuật phân đoạn (Segmentation) là phương pháp quản lý bộ nhớ trong đó không gian địa chỉ logic của chương trình được chia thành các đoạn (segments) có kích thước biến thiên (variable-sized) phù hợp với cấu trúc ngữ nghĩa của chương trình như các hàm, thủ tục hoặc mảng dữ liệu. Ngược lại, kỹ thuật phân trang (Paging) chia bộ nhớ thành các khối có kích thước cố định bằng nhau (pages và frames). Vì câu hỏi yêu cầu kỹ thuật chia bộ nhớ thành các đoạn có kích thước thay đổi, đáp án đúng là Segmentation (B).

\par\vspace{0.3cm}
\hrule
\vspace{0.3cm}
\end{minipage}

\begin{minipage}{\textwidth}
\subsection*{Question 448 \quad \footnotesize{[0448\_q448\_Question\_448.png]}}
\noindent \textbf{The scheme of allocating memory in fixed-sized blocks is known as:}

\begin{enumerate}[label=\Alph*.]
    \item \textbf{Paging} \hfill \textcolor{blue}{\textbf{[Correct]}}
    \item Swapping
    \item Segmentation
    \item Fragmentation
\end{enumerate}

\noindent \textbf{Correct Answer:} \textcolor{blue}{\textbf{A}} \\[2pt]
\noindent \textbf{Explanation:} Cơ chế phân bổ bộ nhớ thành các khối có kích thước cố định bằng nhau được gọi là kỹ thuật phân trang (Paging), trong đó bộ nhớ vật lý được chia thành các khung trang (frames) và bộ nhớ logic được chia thành các trang (pages) có cùng kích cỡ. Việc sử dụng các khối kích thước cố định giúp hệ thống loại bỏ hiện tượng phân mảnh ngoại vi (external fragmentation). Ngược lại, phân đoạn (Segmentation) sử dụng các khối có kích thước biến thiên, còn Swapping là cơ chế hoán chuyển tiến trình với bộ nhớ phụ. Do đó, đáp án đúng là Paging (A).

\par\vspace{0.3cm}
\hrule
\vspace{0.3cm}
\end{minipage}

\begin{minipage}{\textwidth}
\subsection*{Question 449 \quad \footnotesize{[0449\_q449\_Question\_449.png]}}
\noindent \textbf{In a multiprogramming system, how are multiple programs typically managed in memory and executed?}

\begin{enumerate}[label=\Alph*.]
    \item \textbf{They are loaded into memory simultaneously and executed concurrently.} \hfill \textcolor{blue}{\textbf{[Correct]}}
    \item They are loaded one after another sequentially.
    \item They run sequentially, with each program completing before the next one starts.
    \item They are stored in secondary storage until needed for execution.
\end{enumerate}

\noindent \textbf{Correct Answer:} \textcolor{blue}{\textbf{A}} \\[2pt]
\noindent \textbf{Explanation:} Trong hệ thống đa chương (multiprogramming system), nhiều chương trình được nạp đồng thời vào bộ nhớ chính và được CPU thực thi đồng hành/luân phiên (concurrently). Khi một chương trình tạm dừng để chờ các thao tác nhập xuất (I/O) hoàn thành, hệ điều hành sẽ lập tức chuyển CPU sang thực thi chương trình khác nhằm tối đa hóa hiệu suất sử dụng bộ vi xử lý. Các phương án thực thi tuần tự lần lượt từng chương trình (B và C) mô tả hệ thống đơn chương (uniprogramming) truyền thống. Do đó, đáp án A là hoàn toàn chính xác.

\par\vspace{0.3cm}
\hrule
\vspace{0.3cm}
\end{minipage}

\begin{minipage}{\textwidth}
\subsection*{Question 450 \quad \footnotesize{[0450\_q450\_Question\_450.png]}}
\noindent \textbf{What is the primary purpose of a metadata-based file manager?}

\begin{enumerate}[label=\Alph*.]
    \item To organize files based on their content type
    \item \textbf{To store additional information about files, such as author, creation date, and keywords} \hfill \textcolor{blue}{\textbf{[Correct]}}
    \item To encrypt files for enhanced security
    \item To compress files to save disk space
\end{enumerate}

\noindent \textbf{Correct Answer:} \textcolor{blue}{\textbf{B}} \\[2pt]
\noindent \textbf{Explanation:} Siêu dữ liệu (metadata) là dữ liệu cung cấp thông tin mô tả ngữ cảnh về dữ liệu chính. Mục đích trọng tâm của trình quản lý tệp dựa trên siêu dữ liệu (metadata-based file manager) là lưu trữ và quản lý các thuộc tính mở rộng của tệp như tác giả, ngày tạo, thẻ đánh dấu (tags) và từ khóa tìm kiếm, giúp người dùng dễ dàng tìm kiếm và phân loại tệp tin linh hoạt hơn phương pháp tổ chức thư mục phân cấp đơn thuần. Mặc dù hệ thống tệp có thể hỗ trợ nén hoặc mã hóa, đó không phải là mục đích cốt lõi của việc quản lý dựa trên metadata. Vì vậy, đáp án B là chính xác.

\par\vspace{0.3cm}
\hrule
\vspace{0.3cm}
\end{minipage}

\begin{minipage}{\textwidth}
\subsection*{Question 451 \quad \footnotesize{[0451\_q451\_Question\_451.png]}}
\noindent \textbf{Which of the following is a responsibility of the process manager in an operating system?}

\begin{enumerate}[label=\Alph*.]
    \item Managing the communication between hardware and software
    \item \textbf{Scheduling processes for execution} \hfill \textcolor{blue}{\textbf{[Correct]}}
    \item Encrypting and decrypting data files
    \item Loading the operating system during boot
\end{enumerate}

\noindent \textbf{Correct Answer:} \textcolor{blue}{\textbf{B}} \\[2pt]
\noindent \textbf{Explanation:} Bộ quản lý tiến trình (Process Manager) trong hệ điều hành chịu trách nhiệm quản lý vòng đời của các tiến trình, bao gồm việc tạo mới, giám sát trạng thái, kết thúc và đặc biệt là lập lịch điều phối tiến trình (scheduling processes) để phân bổ thời gian sử dụng CPU một cách hợp lý và công bằng. Việc quản lý giao tiếp phần cứng thuộc về trình điều khiển thiết bị (device drivers), còn việc tải hệ điều hành lúc khởi động thuộc về chương trình bootloader. Do đó, lập lịch thực thi tiến trình là trách nhiệm của process manager (đáp án B).

\par\vspace{0.3cm}
\hrule
\vspace{0.3cm}
\end{minipage}

\begin{minipage}{\textwidth}
\subsection*{Question 452 \quad \footnotesize{[0452\_q452\_Question\_452.png]}}
\noindent \textbf{What is the role of data in AI?}

\begin{enumerate}[label=\Alph*.]
    \item To test software bugs
    \item \textbf{To help Al systems learn and improve predictions} \hfill \textcolor{blue}{\textbf{[Correct]}}
    \item To replace programming languages
    \item To make Al systems autonomous without training
\end{enumerate}

\noindent \textbf{Correct Answer:} \textcolor{blue}{\textbf{B}} \\[2pt]
\noindent \textbf{Explanation:} Trong trí tuệ nhân tạo (AI) và học máy (Machine Learning), dữ liệu đóng vai trò là nguồn nguyên liệu cốt lõi giúp các mô hình học hỏi các quy luật, đặc trưng và tri thức để đưa ra dự đoán hoặc quyết định chính xác hơn. Mô hình AI không thể tự động hóa hoặc suy luận tốt nếu không được huấn luyện trên tập dữ liệu chất lượng. Dữ liệu không nhằm mục đích thay thế các ngôn ngữ lập trình hay chỉ dùng để kiểm thử lỗi phần mềm đơn thuần. Do đó, vai trò chính của dữ liệu là giúp hệ thống AI học tập và cải thiện khả năng dự đoán (đáp án B).

\par\vspace{0.3cm}
\hrule
\vspace{0.3cm}
\end{minipage}

\begin{minipage}{\textwidth}
\subsection*{Question 453 \quad \footnotesize{[0453\_q453\_Question\_453.png]}}
\noindent \textbf{To perceive an object in an image, which of the following algorithm can be used?}

\begin{enumerate}[label=\Alph*.]
    \item \textbf{Edge detection} \hfill \textcolor{blue}{\textbf{[Correct]}}
    \item Chatbot
    \item Voice emotional analytic
    \item Text translation
\end{enumerate}

\noindent \textbf{Correct Answer:} \textcolor{blue}{\textbf{A}} \\[2pt]
\noindent \textbf{Explanation:} Phát hiện biên/cạnh (Edge detection) là một kỹ thuật thị giác máy tính và xử lý ảnh kinh điển dùng để xác định các đường ranh giới và đường nét của các đối tượng trong ảnh dựa trên sự biến thiên đột ngột về mức sáng của các pixel, từ đó hỗ trợ nhận dạng và phân biệt các vật thể (đáp án A). Ngược lại, Chatbot và Dịch văn bản thuộc lĩnh vực xử lý ngôn ngữ tự nhiên (NLP), còn Phân tích cảm xúc giọng nói thuộc lĩnh vực xử lý âm thanh. Vì vậy, Edge detection là thuật toán duy nhất trong các lựa chọn dùng để nhận biết đối tượng trong hình ảnh.

\par\vspace{0.3cm}
\hrule
\vspace{0.3cm}
\end{minipage}

\begin{minipage}{\textwidth}
\subsection*{Question 454 \quad \footnotesize{[0454\_q454\_Question\_454.png]}}
\noindent \textbf{What is the correct answer for an expert system?}

\begin{enumerate}[label=\Alph*.]
    \item Machines can learn knowledge without instructions from experts
    \item Experts use machines to talk with remote patients
    \item Machine use searching algorithms on databases to respond to patients
    \item \textbf{Machine learns knowledge given by experts and responds to patients} \hfill \textcolor{blue}{\textbf{[Correct]}}
    \item None of the others
\end{enumerate}

\noindent \textbf{Correct Answer:} \textcolor{blue}{\textbf{D}} \\[2pt]
\noindent \textbf{Explanation:} Hệ chuyên gia (Expert System) là một ứng dụng trí tuệ nhân tạo được xây dựng dựa trên cơ sở tri thức thu thập từ các chuyên gia con người trong một lĩnh vực cụ thể (knowledge base) kết hợp với động cơ suy luận (inference engine) để tư vấn và đưa ra các phản hồi, giải pháp chuyên môn cho người dùng (ví dụ như hỗ trợ bác sĩ chẩn đoán bệnh cho bệnh nhân). Hệ chuyên gia cổ điển không tự học tri thức mà không có sự hướng dẫn của chuyên gia, cũng không phải là phần mềm liên lạc từ xa. Do đó, mô tả đúng nhất là máy móc học/tiếp nhận tri thức do chuyên gia cung cấp và phản hồi cho bệnh nhân (đáp án D).

\par\vspace{0.3cm}
\hrule
\vspace{0.3cm}
\end{minipage}

\begin{minipage}{\textwidth}
\subsection*{Question 455 \quad \footnotesize{[0455\_q455\_Question\_455.png]}}
\noindent \textbf{is an artificial neuron similar to a single biological neuron.}

\begin{enumerate}[label=\Alph*.]
    \item An expert system
    \item An intelligent agent
    \item \textbf{A perceptron} \hfill \textcolor{blue}{\textbf{[Correct]}}
    \item Perception
\end{enumerate}

\noindent \textbf{Correct Answer:} \textcolor{blue}{\textbf{C}} \\[2pt]
\noindent \textbf{Explanation:} Perceptron là mô hình toán học sơ khai của một neuron nhân tạo đơn lẻ, được thiết kế mô phỏng lại cách thức hoạt động của một tế bào thần kinh sinh học (biological neuron). Mô hình này tiếp nhận các tín hiệu đầu vào có gắn trọng số, tính tổng tuyến tính và đi qua hàm kích hoạt để đưa ra tín hiệu đầu ra nhị phân. Các khái niệm như hệ chuyên gia (expert system) hay tác tử thông minh (intelligent agent) là các kiến trúc hệ thống AI bậc cao, không phải là một neuron nhân tạo đơn lẻ. Vì vậy, Perceptron là đáp án chính xác (C).

\par\vspace{0.3cm}
\hrule
\vspace{0.3cm}
\end{minipage}


\end{document}
